\documentclass[10pt]{article}

% Importação do preâmbulo customizado
% --- Pacotes de Geometria e Tipografia ---
\usepackage[a4paper,lmargin=1.5cm,rmargin=1.5cm,tmargin=1.2cm,bmargin=1.5cm]{geometry}
\usepackage[brazilian]{babel}
\usepackage[T1]{fontenc}
\usepackage{microtype}

% --- Matemática e Símbolos ---
\usepackage{amsmath,amsfonts,amssymb}
\usepackage{mathtools}
\usepackage{latexsym}
\usepackage{esint}

% --- Gráficos e Diagramas ---
\usepackage{graphicx}
\usepackage{tikz}
\usetikzlibrary{calc,arrows.meta,shadings}
\usepackage{pgfplots}
\pgfplotsset{compat=1.18}

% --- Elementos Visuais e Layout ---
\usepackage[most]{tcolorbox}
\usepackage{enumitem}
\usepackage[normalem]{ulem}
\usepackage[Smaller]{cancel}
\usepackage{xcolor}
\usepackage[nodisplayskipstretch]{setspace}

% --- Navegação e Links ---
\usepackage[colorlinks=true, linkcolor=black, urlcolor=blue]{hyperref}
\usepackage{bookmark}

% --- Metadados do Documento ---
\hypersetup{
	pdftitle={Solucionário: Variedades Diferenciáveis},
	pdfauthor={Anselmo Ribeiro Lopes},
	pdfkeywords={Variedades Diferenciáveis, Geometria Diferencial, PPGMAT UFCG},
	pdfcreator={TeXstudio},
	pdfproducer={LuaLaTeX com motor LuaHBTeX},
	pdfsubject={Licenciado sob CC BY-NC-SA 4.0. Este material é um registro acadêmico da disciplina de Variedades Diferenciáveis - PPGMAT/UFCG.},
	pdfstartview={FitH},
	pdflang={pt-BR},
	hidelinks,
}

% --- Configurações Globais de Estilo ---
\setstretch{1.3}
\setlength{\parindent}{0pt}
\relpenalty=10000
\binoppenalty=10000
\everymath{\displaystyle}
\pagestyle{empty}
\renewcommand{\CancelColor}{\color{red}}
%\setlength{\overfullrule}{5pt} % Desenha uma barra de 5pt no final de hboxes

% --- Contador e Comandos Customizados do Dr. Geometros ---
\newcounter{numquestao}

% Comando \questao: Enunciado em tcolorbox cinza
% Altere a definição da \questao para isso:
\newcommand{\questao}[1]{%
	\stepcounter{numquestao}%
	% 1. Criamos um destino único (ex: L1.1, L2.1) para a navegação não se perder
	\hypertarget{dest.\listid.\arabic{numquestao}}{}%
	% 2. Criamos o bookmark apontando explicitamente para esse destino
	\bookmark[level=1, dest={dest.\listid.\arabic{numquestao}}]{Questão \arabic{numquestao}}%
	\begin{tcolorbox}[
		enhanced,
		sharp corners,
		boxrule=0.5pt,
		colback=gray!6!white, 
		colframe=black,         
		left=1pt, right=1pt, top=1pt, bottom=1pt,
		width=1\textwidth,
		before skip=3pt,
		after skip=3pt,
		fontupper=\normalfont
		]%
		\mbox{\textbf{Questão \arabic{numquestao}.}}\hspace{0.3em}\ignorespaces#1%
	\end{tcolorbox}%
}

% Comando \resultadoextra: Teoremas ou definições auxiliares
\newcommand{\resultadoextra}[2]{%
	\begin{tcolorbox}[
		enhanced,
		sharp corners,
		boxrule=0.5pt,
		colback=gray!6!white, 
		colframe=black,         
		left=1pt, right=1pt, top=1pt, bottom=1pt,
		width=1\textwidth,
		before skip=3pt,
		after skip=3pt,
		fontupper=\normalfont
		]%
		\mbox{\textbf{#1}}\hspace{0.3em}\ignorespaces#2%
	\end{tcolorbox}%
}

% Comando \solucao: Resolução padrão terminando com quadradinho
\newcommand{\solucao}[1]{%
	\par\vspace{5pt}
	\textbf{\underline{Solução}:} #1 \hfill$\blacksquare$
	\par\vspace{5pt}
}

% Comando \solucaoalt: Soluções alternativas/estruturais
\newcommand{\solucaoalt}[1]{%
	\par\vspace{5pt}
	\textbf{\underline{Solução Alternativa}}: #1 \hfill$\blacksquare$
	\par\vspace{5pt}
}

% --- Operadores Matemáticos e Comandos Rápidos ---
\newcommand{\cqd}{\hfill$\blacksquare$}
\DeclareMathOperator{\supp}{supp}
\DeclareMathOperator{\sen}{sen}
\DeclareMathOperator{\Img}{Im}
\DeclareMathOperator{\dist}{dist}
\newcommand{\subsubset}{\subset\subset}
\DeclareMathOperator{\Inv}{Inv}

\begin{document}
	
	% --- Cabeçalho Institucional ---
	\begin{flushleft}
		\textbf{PROGRAMA DE PÓS-GRADUAÇÃO EM MATEMÁTICA - UFCG\\
			Unidade Acadêmica de Matemática - UAMat}\\
		Variedades Diferenciáveis - Professor Marco Antonio Lázaro Velásquez \hfill Semestre 2026.1 \\
		\textbf{Autor: Anselmo Ribeiro Lopes}
	\end{flushleft}
	
	\vspace{0.5cm}
	
	% --- Prefácio / Notas de Aprimoramento ---
	\begin{tcolorbox}[
		enhanced, sharp corners, boxrule=0.5pt,
		colback=gray!10!white, colframe=black, width=\textwidth,
		fontupper=\normalfont % Fonte normal conforme solicitado
		]
		Este material está sendo desenvolvido como parte do estudo e apoio as atividades da disciplina de \textbf{Variedades Diferenciáveis} no semestre 2026.1, vinculada ao \textbf{Programa de Pós-Graduação em Matemática da UFCG}, sob a orientação do \textbf{Professor Marco Antonio Lázaro Velásquez}. 
		
		O objetivo deste solucionário é reunir ideias, construções geométricas e demonstrações discutidas ao longo do curso. Por se tratar de um registro de acompanhamento acadêmico, as soluções e notas aqui apresentadas estão em \textbf{constante processo de aprimoramento e revisão}, buscando sempre maior rigor, clareza e aprofundamento nos temas da Geometria Diferencial.
	\end{tcolorbox}
	
	\vspace{0.8cm}
	
	\begin{center}
		% Título dentro de um fbox como solicitado
		\fbox{\huge \textbf{Solucionário de Variedades Diferenciáveis}}
	\end{center}
	
	\vspace{1cm}
	
	% --- Sumário ---
	\renewcommand{\contentsname}{Sumário}
	\tableofcontents
	\newpage
	
	% --- Inclusão das Listas ---
	
	% LISTA 01
	\newpage
	\def\listid{L1} % Identificador único para a Lista 1
	\phantomsection 
	\addcontentsline{toc}{section}{Soluções da Lista Exercício 1}
	\setcounter{numquestao}{0} 
	\begin{center}
	\fbox{\textbf{Resolução da Lista 01}}
\end{center}
\begin{tcolorbox}[
	enhanced,
	sharp corners,
	boxrule=0.5pt,
	colback=white, 
	colframe=black,         
	width=1\textwidth,
	%fontupper=\itshape
	]%
	\textbf{Resumo:} Este documento apresenta a resolução detalhada da primeira lista de exercícios da disciplina de Variedades Diferenciáveis. O conteúdo aborda os fundamentos da topologia e da geometria diferencial, incluindo a construção de atlas suaves e maximais, a análise de conectividade em variedades, a estrutura de subvariedades suaves via Teorema do Valor Regular, a geometria de esferas ($\mathbb{S}^n$) e espaços projetivos ($\mathbb{RP}^n$), além de uma introdução formal aos Grupos de Lie e suas propriedades algébricas e topológicas.
\end{tcolorbox}
%%%% QUESTÃO 01 %%%%
\questao{
	Seja $M^{n}$ uma $n$-variedade topológica. Mostre que:
	\begin{enumerate}[label=(\alph*)]
		\item Qualquer atlas suave para $M^{n}$ está contido em um atlas suave maximal (único).
		\item Dois atlas suaves para $M^{n}$ determinam o mesmo atlas suave maximal se, e somente se, a união desses atlas é um atlas suave.
	\end{enumerate}
}
\solucao{
	
	\textbf{Análise do Enunciado}
	\begin{itemize}
		\item \textbf{Hipóteses:} $M^n$ é uma $n$-variedade topológica.
		\item \textbf{Tese (a):} Provar a existência e unicidade de um atlas suave maximal que contenha qualquer atlas suave dado.
		\item \textbf{Tese (b):} Caracterizar a equivalência de atlas suaves através da suavidade de sua união.
	\end{itemize}
	\vspace{5pt}
	\textbf{(a)} Seja $\mathcal{A} = \{(U_\alpha, \varphi_\alpha)\}_{\alpha \in \Lambda}$ um atlas suave para $M^n$. Definimos a coleção $\overline{\mathcal{A}}$ como o conjunto de todas as cartas coordenadas que são suavemente compatíveis com todas as cartas de $\mathcal{A}$:
	\[ \overline{\mathcal{A}} = \{ (V, \psi) : (V, \psi) \text{ é suavemente compatível com toda carta } (U_\alpha, \varphi_\alpha) \in \mathcal{A} \}. \]
	
	Claramente $\mathcal{A} \subseteq \overline{\mathcal{A}}$, pois as cartas de um atlas suave são, por definição, compatíveis entre si. Para mostrar que $\overline{\mathcal{A}}$ é um atlas suave, devemos provar que quaisquer duas cartas $(V_1, \psi_1), (V_2, \psi_2) \in \overline{\mathcal{A}}$ são suavemente compatíveis. 
	
	Seja $p \in V_1 \cap V_2$. Como $\mathcal{A}$ é um atlas, existe uma carta $(U_\alpha, \varphi_\alpha) \in \mathcal{A}$ tal que $p \in U_\alpha$. Definimos o aberto $W = V_1 \cap V_2 \cap U_\alpha$. A aplicação de transição $\psi_2 \circ \psi_1^{-1}$ restrita ao domínio $\psi_1(W) \subset \mathbb{R}^n$ pode ser fatorada como:
	\[ \psi_2 \circ \psi_1^{-1} \Big|_{\psi_1(W)} = \underbrace{(\psi_2 \circ \varphi_\alpha^{-1})}_{\text{suave}} \circ \underbrace{(\varphi_\alpha \circ \psi_1^{-1})}_{\text{suave}}. \]
	
	Visto que ambas as funções no lado direito são suaves (pela definição de $\overline{\mathcal{A}}$), a sua composição também é suave em uma vizinhança de $\psi_1(p)$. Como isso é válido para todo ponto na interseção, concluímos que $\overline{\mathcal{A}}$ é um atlas suave.
	
	\textbf{Maximalidade:} Para mostrar que $\overline{\mathcal{A}}$ é maximal, seja $(V, \psi)$ uma carta suave compatível com todas as cartas de $\overline{\mathcal{A}}$. Como $\mathcal{A} \subseteq \overline{\mathcal{A}}$, em particular, $(V, \psi)$ deve ser compatível com todas as cartas de $\mathcal{A}$. Por definição de $\overline{\mathcal{A}}$, qualquer carta compatível com $\mathcal{A}$ já pertence a $\overline{\mathcal{A}}$. Logo, $(V, \psi) \in \overline{\mathcal{A}}$, o que prova que $\overline{\mathcal{A}}$ é maximal.
	
	\textbf{Unicidade:} Suponha que exista outro atlas maximal $\mathcal{C}$ tal que $\mathcal{A} \subseteq \mathcal{C}$. Como $\mathcal{C}$ é um atlas suave, todas as suas cartas são mutuamente compatíveis. Como $\mathcal{A} \subseteq \mathcal{C}$, todas as cartas de $\mathcal{C}$ são compatíveis com todas as cartas de $\mathcal{A}$. Pela definição de $\overline{\mathcal{A}}$, isto implica $\mathcal{C} \subseteq \overline{\mathcal{A}}$. Por maximalidade de $\mathcal{C}$, temos $\mathcal{C} = \overline{\mathcal{A}}$.
	
	\textbf{(b)} $(\implies)$ Suponha que $\mathcal{A}$ e $\mathcal{B}$ determinam o mesmo atlas suave maximal $\mathcal{M}$. Por definição, $\mathcal{A} \subseteq \mathcal{M}$ e $\mathcal{B} \subseteq \mathcal{M}$. Logo, $\mathcal{A} \cup \mathcal{B} \subseteq \mathcal{M}$. Como subconjuntos de cartas de um atlas suave são mutuamente compatíveis, $\mathcal{A} \cup \mathcal{B}$ é um atlas suave.
	
	$(\impliedby)$ Suponha agora que $\mathcal{A} \cup \mathcal{B}$ seja um atlas suave. Seja $\mathcal{M}$ o atlas suave maximal único que contém a união $\mathcal{A} \cup \mathcal{B}$ (cuja existência foi garantida no item (a)). 
	
	Como $\mathcal{A} \subseteq (\mathcal{A} \cup \mathcal{B}) \subseteq \mathcal{M}$, então $\mathcal{M}$ é um atlas suave maximal que contém $\mathcal{A}$. Pela unicidade, $\mathcal{A}_{max} = \mathcal{M}$. Da mesma forma, como $\mathcal{B} \subseteq (\mathcal{A} \cup \mathcal{B}) \subseteq \mathcal{M}$, temos $\mathcal{B}_{max} = \mathcal{M}$. Portanto, $\mathcal{A}_{max} = \mathcal{B}_{max}$.
}
%%%% QUESTÃO 02 %%%%
\questao{
	Mostre que uma $n$-variedade suave $M^{n}$ é conexa se, e somente se, ela é conexa por caminhos.
}
\solucao{
	
	\textbf{Análise do Enunciado}
	\begin{itemize}
		\item \textbf{Hipóteses:} $M^n$ é uma $n$-variedade suave.
		\item \textbf{Tese:} Provar a equivalência entre a conexidade topológica e a conexidade por caminhos.
	\end{itemize}
	\vspace{5pt}
	\textbf{($\Leftarrow$)} Suponha que $M$ seja conexa por caminhos. Provaremos que $M$ é conexa por contradição. Se $M$ fosse desconexa, existiriam abertos $U, V \subset M$ não vazios e disjuntos tais que $M = U \cup V$. Sejam $p \in U$ e $q \in V$. Por hipótese, existe $\gamma: [0,1] \to M$ contínua com $\gamma(0) = p$ e $\gamma(1) = q$.
	
	Os conjuntos $A = \gamma^{-1}(U)$ e $B = \gamma^{-1}(V)$ seriam abertos em $[0,1]$ pela continuidade de $\gamma$. Teríamos $A \cup B = [0,1]$, $A \cap B = \emptyset$, com $0 \in A$ e $1 \in B$ (logo ambos não vazios). Isso contradiz a conexidade de $[0,1]$. Portanto, $M$ é conexa.
	
	\textbf{($\Rightarrow$)} Suponha $M$ conexa. Fixemos $p \in M$ e definamos o conjunto de pontos que podem ser ligados a $p$ por um caminho contínuo:
	\[ A = \{ q \in M : \exists \, \sigma: [0,1] \to M \text{ contínua com } \sigma(0)=p \text{ e } \sigma(1)=q \}. \]
	
	\textbf{1. $A$ é aberto:} Seja $q \in A$. Existe uma vizinhança coordenada $U$ de $q$ homeomorfa a uma bola em $\mathbb{R}^n$, a qual é conexa por caminhos. Seja $\phi: U \to B^n$ o homeomorfismo. Para qualquer $q' \in U$, definimos o caminho interno $\alpha: [0,1] \to U$ por $\alpha(t) = \phi^{-1}\left((1-t)\phi(q) + t\phi(q')\right)$. Como $q \in A$, existe $\sigma: [0,1] \to M$ de $p$ a $q$. Definimos a concatenação $\beta: [0,1] \to M$:
	\[ \beta(t) = \begin{cases} \sigma(2t), & 0 \leq t \leq 1/2 \\ \alpha(2t-1), & 1/2 < t \leq 1 \end{cases} \]
	$\beta$ é contínua pelo Lema da Colagem ($\sigma(1)=q=\alpha(0)$) e liga $p$ a $q'$. Logo $U \subseteq A$, e $A$ é aberto.
	
	\textbf{2. $A$ é fechado:} Seja $r \in \overline{A}$. Seja $V$ uma vizinhança coordenada de $r$ (homeomorfa a uma bola). Pela definição de fecho, existe $s \in V \cap A$. Como $s \in A$, ele se liga a $p$. Como $s$ e $r$ estão na mesma vizinhança conexa por caminhos $V$, eles se ligam. Por concatenação, $r \in A$. Logo $\overline{A} = A$.
	
	Como $M$ é conexo e $A$ é aberto, fechado e não vazio, $A = M$.
}
%%%% QUESTÃO 03 %%%%
\questao{
	Seja $\mathbb{S}^{n}=\{(x^{1},...,x^{n+1})\in\mathbb{R}^{n+1}; (x^{1})^{2}+\dots+(x^{n+1})^{2}=1\}$ a esfera unitária. Prove que:
	\begin{enumerate}[label=(\alph*)]
		\item $\mathbb{S}^{n}$ é uma $n$-variedade suave, compacta e de Hausdorff com base enumerável exibindo um atlas suave sobre $\mathbb{S}^{n}$.
		\item $\mathbb{S}^{n}$ é a imagem inversa de um valor regular de uma função apropriada.
		\item Mostre que existe um atlas suave de $\mathbb{S}^{n}$ com apenas duas cartas.
		\item $\hat{E}$ possível obtermos um atlas suave para $\mathbb{S}^{n}$ com apenas uma carta?
	\end{enumerate}
}
\solucao{
	
	\textbf{Análise do Enunciado}
	\begin{itemize}
		\item \textbf{Hipóteses:} $\mathbb{S}^n$ definida como a esfera unitária em $\mathbb{R}^{n+1}$.
		\item \textbf{Tese (a):} Provar que $\mathbb{S}^n$ é uma $n$-variedade suave, compacta, Hausdorff e com base enumerável.
		\item \textbf{Tese (b):} Representar $\mathbb{S}^n$ como $f^{-1}(c)$ com $c$ valor regular.
		\item \textbf{Tese (c):} Construir um atlas de duas cartas.
		\item \textbf{Tese (d):} Provar a impossibilidade de um atlas de uma única carta.
	\end{itemize}
	\vspace{5pt}
	\textbf{(a)} Como $\mathbb{S}^n \subset \mathbb{R}^{n+1}$, então pela topologia induzida ela herda as propriedades de \textbf{Hausdorff} e \textbf{base enumerável}. 
	
	\textbf{1. Compacidade}
	
	Seja $f: \mathbb{R}^{n+1} \to \mathbb{R}$ dada por $f(x) = |x|^2$. Como $f$ é contínua e $\mathbb{S}^n = f^{-1}(1)$, a esfera é um conjunto fechado. Sendo também limitada ($|x|=1$), pelo Teorema de Heine-Borel, $\mathbb{S}^n$ é compacta.
	
	\textbf{2. Construção do Atlas das Calotas Polares}
	
	Para cada $i\in\{1,\ldots,n+1\}$, definimos os abertos $U_i^{\pm} = \{x \in \mathbb{S}^n : \pm x^i > 0\}$ e 
	\begin{eqnarray*}
		\varphi_i^{\pm}: U_i^{\pm}&\longrightarrow &\mathbb{R}^n\\
		(x^1, \dots,x^{n+1})&\longmapsto& (x^1, \dots, \widehat{x^i}, \dots, x^{n+1})
	\end{eqnarray*}
	onde, $\widehat{x^i}$ denota que $x_i$ é omitido.
	
	\textbf{Prova de Cobertura}: Seja $x = (x^1, \dots, x^{n+1}) \in \mathbb{S}^n$. Como $|x|^2 = 1$, é impossível que todas as coordenadas de $x$ sejam nulas. Portanto, existe pelo menos um $i \in \{1, \dots, n+1\}$ tal que $x^i \neq 0$. Se $x^i > 0$, então $x \in U_i^+$; se $x^i < 0$, então $x \in U_i^-$.	Isso implica que $\mathbb{S}^n = \bigcup_{i=1}^{n+1} (U_i^+ \cup U_i^-)$.
	
	Notemos que $\varphi_i^{\pm}$ são contínuas e além disso, $\varphi_i^{\pm}( U_i^{\pm})=\mathbb{B}^n$, onde $\mathbb{B}^n=\{x\in\mathbb{R}^n;|x|<1\}$. De fato, dado $x\in U_i^{\pm}$, temos que:
	$$|\varphi_i^{\pm}(x)|^2=(x^1)^2+\cdots+(x^{i-1})^2+(x^{i+1})^2+\cdots+(x^n)^2<(x^1)^2+\cdots+(x^{i-1})^2+(x^i)^2+(x^{i+1})^2+\cdots+(x^n)^2=1.$$ 
	
	\textbf{Cálculo da Inversa}: A inversa $(\varphi_i^{\pm})^{-1}: \mathbb{B}^n\rightarrow U_i^{\pm}$, é obtida isolando $x^i$:
	$$(x^i)^2 = 1 - \sum_{k \neq i} (x^k)^2 \implies x^i = \pm \sqrt{1 - |u|^2},$$ 
	donde
	$$(\varphi_i^{\pm})^{-1}(u)=(u^1, \dots, u^{i-1},\pm\sqrt{1-|u|^2}, u^{i+1}, \dots, u^n),$$ que também são contínuas.
	
	Além disso, temos que $(\varphi_i^{\pm})^{-1}(\mathbb{B}^n)=U_i^{\pm}$, pois
	$$|(\varphi_i^{\pm})^{-1}(u)|^2=(u^1)^2+\cdots+(u^{i-1})^2+1-|u|^2+(u^{i+1})^2+\cdots(u^n)^2=1.$$
	
	Assim, cada uma das $2(n+1)$ cartas $\varphi_i^{\pm}$ é um homeomorfismo de $U_i^{\pm}$ sobre o aberto $\mathbb{B}^n \subset \mathbb{R}^n$, donde $\{(U_i^{\pm}, \varphi_i^{\pm})\}$ constitui um atlas topológico para $\mathbb{S}^n$. Portanto, $\mathbb{S}^n$ é uma $n$-variedade topológica.
	
	\textbf{3. Suavidade das Transições}
	
	A transição $\psi = \varphi_j^{\pm} \circ (\varphi_i^{\pm})^{-1}$ (para $i < j$) em $u \in \varphi_i^+(U_i^+ \cap U_j^+)\subset\mathbb{B}^n$ resulta em:
	$$ \psi(u) = (u^1, \dots, \pm\sqrt{1-|u|^2}, \dots, \widehat{u^j}, \dots, u^n) $$
	Como $|u|^2 < 1$, então a raiz quadrada é suave em $(0, \infty)$, e portanto a transição é suave. De forma similar, vale quando $i>j$. Quando $i=j$, então $\varphi_j^{\pm}\circ(\varphi_i^{\pm})^{-1}=\operatorname{Id}_{\mathbb{B}^n}$ que é suave.
	
	Concluímos que, $\{U_i^{\pm},\varphi_i^{\pm}\}$ é um atlas suave para $\mathbb{S}^n$, e portanto a esfera é uma $n$-variedade suave com essa estrutura.
	
	\textbf{(b)} Consideremos a função diferenciável,
	\begin{eqnarray*}
		f: \mathbb{R}^{n+1}&\longrightarrow &\mathbb{R}\\
		x&\longmapsto& f(x)=\langle x,x\rangle=|x|^2.
	\end{eqnarray*}
	Temos que $\nabla f(x) = 2x$ que é não nulo para todo $x\neq \vec{0}$. Em particular, $1$ é valor regular de $f$, donde $\mathbb{S}^n=f^{-1}(1)$.
	
	\textbf{(c)} Exibiremos o atlas de duas cartas via Projeção Estereográfica. Sejam $U_N = \mathbb{S}^n \setminus \{N\}$ e $U_S = \mathbb{S}^n \setminus \{S\}$, onde $N = (0,\dots,1)$ e $S = (0,\dots,-1)$. Notemos que, $\mathbb{S}^n=U_N\cup U_S$, ou seja, temos uma cobertura.
	
	\textbf{1. Dedução Geométrica e Construção do Atlas Suave}
	
	Para qualquer ponto $x = (x^1, \dots, x^{n+1}) \in U_N$, a projeção estereográfica $\phi_N(x)$ é definida como a interseção da semirreta que parte de $N$ e passa por $x$ com o hiperplano $x_{n+1}=0$ (identificado com $\mathbb{R}^n$).		
	
	A parametrização da reta $Nx$ é dada por $L(t) = N + t(x - N) = (0, \dots, 0, 1) + t(x^1, \dots, x^n, x^{n+1}-1)$ e ela intercepta o hiperplano $x_{n+1}=0$ se, e somente se:
	$$1 + t(x^{n+1} - 1) = 0 \implies t = \frac{1}{1 - x^{n+1}}.$$
	
	\begin{figure}[htbp]
		\centering
		\begin{tikzpicture}[scale=1.8, line join=round, line cap=round]
			% --- Configurações de Estilo ---
			\tikzset{
				projection_inside/.style={thin, dashed, black!70}, 
				projection_outside/.style={thin, black}, 
				sphere_line/.style={thin, black},
				plane_line/.style={thin, blue!40!black}, 
				point/.style={circle, fill=black, inner sep=1pt},
				open_point/.style={circle, draw=black, fill=white, inner sep=0pt, minimum size=3pt},
				color_Int/.style={blue!80!black},
				color_Plane_Fill/.style={fill=blue!5, draw=none} 
			}
			
			% --- 1. O Plano R^n ---
			\fill[color_Plane_Fill] (-4, -0.8) -- (3, -0.8) -- (4, 1.5) -- (-3, 1.5) -- cycle;
			
			% Máscara para limpar o interior da esfera (apenas topo e equador)
			\fill[white] (1,0) arc (0:180:1) arc (180:360:1 and 0.3) -- cycle;
			
			% Rótulo R^n
			\node[color_Int, font=\small] at (-1.8,1.2) {$\text{Hiperplano } \{x_{n+1}=0\}=\mathbb{R}^n$};
			
			% --- 2. A Esfera S^n ---
			% Borda Superior sólida
			\draw[sphere_line] (1,0) arc (0:180:1); 
			
			% Equador (Interseção) em AZUL
			\draw[color_Int, thin] (1,0) arc (0:-180:1 and 0.3); % Frente
			\draw[color_Int, dashed, thin] (1,0) arc (0:180:1 and 0.3); % Atrás
			
			% BORDA INFERIOR DINÂMICA:
			% Parte 1: Tracejada (dentro do plano) de 180° até approx 243° (y = -0.8)
			\draw[dashed, sphere_line] (1,0) arc (0:180:1); % Apenas para referência, ignorar
			\draw[dashed, sphere_line] (-1,0) arc (180:230:1); 
			% Parte 2: SÓLIDA (fora do plano) de 243° até 307° (y = -0.8)
			\draw[sphere_line] (0.6, -0.8) arc (307:234:1); % Ajustado para a base
			% Correção manual do arco inferior para precisão visual:
			%\draw[sphere_line] (-0.6, -0.8) arc (243:307:1);
			% Parte 3: Tracejada (dentro do plano) de 307° até 360°
			\draw[dashed, sphere_line] (0.6, -0.8) arc (307:360:1);
			
			% --- 3. Polos ---
			\node[open_point, label={[font=\small, above right, yshift=-5pt]:$N=(0,\ldots,0,1)$}] (N) at (0,1) {};
			
			% --- 4. Eixo x_{n+1} ---
			\draw[sphere_line, -{stealth}] (0, 1.03) -- (0, 2);       % Sólido superior
			\draw[dashed, sphere_line] (0, 0.95) -- (0, -1);       % Tracejado dentro da esfera
			\draw[sphere_line] (0, -1) -- (0, -1.2);            % Sólido ao sair do plano
			\node[right, font=\small] at (0, 1.9) {$x_{n+1}$};
			
			% --- 5. Geometria da Projeção ---
			\coordinate (P_new) at (-0.7, 0.5);
			\node[point, label={[font=\small, below right]:$x$}] at (P_new) {};
			
			% Ponto phi_N(x)
			\coordinate (U) at (-1.85, -0.32); 
			\node[point, color_Int, label={[font=\small, color=blue!80!black, above left, yshift=-5pt]:$\phi_N(x)$}] at (U) {};
			
			% Reta terminando em y = -1.2
			\coordinate (End) at (-3.08, -1.2);
			
			% Reta: 
			\draw[projection_inside] (N) -- (P_new); 
			\draw[projection_outside] (P_new) -- (U); 
			
			% Divisão da última parte:
			% 1. De phi_N(x) até a borda inferior do plano (-2.52, -0.8) -> Tracejado
			\draw[projection_inside] (U) -- (-2.52, -0.8); 
			% 2. Da borda do plano até o ponto final (End) -> Sólido
			\draw[projection_outside] (-2.52, -0.8) -- (End); 
			
			% --- 6. Rótulos dos Abertos ---
			\node[black!80, font=\small, xshift=-2pt] at (1.05, 0.57) {$U_N$};
		\end{tikzpicture}
		\caption{Projeção Estereográfica a partir do polo norte}
	\end{figure}
	
	Substituindo $t$ nas primeiras $n$ coordenadas de $L(t)$, obtemos a aplicação $\phi_N: U_N \to \mathbb{R}^n$:
	$$\boxed{\phi_N(x^1, \dots, x^{n+1}) = \left( \frac{x^1}{1-x^{n+1}}, \dots, \frac{x^n}{1-x^{n+1}} \right)},$$ que é contínua, pois é uma função racional cujas componentes são contínuas (o denominador $1-x^{n+1}\neq 0$ em $U_N$).
	
	De forma análoga, para $x \in U_S$, a reta $L_S(t) = S + t(x-S) = (0,\dots,-1) + t(x^1, \dots, x^{n+1}+1)$ intercepta $x_{n+1}=0$ quando $$-1+t(x^{n+1}+1)=0 \implies t = \frac{1}{1+x^{n+1}}.$$ 
	Assim, $\phi_S: U_S \to \mathbb{R}^n$ é dada por:
	$$\boxed{\phi_S(x) = \left( \frac{x^1}{1+x^{n+1}}, \dots, \frac{x^n}{1+x^{n+1}} \right),}$$ que é contínua, pois é uma função racional cujas componentes são contínuas (o denominador $1+x^{n+1}\neq 0$ em $U_S$).
	
	Além disso, notemos que $\phi_S(x)=-\phi_N(-x)$. 
	
	\textbf{Cálculo da Inversa $\phi_N^{-1}$}: Seja $u \in \mathbb{R}^n$, queremos $x \in U_N$ tal que $u^i = \frac{x^i}{1-x^{n+1}}$, ou seja, 
	$$\boxed{x^i = u^i(1-x^{n+1}).}\quad (\ast)$$ 
	Elevando ao quadrado e somando de $i=1$ até $n$:
	$$\sum_{i=1}^n (x^i)^2 = \sum_{i=1}^n [u^i(1-x^{n+1})]^2 = (1-x^{n+1})^2 \sum_{i=1}^n (u^i)^2 = |u|^2 (1-x^{n+1})^2$$
	Como $x \in U_N$, temos $\sum_{i=1}^n (x^i)^2 = 1 - (x^{n+1})^2$. Igualando as expressões:
	\begin{align*}
		1 - (x^{n+1})^2 &= |u|^2 (1-x^{n+1})^2 \\
		(1-x^{n+1})(1+x^{n+1}) &= |u|^2 (1-x^{n+1})^2 \\
		1+x^{n+1} &= |u|^2 (1-x^{n+1}) \quad (\text{pois } x \in U_N \implies x^{n+1} \neq 1) \\
		x^{n+1}(1+|u|^2) &= |u|^2 - 1 \implies x^{n+1} = \frac{|u|^2-1}{|u|^2+1}
	\end{align*}
	Substituindo $x^{n+1}$  em $(\ast)$, obtemos que: $$x^i = u^i \left(1 - \frac{|u|^2-1}{|u|^2+1}\right) = u^i \left(\frac{|u|^2+1 - |u|^2 + 1}{|u|^2+1}\right) = \frac{2u^i}{|u|^2+1}.$$
	
	Logo, $\phi_N^{-1}:\mathbb{R}^n\setminus\{0\}\rightarrow U_N$ é dada por:
	$$\phi_N^{-1}(u)=\frac{\left(2u^1,\dots,2u^n,|u|^2-1\right)}{|u|^2+1}$$
	
	\textbf{2. Suavidade das Transições}
	
	A transição $\psi = \phi_S \circ \phi_N^{-1}$ é definida no domínio $\phi_N(U_N \cap U_S) = \mathbb{R}\setminus\{0\}$. Além disso,
	\begin{align*}
		\psi^i(u) = \frac{x^i}{1+x^{n+1}} = \frac{\frac{2u^i}{|u|^2+1}}{1 + \frac{|u|^2-1}{|u|^2+1}} = \frac{2u^i}{|u|^2+1+|u|^2-1} = \frac{2u^i}{2|u|^2} = \frac{u^i}{|u|^2}
	\end{align*}
	A expressão $\psi(u) = \frac{u}{|u|^2}$ é suave em $\mathbb{R}^n \setminus \{0\}$, concluindo que o atlas $\{(U_N,\phi_N),(U_S,\phi_S)\}$ é suave. Portanto, $\mathbb{S}^n$ é uma $n$-variedade com essa estrutura.
	
	\textbf{(d)} Não é possível. Se existisse uma única carta $\phi: \mathbb{S}^n \to \mathbb{R}^n$, ela seria um homeomorfismo sobre um aberto $V \subseteq \mathbb{R}^n$. Pelo item (a), $\mathbb{S}^n$ é compacto, logo $V$ deveria ser um subconjunto compacto e aberto de $\mathbb{R}^n$. Como $\mathbb{R}^n$ é conexo, o único aberto e fechado (compacto) é o conjunto vazio ou o próprio $\mathbb{R}^n$ (que não é compacto). Como $\mathbb{S}^n \neq \emptyset$, chegamos a uma contradição.
}
%%%% QUESTÃO 04 %%%%
\questao{Mostre que o espaço projetivo real $\mathbb{RP}^n$ é uma $n$-variedade suave, compacta e de Hausdorff com base enumerável exibindo um atlas suave sobre $\mathbb{RP}^n$.
}
\solucao{
	
	\textbf{Análise do Enunciado}
	\begin{itemize}
		\item \textbf{Hipóteses:} $\mathbb{RP}^n$ definido como o espaço quociente de $\mathbb{R}^{n+1} \setminus \{0\}$ por $x \sim \lambda x$.
		\item \textbf{Tese:} Provar que $\mathbb{RP}^n$ é uma $n$-variedade suave, compacta, Hausdorff e com base enumerável.
	\end{itemize}
	\vspace{5pt}
	Definimos o espaço projetivo real $\mathbb{RP}^n$ como o espaço quociente de $\mathbb{R}^{n+1} \setminus \{0\}$ pela relação de equivalência $\sim$, onde $x \sim y$ se e somente se $y = \lambda x$ para algum $\lambda \in \mathbb{R} \setminus \{0\}$. Denotamos por $[x]$ a classe de equivalência de $x$.
	
	Consideremos a esfera unitária $\mathbb{S}^n \subset \mathbb{R}^{n+1}$. Definimos a aplicação de projeção canônica $\pi: \mathbb{S}^n \to \mathbb{RP}^n$ por $\pi(x) = [x]$. 
	
	\textbf{Sobrejetividade de $\pi$:} 
	
	Seja $[v] \in \mathbb{RP}^n$ com $v \in \mathbb{R}^{n+1} \setminus \{0\}$ e consideremos o ponto $x = \frac{v}{|v|} \in \mathbb{S}^n$. Como $v = |v|x$ e $|v| \neq 0$, temos que $v \sim x$, donde $\pi(x) = [x] = [v]$. Isso prova que $\pi$ é sobrejetiva. Como a relação $\sim$ em $\mathbb{S}^n$ reduz-se a $x \sim -x$, identificamos $\mathbb{RP}^n$ com o quociente $\mathbb{S}^n / \{x \sim -x\}$.
	
	\textbf{1. Compacidade}
	
	Como provamos na \textbf{Questão 3(a)}, $\mathbb{S}^n$ é compacta e sendo e $\pi$ contínua e sobrejetiva, $\mathbb{RP}^n = \pi(\mathbb{S}^n)$ é a imagem contínua de um compacto, e portanto, compacta.
	
	\textbf{2. Hausdorff e Base Enumerável}
	
	Sabemos que, se $X$ é um espaço compacto Hausdorff e a relação de equivalência $R \subset X \times X$ é um conjunto fechado, então o espaço quociente $X/R$ é Hausdorff. Aqui, $X = \mathbb{S}^n$ e $R = \{(x, -x) : x \in \mathbb{S}^n\}$. A aplicação $\alpha: \mathbb{S}^n \times \mathbb{S}^n \to \mathbb{R}^{n+1}$ dada por $\alpha(x,y) = x+y$ é contínua, e $R = \alpha^{-1}(\{0\})$. Como $\{0\}$ é fechado em $\mathbb{R}^{n+1}$, $R$ é fechado em $\mathbb{S}^n \times \mathbb{S}^n$. Logo, $\mathbb{RP}^n$ é Hausdorff. 
	
	A base enumerável é herdada de $\mathbb{S}^n$ via a aplicação de quociente aberta $\pi$.
	
	\textbf{3. Construção do Atlas Suave}
	
	Para cada $i \in \{1, \dots, n+1\}$, definimos os abertos $V_i = \{ x \in \mathbb{S}^n : x^i \neq 0 \}$. Sejam $U_i = \pi(V_i) \subset \mathbb{RP}^n$. 
	
	\textbf{Prova de Cobertura:} 
	
	Seja $[x] \in \mathbb{RP}^n$. Tomando um representante $x \in \mathbb{S}^n$, temos que $\sum_{j=1}^{n+1} (x^j)^2 = 1$. Isso implica que é impossível que todas as coordenadas $x^j$ sejam nulas simultaneamente. Logo, existe ao menos um índice $i \in \{1, \dots, n+1\}$ tal que $x^i \neq 0$. Assim, $x \in V_i$, o que implica $[x] \in \pi(V_i) = U_i$. Portanto, $\mathbb{RP}^n = \bigcup_{i=1}^{n+1} U_i$.
	
	Definimos as cartas $\phi_i: U_i \to \mathbb{R}^n$ por:
	$$ \boxed{\phi_i([x]) = \left( \frac{x^1}{x^i}, \dots, \frac{x^{i-1}}{x^i}, \frac{x^{i+1}}{x^i}, \dots, \frac{x^{n+1}}{x^i} \right).}$$
	\begin{figure}[htbp]
		\centering
		\begin{tikzpicture}[scale=1.5, line join=round, line cap=round]
			\tikzset{
				point/.style={circle, fill=black, inner sep=1pt},
				arrow/.style={->, >=stealth, thin}
			}
			% Esfera
			\draw[thick] (0,0) circle (1);
			\draw[dashed] (1,0) arc (0:180:1 and 0.3);
			\draw (1,0) arc (0:-180:1 and 0.3);
			
			% Pontos em S^n
			\node[point, label={above:$x$}] (X) at (0.5, 0.86) {};
			\node[point, label={below:$-x$}] (mX) at (-0.5, -0.86) {};
			
			% Espaço RP^n
			\node[point, label={right, xshift=-3.5pt, yshift=7pt:$[x]$}] (Proj) at (2, 0) {};
			
			% Espaço R^n
			\node[point, label={right:$\phi_i([x])$}] (R) at (4, 0) {};
			
			% Setas
			\draw[arrow] (X) -- (Proj) node[midway, below, yshift=-2pt] {$\pi$};
			\draw[arrow] (mX) -- (Proj);
			\draw[arrow] (Proj) -- (R) node[midway, below] {$\phi_i$};
			\draw[arrow] (X) -- (R) node[midway, above] {$\phi_i \circ \pi$};
			
			\node at (0, 1.2) {$\mathbb{S}^n$};
			\node at (2, -0.3) {$\mathbb{RP}^n$};
			\node at (4, -0.3) {$\mathbb{R}^n$};
		\end{tikzpicture}
		\caption{Diagrama de construção dos Atlas $\phi_i$}
	\end{figure}
	
	\textbf{Prova de Boa Definição:} 
	
	Para que $\phi_i$ seja bem definida, deve ter $\phi_i([x]) = \phi_i([y])$ sempre que $[x] = [y]$. Se $[x] = [y]$, então $y = \lambda x$ para algum $\lambda \neq 0$. Então:
	$$ \phi_i([y]) = \left( \frac{\lambda x^1}{\lambda x^i}, \dots, \frac{\lambda x^{n+1}}{\lambda x^i} \right) = \left( \frac{x^1}{x^i}, \dots, \frac{x^{n+1}}{x^i} \right) = \phi_i([x]).$$
	
	\textbf{Cálculo da Inversa $\phi_i^{-1}$:}
	
	Seja $u = (u^1, \dots, u^n) \in \mathbb{R}^n$. Procuramos $x = (x^1, \dots, x^{n+1}) \in \mathbb{S}^n$ tal que $\phi_i([x]) = u$. Isso implica $x^j = u^j x^i$ para todo $j \neq i$. 
	
	Substituindo na equação da esfera:
	\begin{align*}
		\sum_{k=1}^{n+1} (x^k)^2 = 1 &\implies (x^i)^2 + \sum_{j \neq i} (u^j x^i)^2 = 1 \\
		(x^i)^2 (1 + \sum_{j \neq i} (u^j)^2) &= 1 \implies (x^i)^2 (1 + |u|^2) = 1 \\
		x^i &= \frac{1}{\sqrt{1+|u|^2}} \quad (\text{onde escolhemos o representante com } x^i > 0)
	\end{align*}
	As demais coordenadas são $x^j = u^j x^i$. Assim, a inversa $\phi_i^{-1}: \mathbb{R}^n \to U_i$ é:
	$$ \phi_i^{-1}(u) = \left[ \frac{u^1}{\sqrt{1+|u|^2}}, \dots, \frac{1}{\sqrt{1+|u|^2}}, \dots, \frac{u^n}{\sqrt{1+|u|^2}} \right].$$
	Tanto $\phi_i$ quanto $\phi_i^{-1}$ são contínuas, logo $\phi_i$ é um homeomorfismo.
	
	\textbf{4. Suavidade das Transições}
	
	Consideremos $i \neq j$ e para $u \in \phi_i(U_i \cap U_j)$, seja $x = \phi_i^{-1}(u)$. A componente $k$-ésima da transição $\phi_j \circ \phi_i^{-1}$ é dada por:
	$$ 
	(\phi_j \circ \phi_i^{-1}(u))^k = \begin{cases} 
		\frac{x^k}{x^j} = \frac{u^k x^i}{u^j x^i} = \frac{u^k}{u^j}, & \text{se } k \neq i, j \\
		\\ 
		\frac{x^i}{x^j} = \frac{x^i}{u^j x^i} = \frac{1}{u^j}, & \text{se } k = i \\
		\\
		\text{por definição é omitida}, & \text{se } k = j.
	\end{cases} 
	$$
	As componentes são funções racionais com denominador $u^j$. Como $u \in \phi_i(U_i \cap U_j)$, temos que $x^j \neq 0$, o que implica $u^j \neq 0$. Portanto, as transições são suaves ($C^\infty$). Concluímos que $\mathbb{RP}^n$ é uma $n$-variedade suave.
}
%%%% QUESTÃO 05 %%%%
\questao{Mostre que o elipsoide $$\mathbb{E}^{n}=\left\{(x^{1},...,x^{n+1})\in\mathbb{R}^{n+1};\sum_{i=1}^{n+1}\left(\frac{x^{i}}{a^{i}}\right)^{2}=1,a^{i}\ne0,\forall i\right\},$$ é uma $n$-variedade suave difeomorfa a esfera unitária $\mathbb{S}^{n}$.
}
\solucao{
	
	\textbf{Análise do Enunciado}
	\begin{itemize}
		\item \textbf{Hipóteses:} $\mathbb{E}^n$ é definido como um elipsoide em $\mathbb{R}^{n+1}$ com semi-eixos $a^i$.
		\item \textbf{Tese:} Provar que $\mathbb{E}^n$ é uma $n$-variedade suave e que existe um difeomorfismo entre $\mathbb{E}^n$ e $\mathbb{S}^n$.
	\end{itemize}
	\vspace{5pt}
	Para provar que $\mathbb{E}^n$ é uma $n$-variedade suave, demonstraremos que ela é a imagem de $\mathbb{S}^n$ sob um difeomorfismo global de $\mathbb{R}^{n+1}$.
	
	\textbf{1. Construção da Aplicação de Transformação}
	
	Definimos a aplicação linear $F: \mathbb{R}^{n+1} \to \mathbb{R}^{n+1}$ dada por:
	\[ F(x^1, \dots, x^{n+1}) = \left( \frac{x^1}{a^1}, \dots, \frac{x^{n+1}}{a^{n+1}} \right) \]
	Em relação à base canônica de $\mathbb{R}^{n+1}$, a matriz de $F$ é a matriz diagonal $\mathbf{A} = \text{diag}(1/a^1, 1/a^2, \dots, 1/a^{n+1})$:
	
	Visto que $a^i \neq 0$ para todo $i$, o determinante de $F$ é $\det(\mathbf{A}) = \prod_{i=1}^{n+1} \frac{1}{a^i} \neq 0$. Portanto, $F$ é um isomorfismo linear. 
	
	Como toda aplicação linear é suave e $F$ é invertível com inversa $F^{-1}(u^1, \dots, u^{n+1}) = (a^1 u^1, \dots, a^{n+1} u^{n+1})$, cuja matriz é a inversa de $\mathbf{A}$ ($\mathbf{A}^{-1} = \text{diag}(a^1, \dots, a^{n+1})$), concluímos que $F$ é um \textbf{difeomorfismo global} de $\mathbb{R}^{n+1}$.
	
	\textbf{2. Relação entre $\mathbb{E}^n$ e $\mathbb{S}^n$}
	
	Analisemos a imagem de $\mathbb{E}^n$ sob $F$. Um ponto $x = (x^1, \dots, x^{n+1})$ pertence a $\mathbb{E}^n$ se, e somente se, satisfaz a equação:
	\[ \sum_{i=1}^{n+1} \left( \frac{x^i}{a^i} \right)^2 = 1 \]
	Seja $u = F(x)$. As componentes de $u$ são $u^i = \frac{x^i}{a^i}$. A condição acima torna-se:
	\[ \sum_{i=1}^{n+1} (u^i)^2 = 1 \iff |u|^2 = 1 \iff u \in \mathbb{S}^n \]
	Isso prova que $F(\mathbb{E}^n) = \mathbb{S}^n$, ou equivalentemente, $\mathbb{E}^n = F^{-1}(\mathbb{S}^n)$.
	\begin{figure}[htbp]
		\centering
		\begin{tikzpicture}[scale=1.8, line join=round, line cap=round]
			\tikzset{
				sphere_line/.style={thin, black},
				arrow/.style={->, >=stealth, thick, black!70},
				point/.style={circle, fill=black, inner sep=0.8pt}
			}
			% Elipsoide (S^n deformada)
			\draw[sphere_line] (0,0) ellipse (1.4 and 0.7);
			% Equador para efeito 3D
			\draw[dashed, sphere_line] (1.4, 0) arc (0:180:1.4 and 0.3);
			\draw[sphere_line] (1.4, 0) arc (0:-180:1.4 and 0.3); % Aqui corrigi para sphere_line
			\node[font=\small] at (0, -0.5) {$\mathbb{E}^n$};
			
			% Seta de transformação
			\draw[arrow] (1.6, 0) -- (3.8, 0);
			\node[above, font=\small] at (2.7, 0.1) {$F(x) = \left( \frac{x^1}{a^1}, \dots, \frac{x^{n+1}}{a^{n+1}} \right)$};
			
			% Esfera unitária
			\draw[sphere_line] (5,0) circle (0.7);
			% Equador para efeito 3D
			\draw[dashed, sphere_line] (5.7, 0) arc (0:180:0.7 and 0.2);
			\draw[sphere_line] (5.7, 0) arc (0:-180:0.7 and 0.2);
			\node[font=\small] at (5, -0.5) {$\mathbb{S}^n$};
		\end{tikzpicture}
		\caption{Difeomorfismo entre o elipsoide $\mathbb{E}^n$ e a esfera unitária $\mathbb{S}^n$ via a transformação linear de escala $F$.}
	\end{figure}
	
	\textbf{3. Prova de que $\mathbb{E}^n$ é uma Variedade Suave}
	
	Seja $f = F|_{\mathbb{E}^n} : \mathbb{E}^n \to \mathbb{S}^n$ a restrição de $F$ ao elipsoide. 
	\begin{enumerate}[label=(\roman*)]
		\item $f$ é suave, pois é a restrição de uma aplicação linear suave.
		\item $f$ é bijetiva, pois $F$ é um isomorfismo e $F(\mathbb{E}^n) = \mathbb{S}^n$.
		\item A inversa $f^{-1} = F^{-1}|_{\mathbb{S}^n}$ também é suave, pois é a restrição de uma aplicação linear suave.
	\end{enumerate}
	Assim, $f: \mathbb{E}^n \to \mathbb{S}^n$ é um \textbf{difeomorfismo}. 
	
	Como $\mathbb{S}^n$ é uma $n$-variedade suave (provado na Questão 3) e o difeomorfismo preserva a estrutura de variedade suave, concluímos que $\mathbb{E}^n$ também é uma $n$-variedade suave de dimensão $n$. A compacidade, a propriedade de Hausdorff e a base enumerável são herdadas de $\mathbb{S}^n$ via o difeomorfismo $f$.
}
%%%% QUESTÃO 06 %%%%
\questao{Mostre que superfícies regulares de $\mathbb{R}^3$ são 2-variedades suaves. Mais geralmente, mostre que hipersuperfícies regulares de $\mathbb{R}^{n+1}$ são $n$-variedades suaves.
}
\solucao{
	
	\textbf{Análise do Enunciado}
	\begin{itemize}
		\item \textbf{Hipóteses:} Uma hipersuperfície regular $M \subset \mathbb{R}^{n+1}$ definida por $f^{-1}(c)$ com $c$ sendo um valor regular de $f$.
		\item \textbf{Tese:} Provar que $M$ é uma $n$-variedade suave.
	\end{itemize}
	\vspace{5pt}
	\textbf{1. Definição de Hipersuperfície Regular}
	
	Uma hipersuperfície regular $M \subset \mathbb{R}^{n+1}$ é definida como o conjunto de pontos onde uma função suave $f: U \subseteq \mathbb{R}^{n+1} \to \mathbb{R}$ atinge um valor regular. Ou seja, existe um aberto $U \subseteq \mathbb{R}^{n+1}$ e um valor $c \in \mathbb{R}$ tal que:
	\[ M = f^{-1}(c) = \{ p \in U : f(p) = c \} \]
	A condição de ser \textbf{regular} implica que, para todo ponto $p \in M$, o diferencial $df_p : T_p\mathbb{R}^{n+1} \to T_{f(p)}\mathbb{R}$ é sobrejetivo. Como a dimensão do contradomínio é $1$, a sobrejetividade equivale a dizer que o gradiente de $f$ não se anula em $M$:
	\[ \nabla f(p) = \left( \frac{\partial f}{\partial x^1}(p), \dots, \frac{\partial f}{\partial x^{n+1}}(p) \right) \neq \vec{0}, \quad \forall p \in M. \]
	
	\textbf{2. Aplicação do Teorema do Valor Regular}
	
	De acordo com o Teorema do Valor Regular, se $f: U \to \mathbb{R}$ é uma aplicação suave e $c$ é um valor regular de $f$, então o conjunto nível $f^{-1}(c)$ é uma subvariedade suave de $U$ de codimensão $1$. 
	
	A dimensão de $M$ é dada por:
	\[ \dim(M) = \dim(\mathbb{R}^{n+1}) - \text{codim}(M) = (n+1) - 1 = n. \]
	Portanto, $M$ é uma $n$-variedade suave. Para o caso específico de superfícies em $\mathbb{R}^3$ ($n=2$), temos que $M$ é uma 2-variedade suave.
}
%%%% QUESTÃO 07 %%%%
\questao{Seja $\mathbb{B}^n = \left\{(x^1, \dots, x^n) \in \mathbb{R}^n : \sum_{i=1}^n (x^i)^2 < 1 \right\}$ a bola aberta unitária de $\mathbb{R}^n$. Prove que $\mathbb{B}^n$ é uma $n$-variedade suave. Além disso, mostre que $f : \mathbb{B}^n \to \mathbb{R}^n$ definida por $f(x) = \frac{x}{\sqrt{1-|x|^2}}$ é um difeomorfismo.
}
\solucao{
	
	\textbf{Análise do Enunciado}
	\begin{itemize}
		\item \textbf{Hipóteses:} $\mathbb{B}^n$ é a bola aberta unitária em $\mathbb{R}^n$.
		\item \textbf{Tese 1:} Provar que $\mathbb{B}^n$ é uma $n$-variedade suave.
		\item \textbf{Tese 2:} Provar que $f(x) = x/\sqrt{1-|x|^2}$ é um difeomorfismo entre $\mathbb{B}^n$ e $\mathbb{R}^n$.
	\end{itemize}
	\vspace{5pt}
	\textbf{1. Natureza de Variedade Suave de $\mathbb{B}^n$}
	
	Sendo $\mathbb{B}^n$ um subconjunto de $\mathbb{R}^n$, consideramos a topologia induzida. A função norma ao quadrado $g: \mathbb{R}^n \to \mathbb{R}$, dada por $g(x) = |x|^2 = \sum_{i=1}^n (x^i)^2$, é uma função contínua. Notemos que $\mathbb{B}^n = g^{-1}((-\infty, 1))$.
	
	Como $(-\infty, 1)$ é um aberto em $\mathbb{R}$, a imagem inversa $\mathbb{B}^n$ é um aberto de $\mathbb{R}^n$. 
	
	Sabemos que qualquer subconjunto aberto $U \subseteq \mathbb{R}^n$ é, por si só, uma $n$-variedade suave. O atlas suave para $\mathbb{B}^n$ é constituído por uma única carta $(\mathbb{B}^n, \text{Id})$, onde $\text{Id}: \mathbb{B}^n \to \mathbb{B}^n$ é a aplicação identidade. Como a identidade é infinitamente diferenciável, $\mathbb{B}^n$ é uma $n$-variedade suave.
	
	\textbf{2. Prova de que $f$ é um Difeomorfismo}
	
	Para que $f: \mathbb{B}^n \to \mathbb{R}^n$ seja um difeomorfismo, devemos provar que $f$ é bijetiva e que tanto $f$ quanto sua inversa $f^{-1}$ são suaves ($C^\infty$).
	
	\textbf{Suavidade de $f$:} A aplicação $f(x) = \frac{x}{\sqrt{1-|x|^2}}$ é composta por funções polinomiais e a função raiz quadrada. Como $|x|^2 < 1$ para todo $x \in \mathbb{B}^n$, o termo $1-|x|^2$ é estritamente positivo. A função $\sqrt{t}$ é suave para $t > 0$. Portanto, $f$ é a composição de funções suaves, sendo, portanto, suave em todo o seu domínio $\mathbb{B}^n$.
	
	\textbf{Cálculo da Inversa $f^{-1}$:} Seja $y \in \mathbb{R}^n$. Procuramos $x \in \mathbb{B}^n$ tal que $f(x) = y$. 
	Temos a relação:
	\[ y = \frac{x}{\sqrt{1-|x|^2}} \]
	Tomando a norma euclidiana de ambos os lados:
	\[ |y| = \frac{|x|}{\sqrt{1-|x|^2}} \implies |y|^2 = \frac{|x|^2}{1-|x|^2} \]
	Isolamos $|x|^2$:
	$$|y|^2(1-|x|^2)= |x|^2 \implies |y|^2 - |y|^2|x|^2 = |x|^2 \implies |y|^2 = |x|^2(1+|y|^2) \implies |x|^2 = \frac{|y|^2}{1+|y|^2}.$$
	Substituindo na equação original $x = y \sqrt{1-|x|^2}$, obtemos que:
	$$x = y \sqrt{1 - \frac{|y|^2}{1+|y|^2}} = y \sqrt{\frac{1+|y|^2 - |y|^2}{1+|y|^2}} = y \sqrt{\frac{1}{1+|y|^2}}=\frac{y}{\sqrt{1+|y|^2}}.$$
	Assim, a inversa $f^{-1}: \mathbb{R}^n \to \mathbb{B}^n$ é definida por $f^{-1}(y) = \frac{y}{\sqrt{1+|y|^2}}$.
	
	\textbf{Suavidade de $f^{-1}$:} Analogamente a $f$, a aplicação $f^{-1}$ é composta por funções suaves, pois o denominador $\sqrt{1+|y|^2}$ é sempre maior ou igual a $1$, nunca se anulando para qualquer $y \in \mathbb{R}^n$. Logo, $f^{-1}$ é suave.
	
	Como $f$ é uma bijecção suave com inversa suave, $f$ é um difeomorfismo entre a bola aberta $\mathbb{B}^n$ e o espaço euclidiano $\mathbb{R}^n$.
}
%%%% QUESTÃO 08 %%%%
\questao{Seja $GL(n, \mathbb{R}) = \{ A \in M(n, \mathbb{R}) : \det A \neq 0 \}$ o grupo linear geral das matrizes reais de ordem $n$. Mostre que $GL(n, \mathbb{R})$ é uma variedade suave. $GL(n, \mathbb{R})$ é conexo?
}
\solucao{
	
	\textbf{Análise do Enunciado}
	\begin{itemize}
		\item \textbf{Hipóteses:} $GL(n, \mathbb{R})$ é o conjunto das matrizes invertíveis.
		\item \textbf{Tese 1:} Provar que $GL(n, \mathbb{R})$ é uma variedade suave.
		\item \textbf{Tese 2:} Determinar se $GL(n, \mathbb{R})$ é um espaço conexo.
	\end{itemize}
	\vspace{5pt}
	\textbf{1. Prova de que $GL(n, \mathbb{R})$ é uma Variedade Suave}
	
	Consideremos o espaço de todas as matrizes reais de ordem $n$, denotado por $M(n, \mathbb{R})$. Este espaço é isomorfo ao espaço euclidiano $\mathbb{R}^{n^2}$ através da aplicação que mapeia cada matriz às suas componentes em um vetor coluna:
	\[ \Phi: M(n, \mathbb{R}) \to \mathbb{R}^{n^2}, \quad \Phi(A) = (A_{11}, A_{12}, \dots, A_{nn}) \]
	A função determinante $\det: M(n, \mathbb{R}) \to \mathbb{R}$ é um polinômio nas entradas da matriz $A_{ij}$. Como todo polinômio é uma função contínua, a aplicação $\det$ é contínua.
	
	Por definição, o grupo linear geral é o conjunto:
	\[ GL(n, \mathbb{R}) = \{ A \in M(n, \mathbb{R}) : \det A \neq 0 \} = \operatorname{det}^{-1}(\mathbb{R} \setminus \{0\}) \]
	O conjunto $\mathbb{R} \setminus \{0\}$ é um aberto no $\mathbb{R}$. Como a imagem inversa de um aberto por uma aplicação contínua é um aberto, concluímos que $GL(n, \mathbb{R})$ é um \textbf{subconjunto aberto} de $M(n, \mathbb{R}) \cong \mathbb{R}^{n^2}$.
	
	Concluímos que, o atlas suave para $GL(n, \mathbb{R})$ é constituído por uma única carta $(GL(n, \mathbb{R}), \text{Id})$. Portanto, $GL(n, \mathbb{R})$ é uma $n^2$-variedade suave.
	
	\textbf{2. Análise da Conexidade}
	
	Para analisar se $GL(n, \mathbb{R})$ é conexo, consideramos novamente a aplicação determinante $\det: GL(n, \mathbb{R}) \to \mathbb{R} \setminus \{0\}$. Sabemos que:
	\begin{enumerate}[label=(\roman*)]
		\item A aplicação $\det$ é contínua.
		\item O contradomínio $\mathbb{R} \setminus \{0\}$ é a união de dois abertos disjuntos: $(-\infty, 0) \cup (0, \infty)$.
	\end{enumerate}
	Se $GL(n, \mathbb{R})$ fosse conexo, a sua imagem sob uma aplicação contínua também deveria ser conexa. No entanto, a imagem da aplicação $\det$ é todo o conjunto $\mathbb{R} \setminus \{0\}$, que é desconexo. 
	
	Podemos definir dois subconjuntos:
	\[ GL^+(n, \mathbb{R}) = \{ A \in GL(n, \mathbb{R}) : \det A > 0 \} \quad \text{e} \quad GL^-(n, \mathbb{R}) = \{ A \in GL(n, \mathbb{R}) : \det A < 0 \} \]
	Ambos são abertos em $GL(n, \mathbb{R})$ (pois são imagens inversas de abertos do $\mathbb{R}$) e são disjuntos. 
	
	Como $GL(n, \mathbb{R}) = GL^+(n, \mathbb{R}) \cup GL^-(n, \mathbb{R})$, temos que $GL(n, \mathbb{R})$ é a união de dois abertos disjuntos não vazios. Portanto, $GL(n, \mathbb{R})$ é \textbf{desconexo}.
}
%%%% QUESTÃO 09 %%%%
\questao{Seja $O(n) = \{A \in M(n, \mathbb{R}) : AA^t = I\}$ o grupo das matrizes ortogonais reais de ordem $n$. Mostre que $O(n)$ é uma variedade suave, compacta e de Hausdorff com base enumerável de dimensão $\frac{n(n-1)}{2}$.
}
\solucao{
	
	\textbf{Análise do Enunciado}
	\begin{itemize}
		\item \textbf{Hipóteses:} $O(n)$ é o conjunto de matrizes tais que $AA^t = I$.
		\item \textbf{Tese:} Provar que $O(n)$ é uma variedade suave, compacta, Hausdorff, com base enumerável e dimensão $n(n-1)/2$.
	\end{itemize}
	\vspace{5pt}
	\textbf{1. Construção via Teorema do Valor Regular}
	
	Consideremos o espaço de todas as matrizes reais $M(n, \mathbb{R})$, que é isomorfo ao espaço euclidiano $\mathbb{R}^{n^2}$. Definimos o espaço das matrizes simétricas $S(n, \mathbb{R}) = \{ S \in M(n, \mathbb{R}) : S = S^t \}$. 
	
	\textbf{Cálculo da Dimensão de $S(n, \mathbb{R})$:} Uma matriz simétrica é completamente determinada por suas entradas na diagonal principal e por suas entradas acima da diagonal.
	\begin{itemize}
		\item Existem $n$ entradas na diagonal principal: $S_{11}, S_{22}, \dots, S_{nn}$.
		\item Existem $\frac{n(n-1)}{2}$ entradas acima da diagonal principal.
	\end{itemize}
	Assim, a dimensão de $S(n, \mathbb{R})$ é a soma desses valores:
	\[ \dim(S(n, \mathbb{R})) = n + \frac{n(n-1)}{2} = \frac{2n + n^2 - n}{2} = \frac{n(n+1)}{2}. \]
	
	Definimos a aplicação suave $F: M(n, \mathbb{R}) \to S(n, \mathbb{R})$ dada por:
	\[ F(A) = AA^t \]
	Note que $F(A)$ é sempre simétrica, pois $(AA^t)^t = (A^t)^t A^t = AA^t$. O grupo ortogonal é, por definição, a imagem inversa do valor $I \in S(n, \mathbb{R})$:
	\[ O(n) = F^{-1}(I). \]
	
	\textbf{Prova de que $I$ é um Valor Regular:} Para cada $A \in O(n)$, devemos mostrar que o diferencial $$dF_A: T_A M(n, \mathbb{R}) \to T_{F(A)} S(n, \mathbb{R}),$$ é sobrejetivo. Identificando os espaços tangentes com as próprias matrizes, seja $H \in M(n, \mathbb{R})$. O diferencial é calculado via a derivada direcional:
	\begin{align*}
		dF_A(H) &= \lim_{t \to 0} \frac{F(A+tH) - F(A)}{t} \\
		&= \lim_{t \to 0} \frac{(A+tH)(A+tH)^t - AA^t}{t} \\
		&= \lim_{t \to 0} \frac{AA^t + t(AH^t + HA^t) + t^2 HH^t - AA^t}{t} \\
		&= AH^t + HA^t.
	\end{align*}
	Para provar a sobrejetividade, dado qualquer $S \in S(n, \mathbb{R})$, devemos encontrar $H$ tal que $AH^t + HA^t = S$. Escolhemos $H = \frac{1}{2}SA$. Então $H^t = \frac{1}{2}A^t S^t = \frac{1}{2}A^t S$ (pois $S$ é simétrica). Substituindo:
	$$ dF_A(H) = A\left(\frac{1}{2}A^t S\right) + \left(\frac{1}{2}SA\right)A^t = \frac{1}{2}AA^t S + \frac{1}{2}S AA^t=\frac{1}{2}IS + \frac{1}{2}SI = S,$$
	pois $A \in O(n)$.
	
	Como $dF_A$ é sobrejetivo para todo $A \in O(n)$, o valor $I$ é um valor regular. Pelo Teorema do Valor Regular, $O(n)$ é uma subvariedade suave de $M(n, \mathbb{R})$.
	
	\textbf{Cálculo da Dimensão: }A dimensão de $O(n, \mathbb{R})$ é a dimensão do espaço ambiente menos a dimensão do contradomínio da função $F$, ou seja, 
	$$\dim(O(n)) = \dim(M(n, \mathbb{R})) - \dim(S(n, \mathbb{R})) = n^2 - \frac{n(n+1)}{2} = \frac{2n^2 - n^2 - n}{2} = \frac{n(n-1)}{2}.$$
	
	\textbf{2. Propriedades Topológicas}
	
	\textbf{Hausdorff e Base Enumerável:} Como $O(n) \subset M(n, \mathbb{R}) \cong \mathbb{R}^{n^2}$, essas propriedades são herdadas diretamente da topologia induzida do espaço euclidiano.
	
	\textbf{Compacidade:} $O(n)$ é um subconjunto fechado de $M(n, \mathbb{R})$ por ser a imagem inversa do ponto $\{I\}$ sob a aplicação contínua $F$. Para provar que é limitado, observemos que para qualquer $A \in O(n)$, as linhas de $A$ são vetores ortonormais em $\mathbb{R}^n$. Se denotarmos por $r_i$ a $i$-ésima linha de $A$, temos $|r_i|^2 = \sum_{j=1}^n A_{ij}^2 = 1$ para todo $i \in \{1, \dots, n\}$. A soma de todos os quadrados das componentes de $A$ é, portanto:
	\[ \sum_{i=1}^n \sum_{j=1}^n A_{ij}^2 = \sum_{i=1}^n |r_i|^2 = \sum_{i=1}^n 1 = n. \]
	Isso implica que todo $A \in O(n)$ satisfaz $|A|^2 = n$, ou seja, $O(n)$ está contido na esfera de raio $\sqrt{n}$ em $\mathbb{R}^{n^2}$. Como qualquer esfera em $\mathbb{R}^{n^2}$ é um conjunto limitado, $O(n)$ é limitado. Pelo Teorema de Heine-Borel, $O(n)$ é compacto.
}
%%%% QUESTÃO 10 %%%%
\questao{Seja $SL(n, \mathbb{R}) = \{ A \in M(n, \mathbb{R}) : \det A = 1 \}$. Mostre que $SL(n, \mathbb{R})$ é uma variedade suave.
}
\solucao{
	
	\textbf{Análise do Enunciado}
	\begin{itemize}
		\item \textbf{Hipóteses:} $SL(n, \mathbb{R})$ é o conjunto de matrizes com determinante igual a 1.
		\item \textbf{Tese:} Provar que $SL(n, \mathbb{R})$ é uma variedade suave.
	\end{itemize}
	\vspace{5pt}
	\textbf{1. Construção via Teorema do Valor Regular}
	
	Consideremos o espaço de todas as matrizes reais de ordem $n$, $M(n, \mathbb{R})$, que é isomorfo ao espaço euclidiano $\mathbb{R}^{n^2}$. Definimos a aplicação determinante:
	\[ f: M(n, \mathbb{R}) \to \mathbb{R}, \quad f(A) = \det(A) \]
	A função determinante é um polinômio nas entradas da matriz $A_{ij}$, sendo, portanto, uma aplicação suave ($C^\infty$). O Grupo Linear Especial $SL(n, \mathbb{R})$ é, por definição, a imagem inversa do valor $1 \in \mathbb{R}$:
	\[ SL(n, \mathbb{R}) = f^{-1}(1) \]
	
	\textbf{2. Análise do Valor Regular e o Diferencial do Determinante}
	
	Para provar que $SL(n, \mathbb{R})$ é uma variedade suave, devemos mostrar que $1$ é um valor regular de $f$. Isso exige que, para todo $A \in SL(n, \mathbb{R})$, o diferencial $df_A: T_A M(n, \mathbb{R}) \to T_{f(A)} \mathbb{R}$ seja sobrejetivo. 
	
	Identificando $T_A M(n, \mathbb{R}) \cong M(n, \mathbb{R})$, tomemos uma matriz $H \in M(n, \mathbb{R})$. O diferencial do determinante no ponto $A$ é dado pela \textbf{Fórmula de Jacobi}:
	\[ df_A(H) = \text{tr}(\text{adj}(A) H) \]
	onde $\text{adj}(A)$ é a matriz adjunta de $A$. Como $A \in SL(n, \mathbb{R})$, sabemos que $A$ é invertível ($\det A = 1$) e que $\text{adj}(A) = \det(A) A^{-1} = A^{-1}$. Assim, o diferencial simplifica-se para:
	\[ df_A(H) = \text{tr}(A^{-1} H) \]
	
	\textbf{Prova de Sobrejetividade:} Para que $df_A$ seja sobrejetivo, dado qualquer $\lambda \in \mathbb{R}$, devemos encontrar uma matriz $H$ tal que $\text{tr}(A^{-1}H) = \lambda$. Escolhemos a matriz $H = \frac{\lambda}{n} A$. Então:
	\[ df_A(H) = \text{tr}(A^{-1} \cdot \frac{\lambda}{n} A) = \text{tr}(\frac{\lambda}{n} I) = \frac{\lambda}{n} \text{tr}(I) = \frac{\lambda}{n} \cdot n = \lambda \]
	Como $df_A$ é sobrejetivo para todo $A \in SL(n, \mathbb{R})$, $1$ é um valor regular de $f$. Pelo Teorema do Valor Regular, $SL(n, \mathbb{R})$ é uma subvariedade suave de $M(n, \mathbb{R})$.
	
	\textbf{Cálculo da Dimensão:}
	A dimensão de $SL(n, \mathbb{R})$ é a dimensão do espaço ambiente menos a dimensão do contradomínio da função $f$, ou seja, $\dim(SL(n, \mathbb{R})) = \dim(M(n, \mathbb{R})) - \dim(\mathbb{R}) = n^2 - 1$.
}
%%%% QUESTÃO 11 %%%%
\questao{ \textbf{(Construção de variedades suaves)} Seja $M$ um subconjunto. Suponha que são dados uma coleção $\{U_\alpha\}_{\alpha \in \Lambda}$ de subconjuntos de $M$ junto com uma aplicação injetiva $\varphi_\alpha : U_\alpha \to \mathbb{R}^n$, para cada $\alpha \in \Lambda$, tais que as seguintes propriedades são satisfeitas:
	\begin{enumerate}[label=(\alph*)]
		\item Para cada $\alpha \in \Lambda$, $\varphi_\alpha(U_\alpha)$ é um subconjunto aberto de $\mathbb{R}^n$.
		\item Para cada par $\alpha, \beta \in \Lambda$, $\varphi_\alpha(U_\alpha \cap U_\beta)$ e $\varphi_\beta(U_\alpha \cap U_\beta)$ são abertos em $\mathbb{R}^n$.
		\item Quando $U_\alpha \cap U_\beta \neq \emptyset$, $\varphi_\alpha \circ \varphi_\beta^{-1} : \varphi_\beta(U_\alpha \cap U_\beta) \to \varphi_\alpha(U_\alpha \cap U_\beta)$ é um difeomorfismo.
		\item Quantidades contáveis de subconjuntos $U_\alpha$ cobrem $M$.
		\item Quando $p,q$ são pontos distintos de $M$, ou existe algum $U_\alpha$ contendo $p$ e $q$ ou existem conjuntos disjuntos $U_\alpha$ e $U_\beta$ tais que $p \in U_\alpha$ e $q \in U_\beta$.
	\end{enumerate}
	Mostre que $M$ admite uma única estrutura suave de forma que $(U_\alpha, \varphi_\alpha)$ seja uma carta suave.
}
\solucao{
	
	\textbf{Análise do Enunciado}
	\begin{itemize}
		\item \textbf{Hipóteses:} Um conjunto $M$ com uma coleção de aplicações injetivas $\varphi_\alpha$ satisfazendo cinco condições (abertos em $\mathbb{R}^n$, transições difeomórficas, cobertura contável e separabilidade de pontos).
		\item \textbf{Tese:} Provar que $M$ possui uma única estrutura de variedade suave induzida por essas cartas.
	\end{itemize}
	\vspace{5pt}
	\textbf{1. Definição da Topologia em $M$}
	
	Para que $M$ seja uma variedade, devemos primeiro definir uma topologia sobre ele. Definimos que um subconjunto $V \subseteq M$ é \textbf{aberto} se, e somente se, para todo $\alpha \in \Lambda$, a imagem $\varphi_\alpha(V \cap U_\alpha)$ for um aberto em $\mathbb{R}^n$.
	
	Verifiquemos que isso define uma topologia:
	\begin{enumerate}[label=(\roman*)]
		\item $\emptyset$ e $M$ são abertos, pois $\varphi_\alpha(\emptyset) = \emptyset$ e $\varphi_\alpha(M \cap U_\alpha) = \varphi_\alpha(U_\alpha)$, que são abertos em $\mathbb{R}^n$ por hipótese (a).
		\item A união arbitrária de abertos $\bigcup V_i$ é aberta, pois $\varphi_\alpha\left(\left(\bigcup V_i\right) \cap U_\alpha\right) = \bigcup (\varphi_\alpha\left(V_i \cap U_\alpha\right))$, que é união de abertos.
		\item A interseção finita de abertos $V_1 \cap V_2$ é aberta, pois $\varphi_\alpha((V_1 \cap V_2) \cap U_\alpha) = \varphi_\alpha(V_1 \cap U_\alpha) \cap \varphi_\alpha(V_2 \cap U_\alpha)$, que é a interseção de dois abertos em $\mathbb{R}^n$.
	\end{enumerate}
	Assim, $M$ é um espaço topológico e cada $\varphi_\alpha : U_\alpha \to \varphi_\alpha(U_\alpha)$ é, por construção, um homeomorfismo.
	
	\textbf{2. Propriedade de Hausdorff}
	
	Sejam $p, q \in M$ pontos distintos. Devemos encontrar abertos disjuntos $V_p$ e $V_q$ contendo $p$ e $q$, respectivamente. Pela hipótese (e), temos dois casos:
	\begin{itemize}
		\item \textbf{Caso 1:} $p, q$ pertencem ao mesmo $U_\alpha$. Como $\varphi_\alpha(U_\alpha)$ é um aberto de $\mathbb{R}^n$ (Hausdorff), existem abertos disjuntos $W_p, W_q \subset \varphi_\alpha(U_\alpha)$ tais que $\varphi_\alpha(p) \in W_p$ e $\varphi_\alpha(q) \in W_q$. Então, $V_p = \varphi_\alpha^{-1}(W_p)$ e $V_q = \varphi_\alpha^{-1}(W_q)$ são abertos disjuntos em $M$ contendo $p$ e $q$.
		\item \textbf{Caso 2:} Existem $U_\alpha, U_\beta$ disjuntos tais que $p \in U_\alpha$ e $q \in U_\beta$. Como $U_\alpha \cap U_\beta = \emptyset$, estes já são os abertos disjuntos necessários.
	\end{itemize}
	Portanto, $M$ é um espaço de Hausdorff.
	
	\textbf{3. Base Enumerável}
	
	Pela hipótese (d), a coleção $\{U_\alpha\}_{\alpha \in \Lambda}$ é contável. Para cada $\alpha$, a imagem $\varphi_\alpha(U_\alpha)$ é um aberto de $\mathbb{R}^n$, que possui uma base enumerável de bolas abertas $\{B_{\alpha, n}\}_{n \in \mathbb{N}}$. A coleção $\{\varphi_\alpha^{-1}(B_{\alpha, n}) : \alpha \in \Lambda, n \in \mathbb{N}\}$ é uma união contável de coleções enumeráveis, logo é uma base enumerável para a topologia de $M$.
	
	\textbf{4. Estrutura Suave e Unicidade}
	
	Temos agora que $M$ é um espaço de Hausdorff com base enumerável e que a coleção $\mathcal{A} = \{(U_\alpha, \varphi_\alpha)\}_{\alpha \in \Lambda}$ cobre $M$. A hipótese (c) garante que para quaisquer $\alpha, \beta \in \Lambda$, a função de transição $\varphi_\alpha \circ \varphi_\beta^{-1}$ é um difeomorfismo entre abertos de $\mathbb{R}^n$. Portanto, por definição, $\mathcal{A}$ é um atlas suave para $M$.
	
	Para a unicidade da estrutura suave, recorremos ao resultado provado na \textbf{Questão 1}. Vimos que qualquer atlas suave está contido em um único atlas suave maximal. Sendo $\mathcal{A}$ um atlas suave, existe um único atlas maximal $\mathcal{A}_{max}$ tal que $\mathcal{A} \subseteq \mathcal{A}_{max}$. 
	
	Uma estrutura suave em uma variedade é definida precisamente como a classe de equivalência de atlas suaves que compartilham o mesmo atlas maximal. Como a coleção $\{(U_\alpha, \varphi_\alpha)\}_{\alpha \in \Lambda}$ determina um único $\mathcal{A}_{max}$, a estrutura suave induzida em $M$ é única. Concluímos, portanto, que $M$ é uma $n$-variedade suave.
}
%%%% QUESTÃO 12 %%%%
\questao{ \textbf{(Orientabilidade)} Seja $(M^n, \mathcal{A})$ uma $n$-variedade suave. Dizemos que $M^n$ é orientável se $M^n$ pode ser coberta por um atlas $\{U_\alpha, \varphi_\alpha\}_{\alpha \in \Lambda} \subset \mathcal{A}$ tal que para todo par $\alpha, \beta \in \Lambda$ com $U_\alpha \cap U_\beta \neq \emptyset$, a derivada da aplicação transição $\varphi_\beta \circ \varphi_\alpha^{-1}$ tem determinante positivo. Caso tal atlas não exista, dizemos que $M^n$ é não-orientável.
	\begin{enumerate}[label=(\alph*)]
		\item Mostre que se uma $n$-variedade suave $M^n$ pode ser coberta por duas vizinhanças coordenadas $U_1$ e $U_2$ de modo que a interseção $U_1 \cap U_2$ é conexa, então $M^n$ é orientável.
		\item Use item (a) para mostrar que $\mathbb{S}^n$ é orientável.
		\item Mostre diretamente que $\mathbb{RP}^2$ é não-orientável.
\end{enumerate}
}
\solucao{
	
	\textbf{Análise do Enunciado}
	\begin{itemize}
		\item \textbf{Hipóteses:} Definição de orientabilidade via sinal do determinante da Jacobiana das transições.
		\item \textbf{Tese (a):} Provar que a conexidade da interseção de dois abertos de cobertura implica orientabilidade.
		\item \textbf{Tese (b):} Aplicar o resultado (a) à esfera $\mathbb{S}^n$.
		\item \textbf{Tese (c):} Demonstrar a não-orientabilidade de $\mathbb{RP}^2$.
	\end{itemize}
	\vspace{5pt}
	\textbf{(a) Prova da Orientabilidade via Interseção Conexa}
	
	Sejam $(U_1, \varphi_1)$ e $(U_2, \varphi_2)$ as duas cartas que cobrem $M^n$. A aplicação de transição é $\tau = \varphi_2 \circ \varphi_1^{-1}$, definida no aberto $W = \varphi_1(U_1 \cap U_2) \subset \mathbb{R}^n$. 
	
	Consideremos a função determinante do Jacobiano $J_\tau: W \to \mathbb{R}$, dada por $J_\tau(u) = \det(d\tau_u)$. Como $\tau$ é um difeomorfismo, $d\tau_u$ é invertível para todo $u \in W$, logo $J_\tau(u) \neq 0$ para todo $u \in W$.
	
	Por hipótese, $U_1 \cap U_2$ é conexo. Como $\varphi_1$ é um homeomorfismo, a imagem $W = \varphi_1(U_1 \cap U_2)$ também é um conjunto conexo em $\mathbb{R}^n$. A função $J_\tau$ é contínua e nunca se anula em $W$. Pelo Teorema do Valor Intermediário, uma função contínua em um conjunto conexo que nunca assume o valor zero deve manter o mesmo sinal em todo o domínio.
	
	Existem dois casos:
	\begin{enumerate}[label=(\roman*)]
		\item Se $J_\tau(u) > 0$ para todo $u \in W$, o atlas $\{(U_1, \varphi_1), (U_2, \varphi_2)\}$ já é um atlas de orientação.
		\item Se $J_\tau(u) < 0$ para todo $u \in W$, definimos uma nova carta $(U_1, \tilde{\varphi}_1)$ onde $\tilde{\varphi}_1(x) = (-\varphi_1^1(x), \varphi_1^2(x), \dots, \varphi_1^n(x))$. A nova transição será $\tau' = \varphi_2 \circ \tilde{\varphi}_1^{-1}$. A matriz Jacobiana de $\tilde{\varphi}_1$ é $\text{diag}(-1, 1, \dots, 1)$, com determinante $-1$. Pela regra da cadeia:
		\[ \det(d\tau'_u) = \det(d(\varphi_2 \circ \varphi_1^{-1})_u) \cdot \det(d\tilde{\varphi}_1^{-1}) = J_\tau(u) \cdot (-1) > 0. \]
	\end{enumerate}
	Em ambos os casos, construímos um atlas onde a transição tem determinante positivo. Logo, $M^n$ é orientável.
	
	\textbf{(b) Orientabilidade da Esfera $\mathbb{S}^n$}
	
	Utilizaremos a construção da \textbf{Questão 3(c)}, onde $\mathbb{S}^n$ é coberta pelas duas cartas de projeção estereográfica $(U_N, \phi_N)$ e $(U_S, \phi_S)$. 
	
	A interseção dos domínios é dada por $U_N \cap U_S = \mathbb{S}^n \setminus \{N, S\}$. Para analisar a conexidade deste conjunto, observamos que a projeção $\phi_N$ restrita a essa interseção é um homeomorfismo:
	\[ \phi_N : \mathbb{S}^n \setminus \{N, S\} \longrightarrow \mathbb{R}^n \setminus \{0\} \]
	Sendo $\mathbb{R}^n \setminus \{0\}$ homeomorfo ao produto $\mathbb{S}^{n-1} \times (0, \infty)$ via a aplicação $v \mapsto \left( \frac{v}{|v|}, |v| \right)$, e considerando que $(0, \infty) \cong \mathbb{R}$, concluímos que:
	\[ U_N \cap U_S \cong \mathbb{S}^{n-1} \times \mathbb{R} \]
	Para $n \ge 2$, a esfera $\mathbb{S}^{n-1}$ é conexa e, como o produto de espaços conexos é conexo, a interseção $U_N \cap U_S$ é conexa. Sendo a interseção conexa, a condição do item (a) é satisfeita. Portanto, $\mathbb{S}^n$ é orientável.
	
	\textbf{(c) Não-orientabilidade de $\mathbb{RP}^2$}
	
	Utilizaremos as cartas $\phi_i$ definidas na \textbf{Questão 4}. Consideremos as cartas $\phi_1$ e $\phi_2$ nos abertos $U_1$ e $U_2$. Para $n=2$, as coordenadas são $\phi_1([x]) = (x^2/x^1, x^3/x^1)$ e $\phi_2([x]) = (x^1/x^2, x^3/x^2)$.
	
	A transição $\tau = \phi_2 \circ \phi_1^{-1}$ no domínio $u = (u^1, u^2) \in \mathbb{R}^2 \setminus \{0\}$ é:
	\[ \tau(u^1, u^2) = \left( \frac{1}{u^1}, \frac{u^2}{u^1} \right) \]
	Calculamos a matriz Jacobiana $d\tau$:
	\[ d\tau = \begin{bmatrix} 
		\frac{\partial \tau^1}{\partial u^1} & \frac{\partial \tau^1}{\partial u^2} \\ 
		\\
		\frac{\partial \tau^2}{\partial u^1} & \frac{\partial \tau^2}{\partial u^2} \end{bmatrix} = \begin{bmatrix} 
		\frac{-1}{(u^1)^2} & 0 \\
		\\
		\frac{-u^2}{(u^1)^2} & \frac{1}{u^1} \end{bmatrix} \]
	O determinante da transição é:
	\[ \det(d\tau) = \left( -\frac{1}{(u^1)^2} \right) \left( \frac{1}{u^1} \right) - 0 = -\frac{1}{(u^1)^3} \]
	Observe que o sinal do determinante depende do sinal de $u^1$. No domínio $\phi_1(U_1 \cap U_2) = \mathbb{R}^2 \setminus \{0\}$, a coordenada $u^1$ pode assumir tanto valores positivos quanto negativos. 
	
	Se tentarmos mudar a orientação de $\phi_1$ ou $\phi_2$ multiplicando uma coordenada por $-1$, alteraremos o sinal do determinante em todo o domínio. No entanto, como o sinal de $\det(d\tau)$ varia dentro do mesmo domínio de transição (ele é positivo para $u^1 < 0$ e negativo para $u^1 > 0$), não existe escolha de sinais para as cartas que torne o determinante globalmente positivo. Portanto, $\mathbb{RP}^2$ é não-orientável.
}
%%%% QUESTÃO 13 %%%%
\questao{Sejam $M^n$ uma variedade suave e $f : M^n \to \mathbb{R}^k$ uma aplicação suave. Mostre que a representação coordenada $\hat{f}$ de $f$ é suave para qualquer carta suave $(U, \varphi)$ de $M^n$.
}
\solucao{
	
	\textbf{Análise do Enunciado}
	\begin{itemize}
		\item \textbf{Hipóteses:} $f$ é uma aplicação suave de uma variedade $M^n$ para $\mathbb{R}^k$.
		\item \textbf{Tese:} Provar que a representação local $\hat{f} = f \circ \varphi^{-1}$ é suave.
	\end{itemize}
	\vspace{5pt}
	\textbf{1. Definições e Notação}
	
	Seja $(U, \varphi)$ uma carta suave de $M^n$. A representação coordenada de $f$ em relação a esta carta é a aplicação $\hat{f} : \varphi(U) \to \mathbb{R}^k$ definida pela composição:
	\[ \hat{f} = f \circ \varphi^{-1} \]
	onde $\varphi(U)$ é um aberto em $\mathbb{R}^n$ e a aplicação $\varphi^{-1}: \varphi(U) \to U$ é o homeomorfismo inverso da carta.
	
	\textbf{2. Análise da Suavidade}
	
	Para provar que $\hat{f}$ é suave, analisamos a sequência de aplicações:
	\[ \mathbb{R}^n \xrightarrow{\varphi^{-1}} M^n \xrightarrow{f} \mathbb{R}^k \]
	Pela definição de estrutura suave em $M^n$, a aplicação $\varphi^{-1}$ é suave. Como, por hipótese, $f$ é uma aplicação suave de $M^n$ para $\mathbb{R}^k$, a composição $\hat{f} = f \circ \varphi^{-1}$ é a composição de duas aplicações suaves entre variedades. 
	
	Sendo a composição de funções suaves também suave, concluímos que $\hat{f} \in C^\infty(\varphi(U), \mathbb{R}^k)$.
	
	\textbf{3. Independência da Escolha da Carta}
	
	Para garantir a consistência da definição, consideremos outra carta $(V, \psi)$ tal que $U \cap V \neq \emptyset$. A representação de $f$ em relação a $\psi$, denotada por $\tilde{f} = f \circ \psi^{-1}$, relaciona-se com $\hat{f}$ através da função de transição $\tau_{\psi\varphi} = \varphi \circ \psi^{-1}$:
	\[ \tilde{f} = f \circ \psi^{-1} = \underbrace{(f \circ \varphi^{-1})}_{\hat{f}} \circ \underbrace{(\varphi \circ \psi^{-1})}_{\tau_{\psi\varphi}} = \hat{f} \circ \tau_{\psi\varphi} \]
	Como $\tau_{\psi\varphi}$ é um difeomorfismo entre abertos de $\mathbb{R}^n$ e $\hat{f}$ é suave, a composição $\tilde{f}$ é necessariamente suave. Isso prova que a suavidade da representação coordenada é uma propriedade intrínseca da aplicação $f$, independente da carta escolhida.
}
%%%% QUESTÃO 14 %%%%
\questao{
	Sejam $M^{n}$ e $N^{k}$ variedades suaves, e seja $\{U_{\alpha}\}_{\alpha \in \Gamma}$ uma cobertura de $M^{n}$. Suponha que para cada $\alpha \in \Gamma$ temos uma aplicação suave $F_{\alpha} : U_{\alpha} \to N^{k}$ tal que $F_{\alpha} \big|_{U_{\alpha} \cap U_{\beta}} = F_{\beta} \big|_{U_{\alpha} \cap U_{\beta}}$, para quaisquer $\alpha, \beta \in \Gamma$. Mostre que existe uma única aplicação $F : M^{n} \to N^{k}$ tal que $F \big|_{U_{\alpha}} = F_{\alpha}$, para todo $\alpha \in \Gamma$.
}
\solucao{
	
	\textbf{Análise do Enunciado}
	\begin{itemize}
		\item \textbf{Hipóteses:} $\{U_\alpha\}$ cobre $M^n$, e cada $F_\alpha: U_\alpha \to N^k$ é suave e compatível nas interseções.
		\item \textbf{Tese:} Existência e unicidade de uma aplicação global $F: M^n \to N^k$ que seja suave.
	\end{itemize}
	\vspace{5pt}
	\textbf{1. Existência e Boa Definição de $F$:}
	Definimos a aplicação $F: M^n \to N^k$ da seguinte forma: para qualquer ponto $p \in M^n$, como $\{U_\alpha\}_{\alpha \in \Gamma}$ é uma cobertura de $M^n$, existe pelo menos um índice $\alpha \in \Gamma$ tal que $p \in U_\alpha$. Atribuímos:
	\begin{equation*}
		F(p) = F_\alpha(p).
	\end{equation*}
	Para que $F$ esteja bem definida, o valor de $F(p)$ não deve depender da escolha do índice $\alpha$. Seja $p \in U_\alpha \cap U_\beta$. Pela hipótese de compatibilidade, temos que:
	\begin{equation*}
		F_\alpha \big|_{U_\alpha \cap U_\beta} = F_\beta \big|_{U_\alpha \cap U_\beta} \implies F_\alpha(p) = F_\beta(p).
	\end{equation*}
	Portanto, a aplicação $F$ é bem definida em todo o domínio $M^n$.
	
	\textbf{2. Unicidade:}
	Suponha que exista outra aplicação $G: M^n \to N^k$ tal que $G \big|_{U_\alpha} = F_\alpha$ para todo $\alpha \in \Gamma$. Para qualquer $p \in M^n$, escolhemos $\alpha$ tal que $p \in U_\alpha$. Então:
	$$G(p) = G \big|_{U_\alpha}(p) = F_\alpha(p) = F(p).$$
	Como $G(p) = F(p)$ para todo $p \in M^n$, concluímos que $G = F$.
	
	\textbf{3. Suavidade de $F$:}
	Para provar que $F$ é suave, devemos verificar a condição de suavidade em cada ponto $p \in M^n$. 
	Seja $p \in M^n$. Existe um índice $\alpha \in \Gamma$ tal que $p \in U_\alpha$. Por hipótese, a aplicação $F_\alpha: U_\alpha \to N^k$ é suave. 
	
	Pela definição de suavidade em $U_\alpha$, para o ponto $p$, existem cartas coordenadas $(U_p, \varphi)$ em $M^n$ com $p \in U_p \subseteq U_\alpha$ e $(V_{F(p)}, \psi)$ em $N^k$ com $F(p) \in V_{F(p)}$ tais que a representação coordenada de $F_\alpha$ no ponto $p$:
	\begin{equation*}
		\hat{F}_\alpha = \psi \circ F_\alpha \circ \varphi^{-1}
	\end{equation*}
	é uma aplicação suave de $\varphi(U_p) \subseteq \mathbb{R}^n$ em $\psi(V_{F(p)}) \subseteq \mathbb{R}^k$.
	
	Como $F \big|_{U_\alpha} = F_\alpha$, temos que:
	\begin{equation*}
		\psi \circ F \circ \varphi^{-1} = \psi \circ (F \big|_{U_\alpha}) \circ \varphi^{-1} = \psi \circ F_\alpha \circ \varphi^{-1} = \hat{F}_\alpha.
	\end{equation*}
	Assim, $\hat{F}_\alpha$ é suave para todo $p \in M^n$, e portanto a aplicação global $F$ satisfaz a definição de aplicação suave entre variedades.
	
	A aplicação $F$ foi construída a partir de dados locais compatíveis, é única por construção e sua suavidade decorre da suavidade local de cada $F_\alpha$. Assim, $F$ é a única aplicação suave que estende a coleção $\{F_\alpha\}$.
}
%%%% QUESTÃO 15 %%%%
\questao{
	Sejam $M^{n}$ e $N^{k}$ variedades suaves e $F : M^{n} \to N^{k}$ uma aplicação suave. Mostre que a representação coordenada $\hat{F}$ de $F$ é suave para qualquer par de cartas coordenadas suaves $(U, \varphi)$ de $M^{n}$ e $(V, \psi)$ de $N^{k}$, respectivamente.
}
\solucao{
	
	\textbf{Análise do Enunciado}
	\begin{itemize}
		\item \textbf{Hipóteses:} $F$ é uma aplicação suave entre as variedades $M^n$ e $N^k$.
		\item \textbf{Tese:} Provar que a representação coordenada $\hat{F} = \psi \circ F \circ \varphi^{-1}$ é suave para qualquer escolha de cartas.
	\end{itemize}
	\vspace{5pt}
	Seja $\hat{F} = \psi \circ F \circ \varphi^{-1}$ a representação coordenada de $F$ em relação às cartas $(U, \varphi)$ e $(V, \psi)$. O domínio de $\hat{F}$ é o aberto $\varphi(U \cap F^{-1}(V)) \subseteq \mathbb{R}^n$. 
	
	\textbf{1. Existência de representação local suave:}
	Seja $p \in U \cap F^{-1}(V)$. Como $F$ é uma aplicação suave, por definição, existe um par de cartas $(U', \varphi')$ em $M^n$ e $(V', \psi')$ em $N^k$ tais que $p \in U'$, $F(p) \in V'$ e a representação coordenada local
	\begin{equation*}
		\tilde{F} = \psi' \circ F \circ (\varphi')^{-1}
	\end{equation*}
	é suave em uma vizinhança de $\varphi'(p)$.
	
	\textbf{2. Decomposição via funções de transição:}
	Para formalizar o domínio de validade desta composição, definimos o aberto $W = U \cap U' \cap F^{-1}(V \cap V')$. Como $U$ e $U'$ são abertos em $M^n$, $V$ e $V'$ são abertos em $N^k$ e $F$ é contínua, $W$ é uma vizinhança aberta de $p$ em $M^n$. Sobre a imagem $\varphi(W) \subseteq \mathbb{R}^n$, a representação $\hat{F}$ pode ser fatorada da seguinte forma:
	\begin{align*}
		\hat{F} &= \psi \circ F \circ \varphi^{-1} \\
		&= \psi \circ (\psi')^{-1} \circ \psi' \circ F \circ (\varphi')^{-1} \circ \varphi' \circ \varphi^{-1} \\
		&= \underbrace{(\psi \circ (\psi')^{-1})}_{\text{transição em } N^k} \circ \underbrace{\tilde{F}}_{\text{representação local}} \circ \underbrace{(\varphi' \circ \varphi^{-1})}_{\text{transição em } M^n}.
	\end{align*}
	
	O diagrama a seguir ilustra a comutatividade desta composição sobre o domínio $\varphi(W)$:
	\begin{figure}[htbp]
		\centering
		\begin{tikzpicture}[node distance=3cm, auto]
			% Nós
			\node (R1) {$\mathbb{R}^n$};
			\node (R2) [right of=R1] {$\mathbb{R}^n$};
			\node (R3) [right of=R2] {$\mathbb{R}^k$};
			\node (R4) [below of=R3, node distance=1.5cm] {$\mathbb{R}^k$};
			
			% Setas
			\draw[->] (R1) -- node {$\varphi' \circ \varphi^{-1}$} (R2);
			\draw[->] (R2) -- node {$\tilde{F}$} (R3);
			\draw[->] (R3) -- node {$\psi \circ (\psi')^{-1}$} (R4);
			\draw[->] (R1) .. controls (1.5, -1) and (3, -1) .. (R4) node [midway, below] {$\hat{F}$};
		\end{tikzpicture}
		\caption{Diagrama comutativo da representação coordenada.}
	\end{figure}
	
	\textbf{3. Análise da suavidade dos componentes:}
	\begin{itemize}
		\item A aplicação $\varphi' \circ \varphi^{-1}$ é a função de transição entre as cartas $\varphi$ e $\varphi'$ em $M^n$, que é um difeomorfismo, logo é suave.
		\item A aplicação $\tilde{F} = \psi' \circ F \circ (\varphi')^{-1}$ é suave por hipótese, dada a suavidade de $F$.
		\item A aplicação $\psi \circ (\psi')^{-1}$ é a função de transição entre as cartas $\psi'$ e $\psi$ em $N^k$, que é um difeomorfismo, logo é suave.                        
	\end{itemize}
	
	Como $\hat{F}$ é a composição de três aplicações suaves entre abertos de espaços euclidianos, temos que $\hat{F}$ é suave. Como $p$ foi escolhido arbitrariamente, a representação coordenada é suave para qualquer par de cartas suaves.
}
%%%% QUESTÃO 16 %%%%
\questao{
	Mostre que difeomorfismos definem uma relação de equivalência no conjunto de variedades suaves. Além disso, mostre que o conjunto dos difeomorfismos de uma $n$-variedade suave $M^{n}$, munido com a operação da composição, é um grupo.
}
\solucao{
	
	\textbf{Análise do Enunciado}
	\begin{itemize}
		\item \textbf{Hipóteses:} $\mathcal{M}$ é o conjunto de todas as variedades suaves; a relação $\sim$ é definida pela existência de um difeomorfismo entre duas variedades.
		\item \textbf{Tese 1:} Provar que a relação $\sim$ é de equivalência (reflexiva, simétrica e transitiva).
		\item \textbf{Tese 2:} Provar que $\text{Diff}(M^n)$, com a operação de composição, satisfaz os axiomas de grupo.
	\end{itemize}
	\vspace{5pt}
	\textbf{1. Prova da Relação de Equivalência}
	
	Seja $\mathcal{M}$ o conjunto de todas as variedades suaves. Definimos a relação $\sim$ em $\mathcal{M}$ da seguinte forma: para quaisquer variedades suaves $M, N \in \mathcal{M}$, dizemos que $M \sim N$ se, e somente se, existe um difeomorfismo $f: M \to N$. Para que $\sim$ seja uma relação de equivalência, ela deve satisfazer as seguintes propriedades:
	
	\begin{itemize}
		\item \textbf{Reflexividade:} Para qualquer variedade $M$, a aplicação identidade $\text{Id}_M: M \to M$ é um homeomorfismo e é suave (pois a representação coordenada $\text{Id}_{\mathbb{R}^n}$ é suave). Como $(\text{Id}_M)^{-1} = \text{Id}_M$ também é suave, $\text{Id}_M$ é um difeomorfismo. Logo, $M \sim M$.
		
		\item \textbf{Simetria:} Suponha que $M \sim N$. Então existe um difeomorfismo $f: M \to N$. Por definição de difeomorfismo, $f$ é bijetiva e possui uma inversa $f^{-1}: N \to M$ que também é suave. Como $f^{-1}$ é suave e sua inversa $(f^{-1})^{-1} = f$ também é suave, $f^{-1}$ é um difeomorfismo de $N$ para $M$. Logo, $N \sim M$.
		
		\item \textbf{Transitividade:} Suponha que $M \sim N$ e $N \sim L$. Então existem difeomorfismos $f: M \to N$ e $g: N \to L$. Consideremos a composição $g \circ f: M\to P$. 
		Como $f$ e $g$ são bijetivas, $g \circ f$ é bijetiva. A inversa de $g \circ f$ é dada por $(g \circ f)^{-1} = f^{-1} \circ g^{-1}$. 
		Como $f, g, f^{-1}$ e $g^{-1}$ são aplicações suaves, e a composição de aplicações suaves é suave, concluímos que tanto $g \circ f$ quanto $f^{-1} \circ g^{-1}$ são suaves. Portanto, $g \circ f$ é um difeomorfismo de $M$ para $L$, logo $M \sim L$.
	\end{itemize}
	
	Como a relação é reflexiva, simétrica e transitiva, os difeomorfismos definem uma relação de equivalência em $\mathcal{M}$.
	
	\textbf{2. Prova da Estrutura de Grupo de $\text{Diff}(M^n)$}
	
	Seja $\text{Diff}(M^n)$ o conjunto de todos os difeomorfismos de uma $n$-variedade suave $M^n$ para si mesma. Consideremos a operação de composição $\circ$. Verificamos os axiomas de grupo:
	
	\begin{itemize}
		\item \textbf{Fechamento:} Sejam $f, g \in \text{Diff}(M^n)$. Como $f$ e $g$ são difeomorfismos, a composição $f \circ g: M^n \to M^n$ é bijetiva e suave. Sua inversa $(f \circ g)^{-1} = g^{-1} \circ f^{-1}$ também é a composição de difeomorfismos, logo é suave. Portanto, $f \circ g \in \text{Diff}(M^n)$.
		
		\item \textbf{Associatividade:} Para quaisquer $f, g, h \in \text{Diff}(M^n)$ e qualquer $p \in M^n$:
		\begin{equation*}
			(f \circ (g \circ h))(p) = f(g(h(p))) = ((f \circ g) \circ h)(p),
		\end{equation*}
		pois a composição é associativa e composição de difeomorfismos é difeomorfismo.
		
		\item \textbf{Elemento Neutro:} A aplicação identidade $\text{Id}_M: M^n \to M^n$ pertence a $\text{Diff}(M^n)$. Para qualquer $f \in \text{Diff}(M^n)$, temos:
		\begin{equation*}
			f \circ \text{Id}_M = f = \text{Id}_M \circ f.
		\end{equation*}
		
		\item \textbf{Elemento Inverso:} Para qualquer $f \in \text{Diff}(M^n)$, existe por definição a aplicação inversa $f^{-1}: M^n \to M^n$. Como $f$ é um difeomorfismo, $f^{-1}$ também é um difeomorfismo, logo $f^{-1} \in \text{Diff}(M^n)$. Além disso:
		\begin{equation*}
			f \circ f^{-1} = \text{Id}_M = f^{-1} \circ f.
		\end{equation*}
	\end{itemize}
	
	Como todos os axiomas foram satisfeitos,o conjunto $\text{Diff}(M^n)$ munido da composição é um grupo.
}
%%%% QUESTÃO 17 %%%%
\questao{
	Mostre que duas variedades suaves difeomorfas possuem a mesma dimensão.
}
\solucao{
	
	\textbf{Análise do Enunciado}
	\begin{itemize}
		\item \textbf{Hipóteses:} $M$ e $N$ são variedades suaves de dimensões $m$ e $n$, respectivamente, e existe um difeomorfismo $f: M \to N$.
		\item \textbf{Tese:} Provar que $m = n$.
	\end{itemize}
	\vspace{5pt}
	Sejam $M$ e $N$ variedades suaves de dimensões $m$ e $n$, respectivamente. Por hipótese, existe um difeomorfismo $f: M \to N$.
	
	\textbf{1. Representação Coordenada:}
	Seja $p \in M$ e consideremos a representação coordenada $\hat{f} = \psi \circ f \circ \varphi^{-1}$ em relação a cartas suaves $(U, \varphi)$ em $M$ e $(V, \psi)$ em $N$. Como $f$ é um difeomorfismo, $\hat{f}$ é um difeomorfismo entre o aberto $\varphi(U \cap f^{-1}(V)) \subseteq \mathbb{R}^m$ e o aberto $\psi(f(U \cap f^{-1}(V))) \subseteq \mathbb{R}^n$.
	
	\textbf{2. Análise do Diferencial:}
	Para qualquer ponto $x$ no domínio de $\hat{f}$, a aplicação $\hat{f}$ é suave e possui uma inversa suave $\hat{f}^{-1}$. Pela regra da cadeia, a composição $\hat{f}^{-1} \circ \hat{f} = \text{Id}_{\mathbb{R}^m}$ implica que:
	\begin{equation*}
		d(\hat{f}^{-1})_{\hat{f}(x)} \circ d\hat{f}_x = \text{Id}_{\mathbb{R}^m}.
	\end{equation*}
	Analogamente, $\hat{f} \circ \hat{f}^{-1} = \text{Id}_{\mathbb{R}^n}$ implica que:
	\begin{equation*}
		d\hat{f}_x \circ d(\hat{f}^{-1})_{\hat{f}(x)} = \text{Id}_{\mathbb{R}^n}.
	\end{equation*}
	
	\textbf{3. Isomorfismo de Espaços Vetoriais:}
	As equações acima demonstram que o diferencial $d\hat{f}_x : \mathbb{R}^m \to \mathbb{R}^n$ é uma aplicação linear invertível, ou seja, é um isomorfismo de espaços vetoriais. Como a dimensão de um espaço vetorial é preservada por isomorfismos, temos necessariamente que:
	\begin{equation*}
		\dim(\mathbb{R}^m) = \dim(\mathbb{R}^n) \implies m = n.
	\end{equation*}
	Assim, as variedades $M$ e $N$ possuem a mesma dimensão.
}
%%%% QUESTÃO 18 %%%%
\questao{
	Considere variedades suaves $M^{m}$ e $N^{n}$ de dimensão $m$ e $n$, respectivamente. Sejam $\mathcal{U} = \{(U_{i}, \phi_{i})\}_{i \in I}$ e $\mathcal{V} = \{(V_{j}, \psi_{j})\}_{j \in J}$ atlas suaves para $M^{m}$ e $N^{n}$, respectivamente. Mostre que:
	\begin{enumerate}[label=(\alph*)]
		\item $\mathcal{U} \times \mathcal{V} = \{(U_{i} \times V_{j}, \phi_{i} \times \psi_{j})_{i \in I, j \in J}\}$ é um atlas suave em $M^{m} \times N^{n}$.
		\item As projeções $\pi_{1} : M^{m} \times N^{n} \to M^{m}$ e $\pi_{2} : M^{m} \times N^{n} \to N^{n}$ são aplicações suaves.
		\item Se $P^{k}$ é uma outra variedade suave, $g : P^{k} \to M^{m} \times N^{n}$ é suave se, e somente se, $\pi_{i} \circ g$ é suave para todo $i \in \{1,2\}$.
		\item Para $y \in N^{n}$, a aplicação inclusão 
		\begin{eqnarray*}
			\iota : M^{m}& \longrightarrow &M^{m} \times N^{n}\\
			x&\longmapsto& (x,y),
		\end{eqnarray*}
		é suave
		\item $M^{m} \times N^{n}$ é difeomorfo a $N^{n} \times M^{m}$.
	\end{enumerate}
}
\solucao{
	
	\textbf{Análise do Enunciado}
	\begin{itemize}
		\item \textbf{Hipóteses:} $M^m$ e $N^n$ são variedades suaves com atlas $\mathcal{U}$ e $\mathcal{V}$.
		\item \textbf{Tese (a):} Provar que o produto de atlas define uma estrutura suave em $M \times N$.
		\item \textbf{Tese (b):} Provar a suavidade das projeções canônicas $\pi_i$.
		\item \textbf{Tese (c):} Provar a caracterização de suavidade de $g$ via projeções.
		\item \textbf{Tese (d):} Provar que a inclusão de $M$ em $M \times N$ (fixando $y$) é suave.
		\item \textbf{Tese (e):} Provar a existência de um difeomorfismo entre $M \times N$ e $N \times M$.
	\end{itemize}
	\vspace{5pt}
	\textbf{(a) Atlas do Produto}
	
	\textbf{Prova da Cobertura e Estrutura Localmente Euclidiana:}
	
	Primeiro, verificamos que a coleção $\mathcal{U} \times \mathcal{V}$ cobre $M^m \times N^n$. Como $\{U_i\}_{i \in I}$ cobre $M$ e $\{V_j\}_{j \in J}$ cobre $N$, temos:
	\begin{equation*}
		\bigcup_{i \in I, j \in J} (U_i \times V_j) = \left( \bigcup_{i \in I} U_i \right) \times \left( \bigcup_{j \in J} V_j \right) = M^m \times N^n.
	\end{equation*}
	Cada carta $\phi_i \times \psi_j : U_i \times V_j \to \mathbb{R}^{m+n}$ é definida por $(\phi_i \times \psi_j)(p, q) = (\phi_i(p), \psi_j(q))$. Como $\phi_i$ e $\psi_j$ são homeomorfismos sobre abertos de $\mathbb{R}^m$ e $\mathbb{R}^n$, o produto $\phi_i \times \psi_j$ é um homeomorfismo sobre o aberto $\phi_i(U_i) \times \psi_j(V_j) \subseteq \mathbb{R}^{m+n}$.
	
	\textbf{Propriedade de Hausdorff:} Sejam $p = (x_1, y_1)$ e $q = (x_2, y_2)$ pontos distintos em $M^m \times N^n$. Então, ou $x_1 \neq x_2$ ou $y_1 \neq y_2$. Se $x_1 \neq x_2$, como $M^m$ é Hausdorff, existem abertos disjuntos $U_1, U_2 \subset M^m$ tais que $x_1 \in U_1$ e $x_2 \in U_2$. Logo, os conjuntos abertos $U_1 \times N^n$ e $U_2 \times N^n$ são disjuntos em $M^m \times N^n$ e separam $p$ e $q$. O raciocínio é análogo caso $y_1 \neq y_2$. Assim, $M^m \times N^n$ é Hausdorff.
	
	\textbf{Base Enumerável:} Sejam $\mathcal{B}_M = \{U_i\}_{i \in \mathbb{N}}$ e $\mathcal{B}_N = \{V_j\}_{j \in \mathbb{N}}$ bases enumeráveis para $M^m$ e $N^n$, respectivamente. Pela definição de topologia produto, a coleção $\mathcal{B}_{M \times N} = \{ U_i \times V_j : i, j \in \mathbb{N} \}$ constitui uma base para a topologia de $M^m \times N^n$. Como o produto cartesiano $\mathbb{N} \times \mathbb{N}$ é enumerável, a base $\mathcal{B}_{M \times N}$ é enumerável.
	
	\textbf{Suavidade das Transições}
	
	Analisamos a função de transição entre duas cartas $(U_i \times V_j, \phi_i \times \psi_j)$ e $(U_k \times V_l, \phi_k \times \psi_l)$:
	\begin{align*}
		(\phi_k \times \psi_l) \circ (\phi_i \times \psi_j)^{-1}(u, v) &= (\phi_k \times \psi_l)(\phi_i^{-1}(u), \psi_j^{-1}(v)) \\
		&= (\phi_k \circ \phi_i^{-1}(u), \psi_l \circ \psi_j^{-1}(v)).
	\end{align*}
	Como $\phi_k \circ \phi_i^{-1}$ e $\psi_l \circ \psi_j^{-1}$ são funções de transição de atlas suaves, elas são $C^\infty$. Como a aplicação total é definida por componentes suaves, a transição no produto é suave.
	
	Portanto, $M^m \times N^n$ é uma variedade suave de dimensão $m+n$ e $\mathcal{U} \times \mathcal{V}$ é um atlas suave.
	
	\textbf{(b) Suavidade das Projeções}
	
	Analisamos a representação coordenada de $\pi_1$ em relação às cartas $\phi_i \times \psi_j$ e $\phi_i$. Para $(u, v) \in \mathbb{R}^m \times \mathbb{R}^n$:
	\begin{align*}
		\phi_i \circ \pi_1 \circ (\phi_i \times \psi_j)^{-1}(u, v) &= \phi_i(\pi_1(\phi_i^{-1}(u), \psi_j^{-1}(v))) \\
		&= \phi_i(\phi_i^{-1}(u)) = u.
	\end{align*}
	A aplicação $(u, v) \mapsto u$ é a projeção canônica de $\mathbb{R}^{m+n}$ em $\mathbb{R}^m$, que é linear e, portanto, suave. O mesmo raciocínio aplica-se a $\pi_2$.
	
	\textbf{(c) Caracterização da Suavidade de $g$}
	
	$(\implies)$ Se $g: P^k \to M \times N$ é suave, então como $\pi_1$ e $\pi_2$ são suaves (provado acima), as composições $\pi_1 \circ g$ e $\pi_2 \circ g$ são suaves, pois a composição de aplicações suaves é suave.
	
	$(\impliedby)$ Suponha que $g_1 = \pi_1 \circ g$ e $g_2 = \pi_2 \circ g$ sejam suaves. Então $g(p) = (g_1(p), g_2(p))$. Seja $(U, \phi)$ uma carta em $P^k$ e $(U_i \times V_j, \phi_i \times \psi_j)$ uma carta em $M \times N$. A representação coordenada de $g$ é:
	\begin{align*}
		(\phi_i \times \psi_j) \circ g \circ \phi^{-1}(u) &= (\phi_i \circ g_1 \circ \phi^{-1}(u), \psi_j \circ g_2 \circ \phi^{-1}(u)).
	\end{align*}
	Como $g_1$ e $g_2$ são suaves, as representações coordenadas $\phi_i \circ g_1 \circ \phi^{-1}$ e $\psi_j \circ g_2 \circ \phi^{-1}$ são suaves. Como cada componente da representação de $g$ é suave, $g$ é suave.
	
	\textbf{(d) Suavidade da Inclusão $\iota$}
	
	Para um ponto $x \in M^m$ e fixo $y \in N^n$, consideremos a representação coordenada de $\iota$:
	\begin{align*}
		(\phi_i \times \psi_j) \circ \iota \circ \phi_k^{-1}(u) &= (\phi_i \times \psi_j)(\iota(\phi_k^{-1}(u))) \\
		&= (\phi_i \times \psi_j)(\phi_k^{-1}(u), y) \\
		&= (\phi_i \circ \phi_k^{-1}(u), \psi_j(y)).
	\end{align*}
	A primeira componente $\phi_i \circ \phi_k^{-1}$ é uma função de transição suave em $M^m$. A segunda componente $\psi_j(y)$ é uma constante em $\mathbb{R}^n$, logo é suave. Assim, $\iota$ é suave.
	
	\textbf{(e) Difeomorfismo entre $M \times N$ e $N \times M$}
	
	Definimos a aplicação de troca $S: M \times N \to N \times M$ por $S(x, y) = (y, x)$. Claramente $S$ é bijetiva e sua inversa é a própria $S$, logo $S$ é um homeomorfismo.
	
	Analisamos a representação coordenada de $S$ em relação aos atlas produto:
	\begin{align*}
		(\psi_j \times \phi_i) \circ S \circ (\phi_i \times \psi_j)^{-1}(u, v) &= (\psi_j \times \phi_i)(S(\phi_i^{-1}(u), \psi_j^{-1}(v))) \\
		&= (\psi_j \times \phi_i)(\psi_j^{-1}(v), \phi_i^{-1}(u)) \\
		&= (\psi_j(\psi_j^{-1}(v)), \phi_i(\phi_i^{-1}(u))) = (v, u).
	\end{align*}
	A aplicação $(u, v) \mapsto (v, u)$ é uma permutação de coordenadas em $\mathbb{R}^{n+m}$, que é linear e suave. Como $S$ e $S^{-1}$ possuem representações coordenadas suaves, $S$ é um difeomorfismo. Portanto, concluímos que, $M^m \times N^n$ e $N^n \times M^m$ são variedades difeomorfas.
}
%%%% QUESTÃO 19 %%%%
\questao{
	Mostre que a aplicação $F : \mathbb{R}^3 \to \mathbb{R}^6$ definida por $F(x, y, z) = (x^2, y^2, z^2, \sqrt{2}xy, \sqrt{2}xz, \sqrt{2}yz)$ induz uma aplicação $f : \mathbb{S}^2 \to \mathbb{S}^5$.
}
\solucao{
	
	\textbf{Análise do Enunciado}
	\begin{itemize}
		\item \textbf{Hipóteses:} $F$ é uma aplicação de $\mathbb{R}^3$ para $\mathbb{R}^6$.
		\item \textbf{Tese:} Provar que a restrição de $F$ ao domínio $\mathbb{S}^2$ tem imagem contida em $\mathbb{S}^5$.
	\end{itemize}
	\vspace{5pt}
	Para provar que $F$ induz a aplicação $f = F|_{\mathbb{S}^2}$, devemos mostrar que a imagem de qualquer ponto $p = (x, y, z) \in \mathbb{S}^2$ pertence à esfera $\mathbb{S}^5$.
	
	\textbf{1. Verificação da Imagem:}
	Como $p \in \mathbb{S}^2$, temos que $x^2 + y^2 + z^2 = 1$. Calculando a norma euclidiana de $F(p)$ ao quadrado, obtemos:
	\begin{align*}
		\|F(p)\|&=\|F(x, y, z)\|^2 = (x^2)^2 + (y^2)^2 + (z^2)^2 + (\sqrt{2}xy)^2 + (\sqrt{2}xz)^2 + (\sqrt{2}yz)^2 \\
		&= x^4 + y^4 + z^4 + 2x^2y^2 + 2x^2z^2 + 2y^2z^2 \\
		&= (x^2 + y^2 + z^2)^2=(1)^2 = 1,
	\end{align*}
	pois $p \in \mathbb{S}^2$.
	
	Portanto, $F(p) \in \mathbb{S}^5$ para todo $p \in \mathbb{S}^2$, o que valida a existência da aplicação induzida $f$.
	
	\textbf{2. Suavidade:}
	Como $F$ é definida por componentes polinomiais, ela é suave em $\mathbb{R}^3$. Sendo $\mathbb{S}^2$ e $\mathbb{S}^5$ subvariedades suaves, a restrição $f = F|_{\mathbb{S}^2}$ é automaticamente uma aplicação suave entre variedades.	
}
%%%% QUESTÃO 20 %%%%
\questao{
	Mostre que a aplicação $F : \mathbb{R}^3 \to \mathbb{R}^6$ definida por $F(x, y, z) = (x^2, y^2, z^2, \sqrt{2}xy, \sqrt{2}xz, \sqrt{2}yz)$ induz uma aplicação $\bar{f} : \mathbb{RP}^2 \to \mathbb{S}^5$.
}
\solucao{
	
	\textbf{Análise do Enunciado}
	\begin{itemize}
		\item \textbf{Hipóteses:} $F$ é a mesma aplicação da questão anterior.
		\item \textbf{Tese:} Provar que $F$ induz uma aplicação bem definida $\bar{f}$ no espaço projetivo $\mathbb{RP}^2$ com imagem em $\mathbb{S}^5$.
	\end{itemize}
	\vspace{5pt}
	Para que a aplicação $F$ induza uma aplicação $\bar{f}$ no espaço projetivo $\mathbb{RP}^2$, ela deve ser constante sobre cada classe de equivalência da relação $\sim$ em $\mathbb{S}^2$. Como $\mathbb{RP}^2$ é o quociente de $\mathbb{S}^2$ pela identificação $p \sim -p$, a condição necessária e suficiente é que $F(p) = F(-p)$ para todo $p \in \mathbb{S}^2$.
	\begin{figure}[htbp]
		\centering
		\begin{tikzpicture}[scale=1.2, line join=round, line cap=round]
			\tikzset{
				point/.style={circle, fill=black, inner sep=1pt},
				arrow/.style={->, >=stealth, thin}
			}
			
			% Esfera S2
			\draw[thick] (0,0) circle (1);
			\draw[dashed] (1,0) arc (0:180:1 and 0.3);
			\draw (1,0) arc (0:-180:1 and 0.3);
			\node at (0, 1.2) {$\mathbb{S}^2$};
			
			% Pontos p e -p
			\node[point, label={above right:$p$}] (P) at (0.6, 0.8) {};
			\node[point, label={below left:$-p$}] (mP) at (-0.6, -0.8) {};
			
			% Espaco RP2
			\node (RP) at (3,0) {$\mathbb{RP}^2$};
			\node[point, label={below:$[p]$}] (Proj) at (3, -0.5) {};
			
			% Esfera S5
			\draw[thick] (6,0) circle (0.8);
			\draw[dashed] (6.8,0) arc (0:180:0.8 and 0.2);
			\draw (6.8,0) arc (0:-180:0.8 and 0.2);
			\node at (6, 1.1) {$\mathbb{S}^5$};
			\node[point, label={right:$F(p)$}] (S5) at (6.5, 0.3) {};
			
			% Setas
			\draw[arrow] (P) -- node[midway, above] {$\pi$} (Proj);
			\draw[arrow] (mP) -- node[midway, below] {$\pi$} (Proj);
			\draw[arrow] (Proj) -- node[midway, above] {$\bar{f}$} (S5);
			\draw[arrow] (P) .. controls (3, 1.5) and (5, 1.5) .. (S5) node[midway, above] {$F|_{\mathbb{S}^2}$};
			\draw[arrow] (mP) .. controls (3, -2) and (5, -2) .. (S5) node[midway, below] {$F|_{\mathbb{S}^2}$};
			
		\end{tikzpicture}
		\caption{Diagrama de indução da aplicação $\bar{f}$ via projeção canônica $\pi$.}
	\end{figure}
	\textbf{1. Prova de Boa Definição:}
	Seja $p = (x, y, z) \in \mathbb{S}^2$. O ponto equivalente a $p$ é $-p = (-x, -y, -z)$. Aplicamos $F$ ao ponto $-p$:
	\begin{align*}
		F(-x, -y, -z) &= ((-x)^2, (-y)^2, (-z)^2, \sqrt{2}(-x)(-y), \sqrt{2}(-x)(-z), \sqrt{2}(-y)(-z)) \\
		&= (x^2, y^2, z^2, \sqrt{2}xy, \sqrt{2}xz, \sqrt{2}yz) \\
		&= F(x, y, z).
	\end{align*}
	Como $F(p) = F(-p)$ para qualquer $p \in \mathbb{S}^2$, a aplicação $\bar{f} : \mathbb{RP}^2 \to \mathbb{R}^6$ definida por $\bar{f}([p]) = F(p)$ é bem definida, independentemente do representante escolhido para a classe $[p]$.
	
	\textbf{2. Imagem na Esfera $\mathbb{S}^5$:}
	Conforme provamos na \textbf{Questão 19}, para qualquer $p \in \mathbb{S}^2$, a norma $\|F(p)\|^2 = 1$. Portanto, a imagem de qualquer classe $[p] \in \mathbb{RP}^2$ sob $\bar{f}$ pertence obrigatoriamente à esfera $\mathbb{S}^5$, o que valida o codomínio da aplicação $\bar{f} : \mathbb{RP}^2 \to \mathbb{S}^5$.
	
	\textbf{3. Suavidade:}
	Por construção, temos que $F|_{\mathbb{S}^2} = \bar{f} \circ \pi$, onde $\pi : \mathbb{S}^2 \to \mathbb{RP}^2$ é a projeção canônica. Como $\pi$ é um difeomorfismo local, a suavidade de $\bar{f}$ em um ponto $[p]$ é equivalente à suavidade de $F$ em $p$. Como $F$ é definida por componentes polinomiais, a sua restrição a $\mathbb{S}^2$ é suave. Portanto, a aplicação induzida $\bar{f}$ é suave.
}
%%%% QUESTÃO 21 %%%%
\questao{
	Mostre que a aplicação $F : \mathbb{R}^n \to \mathbb{R}^{\frac{n(n+1)}{2}}$ definida por $$F(x^1, \dots, x^n) = ((x^1)^2, \dots, (x^n)^2, \sqrt{2}x^1x^2, \dots, \sqrt{2}x^{n-1}x^n),$$
	induz uma aplicação $f : \mathbb{S}^{n-1} \to \mathbb{S}^{\frac{n(n+1)}{2}-1}$.
}
\solucao{
	
	\textbf{Análise do Enunciado}
	\begin{itemize}
		\item \textbf{Hipóteses:} $F$ é uma aplicação polinomial de $\mathbb{R}^n$ para $\mathbb{R}^{n(n+1)/2}$.
		\item \textbf{Tese:} Provar que a restrição de $F$ à esfera $\mathbb{S}^{n-1}$ tem imagem contida na esfera de dimensão $\frac{n(n+1)}{2}-1$.
	\end{itemize}
	\vspace{5pt}
	Para provar que $F$ induz a aplicação $f = F|_{\mathbb{S}^{n-1}}$, devemos mostrar que, para qualquer ponto $p = (x^1, \dots, x^n) \in \mathbb{S}^{n-1}$, a norma euclidiana de $F(p)$ é igual a $1$.
	
	\textbf{1. Verificação da Imagem:}
	Como $p \in \mathbb{S}^{n-1}$, as coordenadas de $p$ satisfazem a condição $\sum_{i=1}^n (x^i)^2 = 1$. 
	As componentes de $F(p)$ consistem em $n$ termos do tipo $(x^i)^2$ e $\frac{n(n-1)}{2}$ termos do tipo $\sqrt{2}x^i x^j$ (para $i < j$). Calculamos a norma de $F(p)$ ao quadrado:
	$$\|F(p)\|^2 = \sum_{i=1}^n ((x^i)^2)^2 + \sum_{1 \le i < j \le n} (\sqrt{2}x^i x^j)^2 = \sum_{i=1}^n (x^i)^4 + \sum_{1 \le i < j \le n} 2(x^i)^2(x^j)^2=\left( \sum_{i=1}^n (x^i)^2 \right)^2 = (1)^2 = 1,
	$$
	pois $p \in \mathbb{S}^{n-1}$.
	
	Isso prova que $F(p) \in \mathbb{S}^{\frac{n(n+1)}{2}-1}$ para todo $p \in \mathbb{S}^{n-1}$, validando a existência da aplicação induzida $f$.
	
	\textbf{3. Suavidade:}
	Como $F$ é definida por componentes polinomiais, ela é suave em $\mathbb{R}^n$. Sendo $\mathbb{S}^{n-1}$ e $\mathbb{S}^{\frac{n(n+1)}{2}-1}$ subvariedades suaves, a restrição $f = F|_{\mathbb{S}^{n-1}}$ é automaticamente uma aplicação suave entre variedades.
}
%%%% QUESTÃO 22 %%%%
\questao{
	Para cada $x \in \mathbb{S}^n$ defina $\gamma : \mathbb{R} \to \mathbb{S}^n$ por $\gamma(t) = \cos t \, x + \sen t \frac{v}{\|v\|}$, onde $v \in \mathbb{R}^{n+1} \setminus \{0\}$ é perpendicular a $x$. Mostre que $\gamma$ é suave.
}
\solucao{
	
	\textbf{Análise do Enunciado}
	\begin{itemize}
		\item \textbf{Hipóteses:} $\gamma$ é uma curva parametrizada por $t$ em $\mathbb{R}^{n+1}$, com $x$ e $v/\|v\|$ sendo vetores ortonormais.
		\item \textbf{Tese:} Provar que a imagem de $\gamma$ está em $\mathbb{S}^n$ e que $\gamma$ é uma aplicação suave.
	\end{itemize}
	\vspace{5pt}
	Para provar que $\gamma$ é uma aplicação suave para a esfera $\mathbb{S}^n$, devemos primeiro verificar que a imagem de $\gamma$ está contida em $\mathbb{S}^n$ e, em seguida, analisar a suavidade de suas componentes.
	
	\textbf{1. Verificação da Norma:}
	Seja $u = \frac{v}{\|v\|}$ o vetor unitário na direção de $v$. Por hipótese, $\|x\| = 1$, $\|u\| = 1$ e $x \cdot u = 0$ (pois $v \perp x$). Calculamos a norma de $\gamma(t)$ ao quadrado utilizando o produto escalar:
	\begin{align*}
		\|\gamma(t)\|^2 &= \langle \cos t \, x + \sen t \, u, \cos t \, x + \sen t \, u \rangle \\
		&= \cos^2 t \langle x, x \rangle + \sen^2 t \langle u, u \rangle + 2 \cos t \sen t \langle x, u \rangle \\
		&= \cos^2 t (1) + \sen^2 t (1) + 2 \cos t \sen t (0) \\
		&= \cos^2 t + \sen^2 t = 1.
	\end{align*}
	Como $\|\gamma(t)\| = 1$ para todo $t \in \mathbb{R}$, a imagem de $\gamma$ está contida em $\mathbb{S}^n$, logo a aplicação $\gamma : \mathbb{R} \to \mathbb{S}^n$ está bem definida.
	\begin{figure}[htbp]
		\centering
		\begin{tikzpicture}[scale=2, line join=round, line cap=round]
			\tikzset{
				point/.style={circle, fill=black, inner sep=0.8pt},
				vector/.style={->, >=stealth, thick}
			}
			
			% Esfera S^n
			\draw[thick] (0,0) circle (1);
			\draw[dashed] (1,0) arc (0:180:1 and 0.3);
			\draw (1,0) arc (0:-180:1 and 0.3);
			\node at (-1.3, 0) {$\mathbb{S}^n$}; % Rótulo na lateral esquerda
			
			% Vetores Ortogonais
			\draw[vector] (0,0) -- (1,0) node[midway, below] {$\cos t$};
			\draw[vector] (0,0) -- (0,1) node[midway, left] {$\sen t$};
			
			% Pontos x e u
			\node[point, label={right:$x$}] (X) at (1,0) {};
			\node[point, label={above:$u = v/\|v\|$}] (U) at (0,1) {};
			\node[point, label={below:$O$} ] at (0,0) {};
			
			% A curva gamma(t)
			\draw[ultra thick, blue!70!black] (1,0) arc (0:90:1);
			\node[blue!70!black] at (1.1, 0.6) {$\gamma(t)$}; % Deslocados para a direita
			
			% Ângulo t - Agora fechando do eixo X ao eixo Y com seta
			\draw[->] (0.2,0) arc (0:90:0.2);
			\node at (0.21, 0.21) {\small $t$};
			
		\end{tikzpicture}
		\caption{Representação da curva $\gamma(t)$ como um círculo máximo no plano gerado por $\{x, u\}$.}
	\end{figure}
	
	\textbf{2. Análise da Suavidade:}
	Consideremos a aplicação $\gamma$ como uma função indo para $\mathbb{R}^{n+1}$. Se denotarmos as componentes de $x$ por $(x^1, \dots, x^{n+1})$ e as de $u$ por $(u^1, \dots, u^{n+1})$, a $i$-ésima componente de $\gamma(t)$ é dada por:
	\begin{equation*}
		\gamma^i(t) = (\cos t)x^i + (\sen t)u^i.
	\end{equation*}
	Como as funções $\sen t$ e $\cos t$ são infinitamente diferenciáveis em $\mathbb{R}$, e $x^i, u^i$ são constantes, cada componente $\gamma^i$ é uma função suave. 
	
	Sendo $\gamma : \mathbb{R} \to \mathbb{R}^{n+1}$ uma aplicação suave e sua imagem contida na subvariedade suave $\mathbb{S}^n$, a aplicação $\gamma : \mathbb{R} \to \mathbb{S}^n$ é, por definição, suave.
}
%%%% QUESTÃO 23 %%%%
\questao{
	Seja $F : \mathbb{R}^{n+1} \to \mathbb{R}^k$ uma aplicação suave. Seja $f = F|_{\mathbb{S}^n}$. Prove que $f$ é suave.
}
\solucao{
	
	\textbf{Análise do Enunciado}
	\begin{itemize}
		\item \textbf{Hipóteses:} $F$ é suave em $\mathbb{R}^{n+1}$. $f$ é a restrição de $F$ à subvariedade $\mathbb{S}^n$.
		\item \textbf{Tese:} Provar que $f: \mathbb{S}^n \to \mathbb{R}^k$ é uma aplicação suave.
	\end{itemize}
	\vspace{5pt}
	Para provar que a restrição $f$ é suave, devemos analisar a estrutura de $f$ como uma composição de aplicações entre variedades suaves.
	
	\textbf{1. Formalização da Restrição via Inclusão:}
	A restrição da aplicação $F$ ao subconjunto $\mathbb{S}^n$ pode ser expressa formalmente através da aplicação de inclusão $\iota : \mathbb{S}^n \hookrightarrow \mathbb{R}^{n+1}$, definida por $\iota(p) = p$ para todo $p \in \mathbb{S}^n$. Assim, a aplicação $f : \mathbb{S}^n \to \mathbb{R}^k$ é dada pela composição:
	\begin{equation*}
		f = F \circ \iota.
	\end{equation*}
	
	\textbf{2. Suavidade da Aplicação de Inclusão:}
	Sabemos que $\mathbb{S}^n$ é uma subvariedade suave de $\mathbb{R}^{n+1}$. Por definição, a aplicação de inclusão $\iota$ é suave se, para qualquer carta $(U, \varphi)$ de $\mathbb{S}^n$, a representação coordenada $\hat{\iota} = \text{Id}_{\mathbb{R}^{n+1}} \circ \iota \circ \varphi^{-1} = \iota \circ \varphi^{-1}$ for suave.
	
	Como $\mathbb{S}^n$ é uma subvariedade suave, a aplicação $\iota \circ \varphi^{-1}$ é, por construção, uma imersão suave de um aberto de $\mathbb{R}^n$ em $\mathbb{R}^{n+1}$. Portanto, $\iota$ é uma aplicação suave.
	
	\textbf{3. Composição de Aplicações Suaves:}
	Temos agora a seguinte sequência de aplicações:
	\begin{equation*}
		\mathbb{S}^n \xrightarrow{\iota} \mathbb{R}^{n+1} \xrightarrow{F} \mathbb{R}^k.
	\end{equation*}
	Como $\iota$ é uma aplicação suave entre a variedade $\mathbb{S}^n$ e a variedade $\mathbb{R}^{n+1}$, e $F$ é uma aplicação suave entre as variedades $\mathbb{R}^{n+1}$ e $\mathbb{R}^k$, a composição $f = F \circ \iota$ é, necessariamente, uma aplicação suave entre as variedades $\mathbb{S}^n$ e $\mathbb{R}^k$.
	
	Como a composição de funções $C^\infty$ preserva a suavidade, concluímos que a restrição $f = F|_{\mathbb{S}^n}$ é suave.
}
%%%% QUESTÃO 24 %%%%	
\questao{
	Seja $F : \mathbb{S}^n \to \mathbb{R}$ uma função suave tal que $F(x) = F(-x)$. Mostre que $F$ induz uma função $\bar{f}$ no espaço projetivo real $\mathbb{RP}^n$ que também é suave.
}
\solucao{
	
	\textbf{Análise do Enunciado}
	\begin{itemize}
		\item \textbf{Hipóteses:} $F$ é suave em $\mathbb{S}^n$ e possui simetria antípoda ($F(x) = F(-x)$).
		\item \textbf{Tese:} Provar que a aplicação induzida $\bar{f}: \mathbb{RP}^n \to \mathbb{R}$ é bem definida e suave.
	\end{itemize}
	\vspace{5pt}
	Consideremos a estrutura de variedade suave de $\mathbb{RP}^n$ estabelecida na \textbf{Questão 4}, onde $\mathbb{RP}^n$ é definido como o quociente de $\mathbb{S}^n$ pela relação $x \sim -x$.
	
	\textbf{1. Boa Definição de $f$:}
	Por hipótese, temos que $F(x) = F(-x)$ para todo $x \in \mathbb{S}^n$. Assim, definimos a função $f : \mathbb{RP}^n \to \mathbb{R}$ da seguinte forma, $f([x]) = F(x)$, onde $[x]$ denota a classe de equivalência $\{x, -x\}$. Como a imagem de $F$ é a mesma para qualquer representante da classe $[x]$, a função $f$ está bem definida.
	
	\textbf{2. Suavidade de $\bar{f}$:}
	Para provar a suavidade de $\bar{f}$, utilizamos as cartas coordenadas $(U_i, \phi_i)$ definidas na \textbf{Questão 4}. Seja $p = [x] \in U_i$, a representação coordenada de $\bar{f}$ é a aplicação $\bar{f} \circ \phi_i^{-1}$ e a aplicação $\phi_i^{-1} : \phi_i(U_i) \to U_i$ pode ser composta com a projeção $\pi : \mathbb{S}^n \to \mathbb{RP}^n$ para obter uma aplicação suave $\sigma_i : \phi_i(U_i) \to \mathbb{S}^n$, tal que $\pi \circ \sigma_i = \phi_i^{-1}$. Assim, a representação de $\bar{f}$ torna-se, $\bar{f} \circ \phi_i^{-1} = \bar{f} \circ (\pi \circ \sigma_i) = (\bar{f} \circ \pi) \circ \sigma_i$. Por outro lado, da definição de $\bar{f}$, temos que $\bar{f} \circ \pi = F$. Portanto, 
	$$\boxed{\bar{f} \circ \phi_i^{-1} = F \circ \sigma_i}$$
	\begin{figure}[htbp]
		\centering
		\begin{tikzpicture}[scale=1.2, line join=round, line cap=round]
			\tikzset{
				point/.style={circle, fill=black, inner sep=1pt},
				arrow/.style={->, >=stealth, thin}
			}
			
			% S^n
			\draw[thick] (0,0) circle (1);
			\draw[dashed] (1,0) arc (0:180:1 and 0.3);
			\draw (1,0) arc (0:-180:1 and 0.3);
			\node at (0, 1.2) {$\mathbb{S}^n$};
			
			% Pontos x e -x
			\node[point, label={above right:$x$}] (X) at (0.7, 0.7) {};
			\node[point, label={below left:$-x$}] (mX) at (-0.7, -0.7) {};
			
			% RP^n
			\node (RP) at (3,0) {$\mathbb{RP}^n$};
			\node[point, label={below:$[x]$}] (Proj) at (3, -0.5) {};
			
			% R
			\node (R) at (6.1,0) {$\mathbb{R}$};
			\node[point, label={right, xshift=1pt:$F(x)$}] (Val) at (6, 0.3) {};
			
			% Setas
			\draw[arrow] (X) -- node[midway, above] {$\pi$} (Proj);
			\draw[arrow] (mX) -- node[midway, below] {$\pi$} (Proj);
			\draw[arrow] (Proj) -- node[midway, above] {$\bar{f}$} (Val);
			\draw[arrow] (X) .. controls (3, 1.5) and (5, 1.5) .. (Val) node[midway, above] {$F$};
			\draw[arrow] (mX) .. controls (3, -2) and (5, -2) .. (Val) node[midway, below] {$F$};
			
		\end{tikzpicture}
		\caption{Diagrama de comutação - a função $F$ em $\mathbb{S}^n$ induz a função $\bar{f}$ em $\mathbb{RP}^n$ via projeção $\pi$.}
	\end{figure}
	
	Como $F$ é uma função suave em $\mathbb{S}^n$ e $\sigma_i$ é uma aplicação suave entre abertos de $\mathbb{R}^n$ e $\mathbb{S}^n$, a composição $F \circ \sigma_i$ é suave. Como a representação de $\bar{f}$ em qualquer carta da Questão 4 é suave, concluímos que $\bar{f}$ é uma função suave.
}
%%%% QUESTÃO 25 %%%%
\questao{
	Seja $F : \mathbb{S}^n \to \mathbb{S}^n$ definida por $F(x) = -x$. Mostre que $F$ é um difeomorfismo.
}
\solucao{
	
	\textbf{Análise do Enunciado}
	\begin{itemize}
		\item \textbf{Hipóteses:} $F$ é a aplicação antípoda da esfera $\mathbb{S}^n$ para si mesma.
		\item \textbf{Tese:} Provar que $F$ é um difeomorfismo.
	\end{itemize}
	\vspace{5pt}
	Para provar que $F$ é um difeomorfismo, devemos mostrar que ela é bijetiva e que tanto $F$ quanto sua inversa $F^{-1}$ são aplicações suaves.
	
	\textbf{1. Bijetividade e Inversa:}
	Seja $x \in \mathbb{S}^n$. a aplicação $F$ envia cada ponto ao seu ponto antípoda. Verifiquemos a composição de $F$ com ela mesma:
	\begin{equation*}
		(F \circ F)(x) = F(F(x)) = F(-x) = -(-x) = x.
	\end{equation*}
	Como $F \circ F = \text{Id}_{\mathbb{S}^n}$, a aplicação $F$ é a sua própria inversa ($F = F^{-1}$). Toda aplicação que possui inversa é, necessariamente, bijetiva.
	
	\textbf{2. Suavidade de $F$:}
	Consideremos a aplicação linear $L : \mathbb{R}^{n+1} \to \mathbb{R}^{n+1}$ definida por $L(v) = -v$. Como $L$ é uma aplicação linear, ela é suave em todo o seu domínio $\mathbb{R}^{n+1}$.
	
	A aplicação $F$ é precisamente a restrição de $L$ ao subconjunto $\mathbb{S}^n$:
	\begin{equation*}
		F = L|_{\mathbb{S}^n}.
	\end{equation*}
	Sendo $\mathbb{S}^n$ uma subvariedade suave de $\mathbb{R}^{n+1}$ e $L$ uma aplicação suave no espaço ambiente, a restrição $F$ é, por definição, uma aplicação suave entre variedades.
	
	\textbf{3. Suavidade de $F^{-1}$:}
	Visto que $F^{-1} = F$, a suavidade da inversa é idêntica à suavidade da aplicação original. Portanto, $F^{-1}$ também é suave. Como $F$ é bijetiva, suave e possui inversa suave, concluímos que $F$ é um difeomorfismo da esfera $\mathbb{S}^n$ para si mesma.
}
%%%% QUESTÃO 26 %%%%
\questao{
	Seja $F : \mathbb{R}^{n+1} \setminus \{0\} \to \mathbb{R}^{n+1}$ dada por $F(x) = \frac{x}{|x|}$. Mostre que $F$ é uma função suave e $F(\mathbb{R}^{n+1} \setminus \{0\}) = \mathbb{S}^n$.
}
\solucao{
	
	\textbf{Análise do Enunciado}
	\begin{itemize}
		\item \textbf{Hipóteses:} $F$ é a aplicação de normalização de vetores.
		\item \textbf{Tese 1:} Provar a suavidade de $F$.
		\item \textbf{Tese 2:} Provar que a imagem de $F$ é exatamente a esfera $\mathbb{S}^n$.
	\end{itemize}
	\vspace{5pt}
	\textbf{1. Prova da Suavidade de $F$:}
	Analisamos a suavidade de $F$ examinando a expressão de suas componentes. Para qualquer $x = (x^1, \dots, x^{n+1})$ pertencente ao domínio $\mathbb{R}^{n+1} \setminus \{0\}$, a $i$-ésima componente de $F(x)$ é:
	\begin{equation*}
		F_i(x) = \frac{x^i}{\sqrt{(x^1)^2 + \dots + (x^{n+1})^2}}.
	\end{equation*}
	Para provar que $F_i$ é suave, analisamos a composição de funções:
	\begin{itemize}
		\item A função $g: \mathbb{R}^{n+1} \to \mathbb{R}$ dada por $g(x) = \sum_{j=1}^{n+1} (x^j)^2$ é um polinômio, sendo, portanto, suave em todo $\mathbb{R}^{n+1}$.
		\item A função $\sqrt{\cdot}$ é suave no intervalo $(0, \infty)$. Como o domínio de $F$ é precisamente $\mathbb{R}^{n+1} \setminus \{0\}$, temos que $g(x) > 0$ para todo $x$ no domínio. Assim, a composição $\sqrt{g(x)} = |x|$ é suave em $\mathbb{R}^{n+1} \setminus \{0\}$.
		\item O numerador $x^i$ é uma função linear, logo é suave em $\mathbb{R}^{n+1}$.
	\end{itemize}
	Sendo $F_i$ o quociente de duas funções suaves onde o denominador $|x|$ nunca se anula no domínio $\mathbb{R}^{n+1} \setminus \{0\}$, concluímos que cada componente $F_i$ é suave. Portanto, $F$ é uma aplicação suave no aberto $\mathbb{R}^{n+1} \setminus \{0\}$.
	
	\textbf{2. Prova de que a Imagem é $\mathbb{S}^n$:}
	Para mostrar que $F(\mathbb{R}^{n+1} \setminus \{0\}) = \mathbb{S}^n$, provaremos as duas inclusões:
	
	($\subseteq$) Seja $x \in \mathbb{R}^{n+1} \setminus \{0\}$. Calculamos a norma da imagem $F(x)$:
	\begin{align*}
		\|F(x)\| &= \left\| \frac{1}{|x|} \cdot x \right\| = \frac{1}{|x|} \|x\| = \frac{|x|}{|x|} = 1.
	\end{align*}
	Como a norma de $F(x)$ é igual a $1$, temos que $F(x) \in \mathbb{S}^n$ para todo $x$ no domínio. Logo, $F(\mathbb{R}^{n+1} \setminus \{0\}) \subseteq \mathbb{S}^n$.
	
	($\supseteq$) Seja $y \in \mathbb{S}^n$. Por definição, $\|y\| = 1$. Note que $y \neq 0$, logo $y$ pertence ao domínio de $F$. Calculando a imagem de $y$:
	\begin{align*}
		F(y) = \frac{y}{|y|} = \frac{y}{1} = y.
	\end{align*}
	Isso mostra que qualquer ponto $y \in \mathbb{S}^n$ possui um pré-imaginário no domínio de $F$ (o próprio $y$). Logo, $$\mathbb{S}^n \subseteq F(\mathbb{R}^{n+1} \setminus \{0\}).$$
	
	Tendo provadas as duas inclusões, concluímos que $F(\mathbb{R}^{n+1} \setminus \{0\}) = \mathbb{S}^n$.
}
%%%% QUESTÃO 27 %%%%	
\questao{
Mostre que a aplicação $F : \mathbb{R}^3 \to \mathbb{R}^6$ definida por $F(x, y, z) = (x^2, y^2, z^2, \sqrt{2}xy, \sqrt{2}xz, \sqrt{2}yz)$ induz uma aplicação suave injetiva $f : \mathbb{RP}^2 \to \mathbb{S}^5$.
}
\solucao{
	
	\textbf{Análise do Enunciado}
	\begin{itemize}
		\item \textbf{Hipóteses:} $F$ é a aplicação polinomial dada.
		\item \textbf{Tese:} Provar que $F$ induz uma aplicação $f: \mathbb{RP}^2 \to \mathbb{S}^5$ que seja suave e injetiva.
	\end{itemize}
	\vspace{5pt}
	\textbf{1. Indução e Suavidade:}
	Com base nos resultados das \textbf{Questões 19 e 20}, sabemos que $F(\mathbb{S}^2) \subseteq \mathbb{S}^5$ e $F(x) = F(-x)$ para todo $x \in \mathbb{S}^2$. Portanto, a aplicação $f([x]) = F(x)$ está bem definida e é suave, pois é a descida de uma aplicação suave para um espaço quociente via projeção canônica.
	
	\textbf{2. Prova de Injetividade via Sistema de Equações não Linear:}
	Sejam $[x], [y] \in \mathbb{RP}^2$ dois pontos tais que $f([x]) = f([y])$. Isso implica que, para representantes $x, y \in \mathbb{S}^2$, temos a igualdade $F(x) = F(y)$, o que é equivalente ao seguinte sistema de equações:
	\begin{equation*}
		\begin{cases}
			x_1^2 = y_1^2 \\
			x_2^2 = y_2^2 \\
			x_3^2 = y_3^2 \\
			x_1x_2 = y_1y_2 \\
			x_1x_3 = y_1y_3 \\
			x_2x_3 = y_2y_3
		\end{cases}
	\end{equation*}
	
	Como $x \in \mathbb{S}^2$, temos $\sum x_i^2 = 1$, logo $x \neq 0$. Suponhamos, sem perda de generalidade, que $x_1 \neq 0$. A primeira equação do sistema implica que $y_1 = \pm x_1$.
	
	Analisamos a consistência dos sinais nos dois casos possíveis:
	
	\textbf{Caso 1: $y_1 = x_1$}
	Substituindo $y_1 = x_1$ nas equações de produtos mistos:
	\begin{align*}
		x_1x_2 = y_1y_2 &\implies x_1x_2 = x_1y_2 \implies x_1(x_2 - y_2) = 0 \\
		x_1x_3 = y_1y_3 &\implies x_1x_3 = x_1y_3 \implies x_1(x_3 - y_3) = 0.
	\end{align*}
	Como $x_1 \neq 0$, deduzimos que $y_2 = x_2$ e $y_3 = x_3$. Portanto, $y = x$.
	
	\textbf{Caso 2: $y_1 = -x_1$}
	Substituindo $y_1 = -x_1$ nas equações de produtos mistos:
	\begin{align*}
		x_1x_2 = y_1y_2 &\implies x_1x_2 = -x_1y_2 \implies x_1(x_2 + y_2) = 0 \\
		x_1x_3 = y_1y_3 &\implies x_1x_3 = -x_1y_3 \implies x_1(x_3 + y_3) = 0.
	\end{align*}
	Como $x_1 \neq 0$, deduzimos que $y_2 = -x_2$ e $y_3 = -x_3$. Portanto, $y = -x$.
	
	Em ambos os casos, concluímos que $y = \pm x$. Pela definição de equivalência no espaço projetivo real, isso implica que:
	\begin{equation*}
		[x] = [y].
	\end{equation*}
	
	Como a igualdade das imagens $f([x]) = f([y])$ implica a igualdade dos pontos $[x] = [y]$, a aplicação $f$ é injetiva.
	
	\textbf{2* (alternativa). Prova de Injetividade via Operadores Lineares:}
	Para provar a injetividade, utilizaremos a identificação de $F$ com o espaço de matrizes simétricas $S(3, \mathbb{R})$ discutida na \textbf{Questão 9}. Observemos que a aplicação $F(x)$ representa as componentes da matriz simétrica de posto 1 dada pelo produto externo $A_x = x x^T$. Especificamente:
	\begin{equation*}
		A_x = x x^T = \begin{pmatrix} x_1^2 & x_1x_2 & x_1x_3 \\ x_2x_1 & x_2^2 & x_2x_3 \\ x_3x_1 & x_3x_2 & x_3^2 \end{pmatrix}.
	\end{equation*}
	Sejam $[x], [y] \in \mathbb{RP}^2$ tais que $f([x]) = f([y])$. Isso implica que $F(x) = F(y)$, o que equivale a dizer que as matrizes simétricas associadas são idênticas:
	\begin{equation*}
		x x^T = y y^T.
	\end{equation*}
	Como $x, y \in \mathbb{S}^2$, ambos são vetores não nulos. Para qualquer matriz da forma $v v^T$, a imagem do operador linear é o subespaço unidimensional gerado por $v$. Assim, a igualdade de operadores implica a igualdade de suas imagens:
	\begin{equation*}
		\text{Im}(x x^T) = \text{Im}(y y^T) \implies \text{span}\{x\} = \text{span}\{y\}.
	\end{equation*}
	A igualdade de subespaços unidimensionais implica que $x$ e $y$ são colineares, ou seja, existe $\lambda \in \mathbb{R} \setminus \{0\}$ tal que $y = \lambda x$. Como $\|x\| = 1$ e $\|y\| = 1$, temos que $|\lambda| = 1$, logo $\lambda = \pm 1$.
	
	Portanto, $y = x$ ou $y = -x$. Pela definição de equivalência no espaço projetivo, isso significa que $[x] = [y]$. Assim, a aplicação $f$ é injetiva.
	
	Como $f$ é suave e injetiva, a tese está provada.	
}
%%%% QUESTÃO 28 %%%%
\questao{
	Mostre que a aplicação $F : \mathbb{R}^3 \to \mathbb{R}^5$ definida por $$F(x, y, z) = \left(\sqrt{3}xy, \sqrt{3}xz, \sqrt{3}yz, \frac{\sqrt{3}}{2}(x^2 - y^2), \frac{1}{2}x^2 + \frac{1}{2}y^2 - z^2\right),$$ induz uma aplicação suave injetiva $\bar{f} : \mathbb{RP}^2 \to \mathbb{S}^4$.
}
\solucao{
	\textbf{Análise do Enunciado}
	\begin{itemize}
		\item \textbf{Hipóteses:} $F$ é uma aplicação polinomial de $\mathbb{R}^3$ para $\mathbb{R}^5$.
		\item \textbf{Tese 1:} Provar que $F$ induz uma aplicação bem definida $\bar{f}: \mathbb{RP}^2 \to \mathbb{S}^4$ (isso requer que $F$ seja constante nas fibras de $\pi: \mathbb{S}^2 \to \mathbb{RP}^2$ e que a imagem de $\mathbb{S}^2$ esteja em $\mathbb{S}^4$).
		\item \textbf{Tese 2:} Provar que $\bar{f}$ é suave.
		\item \textbf{Tese 3:} Provar que $\bar{f}$ é injetiva.
	\end{itemize}
	\vspace{5pt}
	\textbf{1. Indução e Suavidade:}
	Para que $F$ induza uma aplicação $\bar{f} : \mathbb{RP}^2 \to \mathbb{R}^5$, ela deve ser constante nas fibras da projeção $\pi : \mathbb{S}^2 \to \mathbb{RP}^2$. Seja $p = (x, y, z) \in \mathbb{S}^2$. Calculamos $F(-p)$:
	\begin{align*}
		F(-x, -y, -z) &= \left(\sqrt{3}(-x)(-y), \sqrt{3}(-x)(-z), \sqrt{3}(-y)(-z), \frac{\sqrt{3}}{2}((-x)^2 - (-y)^2), \frac{1}{2}(-x)^2 + \frac{1}{2}(-y)^2 - (-z)^2\right) \\
		&= \left(\sqrt{3}xy, \sqrt{3}xz, \sqrt{3}yz, \frac{\sqrt{3}}{2}(x^2 - y^2), \frac{1}{2}x^2 + \frac{1}{2}y^2 - z^2\right) = F(x, y, z).
	\end{align*}
	Como $F(x) = F(-x)$, a aplicação $\bar{f}([x]) = F(x)$ está bem definida. Além disso, como $F$ é definida por componentes polinomiais, ela é suave em $\mathbb{R}^3$. Sendo $\mathbb{RP}^2$ e $\mathbb{S}^4$ variedades suaves, a aplicação induzida $\bar{f}$ é suave.
	
	\textbf{2. Verificação da Imagem em $\mathbb{S}^4$:}
	Seja $p = (x, y, z) \in \mathbb{S}^2$, logo $x^2 + y^2 + z^2 = 1$. Calculamos a norma de $F(p)$ ao quadrado:
	\begin{align*}
		\|F(p)\|^2 &= (\sqrt{3}xy)^2 + (\sqrt{3}xz)^2 + (\sqrt{3}yz)^2 + \left(\frac{\sqrt{3}}{2}(x^2 - y^2)\right)^2 + \left(\frac{1}{2}x^2 + \frac{1}{2}y^2 - z^2\right)^2 \\
		&= 3x^2y^2 + 3x^2z^2 + 3y^2z^2 + \frac{3}{4}(x^4 + y^4 - 2x^2y^2) + \left(\frac{1}{4}x^4 + \frac{1}{4}y^4 + z^4 + \frac{1}{2}x^2y^2 - x^2z^2 - y^2z^2\right) \\
		&= \left(\frac{3}{4} + \frac{1}{4}\right)x^4 + \left(\frac{3}{4} + \frac{1}{4}\right)y^4 + z^4 + \left(3 - \frac{3}{2} + \frac{1}{2}\right)x^2y^2 + (3 - 1)x^2z^2 + (3 - 1)y^2z^2 \\
		&= x^4 + y^4 + z^4 + 2x^2y^2 + 2x^2z^2 + 2y^2z^2 \\
		&= (x^2 + y^2 + z^2)^2 = (1)^2 = 1.
	\end{align*}
	Portanto, a imagem de $\bar{f}$ está contida em $\mathbb{S}^4$.
	
	\textbf{3. Prova de Injetividade:}
	Sejam $[x], [y] \in \mathbb{RP}^2$ tais que $\bar{f}([x]) = \bar{f}([y])$. Isso implica que $F(x) = F(y)$, resultando no sistema:
	\begin{equation*}
		\begin{cases}
			\sqrt{3}x_1x_2 = \sqrt{3}y_1y_2 \implies x_1x_2 = y_1y_2 \\
			\sqrt{3}x_1x_3 = \sqrt{3}y_1y_3 \implies x_1x_3 = y_1y_3 \\
			\sqrt{3}x_2x_3 = \sqrt{3}y_2y_3 \implies x_2x_3 = y_2y_3 \\
			\frac{\sqrt{3}}{2}(x_1^2 - x_2^2) = \frac{\sqrt{3}}{2}(y_1^2 - y_2^2) \implies x_1^2 - x_2^2 = y_1^2 - y_2^2 \\
			\frac{1}{2}x_1^2 + \frac{1}{2}x_2^2 - x_3^2 = \frac{1}{2}y_1^2 + \frac{1}{2}y_2^2 - y_3^2
		\end{cases}
	\end{equation*}
	Da última equação, e sabendo que $x_1^2 + x_2^2 = 1 - x_3^2$ e $y_1^2 + y_2^2 = 1 - y_3^2$, temos:
	\begin{equation*}
		\frac{1}{2}(1 - x_3^2) - x_3^2 = \frac{1}{2}(1 - y_3^2) - y_3^2 \implies \frac{1}{2} - \frac{3}{2}x_3^2 = \frac{1}{2} - \frac{3}{2}y_3^2 \implies x_3^2 = y_3^2.
	\end{equation*}
	Se $x_3^2 = y_3^2$, então $x_1^2 + x_2^2 = y_1^2 + y_2^2$. Combinando isso com a equação $x_1^2 - x_2^2 = y_1^2 - y_2^2$ através de soma e subtração, obtemos:
	\begin{equation*}
		2x_1^2 = 2y_1^2 \implies x_1^2 = y_1^2 \quad \text{e} \quad 2x_2^2 = 2y_2^2 \implies x_2^2 = y_2^2.
	\end{equation*}
	Agora temos o mesmo sistema da \textbf{Questão 27}: $x_i^2 = y_i^2$ para todo $i$ e $x_i x_j = y_i y_j$ para $i \neq j$. Como demonstrado anteriormente, esse sistema implica necessariamente que $y = x$ ou $y = -x$.
	
	Assim, $[x] = [y]$, o que prova que $\bar{f}$ é injetiva.
}
%%%% RESULTADO EXTRA %%%%
Antes de continuarmos para as próximas questões, é importante estabelecer um resultado fundamental, a saber:
	
\resultadoextra{Teorema da Partição da Unidade.}{
	Seja $M$ uma variedade suave e $\{U_\alpha\}_{\alpha \in \Lambda}$ uma cobertura aberta de $M$. Existe uma família de funções suaves $\{\rho_i\}_{i \in I}$ tal que:
	\begin{enumerate}[label=(\roman*)]
		\item $0 \le \rho_i(p) \le 1$ para todo $p \in M$ e todo $i \in I$.
		\item O suporte de cada $\rho_i$ é compacto e $\text{supp}(\rho_i) \subset U_\alpha$ para algum $\alpha \in \Lambda$.
		\item A família $\{\text{supp}(\rho_i)\}_{i \in I}$ é localmente finita.
		\item $\sum_{i \in I} \rho_i(p) = 1$ para todo $p \in M$.
	\end{enumerate}
}
\solucao{
	
	\textbf{Análise do Enunciado}
	\begin{itemize}
		\item \textbf{Hipóteses:} $M$ é uma variedade suave e $\{U_\alpha\}$ é uma cobertura aberta.
		\item \textbf{Tese:} Provar a existência de uma família de funções suaves com suporte compacto, localmente finita, subordinada à cobertura e com soma unitária.
	\end{itemize}
	\vspace{5pt}
	A demonstração exige a construção de funções suaves com suporte compacto, a refinação da cobertura e a normalização global.
	
	\textbf{1. Construção de Funções Suaves com Suporte Compacto em $\mathbb{R}^n$:}
	Consideremos a função $\psi : \mathbb{R} \to \mathbb{R}$ definida por:
	\begin{equation*}
		\psi(t) = \begin{cases} e^{-1/t} & \text{se } t > 0 \\ 0 & \text{se } t \le 0 \end{cases}
	\end{equation*}
	É um resultado clássico de análise que $\psi$ é suave ($C^\infty$) em todo $\mathbb{R}$, incluindo $t=0$, onde todas as derivadas são nulas. A partir de $\psi$, definimos a função de transição $\eta : \mathbb{R} \to [0, 1]$:
	\begin{equation*}
		\eta(t) = \frac{\psi(t)}{\psi(t) + \psi(1-t)}.
	\end{equation*}
	Note que $\eta(t) = 0$ para $t \le 0$ e $\eta(t) = 1$ para $t \ge 1$. Para qualquer bola aberta $B(x, r) \subset \mathbb{R}^n$, podemos definir uma função suave $\phi_B : \mathbb{R}^n \to [0, 1]$ que é $1$ em $B(x, r/2)$ e $0$ fora de $B(x, r)$, utilizando a composição de $\eta$ com a norma euclidiana:
	\begin{equation*}
		\phi_B(y) = \eta\left( \frac{r/2 - |y-x|}{r/2} \right).
	\end{equation*}
	
	\textbf{2. Refinação da Cobertura e Local Finitude:}
	Como $M$ é uma variedade suave, ela é paracompacta. Portanto, para a cobertura $\{U_\alpha\}$, existe um refinamento aberto $\{V_i\}_{i \in I}$ que é localmente finito e tal que cada $\overline{V}_i$ é compacto e $\overline{V}_i \subset U_\alpha$ para algum $\alpha$. 
	
	Para cada $V_i$, escolhemos uma carta coordenada $(W_i, \varphi_i)$ tal que $V_i \subset W_i$ e a imagem $\varphi_i(V_i)$ seja um aberto em $\mathbb{R}^n$. Podemos então construir, via a função $\phi_B$ definida anteriormente, uma função suave $\psi_i : M \to [0, \infty)$ tal que:
	\begin{itemize}
		\item $\text{supp}(\psi_i) \subset V_i$ (consequentemente, o suporte é compacto e contido em algum $U_\alpha$).
		\item $\psi_i(p) > 0$ para todo $p \in \text{int}(V_i)$.
	\end{itemize}
	
	\textbf{3. Normalização Global:}
	Seja $\Phi : M \to \mathbb{R}$ a soma de todas as funções $\psi_i$:
	\begin{equation*}
		\Phi(p) = \sum_{i \in I} \psi_i(p).
	\end{equation*}
	Como a família $\{\text{supp}(\psi_i)\}$ é localmente finita, para qualquer $p \in M$ existe uma vizinhança onde apenas um número finito de termos da soma é não nulo. Como cada $\psi_i$ é suave, $\Phi$ é suave. Além disso, como $\{V_i\}$ cobre $M$, para cada $p$ existe ao menos um $i$ tal que $\psi_i(p) > 0$, logo $\Phi(p) > 0$ para todo $p \in M$.
	
	Definimos agora as funções $\rho_i : M \to [0, 1]$ por:
	\begin{equation*}
		\rho_i(p) = \frac{\psi_i(p)}{\Phi(p)}.
	\end{equation*}
	
	\textbf{4. Verificação dos Axiomas:}
	\begin{itemize}
		\item \textbf{Suavidade:} $\rho_i$ é o quociente de duas funções suaves com denominador nunca nulo, portanto $\rho_i$ é suave.
		\item \textbf{Suporte:} Como $\text{supp}(\rho_i) = \text{supp}(\psi_i)$, o suporte de $\rho_i$ é compacto e $\text{supp}(\rho_i) \subset U_\alpha$ para algum $\alpha$.
		\item \textbf{Local Finitude:} A local finitude de $\{\rho_i\}$ decorre diretamente da local finitude de $\{\psi_i\}$.
		\item \textbf{Soma Unitária:} Para qualquer $p \in M$:
		\begin{equation*}
			\sum_{i \in I} \rho_i(p) = \sum_{i \in I} \frac{\psi_i(p)}{\Phi(p)} = \frac{1}{\Phi(p)} \sum_{i \in I} \psi_i(p) = \frac{\Phi(p)}{\Phi(p)} = 1.
		\end{equation*}
	\end{itemize}
	
	Assim, concluímos que toda a estrutura necessária para a partição da unidade foi construída.
}
%%%% DEFINIÇÃO %%%%	
\textbf{Definição.} Sejam $M^n$ e $N^k$ variedades suaves e $A \subset M^n$ um subconjunto qualquer. Dizemos que uma função $F : A \to N^k$ é suave se para cada $p \in A$ existe um conjunto aberto $W \subset M^n$ contendo $p$ e uma aplicação suave $\tilde{F} : W \to N^k$ tal que $\tilde{F}|_{W \cap A} = F$.
%%%% QUESTÃO 29 %%%%
\questao{
	\textbf{(Lema de Extensão)} Sejam $M^n$ uma $n$-variedade suave, $A \subset M^n$ um subconjunto fechado e $f : A \to \mathbb{R}^k$ uma função suave. Para qualquer conjunto aberto $U$ contendo $A$, mostre que existe uma função suave $\tilde{f} : M^n \to \mathbb{R}^k$ tal que $\tilde{f}|_A = f$ e $\text{supp } \tilde{f} \subset U$.
}
\solucao{
	
	\textbf{Análise do Enunciado}
	\begin{itemize}
		\item \textbf{Hipóteses:} $A$ é um subconjunto fechado de $M^n$, $f$ é suave em $A$, e $U$ é um aberto tal que $A \subset U$.
		\item \textbf{Tese:} Provar a existência de uma extensão suave $\tilde{f}$ para todo $M^n$ com suporte contido em $U$.
	\end{itemize}
	\vspace{5pt}
	\textbf{1. Construção da Cobertura:}
	Pela definição de suavidade em subconjuntos, para cada $p \in A$, existe um aberto $W_p \subset M^n$ e uma extensão suave $f_p : W_p \to \mathbb{R}^k$ tal que $f_p|_{W_p \cap A} = f$. Definimos $W'_p = W_p \cap U$, garantindo que cada $W'_p$ seja um aberto contendo $p$ e esteja contido em $U$. A coleção $\mathcal{W} = \{W'_p\}_{p \in A} \cup \{M^n \setminus A\}$ constitui uma cobertura aberta de $M^n$.
	
	\textbf{2. Aplicação da Partição da Unidade:}
	Pelo Teorema da Partição da Unidade, existe uma família de funções suaves $\{\rho_i\}_{i \in I}$ subordinada a $\mathcal{W}$. Para cada $\rho_i$, associamos a função $f_i$:
	\begin{itemize}
		\item Se $\text{supp}(\rho_i) \subset W'_p$, definimos $f_i = f_p$.
		\item Se $\text{supp}(\rho_i) \subset M^n \setminus A$, definimos $f_i \equiv 0$.
	\end{itemize}
	Definimos a extensão global $\tilde{f} : M^n \to \mathbb{R}^k$ pela soma localmente finita:
	\begin{equation*}
		\tilde{f}(q) = \sum_{i \in I} \rho_i(q) f_i(q).
	\end{equation*}
	
	\textbf{3. Verificação das Propriedades:}
	Como a soma é localmente finita e cada termo é suave, $\tilde{f}$ é suave. Para qualquer $q \in A$, a função $\rho_i$ associada ao aberto $M^n \setminus A$ é nula, e para todas as demais, $f_i(q) = f(q)$, logo:
	\begin{equation*}
		\tilde{f}(q) = \sum \rho_i(q) f(q) = f(q) \sum \rho_i(q) = f(q).
	\end{equation*}
	Finalmente, se $q \notin U$, então $\rho_i(q) = 0$ para todos os índices onde $f_i \neq 0$ (pois tais suportes estão em $W'_p \subset U$), implicando $\tilde{f}(q) = 0$. Portanto, $\text{supp } \tilde{f} \subset U$.
}
%%%% DEFINIÇÃO %%%%
\textbf{Definição.} Se $M$ é um espaço topológico, uma função de exaustão para $M$ é uma função contínua $f : M \to \mathbb{R}$ tal que o conjunto $M_c = \{x \in M : f(x) \le c\}$ é compacto para cada $c \in \mathbb{R}$.
%%%% QUESTÃO 30 %%%%	
\questao{
	\textbf{(Existência de Funções de Exaustão)} Mostre que qualquer variedade suave admite uma função de exaustão positiva suave.
}
\solucao{
	
	\textbf{Análise do Enunciado}
	\begin{itemize}
		\item \textbf{Hipóteses:} $M$ é uma variedade suave.
		\item \textbf{Tese:} Provar a existência de uma função $f: M \to \mathbb{R}$ que seja suave, positiva e tal que as subníveis $M_c$ sejam compactos.
	\end{itemize}
	\vspace{5pt}
	\textbf{1. Definição e a Propriedade de $\sigma$-compacticidade:}
	Dizemos que um espaço topológico $M$ é $\sigma$-compacto se existe uma família enumerável de subconjuntos compactos $\{K_i\}_{i \in \mathbb{N}}$ tal que $M = \bigcup_{i=1}^\infty K_i$. Como toda variedade suave é, por definição, de Hausdorff e possui base enumerável, ela é necessariamente $\sigma$-compacta.
	
	A partir disso, podemos construir uma \textbf{exaustão por compactos} para $M$. Esta é uma sequência de conjuntos compactos $\{K_i\}_{i=1}^\infty$ tal que:
	\begin{equation*}
		K_1 \subset \text{int}(K_2) \subset K_2 \subset \text{int}(K_3) \subset K_3 \subset \dots \quad \text{e} \quad \bigcup_{i=1}^\infty K_i = M.
	\end{equation*}
	
	\textbf{2. Construção da Função via Partição da Unidade:}
	Para cada $i \ge 1$, definimos o aberto $W_i = \text{int}(K_{i+1}) \setminus K_{i-1}$ (adotando $K_0 = \emptyset$). A coleção $\{W_i\}_{i=1}^\infty$ constitui uma cobertura aberta de $M$. Pelo Teorema da Partição da Unidade, existe uma família de funções suaves $\rho_i : M \to [0, 1]$ tal que $\text{supp}(\rho_i) \subset W_i$ e $\sum_{i=1}^\infty \rho_i(x) = 1$ para todo $x \in M$.
	
	Definimos a função $f : M \to \mathbb{R}$ por:
	\begin{equation*}
		f(x) = \sum_{i=1}^\infty i \cdot \rho_i(x).
	\end{equation*}
	
	\textbf{3. Análise de Suavidade e Positividade:}
	Como a família $\{\text{supp}(\rho_i)\}$ é localmente finita, para qualquer $x \in M$ existe uma vizinhança onde apenas um número finito de termos $\rho_i$ é não nulo. Sendo cada termo $i \cdot \rho_i$ suave, a soma $f$ é suave. Além disso, como $\rho_i(x) \ge 0$ e $\sum \rho_i = 1$, temos $f(x) \ge 1 \cdot \sum \rho_i(x) = 1$. Logo, $f$ é estritamente positiva.
	
	\textbf{4. Prova Algébrica da Propriedade de Exaustão:}
	Seja $c \in \mathbb{R}$ um valor qualquer. Queremos provar que o conjunto $M_c = \{x \in M : f(x) \le c\}$ é compacto. 
	
	Escolhemos um inteiro $N \in \mathbb{N}$ tal que $N > c$. Consideremos agora qualquer ponto $x \in M$ que esteja fora do compacto $K_N$, ou seja, $x \in M \setminus K_N$. 
	
	Para tal ponto $x$, as funções $\rho_i$ com $i < N$ devem ser nulas. De fato, para $i < N$, temos $\text{supp}(\rho_i) \subset \text{int}(K_{i+1}) \subseteq K_N$. Como $x \notin K_N$, segue que $\rho_i(x) = 0$ para todo $i \in \{1, 2, \dots, N-1\}$. 
	
	Calculamos então o valor de $f(x)$ para $x \notin K_N$:
	\begin{align*}
		f(x) &= \sum_{i=1}^\infty i \cdot \rho_i(x) \\
		&= \sum_{i=1}^{N-1} i \cdot 0 + \sum_{i=N}^\infty i \cdot \rho_i(x) \\
		&= \sum_{i=N}^\infty i \cdot \rho_i(x) \ge N \sum_{i=N}^\infty \rho_i(x).
	\end{align*}
	Como $\sum_{i=1}^\infty \rho_i(x) = 1$ e $\rho_i(x) = 0$ para $i < N$, temos que $\sum_{i=N}^\infty \rho_i(x) = 1$. Portanto:
	\begin{equation*}
		f(x) \ge N \cdot 1 = N.
	\end{equation*}
	Dado que escolhemos $N > c$, temos que para todo $x \notin K_N$, $f(x) > c$. 
	
	Isso implica que se $f(x) \le c$, então necessariamente $x \in K_N$. Em termos de conjuntos:
	\begin{equation*}
		M_c \subseteq K_N.
	\end{equation*}
	Como $f$ é contínua, $M_c = f^{-1}((-\infty, c])$ é um subconjunto fechado de $M$. Sendo $M_c$ um fechado contido em um compacto $K_N$, o conjunto $M_c$ é, por consequência, compacto.
	
	Assim, $f$ é uma função de exaustão positiva e suave.
}
%%%% QUESTÃO 31 %%%%
\questao{
	Seja $G$ uma variedade suave munida com uma estrutura de grupo tal que a aplicação $G \times G \to G$, $(g, h) \mapsto gh^{-1}$, é suave. Mostre que $G$ é um grupo de Lie.
}
\solucao{
	
	\textbf{Análise do Enunciado}
	\begin{itemize}
		\item \textbf{Hipóteses:} $G$ é uma variedade suave e um grupo onde a operação combinada $(g,h) \mapsto gh^{-1}$ é suave.
		\item \textbf{Tese:} Provar que $G$ é um grupo de Lie (ou seja, que a multiplicação $m$ e a inversão $i$ são suavemente separadas).
	\end{itemize}
	\vspace{5pt}
	Para provar que $G$ é um grupo de Lie, devemos demonstrar que tanto a aplicação de inversão quanto a aplicação de multiplicação são suaves. Seja $\Phi : G \times G \to G$ a aplicação dada por $\Phi(g, h) = gh^{-1}$, que, por hipótese, é suave.
	
	\textbf{1. Suavidade da Aplicação de Inversão:}
	Consideremos a aplicação de inversão $i : G \to G$ definida por $i(g) = g^{-1}$. 
	Seja $e \in G$ o elemento neutro do grupo. Podemos expressar a aplicação de inversão como uma restrição da aplicação $\Phi$. Fixando o primeiro argumento de $\Phi$ como o elemento neutro $e$, obtemos:
	\begin{equation*}
		\Phi(e, g) = eg^{-1} = g^{-1} = i(g).
	\end{equation*}
	Como $\Phi$ é suave em $G \times G$, a aplicação obtida ao fixar uma das coordenadas (uma aplicação de inclusão $g \mapsto (e, g)$ seguida de $\Phi$) é necessariamente suave. Portanto, a aplicação de inversão $i : G \to G$ é suave.
	
	\textbf{2. Suavidade da Aplicação de Multiplicação:}
	Agora, consideremos a aplicação de multiplicação $m : G \times G \to G$ definida por $m(g, h) = gh$. 
	Podemos expressar a operação de multiplicação utilizando a aplicação $\Phi$ e a aplicação de inversão $i$ que acabamos de provar ser suave. Note que:
	\begin{equation*}
		gh = g(h^{-1})^{-1}.
	\end{equation*}
	Formalmente, definimos a aplicação $\Psi : G \times G \to G \times G$ por $\Psi(g, h) = (g, i(h)) = (g, h^{-1})$. Como $i$ é suave, a aplicação $\Psi$ é suave (pois cada componente é suave). 
	
	A aplicação de multiplicação $m$ pode então ser escrita como a composição:
	\begin{equation*}
		m = \Phi \circ \Psi.
	\end{equation*}
	De fato, para quaisquer $g, h \in G$:
	\begin{align*}
		(\Phi \circ \Psi)(g, h) &= \Phi(\Psi(g, h)) \\
		&= \Phi(g, h^{-1}) \\
		&= g(h^{-1})^{-1} \\
		&= gh.
	\end{align*}
	Como $\Psi$ é suave e $\Phi$ é suave, a composição $\Phi \circ \Psi$ é, necessariamente, suave. Portanto, a aplicação de multiplicação $m$ é suave.
	
	Como a multiplicação e a inversão são aplicações suaves, $G$ satisfaz a definição de um Grupo de Lie.
}
%%%% QUESTÃO 32 %%%%
\questao{
	Mostre que $GL(n, \mathbb{R}) = \{ A \in M(n, \mathbb{R}) : \det A \neq 0 \}$, com a multiplicação usual de matrizes, é um grupo de Lie.
}
\solucao{
	
	\textbf{Análise do Enunciado}
	\begin{itemize}
		\item \textbf{Hipóteses:} $GL(n, \mathbb{R})$ definido como o subconjunto de matrizes invertíveis de $M(n, \mathbb{R})$.
		\item \textbf{Tese:} Provar que $GL(n, \mathbb{R})$ é um grupo de Lie.
	\end{itemize}
	\vspace{5pt}
	\textbf{1. Estrutura de Variedade Suave:}
	Conforme provado na \textbf{Questão 8}, o grupo linear geral $GL(n, \mathbb{R})$ é um subconjunto aberto do espaço de matrizes $M(n, \mathbb{R}) \cong \mathbb{R}^{n^2}$, o que garante que ele seja uma variedade suave de dimensão $n^2$.
	
	\textbf{2. Suavidade da Multiplicação:}
	Seja $m : GL(n, \mathbb{R}) \times GL(n, \mathbb{R}) \to GL(n, \mathbb{R})$ a aplicação de multiplicação $m(A, B) = AB$. A componente $(i, j)$ do produto $AB$ é dada por:
	\begin{equation*}
		(AB)_{ij} = \sum_{k=1}^n A_{ik} B_{kj}.
	\end{equation*}
	Como cada componente $(AB)_{ij}$ é um polinômio nas entradas de $A$ e $B$, a aplicação $m$ é infinitamente diferenciável. Portanto, a multiplicação é suave.
	
	\textbf{3. Suavidade da Inversão:}
	Seja $i : GL(n, \mathbb{R}) \to GL(n, \mathbb{R})$ a aplicação de inversão $i(A) = A^{-1}$. Pela Regra de Cramer, as componentes da matriz inversa são dadas por:
	\begin{equation*}
		(A^{-1})_{ij} = \frac{1}{\det A} (\text{adj } A)_{ji},
	\end{equation*}
	onde $(\text{adj } A)_{ji}$ é o cofator correspondente, que é um polinômio nas entradas de $A$.
	
	Visto que o determinante $\det A$ é uma função suave e nunca se anula no domínio $GL(n, \mathbb{R})$, cada componente $(A^{-1})_{ij}$ é o quociente de duas funções suaves com denominador não nulo. Assim, a aplicação de inversão $i$ é suave.
	
	Tendo sido estabelecida a estrutura de variedade suave e a suavidade das operações de grupo, concluímos que $GL(n, \mathbb{R})$ é um grupo de Lie.
}
%%%% QUESTÃO 33 %%%%
\questao{
	Mostre que $SL(n, \mathbb{R}) = \{ A \in M(n, \mathbb{R}) : \det A = 1 \}$, com a multiplicação usual de matrizes, é um grupo de Lie.
}
\solucao{
	
	\textbf{Análise do Enunciado}
	\begin{itemize}
		\item \textbf{Hipóteses:} $SL(n, \mathbb{R})$ é o subconjunto de matrizes com determinante unitário.
		\item \textbf{Tese:} Provar que $SL(n, \mathbb{R})$ é um grupo de Lie.
	\end{itemize}
	\vspace{5pt}
	\textbf{1. Estrutura de Variedade Suave:}
	Conforme prov도 na \textbf{Questão 10}, a aplicação determinante $\det : M(n, \mathbb{R}) \to \mathbb{R}$ possui o valor $1$ como um valor regular. Pelo Teorema do Valor Regular, o conjunto nível $SL(n, \mathbb{R}) = \det^{-1}(1)$ é uma subvariedade suave de $M(n, \mathbb{R})$ de dimensão $n^2 - 1$.
	
	\textbf{2. Verificação da Estrutura de Subgrupo:}
	Para que $SL(n, \mathbb{R})$ seja um grupo de Lie, as operações de grupo devem estar bem definidas dentro do conjunto. Seja $A, B \in SL(n, \mathbb{R})$, logo $\det A = 1$ e $\det B = 1$.
	
	Pela propriedade multiplicativa do determinante:
	\begin{equation*}
		\det(AB) = \det(A) \cdot \det(B) = 1 \cdot 1 = 1.
	\end{equation*}
	Assim, $AB \in SL(n, \mathbb{R})$, provando que o conjunto é fechado sob a multiplicação. Para a inversão, temos:
	\begin{equation*}
		\det(A^{-1}) = \frac{1}{\det A} = \frac{1}{1} = 1.
	\end{equation*}
	Assim, $A^{-1} \in SL(n, \mathbb{R})$, provando que o conjunto é fechado sob a inversão. Portanto, $SL(n, \mathbb{R})$ é um subgrupo de $GL(n, \mathbb{R})$.
	
	\textbf{3. Suavidade das Operações:}
	Seja $m : GL(n, \mathbb{R}) \times GL(n, \mathbb{R}) \to GL(n, \mathbb{R})$ a aplicação de multiplicação e $i : GL(n, \mathbb{R}) \to GL(n, \mathbb{R})$ a aplicação de inversão. Na \textbf{Questão 32}, provamos que $m$ e $i$ são aplicações suaves.
	
	Como $SL(n, \mathbb{R})$ é uma subvariedade suave de $GL(n, \mathbb{R})$, a restrição da aplicação de multiplicação $m|_{SL \times SL}$ e a restrição da aplicação de inversão $i|_{SL}$ são, por definição, aplicações suaves entre as respectivas variedades.
	
	Como $SL(n, \mathbb{R})$ é uma variedade suave e suas operações de grupo são suaves, concluímos que $SL(n, \mathbb{R})$ é um grupo de Lie.
}
%%%% QUESTÃO 34 %%%%
\questao{
	Mostre que $O(n, \mathbb{R}) = \{ A \in M(n, \mathbb{R}) : AA^T = I \}$, com a multiplicação usual de matrizes, é um grupo de Lie.
}
\solucao{
	
	\textbf{Análise do Enunciado}
	\begin{itemize}
		\item \textbf{Hipóteses:} $O(n, \mathbb{R})$ é o conjunto das matrizes ortogonais.
		\item \textbf{Tese:} Provar que $O(n, \mathbb{R})$ é um grupo de Lie.
	\end{itemize}
	\vspace{5pt}
	\textbf{1. Estrutura de Variedade Suave:}
	Conforme provado detalhadamente na \textbf{Questão 9}, definimos a aplicação $F : M(n, \mathbb{R}) \to S(n, \mathbb{R})$ por $F(A) = AA^T$. Provamos que a identidade $I$ é um valor regular de $F$, o que, pelo Teorema do Valor Regular, implica que o grupo ortogonal $O(n) = F^{-1}(I)$ é uma subvariedade suave de $M(n, \mathbb{R})$ de dimensão $\frac{n(n-1)}{2}$.
	
	\textbf{2. Verificação da Estrutura de Subgrupo:}
	Para que $O(n)$ seja um grupo de Lie, as operações de grupo devem estar bem definidas dentro do conjunto. Sejam $A, B \in O(n)$, logo $AA^T = I$ e $BB^T = I$.
	
	Verifiquemos a multiplicação:
	\begin{equation*}
		(AB)(AB)^T = ABB^T A^T = A I A^T = AA^T = I.
	\end{equation*}
	Assim, $AB \in O(n)$, provando que o conjunto é fechado sob a multiplicação. Para a inversão, note que se $AA^T = I$, então a inversa de $A$ é precisamente a sua transposta, $A^{-1} = A^T$. Verificamos:
	\begin{equation*}
		(A^{-1})(A^{-1})^T = A^T (A^T)^T = A^T A.
	\end{equation*}
	Como $AA^T = I$, a matriz $A$ é invertível e $A^T = A^{-1}$, o que implica que $A^T A = I$. Logo, $A^{-1} \in O(n)$, provando que o conjunto é fechado sob a inversão. Portanto, $O(n)$ é um subgrupo de $GL(n, \mathbb{R})$.
	
	\textbf{3. Suavidade das Operações:}
	Na \textbf{Questão 32}, provamos que o grupo linear geral $GL(n, \mathbb{R})$ é um Grupo de Lie, o que significa que a aplicação de multiplicação $m : GL \times GL \to GL$ e a aplicação de inversão $i : GL \to GL$ são suaves.
	
	Como $O(n)$ é uma subvariedade suave de $GL(n, \mathbb{R})$, a restrição da aplicação de multiplicação $m|_{O \times O}$ e a restrição da aplicação de inversão $i|_{O}$ são, por definição, aplicações suaves entre as respectivas variedades.
	
	Como $O(n)$ é uma variedade suave e suas operações de grupo são suaves, concluímos que $O(n)$ é um grupo de Lie.
}
%%%% QUESTÃO 35 %%%%
\questao{
	Seja $V$ um espaço vetorial de dimensão finita e considere o conjunto $$GL(V) = \{ T : V \to V : T \text{ é transformação linear invertível} \}.$$ Com a composição de aplicações, mostre que $GL(V)$ é um grupo de Lie.
}
\solucao{

	\textbf{Análise do Enunciado}
	\begin{itemize}
		\item \textbf{Hipóteses:} $V$ espaço vetorial de dimensão finita; $GL(V)$ o conjunto de automorfismos de $V$.
		\item \textbf{Tese:} Provar que $GL(V)$ é um grupo de Lie sob composição.
	\end{itemize}
	\vspace{5pt}
	\textbf{1. Estrutura de Variedade Suave:}
	Seja $n = \dim(V)$. Consideremos o espaço de todos os endomorfismos de $V$, denotado por $\text{End}(V)$, que consiste no conjunto de todas as transformações lineares de $V$ em si mesmo ($\text{End}(V) = \{ T : V \to V \mid T \text{ é linear} \}$).
	
	Fixemos uma base $\mathcal{B} = \{e_1, \dots, e_n\}$ de $V$. A aplicação $\Phi : \text{End}(V) \to M(n, \mathbb{R})$, que associa a cada operador linear $T$ a sua matriz representativa $[T]_{\mathcal{B}}$, é um isomorfismo de espaços vetoriais. Como $M(n, \mathbb{R})$ é isomorfo a $\mathbb{R}^{n^2}$, ele é uma variedade suave. 
	
	O grupo $GL(V)$ é o subconjunto de $\text{End}(V)$ composto pelas transformações invertíveis. Através do isomorfismo $\Phi$, temos que:
	\begin{equation*}
		\Phi(GL(V)) = GL(n, \mathbb{R}).
	\end{equation*}
	Conforme provamos na \textbf{Questão 8}, $GL(n, \mathbb{R})$ é um subconjunto aberto de $M(n, \mathbb{R})$, logo é uma variedade suave. Como $\Phi$ é um isomorfismo linear (e, portanto, um difeomorfismo), $GL(V)$ herda a estrutura de variedade suave de $GL(n, \mathbb{R})$, sendo uma variedade de dimensão $n^2$.
	
	\textbf{2. Suavidade da Operação de Grupo:}
	A operação de grupo em $GL(V)$ é a composição de aplicações $\circ$. Para quaisquer $T, S \in GL(V)$, a matriz do operador composto $T \circ S$ em relação à base $\mathcal{B}$ é o produto das matrizes:
	\begin{equation*}
		[T \circ S]_{\mathcal{B}} = [T]_{\mathcal{B}} \cdot [S]_{\mathcal{B}}.
	\end{equation*}
	Vimos na \textbf{Questão 32} que a multiplicação de matrizes em $GL(n, \mathbb{R})$ é uma aplicação suave. Como a composição em $GL(V)$ é isomorfa à multiplicação em $GL(n, \mathbb{R})$ via $\Phi$, a aplicação de multiplicação $m : GL(V) \times GL(V) \to GL(V)$ é suave.
	
	\textbf{3. Suavidade da Inversão:}
	De modo análogo, a matriz do operador inverso $T^{-1}$ é a inversa da matriz de $T$:
	\begin{equation*}
		[T^{-1}]_{\mathcal{B}} = ([T]_{\mathcal{B}})^{-1}.
	\end{equation*}
	Uma vez que a aplicação de inversão de matrizes em $GL(n, \mathbb{R})$ é suave (conforme a \textbf{Questão 32}), a aplicação de inversão $i : GL(V) \to GL(V)$ dada por $i(T) = T^{-1}$ também é suave.
	
	Como $GL(V)$ é uma variedade suave e suas operações de grupo também são suaves, concluímos que $GL(V)$ é um grupo de Lie.
}
%%%% QUESTÃO 36 %%%%
\questao{
	\begin{enumerate}[label=(\alph*)]
		\item Mostre que $\mathbb{R}^n$ é um grupo de Lie com respeito à operação de adição usual.
		\item Mostre que $\mathbb{R}^* = \mathbb{R} \setminus \{0\}$ é um grupo de Lie 1-dimensional com respeito à multiplicação usual de números reais.
	\end{enumerate}
}
\solucao{
	
	\textbf{Análise do Enunciado}
	\begin{itemize}
		\item \textbf{Tese (a):} Provar que $(\mathbb{R}^n, +)$ é um grupo de Lie.
		\item \textbf{Tese (b):} Provar que $(\mathbb{R}^*, \cdot)$ é um grupo de Lie de dimensão 1.
	\end{itemize}
	\vspace{5pt}
	\textbf{(a) O Grupo $(\mathbb{R}^n, +)$:}
	
	\textbf{1. Estrutura de Variedade:}
	A aplicação identidade $\text{Id}: \mathbb{R}^n \to \mathbb{R}^n$ fornece um atlas suave para $\mathbb{R}^n$, tornando-o uma variedade suave de dimensão $n$.
	
	\textbf{2. Suavidade da Operação de Grupo:}
	Neste grupo, a operação de \textbf{multiplicação} é a adição usual. Definimos $m : \mathbb{R}^n \times \mathbb{R}^n \to \mathbb{R}^n$ por $m(x, y) = x + y$. Em termos de componentes, temos:
	\begin{equation*}
		m(x, y) = (x_1 + y_1, x_2 + y_2, \dots, x_n + y_n).
	\end{equation*}
	Como cada componente $m_i(x, y) = x_i + y_i$ é uma função linear (e, portanto, polinomial), a aplicação $m$ é suave.
	
	\textbf{3. Suavidade da Inversão:}
	A aplicação de inversão $i : \mathbb{R}^n \to \mathbb{R}^n$ é dada pelo inverso aditivo: $i(x) = -x$. Componentemente:
	\begin{equation*}
		i(x) = (-x_1, -x_2, \dots, -x_n).
	\end{equation*}
	Sendo a aplicação $i$ linear, ela é infinitamente diferenciável. Portanto, $(\mathbb{R}^n, +)$ é um grupo de Lie.
	
	\vspace{10pt}
	\textbf{(b) O Grupo $(\mathbb{R}^*, \cdot)$:}
	
	\textbf{1. Estrutura de Variedade:}
	O conjunto $\mathbb{R}^* = \mathbb{R} \setminus \{0\}$ é um subconjunto aberto da variedade suave $\mathbb{R}$. Conforme a lógica aplicada na \textbf{Questão 8} e na \textbf{Questão 32}, qualquer subconjunto aberto de uma variedade suave herda a estrutura de variedade suave. Portanto, $\mathbb{R}^*$ é uma variedade suave de dimensão 1.
	
	\textbf{2. Suavidade da Operação de Grupo:}
	A operação de grupo é a multiplicação usual. Definimos $m : \mathbb{R}^* \times \mathbb{R}^* \to \mathbb{R}^*$ por $m(x, y) = xy$. Como a multiplicação de números reais é uma função polinomial, a aplicação $m$ é suave.
	
	\textbf{3. Suavidade da Inversão:}
	A aplicação de inversão $i : \mathbb{R}^* \to \mathbb{R}^*$ é dada por $i(x) = \frac{1}{x}$. 
	Como a função $g(x) = x$ é suave e nunca se anula no domínio $\mathbb{R}^*$, a função recíproca $i(x) = 1/g(x)$ é suave (quociente de funções suaves com denominador não nulo).
	
	Como a estrutura de variedade é suave e as operações de grupo são suaves, $(\mathbb{R}^*, \cdot)$ é um grupo de Lie.
}
%%%% QUESTÃO 37 %%%%
\questao{
	\begin{enumerate}[label=(\alph*)]
		\item Mostre que $\mathbb{C}^* = \mathbb{C} \setminus \{0\}$ é um grupo de Lie 2-dimensional com respeito à multiplicação de números complexos.
		\item Mostre que a 1-esfera $\mathbb{S}^1$, com a estrutura usual de variedade suave, é um grupo de Lie com respeito à multiplicação de números complexos.
	\end{enumerate}
}
\solucao{
	
	\textbf{Análise do Enunciado}
	\begin{itemize}
		\item \textbf{Tese (a):} Provar que $(\mathbb{C}^*, \cdot)$ é um grupo de Lie de dimensão 2.
		\item \textbf{Tese (b):} Provar que $(\mathbb{S}^1, \cdot)$ é um grupo de Lie.
	\end{itemize}
	\vspace{5pt}
	\textbf{(a) O Grupo $(\mathbb{C}^*, \cdot)$:}
	
	\textbf{1. Estrutura de Variedade:}
	Identificamos o plano complexo $\mathbb{C}$ com o espaço euclidiano $\mathbb{R}^2$ via a aplicação $z = x + iy \mapsto (x, y)$. Sob esta identificação, $\mathbb{C}^* = \mathbb{C} \setminus \{0\}$ corresponde ao conjunto $\mathbb{R}^2 \setminus \{(0,0)\}$. Como provamos na \textbf{Questão 8} e na \textbf{Questão 36(b)}, qualquer subconjunto aberto de uma variedade suave é, ele próprio, uma variedade suave. Portanto, $\mathbb{C}^*$ é uma variedade suave de dimensão 2.
	
	\textbf{2. Suavidade da Multiplicação:}
	Sejam $z = x + iy$ e $w = u + iv$ dois elementos de $\mathbb{C}^*$. A multiplicação complexa é dada por:
	\begin{equation*}
		m(z, w) = (x + iy)(u + iv) = (xu - yv) + i(xv + yu).
	\end{equation*}
	Em termos de coordenadas reais, a aplicação $m : \mathbb{R}^2 \times \mathbb{R}^2 \to \mathbb{R}^2$ é:
	\begin{equation*}
		m((x, y), (u, v)) = (xu - yv, xv + yu).
	\end{equation*}
	Como cada componente da aplicação $m$ é um polinômio nas variáveis reais, a multiplicação é suave.
	
	\textbf{3. Suavidade da Inversão:}
	A aplicação de inversão $i : \mathbb{C}^* \to \mathbb{C}^*$ é dada por $i(z) = z^{-1} = \frac{\bar{z}}{|z|^2}$. Em coordenadas reais:
	\begin{equation*}
		i(x, y) = \left( \frac{x}{x^2 + y^2}, \frac{-y}{x^2 + y^2} \right).
	\end{equation*}
	Como o denominador $x^2 + y^2$ é suave e nunca se anula em $\mathbb{C}^*$, a aplicação de inversão é um quociente de funções suaves, logo é suave. Assim, $(\mathbb{C}^*, \cdot)$ é um grupo de Lie.
	
	\vspace{10pt}
	\textbf{(b) O Grupo $(\mathbb{S}^1, \cdot)$:}
	
	\textbf{1. Estrutura de Variedade:}
	A 1-esfera $\mathbb{S}^1 = \{ z \in \mathbb{C} : |z| = 1 \}$ é uma subvariedade suave de $\mathbb{C}$ de dimensão 1, conforme a construção geral da \textbf{Questão 3}.
	
	\textbf{2. Verificação da Operação de Grupo:}
	Sejam $z, w \in \mathbb{S}^1$. Então $|z| = 1$ e $|w| = 1$. Pela propriedade do módulo de números complexos:
	\begin{equation*}
		|zw| = |z||w| = 1 \cdot 1 = 1 \quad \text{e} \quad \dots \quad |z^{-1}| = \frac{1}{|z|} = 1.
	\end{equation*}
	Isso prova que $\mathbb{S}^1$ é um subgrupo de $\mathbb{C}^*$.
	
	\textbf{3. Suavidade das Operações:}
	A aplicação de multiplicação $m_{\mathbb{S}^1}$ e a aplicação de inversão $i_{\mathbb{S}^1}$ são, respectivamente, as restrições das aplicações $m$ e $i$ de $\mathbb{C}^*$ ao subconjunto $\mathbb{S}^1$.
	
	Conforme discutido na \textbf{Questão 23}, a restrição de uma aplicação suave a uma subvariedade suave é automaticamente suave. Como a multiplicação e a inversão são suaves em $\mathbb{C}^*$, suas restrições a $\mathbb{S}^1$ também são suaves. Portanto, $\mathbb{S}^1$ é um grupo de Lie.
}
%%%% QUESTÃO 39 %%%%	
\questao{
	\begin{enumerate}[label=(\alph*)]
		\item Sejam $G_{1}, \dots, G_{k}$ grupos de Lie e considere o produto $$\Pi = G_{1} \times \dots \times G_{k} = \{ g = (g_{1}, \dots, g_{k}) : g_{i} \in G_{i} \text{ para todo } 1 \le i \le k\}.$$ Com a multiplicação $m : \Pi \times \Pi \to \Pi$, definida por $$m(g, h) := (g_{1}h_{1}, \dots, g_{k}h_{k})$$ para todos $g = (g_{1}, \dots, g_{k}), h = (h_{1}, \dots, h_{k}) \in \Pi$, mostre que $\Pi$ é um grupo de Lie.
		\item Mostre que $\mathbb{T}^{n} = \underbrace{\mathbb{S}^{1} \times \dots \times \mathbb{S}^{1}}_{n}$ é um grupo de Lie abeliano.
	\end{enumerate}
}
\solucao{
	
	\textbf{Análise do Enunciado}
	\begin{itemize}
		\item \textbf{Hipóteses:} $G_1, \dots, G_k$ são Grupos de Lie.
		\item \textbf{Tese (a):} Provar que o produto $\Pi = \prod G_i$ é um Grupo de Lie.
		\item \textbf{Tese (b):} Provar que o $n$-toro $\mathbb{T}^n$ é um grupo de Lie abeliano.
	\end{itemize}
	\vspace{5pt}
	\textbf{(a) Prova de que $\Pi$ é um Grupo de Lie:}
	
	\textbf{1. Estrutura de Variedade:}
	Como cada $G_i$ é um Grupo de Lie, cada $G_i$ é, por definição, uma variedade suave. Conforme provamos na \textbf{Questão 18}, o produto cartesiano de variedades suaves $\Pi = G_1 \times \dots \times G_k$ é também uma variedade suave, com dimensão $\dim(\Pi) = \sum_{i=1}^k \dim(G_i)$.
	
	\textbf{2. Suavidade da Multiplicação:}
	A aplicação de multiplicação $m : \Pi \times \Pi \to \Pi$ é definida componente a componente:
	\begin{equation*}
		m(g, h) = (m_1(g_1, h_1), m_2(g_2, h_2), \dots, m_k(g_k, h_k)),
	\end{equation*}
	onde cada $m_i$ é a operação de multiplicação suave do Grupo de Lie $G_i$. Como a aplicação produto de funções suaves é suave, concluímos que $m$ é uma aplicação suave.
	
	\textbf{3. Suavidade da Inversão:}
	A aplicação de inversão $\Inv : \Pi \to \Pi$ é dada por:
	\begin{equation*}
		\Inv(g) = (\Inv_1(g_1), \Inv_2(g_2), \dots, \Inv_k(g_k)),
	\end{equation*}
	onde $\Inv_i$ é a aplicação de inversão suave de $G_i$. Sendo a aplicação produto de funções suaves também suave, $\Inv$ é suave.
	
	Tendo a estrutura de variedade e operações suaves, $\Pi$ é um grupo de Lie.
	
	\vspace{10pt}
	\textbf{(b) O Grupo $\mathbb{T}^n$:}
	
	\textbf{1. Estrutura de Grupo de Lie:}
	Vimos na \textbf{Questão 37} que a 1-esfera $\mathbb{S}^1$ é um Grupo de Lie. O $n$-toro $\mathbb{T}^n$ é o produto de $n$ cópias de $\mathbb{S}^1$, e pelo resultado do item (a), o produto de Grupos de Lie é, necessariamente, um Grupo de Lie.
	
	\textbf{2. Comutatividade da Operação:}
	Dizemos que um grupo é \textbf{abeliano} se a sua operação é comutativa, ou seja, $gh = hg$ para quaisquer elementos do grupo.
	
	A operação em $\mathbb{S}^1$ é a multiplicação de pontos no plano complexo $\mathbb{C}$. Como a multiplicação de coordenadas complexas é comutativa, a operação em $\mathbb{S}^1$ é comutativa. Para quaisquer pontos $g = (g_1, \dots, g_n)$ e $h = (h_1, \dots, h_n)$ em $\mathbb{T}^n$, temos:
	\begin{align*}
		gh &= (g_1 h_1, \dots, g_n h_n) \\
		&= (h_1 g_1, \dots, h_n g_n) \\
		&= hg.
	\end{align*}
	Dessa forma, a comutatividade em cada componente $\mathbb{S}^1$ implica a comutatividade global em $\mathbb{T}^n$. Logo, $\mathbb{T}^n$ é um grupo de Lie abeliano.
}
	
	% LISTA 02
	\newpage
	\def\listid{L2} % Identificador único para a Lista 2
	\phantomsection
	\addcontentsline{toc}{section}{Soluções da Lista Exercício 2}
	\setcounter{numquestao}{0}
	\begin{center}
	\fbox{\textbf{Resolução da Lista 02}}
\end{center}
\begin{tcolorbox}[
	enhanced,
	sharp corners,
	boxrule=0.5pt,
	colback=white, 
	colframe=black,         
	width=1\textwidth,
	%fontupper=\itshape
	]%
	\textbf{Resumo:} Este documento apresenta a resolução detalhada da segunda lista de exercícios da disciplina de Variedades Diferenciáveis. O conteúdo foca na aplicação de conceitos de espaços tangentes, fibrados tangentes e a análise de campos vetoriais em variedades suaves. São explorados a representação coordenada de aplicações em esferas, a estrutura de variedades de espaços projetivos, a construção de atlas para o fibrado tangente $TM$ e o cálculo de colchetes de Lie para campos vetoriais em $\mathbb{R}^3$.
\end{tcolorbox}
%%%% QUESTÃO 01 %%%%
\questao{
	Considere a aplicação antípoda $A : \mathbb{S}^n \to \mathbb{S}^n$ definida por $A(p) = -p$, para cada $p \in \mathbb{S}^n$.
	\begin{enumerate}[label=(\alph*)]
		\item Calcule a representação coordenada $\hat{A}$ de $A$ em coordenadas estereográficas.
		\item Use a representação coordenada $\hat{A}$ de $A$ obtida para mostrar que $A$ é suave.
	\end{enumerate}
}
\solucao{
	
	\textbf{Análise do Enunciado}
	\begin{itemize}
		\item \textbf{Hipóteses:} $A$ é a aplicação antípoda em $\mathbb{S}^n$.
		\item \textbf{Tese (a):} Expressar $\hat{A} = \phi_S \circ A \circ \phi_N^{-1}$ analiticamente.
		\item \textbf{Tese (b):} Confirmar a suavidade de $A$ via representação coordenada.
	\end{itemize}
	
	\textbf{Fundamentação Teórica}
	
	Utilizamos o atlas de projeções estereográficas construído na \textbf{Questão 3 da Lista 01}. Especificamente, as cartas $\phi_N$ (polo norte) e $\phi_S$ (polo sul) e a expressão de $\phi_N^{-1} : \mathbb{R}^n \to \mathbb{S}^n$:
	\begin{equation*}
		\phi_N^{-1}(u) = \frac{\left(2u^1, \dots, 2u^n, |u|^2-1\right)}{|u|^2+1}.
	\end{equation*}
	Embora a suavidade de $A$ tenha sido estabelecida globalmente na \textbf{Questão 20 da Lista 01} (via restrição de aplicação linear no espaço ambiente), a tese aqui exige a verificação local através da representação coordenada $\hat{A}$.
	
	\textbf{Desenvolvimento Formal}
	
	\textbf{(a)} Para $u \in \mathbb{R}^n$, a representação coordenada $\hat{A} = \phi_S \circ A \circ \phi_N^{-1}$ é obtida pela seguinte composição:
	\begin{enumerate}
		\item $\phi_N^{-1}(u) = p = \left( \frac{2u^1}{|u|^2+1}, \dots, \frac{2u^n}{|u|^2+1}, \frac{|u|^2-1}{|u|^2+1} \right)$.
		\item $A(p) = -p = (y^1, \dots, y^{n+1}) = \left( \frac{-2u^1}{|u|^2+1}, \dots, \frac{-2u^n}{|u|^2+1}, \frac{1-|u|^2}{|u|^2+1} \right)$.
	\end{enumerate}
	Aplicando $\phi_S$ ao ponto $A(p)$, temos que $\hat{A}^i(u) = \frac{y^i}{1+y^{n+1}}$ para $i=1,\dots,n$. Calculando o denominador:
	\begin{align*}
		1+y^{n+1} &= 1 + \frac{1-|u|^2}{|u|^2+1} = \frac{|u|^2+1 + 1-|u|^2}{|u|^2+1} = \frac{2}{|u|^2+1}.
	\end{align*}
	Substituindo nas componentes:
	\begin{align*}
		\hat{A}^i(u) &= \frac{\dfrac{-2u^i}{|u|^2+1}}{\dfrac{2}{|u|^2+1}} = -u^i.
	\end{align*}
	Portanto, a representação coordenada é a aplicação linear:
	\begin{equation*}
		\boxed{\hat{A}(u) = -u}.
	\end{equation*}
	
	\textbf{(b)} Como $\hat{A} = -\text{Id}_{\mathbb{R}^n}$ é uma aplicação linear, ela é infinitamente diferenciável ($C^\infty$). Visto que as representações coordenadas de $A$ em qualquer par de cartas do atlas suave de $\mathbb{S}^n$ são composições de $\hat{A}$ com funções de transição suaves, a suavidade de $\hat{A}$ garante que $A$ é uma aplicação suave.
}
%%%% QUESTÃO 02 %%%%
\questao{
	Seja $F : \mathbb{R}^{n+1} \setminus \{0\} \to \mathbb{R}^{m+1} \setminus \{0\}$ uma função suave e homogênea de grau $d$ (isto é, existe $d \in \mathbb{Z}$ tal que $F(\lambda x) = \lambda^d F(x)$ para todo $\lambda \in \mathbb{R} \setminus \{0\}$ e qualquer $x \in \mathbb{R}^{n+1} \setminus \{0\}$). Prove que $f : \mathbb{RP}^n \to \mathbb{RP}^m$ definida por $f([x]) = [F(x)]$ está bem definida e é suave.
}
\solucao{
	
	\textbf{Análise do Enunciado}
	\begin{itemize}
		\item \textbf{Hipóteses:} $F$ é uma aplicação suave entre espaços euclidianos (menos a origem) e satisfaz a condição de homogeneidade $F(\lambda x) = \lambda^d F(x)$.
		\item \textbf{Tese 1:} a aplicação $f$ é bem definida, ou seja, a imagem de uma classe de equivalência $[x] \in \mathbb{RP}^n$ não depende do representante $x$ escolhido.
		\item \textbf{Tese 2:} $f$ é uma aplicação suave entre as variedades $\mathbb{RP}^n$ e $\mathbb{RP}^m$.
	\end{itemize}
	
	\textbf{Fundamentação Teórica}
	
	Recordamos a definição de $\mathbb{RP}^n$ como o espaço quociente $(\mathbb{R}^{n+1} \setminus \{0\}) / \sim$, onde $x \sim y \iff y = \lambda x$ para algum $\lambda \neq 0$. Conforme a \textbf{Questão 4 da Lista 01}, $\mathbb{RP}^n$ é uma variedade suave cujo atlas é composto pelas cartas $\phi_i : U_i \to \mathbb{R}^n$, onde $U_i = \{ [x] \in \mathbb{RP}^n : x^i \neq 0 \}$, dadas por:
	\begin{equation*}
		\phi_i([x]) = \left( \frac{x^1}{x^i}, \dots, \frac{x^{i-1}}{x^i}, \frac{x^{i+1}}{x^i}, \dots, \frac{x^{n+1}}{x^i} \right).
	\end{equation*}
	Uma aplicação $f : \mathbb{RP}^n \to \mathbb{RP}^m$ é suave se, para qualquer par de cartas $\phi_i$ em $\mathbb{RP}^n$ e $\psi_j$ em $\mathbb{RP}^m$, a representação coordenada $\hat{f} = \psi_j \circ f \circ \phi_i^{-1}$ for suave.
	
	\textbf{Desenvolvimento Formal}
	
	\textbf{1. Prova de que $f$ está bem definida:}
	Sejam $x, y \in \mathbb{R}^{n+1} \setminus \{0\}$ representantes da mesma classe de equivalência em $\mathbb{RP}^n$. Então, existe $\lambda \in \mathbb{R} \setminus \{0\}$ tal que $y = \lambda x$. Aplicando a função $F$:
	\begin{align*}
		f([y]) &= [F(y)] \\
		&= [F(\lambda x)] \\
		&= [\lambda^d F(x)] \quad (\text{por homogeneidade de } F).
	\end{align*}
	Como $F(x) \in \mathbb{R}^{m+1} \setminus \{0\}$ e $\lambda \neq 0$, segue que $\lambda^d \neq 0$. Assim, os vetores $F(x)$ e $\lambda^d F(x)$ são linearmente dependentes, o que implica que eles pertencem à mesma classe de equivalência no espaço projetivo $\mathbb{RP}^m$:
	\begin{equation*}
		[\lambda^d F(x)] = [F(x)] \implies f([y]) = f([x]).
	\end{equation*}
	Portanto, a imagem depende apenas da classe $[x]$, e $f$ está bem definida.
	
	\textbf{2. Prova de que $f$ é suave:}
	Seja $[x] \in U_i \subset \mathbb{RP}^n$. Para analisar a suavidade localmente, utilizamos a seção local $\sigma_i : \phi_i(U_i) \to \mathbb{R}^{n+1} \setminus \{0\}$ definida por:
	\begin{equation*}
		\sigma_i(u) = (u^1, \dots, \underbrace{1}_{i\text{-ésima}}, \dots, u^n).
	\end{equation*}
	Note que $\pi \circ \sigma_i = \phi_i^{-1}$, onde $\pi$ é a projeção canônica. A aplicação $f$ pode ser escrita localmente como $f \circ \pi = [F]$. Para qualquer ponto $p \in \mathbb{RP}^n$, existe $j \in \{1, \dots, m+1\}$ tal que a $j$-ésima componente de $F(\sigma_i(u))$, denotada por $F_j(\sigma_i(u))$, é não nula em uma vizinhança de $u$. 
	
	A representação coordenada $\hat{f} = \psi_j \circ f \circ \phi_i^{-1}$ é então:
	\begin{align*}
		\hat{f}(u) &= \psi_j( f(\phi_i^{-1}(u)) ) \\
		&= \psi_j( [F(\sigma_i(u))] ) \\
		&= \left( \frac{F_1(\sigma_i(u))}{F_j(\sigma_i(u))}, \dots, \widehat{\frac{F_j(\sigma_i(u))}{F_j(\sigma_i(u))}}, \dots, \frac{F_{m+1}(\sigma_i(u))}{F_j(\sigma_i(u))} \right).
	\end{align*}
	Como $F$ é uma função suave em $\mathbb{R}^{n+1} \setminus \{0\}$, cada componente $F_k$ é uma função suave. A aplicação $\sigma_i$ também é suave (sendo uma inclusão linear). Portanto, cada componente de $\hat{f}$ é o quociente de duas funções suaves, onde o denominador $F_j(\sigma_i(u))$ é não nulo no domínio considerado.
	
	Pela regra do quociente para funções $C^\infty$, a representação $\hat{f}$ é suave. Como isso vale para qualquer par de cartas, concluímos que $f$ é uma aplicação suave entre variedades.
}
%%%% QUESTÃO 03 %%%%
\questao{
	Em $\mathbb{S}^1$ considere o atlas $\mathcal{U} = \{(U_i^{\pm}, \varphi_i^{\pm}), i=1,2\}$, onde $U_i^{\pm} = \{x \in \mathbb{S}^1 : \pm x_i > 0\}$ e $\varphi_i^{\pm} : U_i^{\pm} \to (-1,1)$ é definida por $\varphi_1^{\pm}(x_1,x_2) = x_2$ e $\varphi_2^{\pm}(x_1,x_2) = x_1$. Encontre o atlas de $T\mathbb{S}^1$ associado ao atlas $\mathcal{U}$ calculando explicitamente as funções de transições.
}
\solucao{
	
	\textbf{Análise do Enunciado}
	\begin{itemize}
		\item \textbf{Hipóteses:} Temos a 1-esfera $\mathbb{S}^1$ com um atlas específico de quatro cartas ($\varphi_1^+, \varphi_1^-, \varphi_2^+, \varphi_2^-$).
		\item \textbf{Tese:} Construir o atlas induzido para o fibrado tangente $T\mathbb{S}^1$ e calcular as funções de transição $\Phi_\beta \circ \Phi_\alpha^{-1}$.
	\end{itemize}
	
	\textbf{Fundamentação Teórica}
	
	De acordo com a teoria de fibrados tangentes (Lee, Capítulo 3), se $\mathcal{U} = \{(U_\alpha, \varphi_\alpha)\}$ é um atlas suave para uma variedade $M$ de dimensão $n$, o atlas induzido $\mathcal{T} = \{(T U_\alpha, \Phi_\alpha)\}$ para $TM$ é definido por:
	\begin{equation*}
		\Phi_\alpha : T U_\alpha \to \mathbb{R}^{2n}, \quad \Phi_\alpha(p, v) = (\varphi_\alpha(p), d(\varphi_\alpha)_p(v)).
	\end{equation*}
	Sejam duas cartas $(U_\alpha, \varphi_\alpha)$ e $(U_\beta, \varphi_\beta)$ com interseção não vazia. A função de transição na base é $\tau = \varphi_\beta \circ \varphi_\alpha^{-1}$. A função de transição no fibrado tangente $\mathbb{R}^{2n} \to \mathbb{R}^{2n}$ é dada por:
	\begin{equation*}
		\Phi_\beta \circ \Phi_\alpha^{-1}(u, w) = (\tau(u), d\tau_u(w)),
	\end{equation*}
	onde $u \in \mathbb{R}^n$ representa a coordenada do ponto e $w \in \mathbb{R}^n$ representa as componentes do vetor tangente. Para $n=1$, $d\tau_u(w) = \tau'(u) \cdot w$.
	
	\textbf{Desenvolvimento Formal}
	
	Consideramos a base $\mathbb{S}^1 \subset \mathbb{R}^2$ definida por $x_1^2 + x_2^2 = 1$. As cartas fornecidas são:
	\begin{itemize}
		\item $\varphi_1^{\pm}(x_1, x_2) = x_2$ (coordenada $u$)
		\item $\varphi_2^{\pm}(x_1, x_2) = x_1$ (coordenada $v$)
	\end{itemize}
	
	Analisaremos as transições entre as cartas $\varphi_1$ e $\varphi_2$. Como as cartas $\varphi_1^+$ e $\varphi_1^-$ (assim como $\varphi_2^+$ e $\varphi_2^-$) possuem domínios disjuntos, as transições não triviais ocorrem entre índices diferentes.
	
	\textbf{Caso 1: Transição entre $\varphi_1^+$ e $\varphi_2^+$}
	A interseção é $U_1^+ \cap U_2^+ = \{ (x_1, x_2) \in \mathbb{S}^1 : x_1 > 0, x_2 > 0 \}$.
	Se $u = \varphi_1^+(x) = x_2$, então como $x_1 > 0$, temos $x_1 = \sqrt{1 - x_2^2}$. Logo:
	\begin{equation*}
		\tau(u) = (\varphi_2^+ \circ (\varphi_1^+)^{-1})(u) = \sqrt{1 - u^2}.
	\end{equation*}
	A derivada é $\tau'(u) = \frac{-u}{\sqrt{1-u^2}}$. Assim, a transição em $T\mathbb{S}^1$ é:
	\begin{equation*}
		\Phi_2^+ \circ (\Phi_1^+)^{-1}(u, w) = \left( \sqrt{1-u^2}, \frac{-u}{\sqrt{1-u^2}}w \right).
	\end{equation*}
	
	\textbf{Caso 2: Transição entre $\varphi_1^+$ e $\varphi_2^-$}
	A interseção é $U_1^+ \cap U_2^- = \{ (x_1, x_2) \in \mathbb{S}^1 : x_1 > 0, x_2 < 0 \}$.
	Novamente, $u = x_2$ e $x_1 = \sqrt{1-x_2^2}$. a transição na base é a mesma:
	\begin{equation*}
		\tau(u) = \varphi_2^- \circ (\varphi_1^+)^{-1}(u) = \sqrt{1 - u^2}.
	\end{equation*}
	A transição no fibrado tangente permanece:
	\begin{equation*}
		\Phi_2^- \circ (\Phi_1^+)^{-1}(u, w) = \left( \sqrt{1-u^2}, \frac{-u}{\sqrt{1-u^2}}w \right).
	\end{equation*}
	
	\textbf{Caso 3: Transição entre $\varphi_1^-$ e $\varphi_2^+$}
	A interseção é $U_1^- \cap U_2^+ = \{ (x_1, x_2) \in \mathbb{S}^1 : x_1 < 0, x_2 > 0 \}$.
	Aqui, $u = x_2$ e, como $x_1 < 0$, temos $x_1 = -\sqrt{1 - x_2^2}$. Logo:
	\begin{equation*}
		\tau(u) = \varphi_2^+ \circ (\varphi_1^-)^{-1}(u) = -\sqrt{1 - u^2}.
	\end{equation*}
	A derivada é $\tau'(u) = \frac{u}{\sqrt{1-u^2}}$. A transição em $T\mathbb{S}^1$ é:
	\begin{equation*}
		\Phi_2^+ \circ (\Phi_1^-)^{-1}(u, w) = \left( -\sqrt{1-u^2}, \frac{u}{\sqrt{1-u^2}}w \right).
	\end{equation*}
	
	\textbf{Caso 4: Transição entre $\varphi_1^-$ e $\varphi_2^-$}
	A interseção é $U_1^- \cap U_2^- = \{ (x_1, x_2) \in \mathbb{S}^1 : x_1 < 0, x_2 < 0 \}$.
	Analogamente ao Caso 3, com $x_1 < 0$:
	\begin{equation*}
		\tau(u) = \varphi_2^- \circ (\varphi_1^-)^{-1}(u) = -\sqrt{1 - u^2} \implies \tau'(u) = \frac{u}{\sqrt{1-u^2}}.
	\end{equation*}
	A transição em $T\mathbb{S}^1$ é:
	\begin{equation*}
		\Phi_2^- \circ (\Phi_1^-)^{-1}(u, w) = \left( -\sqrt{1-u^2}, \frac{u}{\sqrt{1-u^2}}w \right).
	\end{equation*}
	
	Concluímos que o atlas de $T\mathbb{S}^1$ é composto pelas cartas $\Phi_i^{\pm}$ e as funções de transição dependem exclusivamente do sinal de $x_1$ na interseção, assumindo a forma $(\tau(u), \tau'(u)w)$.
}
%%%% QUESTÃO 04 %%%%
\questao{
	Sejam $M^n$ uma $n$-variedade suave, $B \subset M^n$ um subconjunto fechado e $\delta : M^n \to \mathbb{R}$ uma função contínua positiva.
	\begin{enumerate}[label=(\alph*)]
		\item Usando uma partição da unidade, mostre que existe uma função suave $\tilde{\delta} : M^n \to \mathbb{R}$ tal que $0 < \tilde{\delta}(p) < \delta(p)$ para todo $p \in M^n$.
		\item Mostre que existe uma função $\psi : M^n \to \mathbb{R}$ que é suave e positiva em $M^n \setminus B$, zero em $B$ e satisfaz $\psi(x) < \delta(x)$ em toda parte (Sugestão: considere $1/(1 + f)$, onde $f : M^n \setminus B \to \mathbb{R}$ é uma função exaustão positiva).
	\end{enumerate}
}
\solucao{
	
	\textbf{Análise do Enunciado}
	\begin{itemize}
		\item \textbf{Hipóteses:} $M^n$ é uma variedade suave; $B$ é um subconjunto fechado; $\delta$ é uma função contínua tal que $\delta(p) > 0$ para todo $p \in M^n$.
		\item \textbf{Tese (a):} Provar a existência de uma função suave $\tilde{\delta}$ estritamente positiva e estritamente menor que $\delta$.
		\item \textbf{Tese (b):} Provar a existência de uma função suave $\psi$ que seja nula precisamente em $B$ e satisfaça a desigualdade $\psi < \delta$.
	\end{itemize}
	
	\textbf{Fundamentação Teórica}
	Para a parte (a), utilizaremos o **Teorema da Partição da Unidade** (Lista 01), que permite a construção de funções globais suaves a partir de dados locais. 
	
	Para a parte (b), utilizaremos a definição de **Função de Exaustão** (Lista 01): uma função $f : X \to \mathbb{R}$ é de exaustão se, para todo $c \in \mathbb{R}$, o conjunto de subnível $X_c = \{x \in X : f(x) \le c\}$ é compacto. Note que isso implica que, para qualquer compacto $K \subset X$, existe $c$ tal que $K \subset X_c$.
	
	\textbf{Desenvolvimento Formal}
	
	\textbf{(a)} Dado que $\delta$ é contínua e $\delta(p) > 0$ para todo $p \in M^n$, para cada $p \in M^n$ existe uma vizinhança aberta $U_p$ e uma constante $\epsilon_p > 0$ tal que $\delta(q) > \epsilon_p$ para todo $q \in U_p$. 
	
	A coleção $\{U_p\}_{p \in M^n}$ constitui uma cobertura aberta de $M^n$. Pela paracompactidade de $M^n$, existe uma partição da unidade $\{\rho_i\}_{i \in I}$ subordinada a um refinamento $\{V_i\}_{i \in I}$, tal que cada $V_i \subset U_{p(i)}$ para algum $p(i)$. Definimos a função $\tilde{\delta} : M^n \to \mathbb{R}$ por:
	\begin{equation*}
		\tilde{\delta}(q) = \sum_{i \in I} \rho_i(q) \epsilon_{p(i)}.
	\end{equation*}
	Vemos que $\tilde{\delta}$ é suave, pois é uma soma localmente finita de funções suaves. Como $\rho_i(q) \ge 0$ e $\sum \rho_i = 1$ com $\epsilon_{p(i)} > 0$, segue que $\tilde{\delta}(q) > 0$ para todo $q \in M^n$. 
	
	Para qualquer $q \in M^n$, se $\rho_i(q) \neq 0$, então $q \in V_i \subset U_{p(i)}$, o que implica $\epsilon_{p(i)} < \delta(q)$. Assim:
	\begin{equation*}
		\tilde{\delta}(q) = \sum_{i \in I} \rho_i(q) \epsilon_{p(i)} < \sum_{i \in I} \rho_i(q) \delta(q) = \delta(q) \sum_{i \in I} \rho_i(q) = \delta(q).
	\end{equation*}
	Portanto, $0 < \tilde{\delta}(p) < \delta(p)$ para todo $p \in M^n$.
	
	\textbf{(b)} Seja $M^n \setminus B$ a variedade suave aberta resultante. Pelo resultado da Lista 01, existe uma função de exaustão positiva e suave $f : M^n \setminus B \to \mathbb{R}$. Definimos então a função $h : M^n \setminus B \to \mathbb{R}$ por:
	\begin{equation*}
		h(x) = \frac{1}{1 + f(x)}.
	\end{equation*}
	Como $f$ é suave e positiva, $h$ é suave e positiva em $M^n \setminus B$. Pela definição de função de exaustão, para qualquer compacto $K \subset M^n \setminus B$, o conjunto $\{x \in M^n \setminus B : f(x) \le c\}$ é compacto. Isso implica que $f$ é coerciva em $M^n \setminus B$, logo, $h(x) \to 0$ quando $x$ deixa qualquer compacto de $M^n \setminus B$ (o que ocorre quando $x$ tende à fronteira $\partial(M \setminus B) \subset B$ ou quando $x$ deixa de estar em qualquer compacto de $M^n$).
	
	Para estender a função a todo $M^n$, utilizamos a função $\tilde{\delta}$ construída no item (a) e a existência de uma função suave $\rho : M^n \to [0, 1]$ tal que $\rho^{-1}(0) = B$. 
	
	Definimos a função $\psi : M^n \to \mathbb{R}$ por:
	\begin{equation*}
		\psi(x) = \rho(x) \tilde{\delta}(x).
	\end{equation*}
	Verificamos as propriedades de $\psi$:
	\begin{enumerate}
		\item \textbf{Suavidade:} $\psi$ é o produto de duas funções suaves, logo é suave em $M^n$.
		\item \textbf{Aniquilação em $B$:} Para todo $x \in B$, $\rho(x) = 0$, implicando $\psi(x) = 0$.
		\item \textbf{Positividade em $M^n \setminus B$:} Para todo $x \in M^n \setminus B$, $\rho(x) > 0$ e $\tilde{\delta}(x) > 0$, logo $\psi(x) > 0$.
		\item \textbf{Limitante superior:} Dado que $0 \le \rho(x) \le 1$ e $0 < \tilde{\delta}(x) < \delta(x)$, temos:
		\begin{equation*}
			\psi(x) = \rho(x) \tilde{\delta}(x) \le 1 \cdot \tilde{\delta}(x) < \delta(x).
		\end{equation*}
	\end{enumerate}
	Assim, a função $\psi$ satisfaz todas as condições requeridas.
}
%%%% QUESTÃO 05 %%%%
\questao{
	Sejam $M^n$ uma $n$-variedade suave, $p \in M^n$ e $X \in T_p M$. Mostre que se $f,g \in C^\infty(M)$ são funções que coincidem em alguma vizinhança de $p$, então $X f = Xg$. Conclua que a noção de espaço tangente em uma variedade suave é realmente uma construção local.
}
\solucao{
	
	\textbf{Análise do Enunciado}
	\begin{itemize}
		\item \textbf{Hipóteses:} $M^n$ é uma variedade suave, $p \in M^n$ e $X$ é um vetor tangente em $p$. As funções $f, g \in C^\infty(M)$ coincidem em um aberto $U$ tal que $p \in U$.
		\item \textbf{Tese 1:} Provar que $X f = X g$.
		\item \textbf{Tese 2:} Concluir que a definição de $T_p M$ é local.
	\end{itemize}
	
	\textbf{Fundamentação Teórica}
	Conforme a definição de Lee (Capítulo 3), um vetor tangente $X \in T_p M$ é uma \textbf{derivação} em $p$: uma aplicação linear $X : C^\infty(M) \to \mathbb{R}$ que satisfaz a regra de Leibniz:
	\begin{equation*}
		X(fg) = f(p)Xg + g(p)Xf, \quad \forall f, g \in C^\infty(M).
	\end{equation*}
	Para a construção da prova, utilizaremos a técnica de funções de corte (bump functions). A existência de tais funções é garantida pelo **Teorema da Partição da Unidade** e detalhada na prova do **Lema de Extensão (Questão 27 da Lista 01)**, que assegura a possibilidade de construir funções suaves com suporte compacto contidos em abertos arbitrários.
	
	\textbf{Desenvolvimento Formal}
	
	Seja $h = f - g$. Pela linearidade do operador $X$, temos que $X f = X g$ se, e somente se, $X h = 0$. Por hipótese, $f$ e $g$ coincidem em uma vizinhança aberta $U$ de $p$, logo $h(q) = 0$ para todo $q \in U$.
	
	Com base no resultado da **Questão 27 da Lista 01**, existe uma função suave $\rho \in C^\infty(M)$ (uma bump function) tal que:
	\begin{enumerate}
		\item $0 \le \rho(q) \le 1$ para todo $q \in M^n$.
		\item $\rho(p) = 1$.
		\item $\text{supp}(\rho) \subset U$.
	\end{enumerate}
	
	Consideremos a decomposição da função $h$:
	\begin{equation*}
		h = \rho h + (1 - \rho) h.
	\end{equation*}
	Aplicando a derivação $X$:
	\begin{equation*}
		X h = X(\rho h) + X((1 - \rho) h).
	\end{equation*}
	Analisamos os termos:
	\begin{enumerate}
		\item \textbf{Termo $\rho h$:} Como $\text{supp}(\rho) \subset U$ e $h|_U \equiv 0$, o produto $\rho h$ é a função identicamente nula em todo $M^n$. Portanto, $X(\rho h) = 0$.
		\item \textbf{Termo $(1 - \rho) h$:} Pela regra de Leibniz:
		\begin{equation*}
			X((1 - \rho) h) = (1 - \rho)(p) X h + h(p) X(1 - \rho).
		\end{equation*}
		Visto que $\rho(p) = 1$ e $h(p) = 0$, temos $(1 - \rho)(p) = 0$ e $h(p) = 0$. Logo, $X((1 - \rho) h) = 0$.
	\end{enumerate}
	Consequentemente, $X h = 0$, o que implica $X f = X g$.
	
	\textbf{Conclusão sobre a Localidade:}
	Este resultado demonstra que a ação de um vetor tangente $X \in T_p M$ sobre uma função $f$ depende exclusivamente dos valores de $f$ em uma vizinhança arbitrária de $p$. Alterações em $f$ fora de $U$ não afetam o valor de $X f$, provando que o espaço tangente é, fundamentalmente, uma \textbf{construção local}.
}
%%%% QUESTÃO 06 %%%%
\questao{
	Sejam $M^n$ uma $n$-variedade suave, $U \subset M^n$ uma subvariedade aberta e $\iota : U \hookrightarrow M^n$ a aplicação inclusão. Para cada $p \in U$, mostre que $\iota_* : T_p U \to T_p M$ é um isomorfismo. Conclua que qualquer vetor tangente $X \in T_p M$ pode ser aplicado as funções definidas apenas em uma vizinhança de $p \in M^n$.
}
\solucao{
	
	\textbf{Análise do Enunciado}
	\begin{itemize}
		\item \textbf{Hipóteses:} $M^n$ é uma variedade suave, $U$ é um aberto de $M^n$ (subvariedade aberta) e $\iota$ é a inclusão canônica.
		\item \textbf{Tese 1:} Provar que a aplicação diferencial (push-forward) $\iota_*$ é um isomorfismo de espaços vetoriais entre $T_p U$ e $T_p M$.
		\item \textbf{Tese 2:} Concluir a aplicabilidade de $X \in T_p M$ a funções locais.
	\end{itemize}
	
	\textbf{Fundamentação Teórica}
	Conforme a definição de Lee (Capítulo 3), dado um mapa suave $F : M \to N$, a aplicação diferencial $F_* : T_p M \to T_{F(p)} N$ é definida para qualquer $v \in T_p M$ e $f \in C^\infty(N)$ por:
	\begin{equation*}
		(F_* v)(f) = v(f \circ F).
	\end{equation*}
	No caso da inclusão $\iota : U \hookrightarrow M^n$, temos $\iota(p) = p$ e $f \circ \iota = f|_U$ (a restrição da função $f$ ao aberto $U$). Assim, para $v \in T_p U$:
	\begin{equation*}
		(\iota_* v)(f) = v(f|_U).
	\end{equation*}
	Para provar que $\iota_*$ é um isomorfismo, devemos mostrar que ela é injetiva e sobrejetiva.
	
	\textbf{Desenvolvimento Formal}
	
	\textbf{1. Injetividade de $\iota_*$:}
	Suponha que $\iota_* v = 0$ para algum $v \in T_p U$. Isso implica que $(\iota_* v)(f) = 0$ para toda função $f \in C^\infty(M)$, ou seja:
	\begin{equation*}
		v(f|_U) = 0, \quad \forall f \in C^\infty(M).
	\end{equation*}
	Agora, seja $g \in C^\infty(U)$ qualquer função suave definida em $U$. Pelo \textbf{Lema de Extensão (Questão 27 da Lista 01)}, existe uma função $\tilde{g} \in C^\infty(M)$ tal que $\tilde{g}|_U = g$. Como $\iota_* v = 0$, temos:
	\begin{equation*}
		v(g) = v(\tilde{g}|_U) = (\iota_* v)(\tilde{g}) = 0.
	\end{equation*}
	Como $v(g) = 0$ para toda função $g \in C^\infty(U)$, concluímos que $v = 0$. Logo, $\iota_*$ é injetiva.
	
	\textbf{2. Sobrejetividade de $\iota_*$:}
	Seja $X \in T_p M$ um vetor tangente em $p$. Queremos encontrar um $v \in T_p U$ tal que $\iota_* v = X$. Definimos a aplicação $v : C^\infty(U) \to \mathbb{R}$ por:
	\begin{equation*}
		v(g) = X(\tilde{g}),
	\end{equation*}
	onde $\tilde{g} \in C^\infty(M)$ é qualquer extensão suave de $g$ para $M$. 
	
	Pelo resultado da \textbf{Questão 05 desta Lista 02}, sabemos que se duas funções $\tilde{g}_1, \tilde{g}_2 \in C^\infty(M)$ coincidem em uma vizinhança de $p$ (neste caso, ambas coincidem com $g$ em $U$), então $X(\tilde{g}_1) = X(\tilde{g}_2)$. Portanto, a aplicação $v$ está bem definida e independe da escolha da extensão. É trivial verificar que $v$ é linear e satisfaz a regra de Leibniz, logo $v \in T_p U$. Por construção:
	\begin{equation*}
		(\iota_* v)(f) = v(f|_U) = X(\tilde{f}) = X(f), \quad \forall f \in C^\infty(M),
	\end{equation*}
	onde $\tilde{f}$ é a extensão de $f|_U$ (que é a própria $f$). Assim, $\iota_* v = X$, provando que $\iota_*$ é sobrejetiva.
	
	Como $\iota_*$ é linear, injetiva e sobrejetiva, ela é um isomorfismo de espaços vetoriais.
	
	\textbf{Conclusão Final:}
	O fato de $\iota_*$ ser um isomorfismo implica que existe uma correspondência biunívoca entre as derivações em $p$ considerando apenas funções locais ($T_p U$) e derivações considerando funções globais ($T_p M$). Como provamos na \textbf{Questão 05}, a ação de $X \in T_p M$ depende apenas dos valores de $f$ em uma vizinhança de $p$. Portanto, qualquer vetor tangente $X \in T_p M$ pode ser aplicado a qualquer função $g \in C^\infty(U)$ via a identificação $X(g) = X(\tilde{g})$, validando que o espaço tangente é uma construção local.
}
%%%% QUESTÃO 07 %%%%
\questao{
	Seja $F : M^n \to N^k$ um difeomorfismo local entre variedades suaves. Para cada $p \in M^n$, mostre que $F_* : T_p M \to T_{F(p)} N$ é um isomorfismo.
}
\solucao{
	
	\textbf{Análise do Enunciado}
	\begin{itemize}
		\item \textbf{Hipóteses:} $F$ é um difeomorfismo local. Isso significa que, para cada $p \in M^n$, existe uma vizinhança aberta $U \subset M^n$ de $p$ tal que a restrição $F|_U : U \to F(U)$ é um difeomorfismo sobre a sua imagem $V = F(U) \subset N^k$.
		\item \textbf{Tese:} Provar que a aplicação diferencial (push-forward) $F_* : T_p M \to T_{F(p)} N$ é um isomorfismo de espaços vetoriais.
	\end{itemize}
	
	\textbf{Fundamentação Teórica}
	Utilizaremos a definição de aplicação diferencial de John M. Lee (Capítulo 3). Se $F : M \to N$ é uma aplicação suave, seu diferencial $F_* : T_p M \to T_{F(p)} N$ é a aplicação linear definida por:
	\begin{equation*}
		(F_* v)(f) = v(f \circ F), \quad \forall v \in T_p M, \forall f \in C^\infty(N).
	\end{equation*}
	A prova baseia-se em dois resultados fundamentais:
	\begin{enumerate}
		\item \textbf{Regra da Cadeia:} Se $F : M \to N$ e $G : N \to P$ são aplicações suaves, então $(G \circ F)_* = G_* \circ F_*$.
		\item \textbf{Diferencial da Identidade:} Se $\text{Id}_M : M \to M$ é a aplicação identidade, então $(\text{Id}_M)_* : T_p M \to T_p M$ é a aplicação identidade $\text{Id}_{T_p M}$.
	\end{enumerate}
	
	\textbf{Desenvolvimento Formal}
	
	Seja $p \in M^n$. Por hipótese, $F$ é um difeomorfismo local, logo existe um aberto $U \subset M^n$ contendo $p$ tal que a aplicação $\Phi = F|_U : U \to V$ é um difeomorfismo, onde $V = F(U)$ é um aberto em $N^k$.
	
	Como $\Phi$ é um difeomorfismo, existe uma aplicação inversa suave $\Phi^{-1} : V \to U$. Consideremos a aplicação diferencial desta inversa, $(\Phi^{-1})_* : T_{F(p)} N \to T_p U$. 
	
	Pela **Questão 06 desta Lista 02**, sabemos que a inclusão $\iota : U \hookrightarrow M^n$ induz um isomorfismo $\iota_* : T_p U \to T_p M$. Portanto, podemos identificar $T_p U$ com $T_p M$ e tratar $\Phi_*$ como a própria aplicação $F_*$ no ponto $p$.
	
	Analisamos agora as composições das aplicações diferenciais:
	\begin{enumerate}
		\item A composição $\Phi^{-1} \circ \Phi$ é a identidade $\text{Id}_U$ em $U$. Pela regra da cadeia:
		\begin{equation*}
			(\Phi^{-1} \circ \Phi)_* = (\Phi^{-1})_* \circ \Phi_* = (\text{Id}_U)_* = \text{Id}_{T_p M}.
		\end{equation*}
		\item A composição $\Phi \circ \Phi^{-1}$ é a identidade $\text{Id}_V$ em $V$. Pela regra da cadeia:
		\begin{equation*}
			(\Phi \circ \Phi^{-1})_* = \Phi_* \circ (\Phi^{-1})_* = (\text{Id}_V)_* = \text{Id}_{T_{F(p)} N}.
		\end{equation*}
	\end{enumerate}
	
	A existência de uma aplicação linear $(\Phi^{-1})_*$ tal que a composição com $\Phi_*$ resulta na identidade em ambos os lados prova que $\Phi_*$ é um isomorfismo de espaços vetoriais. Como $\Phi_*$ coincide com $F_*$ no ponto $p$, concluímos que $F_* : T_p M \to T_{F(p)} N$ é um isomorfismo.
	
	Um corolário imediato desta prova é que, para $F_*$ ser um isomorfismo, as dimensões dos espaços tangentes devem ser iguais, logo $n = k$.
}
%%%% QUESTÃO 08 %%%%
\questao{
	Para cada espaço vetorial $V$ (de dimensão finita) e cada $a \in V$, mostre que existe um isomorfismo natural $\Psi_V : V \to T_a V$ tal que para qualquer aplicação linear $L : V \to W$ o seguinte diagrama comuta:
	\begin{center}
		\begin{tikzpicture}[node distance=2.5cm, auto]
			\node (V) {$V$};
			\node (TaV) [right of=V] {$T_a V$};
			\node (W) [below of=V] {$W$};
			\node (TaW) [right of=W] {$T_{L(a)} W$};
			
			\draw[->] (V) -- node {$\Psi_V$} (TaV);
			\draw[->] (V) -- node [left] {$L$} (W);
			\draw[->] (TaV) -- node [right] {$L_*$} (TaW);
			\draw[->] (W) -- node {$\Psi_W$} (TaW);
		\end{tikzpicture}
	\end{center}
}
\solucao{
	
	\textbf{Análise do Enunciado}
	\begin{itemize}
		\item \textbf{Hipóteses:} $V$ e $W$ são espaços vetoriais de dimensão finita; $a \in V$; $L : V \to W$ é uma aplicação linear.
		\item \textbf{Tese 1:} Construir um isomorfismo $\Psi_V : V \to T_a V$.
		\item \textbf{Tese 2:} Provar que $\Psi_W \circ L = L_* \circ \Psi_V$, ou seja, que o diagrama comutativo é válido.
	\end{itemize}
	
	\textbf{Fundamentação Teórica}
	Um espaço vetorial $V$ de dimensão $n$ é uma variedade suave com um atlas global composto por uma única carta $\varphi : V \to \mathbb{R}^n$ (qualquer isomorfismo de $V$ para $\mathbb{R}^n$ serve). 
	
	De acordo com a definição de Lee (Capítulo 3), o espaço tangente $T_a V$ consiste em todas as derivações em $a$. Para qualquer vetor $v \in V$, podemos definir uma derivação $\Psi_V(v) \in T_a V$ através da derivada direcional:
	\begin{equation*}
		(\Psi_V(v)) f = df_a(v) = \lim_{t \to 0} \frac{f(a + tv) - f(a)}{t}, \quad \forall f \in C^\infty(V).
	\end{equation*}
	A aplicação diferencial $L_* : T_a V \to T_{L(a)} W$ (push-forward) é definida por $(L_* X) h = X(h \circ L)$ para todo $X \in T_a V$ e $h \in C^\infty(W)$.
	
	\textbf{Desenvolvimento Formal}
	
	\textbf{1. Prova de que $\Psi_V$ é um isomorfismo:}
	A aplicação $\Psi_V : V \to T_a V$ é linear, pois o limite da soma é a soma dos limites e o limite do produto por escalar é o produto pelo escalar. 
	
	Para provar que $\Psi_V$ é um isomorfismo, observamos que $\dim(V) = n$ e, como $V$ é uma variedade suave de dimensão $n$, temos $\dim(T_a V) = n$. Portanto, basta provar que $\Psi_V$ é injetiva. 
	
	Suponha $\Psi_V(v) = 0$. Então, para toda função $f \in C^\infty(V)$, temos $df_a(v) = 0$. Escolhemos a função linear $f(x) = \ell(x)$ para algum $\ell \in V^*$ (o espaço dual de $V$). Para tal função, temos:
	\begin{equation*}
		df_a(v) = \lim_{t \to 0} \frac{\ell(a+tv) - \ell(a)}{t} = \lim_{t \to 0} \frac{\ell(a) + t\ell(v) - \ell(a)}{t} = \ell(v).
	\end{equation*}
	Como $\ell(v) = 0$ para toda $\ell \in V^*$, segue necessariamente que $v = 0$. Logo, $\Psi_V$ é injetiva e, consequentemente, um isomorfismo.
	
	\textbf{2. Prova da Comutatividade do Diagrama:}
	Desejamos mostrar que $\Psi_W \circ L = L_* \circ \Psi_V$. Seja $v \in V$. Aplicamos ambas as composições a uma função arbitrária $h \in C^\infty(W)$.
	
	Primeiro, calculamos a ação de $\Psi_W(L(v))$ sobre $h$:
	\begin{equation*}
		(\Psi_W(L(v))) h = dh_{L(a)}(L(v)) = \lim_{t \to 0} \frac{h(L(a) + tL(v)) - h(L(a))}{t}.
	\end{equation*}
	Como $L$ é linear, temos $L(a) + tL(v) = L(a + tv)$. Assim:
	\begin{equation*}
		(\Psi_W(L(v))) h = \lim_{t \to 0} \frac{h(L(a + tv)) - h(L(a))}{t}.
	\end{equation*}
	
	Agora, calculamos a ação de $L_* (\Psi_V(v))$ sobre $h$:
	\begin{align*}
		(L_* (\Psi_V(v))) h &= (\Psi_V(v))(h \circ L) \quad (\text{pela definição de } L_*) \\
		&= d(h \circ L)_a(v) \quad (\text{pela definição de } \Psi_V) \\
		&= \lim_{t \to 0} \frac{(h \circ L)(a + tv) - (h \circ L)(a)}{t} \\
		&= \lim_{t \to 0} \frac{h(L(a + tv)) - h(L(a))}{t}.
	\end{align*}
	
	Comparando os dois resultados, vemos que:
	\begin{equation*}
		(\Psi_W \circ L)(v) = (L_* \circ \Psi_V)(v).
	\end{equation*}
	Como isso é válido para todo $v \in V$, a aplicação $\Psi_W \circ L$ é igual a $L_* \circ \Psi_V$. Portanto, o diagrama comuta.
}
%%%% QUESTÃO 09 %%%%
\questao{
	Sejam $M^{n}$ e $N^{k}$ variedades suaves, com $M^{n}$ conexa, e $F : M^{n} \to N^{k}$ uma aplicação suave tal que $F_* : T_p M \to T_{F(p)} N$ é a aplicação nula, para todo $p \in M^{n}$. Mostre que $F$ é constante.
}
\solucao{
	
	\textbf{Análise do Enunciado}
	\begin{itemize}
		\item \textbf{Hipóteses:} $M^n$ é uma variedade suave e conexa; $N^k$ é uma variedade suave; $F : M^n \to N^k$ é suave; para todo $p \in M^n$, o diferencial $F_{*,p}$ é o operador nulo ($F_{*,p}(v) = 0$ para todo $v \in T_p M$).
		\item \textbf{Tese:} Provar que $F(p) = F(q)$ para quaisquer $p, q \in M^n$.
	\end{itemize}
	
	\textbf{Fundamentação Teórica}
	Para a demonstração, utilizaremos os seguintes conceitos e resultados:
	\begin{enumerate}
		\item \textbf{Conexidade por Caminhos:} Como demonstrado na \textbf{Questão 02 da Lista 01}, para uma variedade suave, a conexidade topológica é equivalente à conexidade por caminhos. Assim, a hipótese de que $M^n$ é conexa garante que, para quaisquer $p, q \in M^n$, existe uma curva suave $\gamma : [0, 1] \to M^n$ tal que $\gamma(0) = p$ e $\gamma(1) = q$.
		\item \textbf{Regra da Cadeia:} Se $\gamma : [0, 1] \to M^n$ é uma curva suave, a derivada de $\gamma$ no tempo $t$ é um vetor tangente $\gamma'(t) \in T_{\gamma(t)} M$. A aplicação composta $F \circ \gamma : [0, 1] \to N^k$ é uma curva suave em $N^k$ cujo vetor tangente é dado por:
		\begin{equation*}
			(F \circ \gamma)'(t) = F_*(\gamma'(t)).
		\end{equation*}
		\item \textbf{Teorema do Valor Médio (ou Fundamental do Cálculo):} Em $\mathbb{R}^k$, se a derivada de uma função é nula em todo o intervalo, a função é constante.
	\end{enumerate}
	
	\textbf{Desenvolvimento Formal}
	
	Sejam $p, q$ dois pontos quaisquer em $M^n$. Pelo resultado citado anteriormente, existe uma curva suave $\gamma : [0, 1] \to M^n$ tal que $\gamma(0) = p$ e $\gamma(1) = q$. Consideremos a aplicação composta:
	\begin{equation*}
		\sigma : [0, 1] \to N^k, \quad \sigma(t) = (F \circ \gamma)(t).
	\end{equation*}
	A aplicação $\sigma$ é uma curva suave em $N^k$. Para analisar o comportamento de $\sigma$, escolhemos um atlas suave para $N^k$. Para cada $t \in [0, 1]$, seja $(V, \psi)$ uma carta coordenada em $N^k$ tal que $\sigma(t) \in V$. A representação coordenada da curva, $\hat{\sigma} = \psi \circ \sigma$, é uma aplicação de um intervalo para $\mathbb{R}^k$.
	
	A derivada de $\hat{\sigma}$ no tempo $t$ é dada por:
	\begin{align*}
		\frac{d}{dt} (\psi \circ F \circ \gamma)(t) &= d\psi_{\sigma(t)} \left( (F \circ \gamma)'(t) \right) \\
		&= d\psi_{\sigma(t)} \left( F_*(\gamma'(t)) \right).
	\end{align*}
	Por hipótese, o diferencial $F_*$ é a aplicação nula para todo ponto de $M^n$. Como $\gamma'(t) \in T_{\gamma(t)} M$, temos que $F_*(\gamma'(t)) = 0$ para todo $t \in [0, 1]$. Substituindo na expressão anterior:
	\begin{equation*}
		\frac{d}{dt} (\psi \circ F \circ \gamma)(t) = d\psi_{\sigma(t)}(0) = 0.
	\end{equation*}
	Sendo a derivada de $\hat{\sigma}$ nula para todo $t$ em seu domínio, a função $\hat{\sigma}$ é constante em cada componente do intervalo $[0,1]$. Como a representação coordenada é constante e $\psi$ é um homeomorfismo, a curva $\sigma(t)$ é constante em $N^k$.
	
	Portanto, temos que:
	\begin{equation*}
		\sigma(0) = \sigma(1) \implies F(\gamma(0)) = F(\gamma(1)) \implies F(p) = F(q).
	\end{equation*}
	Como $p$ e $q$ foram escolhidos arbitrariamente em $M^n$, concluímos que $F$ assume um único valor em todo o domínio. Logo, $F$ é constante.
}
%%%% QUESTÃO 10 %%%%
\questao{
	Mostre que se uma $n$-variedade suave (não vazia) é difeomorfa a uma $k$-variedade suave, então $n = k$.
}
\solucao{
	
	\textbf{Análise do Enunciado}
	\begin{itemize}
		\item \textbf{Hipóteses:} $M^n$ e $N^k$ são variedades suaves não vazias e existe um difeomorfismo $f : M^n \to N^k$.
		\item \textbf{Tese:} Provar que as dimensões são iguais, ou seja, $n = k$.
	\end{itemize}
	
	\textbf{Fundamentação Teórica}
	Um difeomorfismo $f : M^n \to N^k$ é, por definição, uma aplicação suave que possui uma inversa suave $f^{-1}$. Consequentemente, $f$ é, em particular, um \textbf{difeomorfismo local} em cada ponto $p \in M^n$.
	
	Utilizaremos o resultado obtido na \textbf{Questão 07 desta Lista 02}, que estabelece que, se $f$ é um difeomorfismo local, então a aplicação diferencial $f_* : T_p M \to T_{f(p)} N$ é um isomorfismo de espaços vetoriais para todo $p \in M^n$.
	
	\textbf{Desenvolvimento Formal}
	
	Seja $f : M^n \to N^k$ um difeomorfismo. Fixemos um ponto $p \in M^n$. Como $f$ é um difeomorfismo, ele é um difeomorfismo local em $p$. Pelo resultado da \textbf{Questão 07}, a aplicação diferencial:
	\begin{equation*}
		f_* : T_p M \longrightarrow T_{f(p)} N
	\end{equation*}
	é um isomorfismo de espaços vetoriais. 
	
	De acordo com a teoria fundamental de álgebra linear, dois espaços vetoriais são isomorfos se, e somente se, possuem a mesma dimensão. Sabemos que:
	\begin{enumerate}
		\item A dimensão do espaço tangente $T_p M$ é igual à dimensão da variedade $M$, logo $\dim(T_p M) = n$.
		\item A dimensão do espaço tangente $T_{f(p)} N$ é igual à dimensão da variedade $N$, logo $\dim(T_{f(p)} N) = k$.
	\end{enumerate}
	
	Sendo $f_*$ um isomorfismo, segue necessariamente que:
	\begin{equation*}
		\dim(T_p M) = \dim(T_{f(p)} N) \implies n = k.
	\end{equation*}
	Dessa forma, a dimensão de uma variedade suave é um invariante sob difeomorfismos.
}
%%%% QUESTÃO 11 %%%%
\questao{
	Mostre que a aplicação $F : \mathbb{R}^3 \to \mathbb{R}^6$ definida por $$F(x, y, z) = (x^2, y^2, z^2, \sqrt{2}xy, \sqrt{2}xz, \sqrt{2}yz),$$ induz uma aplicação suave injetiva $f : \mathbb{RP}^2 \to \mathbb{S}^5$ cuja derivada $Df$ também é injetiva.
}
\solucao{
	\textbf{Análise do Enunciado}
	\begin{itemize}
		\item \textbf{Hipóteses:} $F$ é a aplicação polinomial dada de $\mathbb{R}^3$ para $\mathbb{R}^6$.
		\item \textbf{Tese 1:} $F$ induz uma aplicação suave $f: \mathbb{RP}^2 \to \mathbb{S}^5$.
		\item \textbf{Tese 2:} A aplicação $f$ é injetiva.
		\item \textbf{Tese 3:} A derivada (diferencial) $Df$ é injetiva em cada ponto (ou seja, $f$ é uma imersão).
	\end{itemize}
	
	\textbf{Fundamentação Teórica}
	Conforme demonstrado nas \textbf{Questões 19 e 20 da Lista 01}, a aplicação $F$ satisfaz $F(\mathbb{S}^2) \subset \mathbb{S}^5$ e $F(x) = F(-x)$ para todo $x \in \mathbb{S}^2$. Pela teoria de variedades quocientes, isso garante a existência de uma aplicação suave $\bar{f} : \mathbb{RP}^2 \to \mathbb{S}^5$ tal que $F|_{\mathbb{S}^2} = \bar{f} \circ \pi$, onde $\pi : \mathbb{S}^2 \to \mathbb{RP}^2$ é a projeção canônica.
	
	A injetividade de $f$ foi rigorosamente provada na \textbf{Questão 27 da Lista 01}, onde se demonstrou que $$F(x) = F(y) \implies x = \pm ,y$$ o que implica $[x] = [y]$ em $\mathbb{RP}^2$.
	
	Para a injetividade de $Df$, utilizaremos a caracterização do espaço tangente à esfera. Na \textbf{Questão 3(b) da Lista 01}, definimos $\mathbb{S}^2$ como o conjunto nível $f^{-1}(1)$ de $f(x) = |x|^2$. Pelo Teorema do Valor Regular, o espaço tangente $T_p \mathbb{S}^2$ é o núcleo do diferencial $df_p$. Como $\nabla f(p) = 2p$, temos que um vetor $v \in \mathbb{R}^3$ pertence a $T_p \mathbb{S}^2$ se, e somente se:
	\begin{equation*}
		df_p(v) = \langle \nabla f(p), v \rangle = \langle 2p, v \rangle = 0 \iff \langle p, v \rangle = 0.
	\end{equation*}
	Assim, $T_p \mathbb{S}^2$ é o hiperplano ortogonal ao vetor posição $p$.
	
	\textbf{Desenvolvimento Formal}
	
	\textbf{1. Indução e Injetividade de $f$:}
	As teses de que $f$ está bem definida, é suave e é injetiva seguem diretamente dos resultados obtidos nas \textbf{Questões 19, 20 e 27 da Lista 01}.
	
	\textbf{2. Injetividade da Derivada $Df$:}
	Para provar que $df_{[x]}$ é injetivo, analisamos o diferencial de $F$ em um ponto $p = (x, y, z) \in \mathbb{S}^2$. O diferencial $DF_p$ é representado pela matriz Jacobiana de $F$:
	\begin{equation*}
		DF_p = \begin{bmatrix} 
			2x & 0 & 0 \\ 
			0 & 2y & 0 \\ 
			0 & 0 & 2z \\ 
			\sqrt{2}y & \sqrt{2}x & 0 \\ 
			\sqrt{2}z & 0 & \sqrt{2}x \\ 
			0 & \sqrt{2}z & \sqrt{2}y 
		\end{bmatrix}.
	\end{equation*}
	Seja $v = (v_1, v_2, v_3) \in T_p \mathbb{S}^2$. Pela fundamentação acima, temos a restrição $\langle p, v \rangle = xv_1 + yv_2 + zv_3 = 0$.
	
	Suponhamos que $DF_p(v) = 0$. Isso implica que todas as componentes do vetor resultante sejam nulas:
	\begin{align*}
		2xv_1 = 0, \quad 2yv_2 = 0, \quad 2zv_3 = 0, \quad \sqrt{2}(yv_1 + xv_2) = 0, \quad \sqrt{2}(zv_1 + xv_3) = 0, \quad \text{e} \quad \sqrt{2}(zv_2 + yv_3) = 0.
	\end{align*}
	
	Analisamos a consistência desse sistema:
	\begin{enumerate}
		\item Se $x, y, z \neq 0$, as três primeiras equações implicam imediatamente que $v_1 = v_2 = v_3 = 0$, logo $v = 0$.
		\item Se alguma coordenada for nula, por exemplo, $x = 0$, então $p = (0, y, z)$ com $y^2 + z^2 = 1$. As equações tornam-se:
		\begin{equation*}
			0 = 0, \quad 2yv_2 = 0, \quad 2zv_3 = 0, \quad \sqrt{2}yv_1 = 0, \quad \sqrt{2}zv_1 = 0, \quad \sqrt{2}(zv_2 + yv_3) = 0.
		\end{equation*}
		Como $y$ e $z$ não podem ser ambos nulos, a quarta e quinta equações ($\sqrt{2}yv_1 = 0$ e $\sqrt{2}zv_1 = 0$) forçam $v_1 = 0$. 
		
		As equações $2yv_2 = 0$ e $2zv_3 = 0$ forçam $v_2 = 0$ e $v_3 = 0$ respectivamente (pois se $y=0$, então $z \neq 0$ e vice-versa).
	\end{enumerate}
	Em todos os casos, $DF_p(v) = 0 \implies v = 0$ para qualquer $v \in T_p \mathbb{S}^2$. Assim, o diferencial $DF_p$ é injetivo quando restrito ao espaço tangente da esfera.
	
	Como $F|_{\mathbb{S}^2} = f \circ \pi$, pela regra da cadeia temos $(F|_{\mathbb{S}^2})_* = f_* \circ \pi_*$. Visto que $\pi_*$ é um isomorfismo (pois é difeomorfismo local), a injetividade de $(F|_{\mathbb{S}^2})_*$ implica a injetividade de $f_*$ (ou $Df$).
}
%%%% QUESTÃO 12 %%%%
\questao{
	Mostre que a aplicação $F : \mathbb{R}^3 \to \mathbb{R}^5$ definida por 
	$$F(x, y, z) = \left(\sqrt{3}xy, \sqrt{3}xz, \sqrt{3}yz, \frac{\sqrt{3}}{2}(x^2 - y^2), \frac{1}{2}x^2 + \frac{1}{2}y^2 - z^2\right),$$ 
	induz uma aplicação suave injetiva $\bar{f} : \mathbb{RP}^2 \to \mathbb{S}^4$ cuja derivada $D\bar{f}$ também é injetiva.
}
\solucao{
	
	\textbf{Análise do Enunciado}
	\begin{itemize}
		\item \textbf{Hipóteses:} $F$ é uma aplicação polinomial de $\mathbb{R}^3$ para $\mathbb{R}^5$.
		\item \textbf{Tese 1:} Provar que $F$ induz uma aplicação $\bar{f} : \mathbb{RP}^2 \to \mathbb{S}^4$ bem definida e suave.
		\item \textbf{Tese 2:} Provar que $\bar{f}$ é injetiva.
		\item \textbf{Tese 3:} Provar que a derivada $D\bar{f}$ é injetiva em cada ponto (ou seja, $\bar{f}$ é uma imersão).
	\end{itemize}
	
	\textbf{Fundamentação Teórica}
	\begin{enumerate}
		\item \textbf{Mapas Induzidos em Espaços Projetivos:} De acordo com a teoria de variedades quocientes (Lee, Cap. 4), uma aplicação $F: \mathbb{R}^{n+1}\setminus\{0\} \to \mathbb{R}^{m+1}\setminus\{0\}$ induz um mapa suave $\bar{f}: \mathbb{RP}^n \to \mathbb{RP}^m$ se, e somente se, $F$ for constante nas fibras da projeção $\pi$, o que é garantido se $F$ for homogênea de algum grau $d$. Este princípio foi aplicado na \textbf{Questão 28 da Lista 01}.
		\item \textbf{Subvariedades e a Esfera $\mathbb{S}^n$:} Para que a imagem de $\bar{f}$ esteja em $\mathbb{S}^4$, devemos provar que a norma euclidiana $\|F(p)\|$ é unitária para todo $p \in \mathbb{S}^2$.
		\item \textbf{Injetividade e a Relação de Equivalência:} A injetividade de $\bar{f}$ requer que $F(x) = F(y) \implies x = \pm y$. A lógica de resolução de sistemas de equações quadráticas para este fim foi estabelecida na \textbf{Questão 27 da Lista 01}.
		\item \textbf{Imersões e o Diferencial:} Uma aplicação suave $f: M \to N$ é uma imersão se seu diferencial $df_p: T_p M \to T_{f(p)} N$ for injetivo para todo $p \in M$. Para $M = \mathbb{S}^2$, o espaço tangente $T_p \mathbb{S}^2$ é caracterizado por $\{v \in \mathbb{R}^3 : \langle p, v \rangle = 0\}$ (Lee, Cap. 3; \textbf{Questão 3b da Lista 01}).
	\end{enumerate}
	
	\textbf{Desenvolvimento Formal}
	
	\textbf{1. Indução e Suavidade da Aplicação $\bar{f}$}
	
	Seja $p = (x, y, z) \in \mathbb{S}^2$. Note que $F$ é homogênea de grau 2, pois para qualquer $\lambda \in \mathbb{R}\setminus\{0\}$:
	\begin{align*}
		F(\lambda x, \lambda y, \lambda z) &= \left(\sqrt{3}\lambda^2 xy, \sqrt{3}\lambda^2 xz, \sqrt{3}\lambda^2 yz, \frac{\sqrt{3}}{2}(\lambda^2 x^2 - \lambda^2 y^2), \frac{1}{2}\lambda^2 x^2 + \frac{1}{2}\lambda^2 y^2 - \lambda^2 z^2\right) \\
		&= \lambda^2 F(x, y, z).
	\end{align*}
	Pela fundamentação citada, $F$ induz um mapa $\bar{f}: \mathbb{RP}^2 \to \mathbb{RP}^4$ dado por $\bar{f}([p]) = [F(p)]$. Como $\|F(p)\|^2 = 1$ para $p \in \mathbb{S}^2$ (verificação abaixo), a imagem de cada classe $[p]$ pode ser representada por um único vetor unitário em $\mathbb{S}^4$, permitindo a definição $\bar{f}: \mathbb{RP}^2 \to \mathbb{S}^4$.
	
	Verificação da norma para $p \in \mathbb{S}^2$ ($x^2+y^2+z^2=1$):
	\begin{align*}
		\|F(p)\|^2 &= 3(x^2y^2 + x^2z^2 + y^2z^2) + \frac{3}{4}(x^4 + y^4 - 2x^2y^2) + \left(\frac{1}{4}x^4 + \frac{1}{4}y^4 + z^4 + \frac{1}{2}x^2y^2 - x^2z^2 - y^2z^2\right) \\
		&= x^4 + y^4 + z^4 + 2x^2y^2 + 2x^2z^2 + 2y^2z^2 = (x^2 + y^2 + z^2)^2 = 1.
	\end{align*}
	
	\textbf{2. Prova de Injetividade}
	
	Sejam $[x], [y] \in \mathbb{RP}^2$ tais que $\bar{f}([x]) = \bar{f}([y])$. Para representantes $x, y \in \mathbb{S}^2$, temos $F(x) = F(y)$, resultando no sistema:
	\begin{enumerate}[label=(\roman*)]
		\item $x_1x_2 = y_1y_2$; (ii) $x_1x_3 = y_1y_3$; (iii) $x_2x_3 = y_2y_3$; (iv) $x_1^2 - x_2^2 = y_1^2 - y_2^2$; (v) $\frac{1}{2}x_1^2 + \frac{1}{2}x_2^2 - x_3^2 = \frac{1}{2}y_1^2 + \frac{1}{2}y_2^2 - y_3^2$.
	\end{enumerate}
	Da equação (v), e utilizando $x_1^2+x_2^2 = 1-x_3^2$, temos:
	\begin{align*}
		\frac{1}{2}(1-x_3^2) - x_3^2 &= \frac{1}{2}(1-y_3^2) - y_3^2 \implies \frac{1}{2} - \frac{3}{2}x_3^2 = \frac{1}{2} - \frac{3}{2}y_3^2 \implies x_3^2 = y_3^2.
	\end{align*}
	Combinando $x_1^2+x_2^2 = y_1^2+y_2^2$ com a equação (iv) ($x_1^2-x_2^2 = y_1^2-y_2^2$), obtemos por soma e subtração:
	\begin{align*}
		2x_1^2 = 2y_1^2 \implies x_1^2 = y_1^2 \quad \text{e} \quad 2x_2^2 = 2y_2^2 \implies x_2^2 = y_2^2.
	\end{align*}
	Como $x_i^2 = y_i^2$ e $x_ix_j = y_iy_j$, a lógica da \textbf{Questão 27 da Lista 01} implica que $y = \pm x$. Logo, $[x] = [y]$.
	
	\textbf{3. Injetividade da Derivada $D\bar{f}$}
	
	A matriz Jacobiana $DF_p$ é:
	\begin{equation*}
		DF_p = \begin{bmatrix} 
			\sqrt{3}y & \sqrt{3}x & 0 \\ \sqrt{3}z & 0 & \sqrt{3}x \\ 0 & \sqrt{3}z & \sqrt{3}y \\ \sqrt{3}x & -\sqrt{3}y & 0 \\ x & y & -2z 
		\end{bmatrix}.
	\end{equation*}
	Seja $v \in T_p\mathbb{S}^2$, logo $xv_1 + yv_2 + zv_3 = 0$. Suponhamos $DF_p(v) = 0$:
	\begin{align*}
		\sqrt{3}(yv_1 + xv_2) &= 0 \quad (1) \\
		\sqrt{3}(xv_1 - yv_2) &= 0 \quad (2) \\
		\sqrt{3}(zv_1 + xv_3) &= 0 \quad (3)
	\end{align*}
	Das equações (1) e (2), o determinante do sistema em $v_1, v_2$ é $-(y^2+x^2)$. 
	Se $x^2+y^2 \neq 0$, então $v_1=v_2=0$, e por $\langle p, v \rangle = 0$, temos $zv_3=0 \implies v_3=0$.
	Se $x^2+y^2=0$, então $x=0, y=0 \implies z=\pm 1$. A condição $\langle p, v \rangle = 0$ implica $v_3=0$. Substituindo na Jacobiana, a linha 2 dá $\pm\sqrt{3}v_1=0$ e a linha 3 dá $\pm\sqrt{3}v_2=0$. 
	Em todos os casos, $v=0$. O núcleo de $DF_p|_{T_p\mathbb{S}^2}$ é trivial. Como $\pi_*$ é um isomorfismo, a derivada $D\bar{f}$ é injetiva.
}
%%%% QUESTÃO 13 %%%%
\questao{
	Seja $M^n$ uma $n$-variedade suave. Mostre que o fibrado tangente $TM$ admite uma topologia e uma estrutura suave de tal forma que $TM$ é uma $2n$-variedade suave.
}
\solucao{
	
	\textbf{Análise do Enunciado}
	\begin{itemize}
		\item \textbf{Hipóteses:} $M^n$ é uma $n$-variedade suave.
		\item \textbf{Tese:} Provar que o conjunto $TM = \bigcup_{p \in M} T_p M$ possui uma topologia e um atlas suave que o tornam uma variedade suave de dimensão $2n$.
	\end{itemize}
	
	\textbf{Fundamentação Teórica}
	\begin{enumerate}
		\item \textbf{Definição do Espaço Tangente:} Como estabelecido na \textbf{Lista 01}, para cada $p \in M$, $T_p M$ é um espaço vetorial de dimensão $n$. O fibrado tangente $TM$ é a união disjunta de todos esses espaços.
		\item \textbf{Construção de Atlas Induzidos:} Se $\mathcal{U} = \{(U_\alpha, \phi_\alpha)\}_{\alpha \in \Lambda}$ é um atlas suave para $M$, podemos definir a \textbf{extensão} desse atlas para $TM$. Para cada $\alpha \in \Lambda$, definimos o conjunto $T U_\alpha = \bigcup_{p \in U_\alpha} T_p M$, que é um subconjunto de $TM$.
		\item \textbf{Mapeamentos de Coordenadas:} A aplicação $\Phi_\alpha : T U_\alpha \to \mathbb{R}^{2n}$ é definida enviando um vetor $(p, v)$ para a coordenada do ponto $p$ e para as componentes do vetor $v$ na base induzida pela carta $\phi_\alpha$.
		\item \textbf{Axiomas de Variedade:} Para que $TM$ seja uma variedade, devemos garantir que a topologia induzida seja de Hausdorff, possua base enumerável e que as funções de transição sejam $C^\infty$.
	\end{enumerate}
	
	\textbf{Desenvolvimento Formal}
	
	\textbf{1. Definição das Cartas de $TM$}
	
	Seja $\phi_\alpha : U_\alpha \to \mathbb{R}^n$ uma carta de $M$. Para qualquer $p \in U_\alpha$, o diferencial da carta no ponto $p$ é uma aplicação linear:
	\begin{equation*}
		d(\phi_\alpha)_p : T_p M \to T_{\phi_\alpha(p)} \mathbb{R}^n \cong \mathbb{R}^n.
	\end{equation*}
	Como $\phi_\alpha$ é um homeomorfismo, $d(\phi_\alpha)_p$ é um isomorfismo de espaços vetoriais. Definimos a aplicação $\Phi_\alpha : T U_\alpha \to \mathbb{R}^{2n}$ por:
	\begin{equation*}
		\Phi_\alpha(p, v) = (\phi_\alpha(p), d(\phi_\alpha)_p(v)).
	\end{equation*}
	Seja $u = \phi_\alpha(p) \in \mathbb{R}^n$ e $\xi = d(\phi_\alpha)_p(v) \in \mathbb{R}^n$. A inversa $\Phi_\alpha^{-1} : \phi_\alpha(U_\alpha) \times \mathbb{R}^n \to T U_\alpha$ é dada por:
	\begin{equation*}
		\Phi_\alpha^{-1}(u, \xi) = (\phi_\alpha^{-1}(u), \sum_{i=1}^n \xi^i \frac{\partial}{\partial x^i} \Big|_{\phi_\alpha^{-1}(u)}).
	\end{equation*}
	
	\textbf{2. Construção da Topologia}
	
	Definimos a topologia em $TM$ declarando que um conjunto $W \subset TM$ é aberto se, e somente se, $\Phi_\alpha(W \cap T U_\alpha)$ é um aberto de $\mathbb{R}^{2n}$ para cada $\alpha \in \Lambda$. 
	
	Como as cartas $\Phi_\alpha$ são homeomorfismos sobre abertos de $\mathbb{R}^{2n}$, e $M$ é de Hausdorff e possui base enumerável, as mesmas propriedades são herdadas por $TM$. De fato, se $p, q \in M$ são distintos, eles podem ser separados em $M$, e consequentemente seus fibrados tangentes podem ser separados em $TM$. Se $\mathcal{B}$ é uma base enumerável para $M$, então $\mathcal{B} \times \mathbb{Q}^n$ (ou qualquer base enumerável de $\mathbb{R}^n$) induz uma base enumerável para $TM$.
	
	\textbf{3. Suavidade das Funções de Transição}
	
	Sejam $(U_\alpha, \phi_\alpha)$ e $(U_\beta, \phi_\beta)$ duas cartas de $M$ com $U_\alpha \cap U_\beta \neq \emptyset$. A função de transição na base $M$ é $$\tau = \phi_\beta \circ \phi_\alpha^{-1} : \phi_\alpha(U_\alpha \cap U_\beta) \to \phi_\beta(U_\alpha \cap U_\beta).$$
	
	A função de transição em $TM$ é $\Psi = \Phi_\beta \circ \Phi_\alpha^{-1}$. Para $(u, \xi) \in \mathbb{R}^n \times \mathbb{R}^n$:
	\begin{align*}
		\Psi(u, \xi) &= \Phi_\beta(\Phi_\alpha^{-1}(u, \xi)) \\
		&= \Phi_\beta(p, v) \quad \text{onde } p = \phi_\alpha^{-1}(u) \text{ e } v = \sum \xi^i \frac{\partial}{\partial x^i} \Big|_p \\
		&= (\phi_\beta(p), d(\phi_\beta)_p(v)) \\
		&= (\phi_\beta(\phi_\alpha^{-1}(u)), d(\phi_\beta)_p(d(\phi_\alpha)^{-1}_p(\xi))).
	\end{align*}
	Pela regra da cadeia para diferenciais, temos $d(\phi_\beta)_p \circ d(\phi_\alpha)^{-1}_p = d(\phi_\beta \circ \phi_\alpha^{-1})_{\phi_\alpha(p)}$. Portanto:
	\begin{equation*}
		\Psi(u, \xi) = (\tau(u), d\tau_u(\xi)).
	\end{equation*}
	Como $M$ é uma variedade suave, $\tau$ é uma aplicação $C^\infty$. O diferencial $d\tau_u$ é representado pela matriz Jacobiana $J_\tau(u)$, cujas entradas são as derivadas parciais de $\tau$. Como $\tau$ é $C^\infty$, suas derivadas parciais também são $C^\infty$. 
	
	Logo, a aplicação $(u, \xi) \mapsto (\tau(u), J_\tau(u) \cdot \xi)$ é a composição de funções suaves em $\mathbb{R}^{2n}$, sendo, portanto, suave.
	
	\textbf{Conclusão}
	
	Tendo definido um atlas $\mathcal{T} = \{(T U_\alpha, \Phi_\alpha)\}$ com funções de transição suaves, e tendo verificado as propriedades topológicas de Hausdorff e base enumerável, concluímos que $TM$ é uma variedade suave de dimensão $n + n = 2n$.
}
%%%% RESULTADO EXTRA %%%%
Abaixo, estabelecemos um resultado fundamental da teoria de Grupos de Lie. Este teorema demonstra que a estrutura algébrica de um grupo, quando compatível com a estrutura de variedade suave, impõe uma rigidez geométrica ao fibrado tangente, tornando-o globalmente trivial. Este resultado será utilizado como base para as \textbf{Soluções Alternativas} das Questões 14 e 15, permitindo que a prova de trivialidade do fibrado tangente seja reduzida a uma propriedade de simetria do grupo.

\resultadoextra{Trivialidade do Fibrado Tangente de Grupos de Lie}{
	Seja $G$ um Grupo de Lie de dimensão $n$. Então $TG$ é difeomorfo a $G \times \mathbb{R}^n$.
}
\solucao{

	Seja $e \in G$ o elemento neutro do grupo. Para cada $g \in G$, definimos a aplicação de translação à esquerda $L_g : G \to G$ por $L_g(h) = gh$. Por definição de Grupo de Lie, a multiplicação é suave, logo $L_g$ é uma aplicação suave. Como $L_g$ possui inversa suave $L_{g^{-1}}$, então $L_g$ é um difeomorfismo de $G$ em si mesmo.
	
	O diferencial de $L_g$ no ponto $e$, denotado por $(L_g)_* : T_e G \to T_g G$, é um isomorfismo de espaços vetoriais, pois é o diferencial de um difeomorfismo. 
	
	Definimos a aplicação $\Psi : G \times T_e G \to TG$ por:
	\begin{equation*}
		\Psi(g, v) = (g, (L_g)_* v).
	\end{equation*}
	
	\textbf{1. Bijetividade:} Para qualquer vetor tangente $(p, w) \in TG$, onde $p \in G$ e $w \in T_p G$, existe um único $v \in T_e G$ tal que $(L_p)_* v = w$, dado que $v = (L_{p^{-1}})_* w$. Assim, $\Psi(p, v) = (p, w)$, provando que $\Psi$ é bijetiva.
	
	\textbf{2. Suavidade:} A suavidade de $\Psi$ decorre da suavidade da operação de multiplicação do grupo. Em coordenadas locais, a aplicação $(g, v) \mapsto (L_g)_* v$ é representada por funções suaves nas entradas de $g$ e $v$.
	
	\textbf{3. Suavidade da Inversa:} A inversa $\Psi^{-1} : TG \to G \times T_e G$ é dada por:
	\begin{equation*}
		\Psi^{-1}(p, w) = (p, (L_{p^{-1}})_* w).
	\end{equation*}
	Visto que a inversão $p \mapsto p^{-1}$ e a translação são suaves, $\Psi^{-1}$ também é suave.
	
	Como $T_e G$ é um espaço vetorial de dimensão $n$, ele é isomorfo a $\mathbb{R}^n$. Logo, temos o difeomorfismo $TG \cong G \times \mathbb{R}^n$.
}
%%%% QUESTÃO 14 %%%%
\questao{
	Mostre que $T\mathbb{R}^n$ é difeomorfo a $\mathbb{R}^{2n}$.
}
\solucao{
	
	\textbf{Análise do Enunciado}
	\begin{itemize}
		\item \textbf{Hipóteses:} O espaço base é a variedade suave $\mathbb{R}^n$.
		\item \textbf{Tese:} Provar a existência de um difeomorfismo entre o fibrado tangente $T\mathbb{R}^n$ e o espaço euclidiano $\mathbb{R}^{2n}$.
	\end{itemize}
	
	\textbf{Fundamentação Teórica}
	Utilizamos a construção de atlas induzidos para o fibrado tangente (Lee, Cap. 3). Para uma variedade que admite uma carta global $(M, \phi)$, o fibrado tangente é trivializado pela aplicação $\Phi(p, v) = (\phi(p), d\phi_p(v))$.
	
	\textbf{Desenvolvimento Formal}
	
	Seja $\phi = \text{Id}_{\mathbb{R}^n}$ a carta global para $\mathbb{R}^n$. Para qualquer ponto $p \in \mathbb{R}^n$, o diferencial da identidade $$d(\text{Id}_{\mathbb{R}^n})_p : T_p\mathbb{R}^n \to T_p\mathbb{R}^n \cong \mathbb{R}^n $$ é a aplicação identidade nos espaços tangentes. Definimos $\Psi : T\mathbb{R}^n \to \mathbb{R}^{2n}$ por:
	\begin{equation*}
		\Psi(p, v) = (p, v).
	\end{equation*}
	Onde $v$ é representado por suas componentes na base canônica $\left\{\frac{\partial}{\partial x^i}\right\}$.
	
	\textbf{1. Bijetividade:} Para qualquer par $(u, \xi) \in \mathbb{R}^n \times \mathbb{R}^n$, existe um único ponto $p = u \in \mathbb{R}^n$ e um único vetor $v = \sum_{i=1}^n \xi^i \frac{\partial}{\partial x^i} \Big|_p \in T_p\mathbb{R}^n$ tal que $\Psi(p, v) = (u, \xi)$. Logo, $\Psi$ é bijetiva.
	
	\textbf{2. Suavidade de $\Psi$ e $\Psi^{-1}$:}
	A aplicação $\Psi$ é, por construção, a própria carta $\Phi$ definida na Questão 13 para a identidade. Como a função de transição para uma única carta é a identidade $\text{Id}_{\mathbb{R}^{2n}}$, que é infinitamente diferenciável, $\Psi$ é suave.
	
	A inversa $\Psi^{-1} : \mathbb{R}^{2n} \to T\mathbb{R}^n$ é dada por:
	\begin{equation*}
		\Psi^{-1}(u, \xi) = \left( u, \sum_{i=1}^n \xi^i \frac{\partial}{\partial x^i} \Big|_u \right).
	\end{equation*}
	Sendo $\Psi^{-1}$ a inversa de um homeomorfismo que define a estrutura suave de $T\mathbb{R}^n$, ela é, por definição, suave.
	
	\textbf{Conclusão}
	Como $\Psi$ é uma bijecção suave com inversa suave, ela é um difeomorfismo. Portanto, $T\mathbb{R}^n \cong \mathbb{R}^{2n}$.
}
\solucaoalt{
	
	\textbf{Análise via Grupos de Lie}
	
	Como provamos na \textbf{Questão 36(a) da Lista 01}, $\mathbb{R}^n$ é um Grupo de Lie sob a operação de adição usual. 
	
	De acordo com o resultado extra demonstrado anteriormente, o fibrado tangente de qualquer Grupo de Lie $G$ é trivial, sendo difeomorfo a $G \times T_e G$. No caso de $\mathbb{R}^n$, o elemento neutro é o vetor nulo $0$. A translação à esquerda $L_g: \mathbb{R}^n \to \mathbb{R}^n$ é dada por $L_g(h) = g+h$.
	
	O diferencial $(L_g)_* : T_0\mathbb{R}^n \to T_g\mathbb{R}^n$ é um isomorfismo. Definimos a aplicação:
	\begin{equation*}
		\Psi : \mathbb{R}^n \times T_0\mathbb{R}^n \to T\mathbb{R}^n, \quad \Psi(g, v) = (g, (L_g)_* v).
	\end{equation*}
	Visto que $T_0\mathbb{R}^n$ é um espaço vetorial de dimensão $n$, temos $T_0\mathbb{R}^n \cong \mathbb{R}^n$. Assim, $\Psi$ estabelece um difeomorfismo entre $\mathbb{R}^n \times \mathbb{R}^n$ e $T\mathbb{R}^n$. Como $\mathbb{R}^n \times \mathbb{R}^n$ é canonicamente difeomorfo a $\mathbb{R}^{2n}$, segue que $T\mathbb{R}^n \cong \mathbb{R}^{2n}$.
}
%%%% QUESTÃO 15 %%%%
\questao{
	Mostre que $T\mathbb{S}^1$ é difeomorfo a $\mathbb{S}^1 \times \mathbb{R}$.
}
\solucao{
	
	\textbf{Análise do Enunciado}
	\begin{itemize}
		\item \textbf{Hipóteses:} O espaço base é a 1-esfera $\mathbb{S}^1 \subset \mathbb{R}^2$.
		\item \textbf{Tese:} Provar que $T\mathbb{S}^1 \cong \mathbb{S}^1 \times \mathbb{R}$.
	\end{itemize}
	
	\textbf{Fundamentação Teórica}
	Uma variedade $M^n$ é paralelizável se admite $n$ campos de vetores suaves que formam uma base de $T_p M$ para todo $p \in M$. Se $M$ é paralelizável, então $TM \cong M \times \mathbb{R}^n$. Para $\mathbb{S}^1 \subset \mathbb{R}^2$, o espaço tangente em $p = (x, y)$ é $T_p\mathbb{S}^1 = \{v \in \mathbb{R}^2 : \langle p, v \rangle = 0\}$.
	
	\textbf{Desenvolvimento Formal}
	
	\textbf{1. Construção de um Campo de Vetores Global:}
	Definimos o campo de vetores $X$ em $\mathbb{S}^1$ por $X(p) = (-y, x)$ para $p=(x,y)$.
	\begin{itemize}
		\item \textbf{Tangencialidade:} $\langle p, X(p) \rangle = \langle (x, y), (-y, x) \rangle = -xy + yx = 0$. Logo, $X(p) \in T_p\mathbb{S}^1$.
		\item \textbf{Não-nulidade:} $\|X(p)\|^2 = (-y)^2 + x^2 = 1$. Logo, $X(p)$ nunca é nulo.
		\item \textbf{Suavidade:} As componentes são lineares, logo $X$ é suave.
	\end{itemize}
	Como $\dim(T_p\mathbb{S}^1)=1$ e $X(p) \neq 0$, o conjunto $\{X(p)\}$ é uma base de $T_p\mathbb{S}^1$ para todo $p$.
	
	\textbf{2. Definição da Aplicação $\Psi$:}
	Definimos $\Psi : \mathbb{S}^1 \times \mathbb{R} \to T\mathbb{S}^1$ por $\Psi(p, \lambda) = (p, \lambda X(p))$.
	
	\textbf{3. Prova de que $\Psi$ é um Difeomorfismo}
	
	\textbf{A. Bijetividade e a Decomposição de $v$:}
	Seja $(p, v) \in T\mathbb{S}^1$. Como $\dim(T_p\mathbb{S}^1) = 1$ e o vetor $X(p) = (-y, x)$ é não nulo ($\|X(p)\| = 1$), o conjunto $\{X(p)\}$ constitui uma base para o espaço tangente em $p$. Portanto, qualquer vetor $v \in T_p\mathbb{S}^1$ pode ser escrito unicamente como um múltiplo escalar de $X(p)$:
	\begin{equation*}
		v = \lambda X(p), \quad \text{para algum } \lambda \in \mathbb{R}.
	\end{equation*}
	Para determinar $\lambda$ explicitamente, aplicamos o produto escalar em ambos os lados da equação com o vetor $X(p)$:
	\begin{align*}
		\langle v, X(p) \rangle &= \langle \lambda X(p), X(p) \rangle \\
		\langle v, X(p) \rangle &= \lambda \langle X(p), X(p) \rangle \\
		\langle v, X(p) \rangle &= \lambda \|X(p)\|^2 = \lambda (1) = \lambda.
	\end{align*}
	Assim, temos a representação única $v = \langle v, X(p) \rangle X(p)$. Isso prova que para cada $(p, v) \in T\mathbb{S}^1$, existe um único $\lambda = \langle v, X(p) \rangle$ tal que $\Psi(p, \lambda) = (p, v)$, logo $\Psi$ é bijetiva.
	
	\textbf{B. Suavidade de $\Psi$ via Coordenadas:}
	Para provar que $\Psi$ é suave, devemos mostrar que sua representação coordenada $\Phi \circ \Psi$ é suave, onde $\Phi : T U \to \mathbb{R}^2$ é a carta do fibrado tangente definida na \textbf{Questão 13}.
	\begin{align*}
		(\Phi \circ \Psi)(p, \lambda) &= \Phi(p, \lambda X(p)) \\
		&= (\phi(p), d\phi_p(\lambda X(p))) \\
		&= (\phi(p), \lambda \cdot d\phi_p(X(p))).
	\end{align*}
	Seja $u = \phi(p)$. A componente $d\phi_p(X(p))$ é a representação coordenada do campo de vetores $X$ na carta $\phi$, denotada por $\hat{X}(u)$. Como $X$ é um campo de vetores suave, $\hat{X} : \phi(U) \to \mathbb{R}$ é uma função suave. A aplicação $\Phi \circ \Psi(u, \lambda) = (u, \lambda \hat{X}(u))$ é a composição de funções suaves (multiplicação e projeção), sendo, portanto, suave.
	
	\textbf{C. Suavidade de $\Psi^{-1}$:}
	A inversa $\Psi^{-1} : T\mathbb{S}^1 \to \mathbb{S}^1 \times \mathbb{R}$ é definida por:
	\begin{equation*}
		\Psi^{-1}(p, v) = (p, \langle v, X(p) \rangle).
	\end{equation*}
	Analisamos a suavidade da segunda componente $\lambda(p, v) = \langle v, X(p) \rangle$ em coordenadas locais. Seja $\Phi(p, v) = (u, \xi)$, onde $\xi$ são as componentes de $v$ na base $\{\frac{\partial}{\partial x}\}$. Então:
	\begin{align*}
		\lambda(u, \xi) &= \langle \Phi^{-1}(u, \xi), X(\Phi^{-1}(u)) \rangle \\
		&= \sum_{i=1}^1 \xi^i \langle \frac{\partial}{\partial x^i} \Big|_p, X(p) \rangle.
	\end{align*}
	Como $X(p)$ é um campo suave, a função $u \mapsto \langle \frac{\partial}{\partial x^i} \Big|_p, X(p) \rangle$ é suave. A aplicação $\lambda(u, \xi)$ é, portanto, uma soma de produtos de funções suaves em $u$ e $\xi$. Visto que a primeira componente de $\Psi^{-1}$ é a projeção canônica $\pi$, que é suave por definição, $\Psi^{-1}$ é suave.
	
	\textbf{Conclusão:}
	Como $\Psi$ é uma bijecção suave com inversa suave, ela é um difeomorfismo. Logo, $T\mathbb{S}^1 \cong \mathbb{S}^1 \times \mathbb{R}$.
}
\solucaoalt{
	
	\textbf{Análise via Grupos de Lie}
	
	Como provamos na \textbf{Questão 37(b) da Lista 01}, a 1-esfera $\mathbb{S}^1$ é um Grupo de Lie sob a multiplicação de números complexos.
	
	De acordo com o resultado extra demonstrado anteriormente, todo Grupo de Lie $G$ possui fibrado tangente trivial, sendo difeomorfo a $G \times T_e G$. No caso de $\mathbb{S}^1$, o elemento neutro é $e=1$. A translação à esquerda $L_g: \mathbb{S}^1 \to \mathbb{S}^1$ é dada por $L_g(h) = gh$.
	
	O diferencial $(L_g)_* : T_1\mathbb{S}^1 \to T_g\mathbb{S}^1$ é um isomorfismo para todo $g \in \mathbb{S}^1$. Definimos a aplicação:
	\begin{equation*}
		\Psi : \mathbb{S}^1 \times T_1\mathbb{S}^1 \to T\mathbb{S}^1, \quad \Psi(g, v) = (g, (L_g)_* v).
	\end{equation*}
	Visto que $T_1\mathbb{S}^1$ é um espaço vetorial de dimensão 1, temos $T_1\mathbb{S}^1 \cong \mathbb{R}$. Pela teoria de Grupos de Lie, a aplicação $\Psi$ é um difeomorfismo. Portanto, $T\mathbb{S}^1 \cong \mathbb{S}^1 \times \mathbb{R}$.
}
%%%% QUESTÃO 16 %%%%
\questao{
	Seja $F : \mathbb{R}^n \times \mathbb{R}^n \to \mathbb{R}^2$ a função suave dada por $F(x,v) = (|x|^2 - 1, \langle x, v \rangle)$. Mostre que $F^{-1}(0,0)$ é difeomorfo ao fibrado tangente de $\mathbb{S}^{n-1}$.
}
\solucao{
	
	\textbf{Análise do Enunciado}
	\begin{itemize}
		\item \textbf{Hipóteses:} $F$ é uma aplicação de $\mathbb{R}^{2n}$ para $\mathbb{R}^2$ definida por $F(x,v) = (f_1(x,v), f_2(x,v))$, onde $f_1(x,v) = |x|^2 - 1$ e $f_2(x,v) = \langle x, v \rangle$.
		\item \textbf{Tese:} Provar que o conjunto nível $M = F^{-1}(0,0)$ é difeomorfo ao fibrado tangente $T\mathbb{S}^{n-1}$.
	\end{itemize}
	
	\textbf{Fundamentação Teórica}
	\begin{enumerate}
		\item \textbf{Caracterização do Espaço Tangente da Esfera:} Conforme provamos na \textbf{Questão 3(b) da Lista 01}, para qualquer $x \in \mathbb{S}^{n-1}$, o espaço tangente é dado por $T_x \mathbb{S}^{n-1} = \{ v \in \mathbb{R}^n : \langle x, v \rangle = 0 \}$.
		\item \textbf{Definição de Fibrado Tangente:} Por definição (Lee, Cap. 3), o fibrado tangente de $\mathbb{S}^{n-1}$ é a união disjunta de seus espaços tangentes:
		\begin{equation*}
			T\mathbb{S}^{n-1} = \bigcup_{x \in \mathbb{S}^{n-1}} \{x\} \times T_x \mathbb{S}^{n-1} = \{ (x, v) \in \mathbb{R}^n \times \mathbb{R}^n : |x|^2 = 1 \text{ e } \langle x, v \rangle = 0 \}.
		\end{equation*}
		\item \textbf{Teorema do Valor Regular:} Se $c$ é um valor regular de uma aplicação suave $F$, então $F^{-1}(c)$ é uma subvariedade suave de codimensão igual à dimensão do contradomínio.
	\end{enumerate}
	
	\textbf{Desenvolvimento Formal}
	
	\textbf{1. Igualdade de Conjuntos}
	
	Analisemos o conjunto $M = F^{-1}(0,0)$. Por definição:
	\begin{align*}
		M &= \{ (x, v) \in \mathbb{R}^n \times \mathbb{R}^n : F(x, v) = (0, 0) \} \\
		&= \{ (x, v) \in \mathbb{R}^n \times \mathbb{R}^n : |x|^2 - 1 = 0 \text{ e } \langle x, v \rangle = 0 \} \\
		&= \{ (x, v) \in \mathbb{R}^n \times \mathbb{R}^n : |x|^2 = 1 \text{ e } \langle x, v \rangle = 0 \}.
	\end{align*}
	Comparando com a definição de $T\mathbb{S}^{n-1}$ exposta na fundamentação, vemos que $M$ e $T\mathbb{S}^{n-1}$ são idênticos como subconjuntos de $\mathbb{R}^{2n}$.
	
	\textbf{2. Prova de que $(0,0)$ é Valor Regular}
	
	Para garantir que $M$ é uma variedade suave, devemos mostrar que o diferencial $DF_{(x,v)} : \mathbb{R}^{2n} \to \mathbb{R}^2$ é sobrejetivo para todo $(x, v) \in M$. A matriz Jacobiana de $F$ é dada por:
	\begin{equation*}
		DF_{(x,v)} = \begin{bmatrix} 
			\frac{\partial f_1}{\partial x_1} & \dots & \frac{\partial f_1}{\partial x_n} & \frac{\partial f_1}{\partial v_1} & \dots & \frac{\partial f_1}{\partial v_n} \\
			\frac{\partial f_2}{\partial x_1} & \dots & \frac{\partial f_2}{\partial x_n} & \frac{\partial f_2}{\partial v_1} & \dots & \frac{\partial f_2}{\partial v_n}
		\end{bmatrix} = \begin{bmatrix} 
			2x_1 & \dots & 2x_n & 0 & \dots & 0 \\
			v_1 & \dots & v_n & x_1 & \dots & x_n
		\end{bmatrix}.
	\end{equation*}
	
	O diferencial é sobrejetivo se, e somente se, as duas linhas da matriz Jacobiana forem linearmente independentes. Suponhamos que existam $\alpha, \beta \in \mathbb{R}$ tais que:
	\begin{equation*}
		\alpha (2x_1, \dots, 2x_n, 0, \dots, 0) + \beta (v_1, \dots, v_n, x_1, \dots, x_n) = (0, \dots, 0).
	\end{equation*}
	Isso nos dá dois sistemas de equações:
	\begin{enumerate}
		\item Para as últimas $n$ componentes: $\beta x_i = 0$ para todo $i \in \{1, \dots, n\}$. Como $(x, v) \in M$, temos $|x|^2 = 1$, logo $x \neq 0$. Isso força $\beta = 0$.
		\item Para as primeiras $n$ componentes: $\alpha (2x_i) + \beta v_i = 0$. Substituindo $\beta = 0$, temos $2\alpha x_i = 0$. Novamente, como $x \neq 0$, temos $\alpha = 0$.
	\end{enumerate}
	Como $\alpha = \beta = 0$ é a única solução, as linhas são linearmente independentes. Portanto, $(0,0)$ é um valor regular de $F$, e $M$ é uma subvariedade suave de $\mathbb{R}^{2n}$ de dimensão $2n - 2$.
	
	\textbf{3. Difeomorfismo com $T\mathbb{S}^{n-1}$}
	
	Definimos a aplicação $\Psi : M \to T\mathbb{S}^{n-1}$ como a aplicação identidade $\Psi(x, v) = (x, v)$.
	\begin{itemize}
		\item \textbf{Bijetividade:} Já provamos que $M = T\mathbb{S}^{n-1}$ como conjuntos.
		\item \textbf{Suavidade:} $\Psi$ é a restrição da identidade $\text{Id}_{\mathbb{R}^{2n}}$ a uma subvariedade, logo é suave.
		\item \textbf{Inversa:} A inversa $\Psi^{-1}$ também é a restrição da identidade, logo é suave.
	\end{itemize}
	
	Visto que $T\mathbb{S}^{n-1}$ possui a estrutura de variedade suave induzida por seu atlas (Questão 13), e $M$ possui a estrutura de subvariedade induzida pelo Teorema do Valor Regular, a aplicação $\Psi$ é um difeomorfismo.
	
	\textbf{Conclusão}
	O conjunto nível $F^{-1}(0,0)$ é a representação extrínseca do fibrado tangente da esfera, e ambos são difeomorfos.
}
%%%% QUESTÃO 17 %%%%
\questao{
	Sejam $M^n$ uma $n$-variedade suave e $Y : M^n \to TM$ um campo vetorial áspero. Mostre que $Y$ é suave se, e somente se, para qualquer aberto $U \subset M^n$ e toda $f \in C^\infty(U)$, a função $Yf : M^n \to \mathbb{R}$, definida por $(Yf)(p) = Y_p(f)$ para cada $p \in U$, é suave.
}
\solucao{
	
	\textbf{Análise do Enunciado}
	\begin{itemize}
		\item \textbf{Hipóteses:} $M^n$ é uma variedade suave; $Y$ é um campo vetorial áspero (ou seja, uma seção do fibrado tangente $T M$, onde para cada $p \in M$, $Y_p \in T_p M$).
		\item \textbf{Tese:} Provar a equivalência: $Y$ é uma aplicação suave $M \to TM \iff$ a aplicação $f \mapsto Yf$ preserva a suavidade (transforma funções suaves em funções suaves).
	\end{itemize}
	
	\textbf{Fundamentação Teórica}
	\begin{enumerate}
		\item \textbf{Suavidade de Campos Vetoriais:} Segundo a definição de Lee (Cap. 3), um campo de vetores $Y$ é suave se a aplicação $Y : M \to TM$ for suave. Em coordenadas locais $(U, \phi)$ com $\phi(p) = (x^1, \dots, x^n)$, $Y$ pode ser escrito como:
		\begin{equation*}
			Y_p = \sum_{i=1}^n Y^i(p) \frac{\partial}{\partial x^i} \Big|_p,
		\end{equation*}
		onde as funções $Y^i : U \to \mathbb{R}$ são as componentes do campo. $Y$ é suave se, e somente se, cada função componente $Y^i$ for suave ($C^\infty$).
		\item \textbf{Ação de Vetores Tangentes:} Um vetor $Y_p \in T_p M$ atua como uma derivação sobre funções suaves. Para $f \in C^\infty(U)$, a ação é dada por:
		\begin{equation*}
			(Yf)(p) = Y_p(f) = \sum_{i=1}^n Y^i(p) \frac{\partial f}{\partial x^i}(p).
		\end{equation*}
	\end{enumerate}
	
	\textbf{Desenvolvimento Formal}
	
	\textbf{1. Direção ($\implies$): Suponha que $Y$ seja suave.}
	
	Se $Y$ é suave, então para qualquer carta local $(U, \phi)$, as funções componentes $Y^i : U \to \mathbb{R}$ são suaves. Seja $f \in C^\infty(U)$ uma função suave. A função $(Yf)$ é definida localmente por:
	\begin{equation*}
		(Yf)(p) = \sum_{i=1}^n Y^i(p) \frac{\partial f}{\partial x^i}(p).
	\end{equation*}
	Como $f$ é suave, suas derivadas parciais $\frac{\partial f}{\partial x^i}$ são também funções suaves em $U$. A expressão de $(Yf)(p)$ é uma soma de produtos de funções suaves ($Y^i$ e $\frac{\partial f}{\partial x^i}$). Como o espaço $C^\infty(U)$ é um anel (estável sob soma e produto), a função $(Yf)$ é necessariamente suave em $U$.
	
	\textbf{2. Direção ($\impliedby$): Suponha que $Yf$ seja suave para toda $f \in C^\infty(U)$.}
	
	Para provar que $Y$ é suave, devemos mostrar que suas componentes $Y^i$ são suaves para qualquer carta local $(U, \phi)$. 
	
	Consideremos as funções de coordenada $x^j : U \to \mathbb{R}$ para $j \in \{1, \dots, n\}$. Estas funções são, por definição de variedade suave, funções suaves ($x^j \in C^\infty(U)$). Pela nossa hipótese, a aplicação de $Y$ sobre qualquer função suave resulta em uma função suave. Portanto, a função $(Y x^j)$ deve ser suave para cada $j$.
	
	Calculamos a ação de $Y$ sobre a função de coordenada $x^j$:
	\begin{align*}
		(Y x^j)(p) &= Y_p(x^j) \\
		&= \sum_{i=1}^n Y^i(p) \frac{\partial x^j}{\partial x^i} \Big|_p.
	\end{align*}
	Visto que $\frac{\partial x^j}{\partial x^i} = \delta^j_i$ (o delta de Kronecker), a soma simplifica-se para:
	\begin{equation*}
		(Y x^j)(p) = \sum_{i=1}^n Y^i(p) \delta^j_i = Y^j(p).
	\end{equation*}
	Como assumimos que $(Y f)$ é suave para toda $f \in C^\infty(U)$, e $x^j$ é tal função, concluímos que $Y^j$ é uma função suave para todo $j \in \{1, \dots, n\}$.
	
	Sendo todas as funções componentes $Y^i$ suaves, a aplicação $Y : M \to TM$ é suave por definição.
	
	\textbf{Conclusão:}
	A suavidade de um campo de vetores é equivalente à sua capacidade de mapear funções suaves em funções suaves. Isso prova que a definição geométrica e a definição analítica de campos vetoriais suaves são equivalentes.
}
%%%% QUESTÃO 18 %%%%
\questao{
	Seja $M^n$ uma $n$-variedade suave. Mostre que uma aplicação $\mathcal{J} : C^\infty(M) \to C^\infty(M)$ é uma derivação se, e somente se, $\mathcal{J}f = Yf$ para algum campo de vetores suave $Y \in \mathfrak{X}(M)$.
}
\solucao{
	
	\textbf{Análise do Enunciado}
	\begin{itemize}
		\item \textbf{Hipóteses:} $\mathcal{J}$ é uma aplicação linear de $C^\infty(M)$ em si mesma que satisfaz a regra de Leibniz: $\mathcal{J}(fg) = f\mathcal{J}g + g\mathcal{J}f$ para todo $f, g \in C^\infty(M)$.
		\item \textbf{Tese:} Provar que existe um único campo de vetores suave $Y$ tal que, para toda função $f$, a função $\mathcal{J}f$ é idêntica à função $Yf$.
	\end{itemize}
	
	\textbf{Fundamentação Teórica}
	\begin{enumerate}
		\item \textbf{Derivação em um Ponto:} Um vetor tangente $Y_p \in T_p M$ é definido como uma derivação no ponto $p$.
		\item \textbf{Campos Vetoriais:} Um campo de vetores $Y$ é suave se, para toda $f \in C^\infty(M)$, a função $Yf : p \mapsto Y_p(f)$ é suave (resultado provado na \textbf{Questão 17 desta Lista 02}).
		\item \textbf{Localidade e Lema de Extensão:} Para que $\mathcal{J}$ defina um campo de vetores, a operação $\mathcal{J}f(p)$ deve depender apenas dos valores de $f$ em uma vizinhança de $p$. O \textbf{Lema de Extensão (Questão 27 da Lista 01)} garante que qualquer função suave local pode ser estendida para uma função suave global, permitindo a identificação de $\mathcal{J}$ com a ação de vetores tangentes.
	\end{enumerate}
	
	\textbf{Desenvolvimento Formal}
	
	\textbf{1. Construção do Campo de Vetores $Y$}
	
	Para cada ponto $p \in M$, definimos a aplicação $Y_p : C^\infty(M) \to \mathbb{R}$ por:
	\begin{equation*}
		Y_p(f) = (\mathcal{J}f)(p).
	\end{equation*}
	Verificamos que $Y_p$ é um vetor tangente em $p$ (uma derivação):
	\begin{itemize}
		\item \textbf{Linearidade:} Como $\mathcal{J}$ é linear, $Y_p(af+bg) = (\mathcal{J}(af+bg))(p) = (a\mathcal{J}f + b\mathcal{J}g)(p) = a Y_p(f) + b Y_p(g)$.
		\item \textbf{Regra de Leibniz:} $Y_p(fg) = (\mathcal{J}(fg))(p) = (f\mathcal{J}g + g\mathcal{J}f)(p) = f(p)Y_p(g) + g(p)Y_p(f)$.
	\end{itemize}
	Portanto, $Y_p \in T_p M$ para todo $p \in M$, o que define um campo vetorial áspero $Y$.
	
	\textbf{2. Prova de que $Y$ é Suave}
	
	Pela fundamentação da \textbf{Questão 17}, $Y$ é um campo de vetores suave se, e somente se, para toda $f \in C^\infty(M)$, a função $Yf$ é suave. 
	
	Por definição de nossa construção:
	\begin{equation*}
		(Yf)(p) = Y_p(f) = (\mathcal{J}f)(p).
	\end{equation*}
	Logo, a função $Yf$ é exatamente a função $\mathcal{J}f$. Como, por hipótese, $\mathcal{J}$ mapeia $C^\infty(M)$ em $C^\infty(M)$, a função $\mathcal{J}f$ é suave para qualquer $f$ suave. Portanto, $Y$ é um campo de vetores suave ($Y \in \mathfrak{X}(M)$).
	
	\textbf{3. Unicidade e Direção Inversa}
	
	\begin{itemize}
		\item \textbf{Unicidade:} Se existir outro campo $Z$ tal que $Zf = \mathcal{J}f$, então para todo $p$, $Z_p(f) = (\mathcal{J}f)(p) = Y_p(f)$ para toda $f$. Como as funções suaves separam os pontos do espaço tangente, temos $Z_p = Y_p$ para todo $p$, logo $Z = Y$.
		\item \textbf{Inversa ($\impliedby$):} Seja $Y \in \mathfrak{X}(M)$. Definimos $\mathcal{J}f = Yf$. Como $Y$ é suave, $\mathcal{J}f \in C^\infty(M)$. A linearidade de $\mathcal{J}$ segue da linearidade de $Y_p$ em cada ponto, e a regra de Leibniz segue da definição de vetor tangente. Assim, $\mathcal{J}$ é uma derivação.
	\end{itemize}
	
	\textbf{Conclusão:}
	A correspondência $\mathcal{J} \longleftrightarrow Y$ é um isomorfismo entre o espaço de derivações de $C^\infty(M)$ e o espaço de campos vetoriais suaves $\mathfrak{X}(M)$.
}
%%%% QUESTÃO 19 %%%%
\questao{
	Seja $F : \mathbb{R}^4 \to \mathbb{R}^4$ o campo de vetores suave definido por $F(x_1, x_2, x_3, x_4) = (x_2, -x_1, x_4, -x_3)$. Mostre que $X = F|_{\mathbb{S}^3}$ é um campo de vetores em $\mathbb{S}^3$ sem singularidades.
}
\solucao{
	
	\textbf{Análise do Enunciado}
	\begin{itemize}
		\item \textbf{Hipóteses:} $F$ é uma aplicação suave de $\mathbb{R}^4$ em $\mathbb{R}^4$. $\mathbb{S}^3$ é a esfera unitária em $\mathbb{R}^4$. $X$ é a restrição de $F$ a $\mathbb{S}^3$.
		\item \textbf{Tese 1:} Provar que $X$ é um campo de vetores em $\mathbb{S}^3$ (ou seja, $F(p) \in T_p\mathbb{S}^3$ para todo $p \in \mathbb{S}^3$).
		\item \textbf{Tese 2:} Provar que $X$ não possui singularidades (ou seja, $X(p) \neq 0$ para todo $p \in \mathbb{S}^3$).
	\end{itemize}
	
	\textbf{Fundamentação Teórica}
	\begin{enumerate}
		\item \textbf{Espaço Tangente da Esfera:} Como provamos na \textbf{Questão 3(b) da Lista 01}, para qualquer ponto $p \in \mathbb{S}^n$, o espaço tangente $T_p\mathbb{S}^n$ é o subespaço de $\mathbb{R}^{n+1}$ ortogonal ao vetor posição $p$. Portanto, um vetor $v \in \mathbb{R}^{n+1}$ pertence a $T_p\mathbb{S}^n$ se, e somente se, $\langle p, v \rangle = 0$.
		\item \textbf{Singularidade:} Um ponto $p$ é dito ser uma singularidade de um campo de vetores $X$ se $X(p) = 0$ (o vetor nulo).
	\end{enumerate}
	
	\textbf{Desenvolvimento Formal}
	
	\textbf{1. Prova de que $X$ é um campo de vetores em $\mathbb{S}^3$}
	
	Seja $p = (x_1, x_2, x_3, x_4) \in \mathbb{S}^3$. O vetor $F(p)$ é dado por $(x_2, -x_1, x_4, -x_3)$. Calculamos o produto escalar entre o vetor posição $p$ e o vetor $F(p)$:
	\begin{align*}
		\langle p, F(p) \rangle &= \langle (x_1, x_2, x_3, x_4), (x_2, -x_1, x_4, -x_3) \rangle \\
		&= x_1(x_2) + x_2(-x_1) + x_3(x_4) + x_4(-x_3) \\
		&= x_1x_2 - x_1x_2 + x_3x_4 - x_3x_4 \\
		&= 0.
	\end{align*}
	Como o produto escalar é nulo para todo $p \in \mathbb{S}^3$, concluímos que $F(p) \in T_p\mathbb{S}^3$. Portanto, a restrição $X = F|_{\mathbb{S}^3}$ é, por definição, um campo de vetores em $\mathbb{S}^3$.
	
	\textbf{2. Prova da ausência de singularidades}
	
	Calculamos a norma euclidiana de $X(p)$ para qualquer $p \in \mathbb{S}^3$:
	\begin{align*}
		\|X(p)\|^2 &= \|F(p)\|^2 = (x_2)^2 + (-x_1)^2 + (x_4)^2 + (-x_3)^2 \\
		&= x_1^2 + x_2^2 + x_3^2 + x_4^2.
	\end{align*}
	Visto que $p \in \mathbb{S}^3$, sabemos que $\|p\|^2 = 1$, logo:
	\begin{equation*}
		\|X(p)\|^2 = 1 \implies \|X(p)\| = 1 \neq 0.
	\end{equation*}
	Como a norma de $X(p)$ é constante e igual a $1$ para todo $p \in \mathbb{S}^3$, o campo de vetores $X$ nunca se anula. Portanto, $X$ não possui singularidades.
}
%%%% QUESTÃO 20 %%%%
\questao{
	Seja $F : \mathbb{R}^{2n} \to \mathbb{R}^{2n}$ o campo de vetores suave definido por $$F(x_1, x_2, x_3, x_4, \dots, x_{2n-1}, x_{2n}) = (x_2, -x_1, x_4, -x_3, \dots, x_{2n}, -x_{2n-1}).$$ 
	Mostre que $X = F|_{\mathbb{S}^{2n-1}}$ é um campo de vetores em $\mathbb{S}^{2n-1}$ sem singularidades.
}
\solucao{
	
	\textbf{Análise do Enunciado}
	\begin{itemize}
		\item \textbf{Hipóteses:} $F$ é uma aplicação suave de $\mathbb{R}^{2n}$ em $\mathbb{R}^{2n}$. $\mathbb{S}^{2n-1}$ é a esfera unitária em $\mathbb{R}^{2n}$.
		\item \textbf{Tese:} Provar que a restrição $X = F|_{\mathbb{S}^{2n-1}}$ é um campo de vetores sem singularidades.
	\end{itemize}
	
	\textbf{Fundamentação Teórica}
	Utilizamos os mesmos fundamentos da \textbf{Questão 19}: a condição de tangencialidade $\langle p, X(p) \rangle = 0$ e a verificação de que $\|X(p)\| \neq 0$ para todo $p$.
	
	\textbf{Desenvolvimento Formal}
	
	\textbf{1. Prova de Tangencialidade}
	Seja $p = (x_1, x_2, \dots, x_{2n-1}, x_{2n}) \in \mathbb{S}^{2n-1}$. O vetor $F(p)$ é dado por $$(x_2, -x_1, \dots, x_{2n}, -x_{2n-1}).$$ O produto escalar é:
	\begin{align*}
		\langle p, F(p) \rangle &= \sum_{k=1}^{n} (x_{2k-1} x_{2k} + x_{2k} (-x_{2k-1})) \\
		&= \sum_{k=1}^{n} (x_{2k-1} x_{2k} - x_{2k} x_{2k-1}) \\
		&= \sum_{k=1}^{n} 0 = 0.
	\end{align*}
	Como $\langle p, F(p) \rangle = 0$ para todo $p \in \mathbb{S}^{2n-1}$, $X$ é um campo de vetores em $\mathbb{S}^{2n-1}$.
	
	\textbf{2. Prova da ausência de singularidades}
	A norma de $X(p)$ ao quadrado é:
	\begin{align*}
		\|X(p)\|^2 &= \sum_{k=1}^n (x_{2k}^2 + (-x_{2k-1})^2) \\
		&= \sum_{i=1}^{2n} x_i^2 = \|p\|^2.
	\end{align*}
	Dado que $p \in \mathbb{S}^{2n-1}$, temos $\|p\|^2 = 1$. Logo, $\|X(p)\| = 1$ para todo $p \in \mathbb{S}^{2n-1}$, o que implica que $X$ não possui singularidades.
}
%%%% QUESTÃO 21 %%%%
\questao{
	Seja $F : \mathbb{R}^{n+1} \setminus \{0\} \to \mathbb{R}^{n+1}$ o campo de vetores suave definido por 
	$$F(x_1, \dots, x_{n+1}) = \frac{1}{x_1^2 + \dots + x_{n+1}^2} (-x_1 x_{n+1}, \dots, -x_n x_{n+1}, 1 - x_{n+1}^2).$$
	Mostre que $X = F|_{\mathbb{S}^n}$ é um campo de vetores em $\mathbb{S}^n$ com apenas duas singularidades.
}
\solucao{
	
	\textbf{Análise do Enunciado}
	\begin{itemize}
		\item \textbf{Hipóteses:} $F$ é uma aplicação suave em $\mathbb{R}^{n+1} \setminus \{0\}$. $X$ é a restrição de $F$ à esfera $\mathbb{S}^n$.
		\item \textbf{Tese 1:} Provar que $X$ é um campo de vetores em $\mathbb{S}^n$.
		\item \textbf{Tese 2:} Provar que $X(p) = 0$ se, e somente se, $p$ for um dos dois polos da esfera.
	\end{itemize}
	
	\textbf{Fundamentação Teórica}
	Como nas questões anteriores, utilizaremos a condição $\langle p, X(p) \rangle = 0$ para tangencialidade. Para a análise de singularidades, resolveremos o sistema de equações $F(p) = 0$ sob a restrição $\|p\|^2 = 1$.
	
	\textbf{Desenvolvimento Formal}
	
	\textbf{1. Prova de Tangencialidade}
	Seja $p = (x_1, \dots, x_{n+1}) \in \mathbb{S}^n$. Para este ponto, $\|p\|^2 = \sum_{i=1}^{n+1} x_i^2 = 1$. O vetor $F(p)$ simplifica-se para:
	\begin{equation*}
		F(p) = (-x_1 x_{n+1}, \dots, -x_n x_{n+1}, 1 - x_{n+1}^2).
	\end{equation*}
	Calculamos o produto escalar $\langle p, F(p) \rangle$:
	\begin{align*}
		\langle p, F(p) \rangle &= \left( \sum_{i=1}^n x_i (-x_i x_{n+1}) \right) + x_{n+1}(1 - x_{n+1}^2) \\
		&= -x_{n+1} \left( \sum_{i=1}^n x_i^2 \right) + x_{n+1} - x_{n+1}^3.
	\end{align*}
	Como $\sum_{i=1}^n x_i^2 = 1 - x_{n+1}^2$, substituímos:
	\begin{align*}
		\langle p, F(p) \rangle &= -x_{n+1}(1 - x_{n+1}^2) + x_{n+1} - x_{n+1}^3 \\
		&= -x_{n+1} + x_{n+1}^3 + x_{n+1} - x_{n+1}^3 \\
		&= 0.
	\end{align*}
	Portanto, $X$ é um campo de vetores em $\mathbb{S}^n$.
	
	\textbf{2. Análise das Singularidades}
	
	Um ponto $p \in \mathbb{S}^n$ é uma singularidade se $F(p) = 0$. Isso implica que todas as componentes de $F(p)$ devem ser nulas:
	\begin{enumerate}
		\item $-x_i x_{n+1} = 0$ para todo $i \in \{1, \dots, n\}$.
		\item $1 - x_{n+1}^2 = 0$.
	\end{enumerate}
	Da equação (2), temos que $x_{n+1}^2 = 1$, logo $x_{n+1} = 1$ ou $x_{n+1} = -1$.
	
	\textbf{Caso I: $x_{n+1} = 1$}
	Substituindo na equação (1), temos $-x_i(1) = 0 \implies x_i = 0$ para todo $i \in \{1, \dots, n\}$. O ponto é $p = (0, 0, \dots, 0, 1)$, que é o \textbf{polo norte} $N$.
	
	\textbf{Caso II: $x_{n+1} = -1$}
	Substituindo na equação (1), temos $-x_i(-1) = 0 \implies x_i = 0$ para todo $i \in \{1, \dots, n\}$. O ponto é $p = (0, 0, \dots, 0, -1)$, que é o \textbf{polo sul} $S$.
	
	Como esses são os únicos pontos que satisfazem o sistema, o campo de vetores $X$ possui exatamente duas singularidades: $\{N, S\}$.
}
%%%% QUESTÃO 22 %%%%
\questao{
	Sejam $F : M^n \to N^k$ uma aplicação suave, $V_1, V_2 \in \mathfrak{X}(M)$ e $W_1, W_2 \in \mathfrak{X}(N)$. Suponha que $V_i$ é $F$-relacionado com $W_i$ para $i \in \{1,2\}$. Mostre que $[V_1, V_2]$ é $F$-relacionado a $[W_1, W_2]$.
}
\solucao{
	\textbf{Análise do Enunciado}
	\begin{itemize}
		\item \textbf{Hipóteses:} $F$ é suave. $V_1, V_2$ são campos em $M$; $W_1, W_2$ são campos em $N$. A relação de $F$-relacionamento implica que para todo $p \in M$, $F_{*,p}(V_{1,p}) = W_{1,F(p)}$ e $F_{*,p}(V_{2,p}) = W_{2,F(p)}$.
		\item \textbf{Tese:} Provar que o colchete de Lie $[V_1, V_2]$ é $F$-relacionado ao colchete $[W_1, W_2]$.
	\end{itemize}
	
	\textbf{Fundamentação Teórica}
	\begin{enumerate}
		\item \textbf{Caracterização de Campos $F$-relacionados:} De acordo com Lee (Cap. 9), um campo $V \in \mathfrak{X}(M)$ é $F$-relacionado a $W \in \mathfrak{X}(N)$ se, e somente se, para toda função $g \in C^\infty(N)$, temos a identidade:
		\begin{equation*}
			V(g \circ F) = (Wg) \circ F.
		\end{equation*}
		\item \textbf{Definição do Colchete de Lie:} O colchete de Lie $[V, W]$ é o único campo vetorial que atua sobre funções $f \in C^\infty(M)$ como a comutador de derivações:
		\begin{equation*}
			[V, W]f = V(Wf) - W(Vf).
		\end{equation*}
	\end{enumerate}
	
	\textbf{Desenvolvimento Formal}
	
	Para provar que $[V_1, V_2]$ é $F$-relacionado a $[W_1, W_2]$, devemos mostrar que para qualquer função de teste $g \in C^\infty(N)$, a seguinte igualdade:
	\begin{align*}
		[V_1, V_2](g \circ F) = ([W_1, W_2]g) \circ F.
	\end{align*}
	
	Partindo do lado esquerdo e aplicando a definição de colchete de Lie em $M$:
	\begin{align*}
		[V_1, V_2](g \circ F) &= V_1(V_2(g \circ F)) - V_2(V_1(g \circ F)).
	\end{align*}
	Como $V_2$ é $F$-relacionado a $W_2$, temos $V_2(g \circ F) = (W_2 g) \circ F$. Substituindo:
	\begin{align*}
		[V_1, V_2](g \circ F) &= V_1((W_2 g) \circ F) - V_2((W_1 g) \circ F).
	\end{align*}
	Agora, aplicamos novamente a propriedade de $F$-relacionamento, desta vez para $V_1$ e $W_1$, e para $V_2$ e $W_2$:
	\begin{align*}
		V_1((W_2 g) \circ F) &= (W_1(W_2 g)) \circ F \\
		V_2((W_1 g) \circ F) &= (W_2(W_1 g)) \circ F.
	\end{align*}
	Substituindo estas expressões de volta na equação original:
	\begin{align*}
		[V_1, V_2](g \circ F) &= (W_1(W_2 g)) \circ F - (W_2(W_1 g)) \circ F \\
		&= (W_1(W_2 g) - W_2(W_1 g)) \circ F \\
		&= ([W_1, W_2]g) \circ F.
	\end{align*}
	A igualdade foi provada. Portanto, $[V_1, V_2]$ é $F$-relacionado a $[W_1, W_2]$.
}
%%%% QUESTÃO 23 %%%%
\questao{
	Seja $F : M^n \to N^k$ um difeomorfismo e considere $V_1, V_2 \in \mathfrak{X}(M)$. Mostre que $F_*[V_1, V_2] = [F_*V_1, F_*V_2]$.
}
\solucao{
	\textbf{Análise do Enunciado}
	\begin{itemize}
		\item \textbf{Hipóteses:} $F$ é um difeomorfismo. $V_1, V_2$ são campos suaves em $M$.
		\item \textbf{Tese:} O push-forward do colchete é igual ao colchete dos push-forwards.
	\end{itemize}
	
	\textbf{Fundamentação Teórica}
	\begin{enumerate}
		\item \textbf{Push-forward de Campos via Difeomorfismos:} Se $F$ é um difeomorfismo, para qualquer $V \in \mathfrak{X}(M)$, o campo $F_* V \in \mathfrak{X}(N)$ é definido por $(F_* V)_{F(p)} = F_{*,p}(V_p)$. Por construção, $V$ é $F$-relacionado a $F_* V$.
		\item \textbf{Resultado Anterior:} Na \textbf{Questão 22 desta Lista 02}, provamos que se $V_i$ é $F$-relacionado a $W_i$, então $[V_1, V_2]$ é $F$-relacionado a $[W_1, W_2]$.
	\end{enumerate}
	
	\textbf{Desenvolvimento Formal}
	
	Sejam $W_1 = F_* V_1$ e $W_2 = F_* V_2$. Por definição de push-forward via difeomorfismo, $V_1$ é $F$-relacionado a $W_1$ e $V_2$ é $F$-relacionado a $W_2$.
	
	Aplicando o resultado da \textbf{Questão 22}, concluímos que o campo $[V_1, V_2]$ é $F$-relacionado ao campo $[W_1, W_2]$. 
	
	No entanto, o push-forward do colchete $[V_1, V_2]$ por $F$, denotado por $F_*[V_1, V_2]$, é, por definição, o único campo em $N$ que é $F$-relacionado a $[V_1, V_2]$. Como $[W_1, W_2]$ também é $F$-relacionado a $[V_1, V_2]$, a unicidade do campo $F$-relacionado implica que:
	\begin{equation*}
		F_*[V_1, V_2] = [W_1, W_2].
	\end{equation*}
	Substituindo $W_1$ e $W_2$ pelas suas definições:
	\begin{equation*}
		F_*[V_1, V_2] = [F_*V_1, F_*V_2].
	\end{equation*}
}
%%%% QUESTÃO 24 %%%%
\questao{
	Sejam $M^n$ e $N^k$ variedades suaves e $f : M^n \to N^k$ uma aplicação suave. Defina a aplicação $F : M^n \to M^n \times N^k$ por $F(x) = (x, f(x))$. Mostre que para cada $V \in \mathfrak{X}(M)$, existe um campo vetorial suave em $M^n \times N^k$ que é $F$-relacionado com $V$.
}
\solucao{
	\textbf{Análise do Enunciado}
	\begin{itemize}
		\item \textbf{Hipóteses:} $f: M \to N$ é suave. $F(x) = (x, f(x))$ é a aplicação grafo.
		\item \textbf{Tese:} Para qualquer $V \in \mathfrak{X}(M)$, existe $\mathcal{V} \in \mathfrak{X}(M \times N)$ tal que $\mathcal{V}_{F(p)} = F_{*,p}(V_p)$.
	\end{itemize}
	
	\textbf{Fundamentação Teórica}
	\begin{enumerate}
		\item \textbf{Diferencial de Aplicações Produto:} Se $F(x) = (g(x), h(x))$, então $F_* V = (g_* V, h_* V)$. No nosso caso, $g = \text{Id}_M$ e $h = f$. Logo:
		\begin{equation*}
			F_{*,p}(V_p) = (V_p, f_{*,p}(V_p)) \in T_p M \times T_{f(p)} N.
		\end{equation*}
		\item \textbf{Soma Direta do Espaço Tangente:} Como visto anteriormente, $T_{(p,q)}(M \times N) \cong T_p M \oplus T_q N$. Um campo vetorial em $M \times N$ é suave se suas componentes nos fatores são suaves.
	\end{enumerate}
	
	\textbf{Desenvolvimento Formal}
	
	Queremos construir um campo $\mathcal{V} \in \mathfrak{X}(M \times N)$ tal que $\mathcal{V}_{F(p)} = (V_p, f_{*,p}(V_p))$. 
	
	Definimos o campo $\mathcal{V}$ em $M \times N$ da seguinte forma: para cada ponto $(p, q) \in M \times N$, associamos o vetor:
	\begin{equation*}
		\mathcal{V}_{(p,q)} = (V_p, (f_{*,p} V_p)).
	\end{equation*}
	Note que a segunda componente depende apenas de $p$ e do campo $V$, sendo constante em relação a $q$.
	
	\textbf{1. Prova de Suavidade:}
	Seja $(U, \phi)$ uma carta em $M$ e $(W, \psi)$ uma carta em $N$. O atlas produto fornece coordenadas $(x^i, y^j)$. As componentes de $\mathcal{V}$ em relação a este atlas são:
	\begin{itemize}
		\item Componentes em $M$: $V^i(p)$, que são suaves pois $V \in \mathfrak{X}(M)$.
		\item Componentes em $N$: $\sum_{i=1}^n V^i(p) \frac{\partial f^j}{\partial x^i}(p)$. Como $V^i$ e as derivadas de $f$ são suaves, essas componentes são suaves em $p$ e constantes em $q$.
	\end{itemize}
	Logo, $\mathcal{V}$ é um campo vetorial suave em $M \times N$.
	
	\textbf{2. Verificação do $F$-relacionamento:}
	Para qualquer $p \in M$:
	\begin{align*}
		\mathcal{V}_{F(p)} &= \mathcal{V}_{(p, f(p))} \\
		&= (V_p, f_{*,p}(V_p)) \\
		&= F_{*,p}(V_p).
	\end{align*}
	
	Isso prova que $\mathcal{V}$ é $F$-relacionado com $V$.
}
%%%% QUESTÃO 25 %%%%
\questao{
	Sejam $M_1^{n_1}, \dots, M_k^{n_k}$ variedades suaves. Para cada $i \in \{1, \dots, k\}$, considere a projeção $$\pi_i : M_1 \times \dots \times M_k \to M_i$$ sobre o $i$-ésimo fator. Mostre que para cada $X \in \mathfrak{X}(M_i)$, existe um campo vetorial suave sobre $M_1 \times \dots \times M_k$ que é $\pi_i$-relacionado com $X$.
}
\solucao{
	\textbf{Análise do Enunciado}
	\begin{itemize}
		\item \textbf{Hipóteses:} $M = \prod M_j$ é a variedade produto. $\pi_i$ é a projeção no fator $i$. $X$ é um campo suave em $M_i$.
		\item \textbf{Tese:} Existe $\mathcal{X} \in \mathfrak{X}(M)$ tal que $(\pi_i)_* \mathcal{X} = X \circ \pi_i$.
	\end{itemize}
	
	\textbf{Fundamentação Teórica}
	O espaço tangente de um produto é a soma direta dos espaços tangentes: $$T_{(p_1, \dots, p_k)} M \cong \bigoplus_{j=1}^k T_{p_j} M_j.$$ O diferencial da projeção $\pi_i$ é a projeção linear sobre a $i$-ésima componente:
	\begin{equation*}
		(\pi_i)_*(v_1, \dots, v_k) = v_i.
	\end{equation*}
	
	\textbf{Desenvolvimento Formal}
	
	Definimos o campo vetorial $\mathcal{X}$ em $M$ da seguinte forma: para cada ponto $p = (p_1, \dots, p_k) \in M$, associamos o vetor:
	\begin{equation*}
		\mathcal{X}_p = (0, \dots, 0, X_{p_i}, 0, \dots, 0),
	\end{equation*}
	onde o vetor $X_{p_i} \in T_{p_i} M_i$ ocupa a $i$-ésima posição e todas as outras componentes são o vetor nulo nos respectivos espaços tangentes $T_{p_j} M_j$.
	
	\textbf{1. Suavidade:} 
	Seja $(U_j, \phi_j)$ um atlas para cada $M_j$. O atlas produto $\Phi = (\phi_1, \dots, \phi_k)$ fornece coordenadas para $M$. As componentes de $\mathcal{X}$ em $\Phi$ são nulas para todas as coordenadas exceto as que correspondem ao fator $M_i$, cujas componentes são as de $X$ em $\phi_i$. Como $X$ é suave, $\mathcal{X}$ é suave.
	
	\textbf{2. Verificação do $\pi_i$-relacionamento:}
	Para qualquer $p \in M$, temos:
	\begin{align*}
		(\pi_i)_* \mathcal{X}_p &= (\pi_i)_* (0, \dots, X_{p_i}, \dots, 0) \\
		&= X_{p_i} \\
		&= X_{\pi_i(p)}.
	\end{align*}
	Isso prova que $\mathcal{X}$ é $\pi_i$-relacionado com $X$.
}
%%%% QUESTÃO 26 %%%%
\questao{
	Seja $F : M \to N$ um difeomorfismo local. Mostre que para cada $Y \in \mathfrak{X}(N)$, existe um único campo vetorial suave em $M$ que é $F$-relacionado com $Y$.
}
\solucao{
	\textbf{Análise do Enunciado}
	\begin{itemize}
		\item \textbf{Hipóteses:} $F$ é um difeomorfismo local. $Y$ é um campo suave em $N$.
		\item \textbf{Tese:} Existe um único $X \in \mathfrak{X}(M)$ tal que $F_* X = Y \circ F$.
	\end{itemize}
	
	\textbf{Fundamentação Teórica}
	\begin{enumerate}
		\item \textbf{Difeomorfismo Local:} Para cada $p \in M$, existe aberto $U$ tal que $F|_U : U \to F(U)$ é um difeomorfismo.
		\item \textbf{Isomorfismo do Diferencial:} Como provamos na \textbf{Questão 07 desta Lista 02}, se $F$ é um difeomorfismo local, então $F_{*,p} : T_p M \to T_{F(p)} N$ é um isomorfismo de espaços vetoriais para todo $p \in M$.
	\end{enumerate}
	
	\textbf{Desenvolvimento Formal}
	
	\textbf{1. Existência e Unicidade Pontual}
	Para cada $p \in M$, queremos encontrar $X_p \in T_p M$ tal que $F_{*,p}(X_p) = Y_{F(p)}$. Como $F_{*,p}$ é um isomorfismo, existe um único vetor $X_p$ dado por:
	\begin{equation*}
		X_p = (F_{*,p})^{-1}(Y_{F(p)}).
	\end{equation*}
	Isso define unicamente o campo vetorial áspero $X$.
	
	\textbf{2. Prova de Suavidade e a Matriz Jacobiana}
	Seja $(U, \phi)$ uma carta em $M$ e $(V, \psi)$ uma carta em $N$ tais que $F(U) \subset V$. A representação coordenada de $F$ é $\hat{F} = \psi \circ F \circ \phi^{-1}$. Definimos a \textbf{Matriz Jacobiana} de $F$ no ponto $p$, denotada por $J_F(p)$, como a matriz cujas entradas são as derivadas parciais das componentes de $\hat{F}$:
	\begin{equation*}
		(J_F(p))_{ji} = \frac{\partial \hat{F}^j}{\partial x^i}(\phi(p)), \quad \text{para } i=1, \dots, n \text{ e } j=1, \dots, k.
	\end{equation*}
	A representação coordenada de $X$ é dada por $X^i(p)$. A condição $F_* X = Y \circ F$ traduz-se, em coordenadas, no sistema linear:
	\begin{equation*}
		\sum_{i=1}^n (J_F(p))_{ji} X^i(p) = Y^j(F(p)).
	\end{equation*}
	Em notação matricial, isso é $J_F(p) \cdot \mathbf{X}(p) = \mathbf{Y}(F(p))$. Como $F$ é um difeomorfismo local, $\det(J_F(p)) \neq 0$, logo a matriz é invertível:
	\begin{equation*}
		\begin{bmatrix} X^1(p) \\ \vdots \\ X^n(p) \end{bmatrix} = (J_F(p))^{-1} \begin{bmatrix} Y^1(F(p)) \\ \vdots \\ Y^k(F(p)) \end{bmatrix}.
	\end{equation*}
	Como as funções $\hat{F}^j$ são suaves, as entradas de $J_F(p)$ são suaves. Pela fórmula de Cramer, a inversa de uma matriz cujas entradas são suaves também possui entradas suaves (sendo o quociente do determinante por cofatores). Como $Y^j$ e $F$ são suaves, o produto matricial resulta em funções $X^i$ que são suaves.
	
	\textbf{Conclusão}
	O campo $X$ é suave e é o único campo $F$-relacionado a $Y$.
}
% =============================================================================
% SEÇÃO DE COLCHETES DE LIE COM NOTAÇÃO DE EINSTEIN
% =============================================================================
%%%% RESULTADO EXTRA %%%%
\resultadoextra{Convenção de Notação: Soma de Einstein}{
	Para simplificar as expressões algébricas em cálculos de alta densidade, adotamos a \textbf{Notação de Einstein}. Nesta convenção, sempre que um índice aparece repetido em um único termo (uma vez como sobrescrito e outra como subscrito), assume-se a soma automática sobre a dimensão da variedade. 
	
	Exemplo: $V^i \partial_i f \equiv \sum_{i=1}^n V^i \frac{\partial f}{\partial x^i}$.
}
%%%% QUESTÃO 27 %%%%
\questao{
	Dados os seguintes campos vetoriais $V, W \in \mathfrak{X}(\mathbb{R}^3)$, calcule o colchete de Lie $[V, W]$:
	\begin{enumerate}[label=(\alph*)]
		\item $V = y \frac{\partial}{\partial z} - 2xy^2 \frac{\partial}{\partial y}, \quad W = \frac{\partial}{\partial y}$.
		\item $V = x \frac{\partial}{\partial y} - y \frac{\partial}{\partial x}, \quad W = y \frac{\partial}{\partial z} - z \frac{\partial}{\partial y}$.
		\item $V = x \frac{\partial}{\partial y} - y \frac{\partial}{\partial x}, \quad W = x \frac{\partial}{\partial y} + y \frac{\partial}{\partial x}$.
	\end{enumerate}
}
\solucao{
	\textbf{Análise do Enunciado}
	\begin{itemize}
		\item \textbf{Hipóteses:} $V$ e $W$ são campos vetoriais suaves em $\mathbb{R}^3$.
		\item \textbf{Tese:} Calcular a aplicação $[V, W]$.
	\end{itemize}
	
	\textbf{Fundamentação Teórica}
	Conforme a teoria de campos vetoriais, a suavidade de $[V, W]$ é garantida pois as componentes de $V$ e $W$ são suaves, e o colchete é formado por derivadas parciais dessas componentes (aplicando a lógica de suavidade da \textbf{Questão 26}). 
	
	Utilizando a notação de Einstein, as componentes do colchete $[V, W]^k$ são dadas por:
	\begin{equation*}
		[V, W]^k = V^i \partial_i W^k - W^i \partial_i V^k.
	\end{equation*}
	
	\textbf{Desenvolvimento Formal}
	
	\textbf{(a)} $V = -2xy^2 \partial_y + y \partial_z$ e $W = \partial_y$.
	
	Componentes: $V^x = 0, V^y = -2xy^2, V^z = y$ e $W^x = 0, W^y = 1, W^z = 0$.
	\begin{align*}
		[V, W]^x &= V^i \partial_i W^x - W^i \partial_i V^x = 0 - 0 = 0. \\
		[V, W]^y &= V^i \partial_i W^y - W^i \partial_i V^y = 0 - (1 \cdot \partial_y(-2xy^2)) = -(-4xy) = 4xy. \\
		[V, W]^z &= V^i \partial_i W^z - W^i \partial_i V^z = 0 - (1 \cdot \partial_y(y)) = -1.
	\end{align*}
	Logo: $\boxed{[V, W] = 4xy \frac{\partial}{\partial y} - \frac{\partial}{\partial z}}$.
	
	\textbf{(b)} $V = -y \partial_x + x \partial_y + 0 \partial_z$ e $W = 0 \partial_x - z \partial_y + y \partial_z$.
	
	Componentes: $V^x = -y, V^y = x, V^z = 0$ e $W^x = 0, W^y = -z, W^z = y$.
	\begin{align*}
		[V, W]^x &= V^i \partial_i W^x - W^i \partial_i V^x = 0 - (W^y \partial_y V^x + W^z \partial_z V^x) = -((-z)(-1) + y(0)) = -z. \\
		[V, W]^y &= V^i \partial_i W^y - W^i \partial_i V^y = (V^x \partial_x W^y + V^y \partial_y W^y) - (W^y \partial_y V^y + W^z \partial_z V^y) \\
		&= (-y(0) + x(0)) - (-z(0) + y(0)) = 0. \\
		[V, W]^z &= V^i \partial_i W^z - W^i \partial_i V^z = (V^x \partial_x W^z + V^y \partial_y W^z) - (W^y \partial_y V^z + W^z \partial_z V^z) \\
		&= (-y(0) + x(1)) - (0) = x.
	\end{align*}
	Logo: $\boxed{[V, W] = -z \frac{\partial}{\partial x} + x \frac{\partial}{\partial z}}$.
	
	\textbf{(c)} $V = -y \partial_x + x \partial_y + 0 \partial_z$ e $W = y \partial_x + x \partial_y + 0 \partial_z$.
	
	Componentes: $V^x = -y, V^y = x, V^z = 0$ e $W^x = y, W^y = x, W^z = 0$.
	\begin{align*}
		[V, W]^x &= V^i \partial_i W^x - W^i \partial_i V^x = (V^x \partial_x W^x + V^y \partial_y W^x) - (W^x \partial_x V^x + W^y \partial_y V^x) \\
		&= (-y(0) + x(1)) - (y(0) + x(-1)) = x - (-x) = 2x. \\
		[V, W]^y &= V^i \partial_i W^y - W^i \partial_i V^y = (V^x \partial_x W^y + V^y \partial_y W^y) - (W^x \partial_x V^y + W^y \partial_y V^y) \\
		&= (-y(1) + x(0)) - (y(1) + x(0)) = -y - y = -2y. \\
		[V, W]^z &= 0 - 0 = 0.
	\end{align*}
	Logo: $\boxed{[V, W] = 2x \frac{\partial}{\partial x} - 2y \frac{\partial}{\partial y}}$.
}
	
	% LISTA 03
	\newpage
	\def\listid{L3} % Identificador único para a Lista 3
	\phantomsection
	\addcontentsline{toc}{section}{Soluções da Lista Exercício 3}
	\setcounter{numquestao}{0}
	\begin{center}
	\fbox{\textbf{Resolução da Lista 03}}
\end{center}

\begin{tcolorbox}[
	enhanced,
	sharp corners,
	boxrule=0.5pt,
	colback=white, 
	colframe=black,         
	width=1\textwidth,
	]%
	\textbf{Resumo:} Este documento apresenta a resolução detalhada da terceira lista de exercícios da disciplina de Variedades Diferenciáveis. O conteúdo aborda a teoria de fibrados vetoriais, incluindo a construção de estruturas suaves em fibrados, a análise de seções locais, a estrutura do fibrado cotangente e a teoria de integração de 1-formas em curvas.
\end{tcolorbox}

%%%% QUESTÃO 01 %%%%
\questao{
	Seja $M^{n}$ uma $n$-variedade suave. Suponha que
	\begin{enumerate}[label=(\roman*)]
		\item para cada $p \in M^{n}$, temos associado um $k$-espaço vetorial real $E_p$.
	\end{enumerate}
	Sejam $E = \coprod_{p \in M^n} E_p$ e $\pi : E \to M^n$ uma aplicação que faz corresponder cada elemento de $E_p$ ao ponto $p$. Suponha também que nos é dado:
	\begin{enumerate}[label=(\roman*), resume]
		\item uma cobertura aberta $\{U_\alpha\}_{\alpha \in \Gamma}$ de $M^n$;
		\item para cada $\alpha \in \Gamma$, uma aplicação bijetora $\Phi_\alpha : \pi^{-1}(U_\alpha) \to U_\alpha \times \mathbb{R}^k$ cuja restrição a cada $E_p$ é um isomorfismo linear de $E_p$ em $\{p\} \times \mathbb{R}^k \cong \mathbb{R}^k$;
		\item para cada $\alpha, \beta \in \Gamma$ tal que $U_\alpha \cap U_\beta \neq \emptyset$, uma aplicação suave $\tau_{\alpha\beta} : U_\alpha \cap U_\beta \to M(k, \mathbb{R})$ tal que a composta $\Phi_\alpha \circ \Phi_\beta^{-1} : (U_\alpha \cap U_\beta) \times \mathbb{R}^k \to (U_\alpha \cap U_\beta) \times \mathbb{R}^k$ é dada por $\Phi_\alpha \circ \Phi_\beta^{-1}(p, v) = (p, \tau_{\alpha\beta}(p)v)$.
	\end{enumerate}
	Mostre que $E$ admite uma única estrutura de variedade suave, tornando-se um fibrado vetorial suave de posto $k$ sobre $M^n$, com $\pi$ como projeção e as aplicações $\Phi_\alpha$ como trivializações locais suaves.
}
\solucao{
	
	\textbf{Análise do Enunciado}
	\begin{itemize}
		\item \textbf{Hipóteses:} $M^n$ é uma variedade suave; $E$ é a união disjunta de espaços vetoriais $E_p$; $\pi$ é a projeção canônica; $\Phi_\alpha$ são trivializações locais bijetivas; $\tau_{\alpha\beta}$ são funções de transição suaves com valores na álgebra de matrizes $M(k, \mathbb{R})$.
		\item \textbf{Tese:} Provar que $E$ possui uma única estrutura de variedade suave de dimensão $n+k$ que torna $\pi$ suave e as $\Phi_\alpha$ trivializações suaves.
	\end{itemize}
	
	\textbf{Fundamentação Teórica}
	\begin{enumerate}
		\item \textbf{Lema de Construção de Variedades (Questão 11, Lista 01):} Para dotar um conjunto $E$ de uma estrutura de variedade suave, é suficiente fornecer uma cobertura aberta $\{W_\alpha\}$ e bijeções $\Psi_\alpha : W_\alpha \to V_\alpha \subset \mathbb{R}^m$ tais que as funções de transição $\Psi_\alpha \circ \Psi_\beta^{-1}$ sejam difeomorfismos entre abertos de $\mathbb{R}^m$. 
		\item \textbf{Fibrados Vetoriais (Lee, Cap. 5):} Um fibrado vetorial suave de posto $k$ é uma variedade suave $E$ munida de uma projeção suave $\pi: E \to M$ e trivializações locais $\Phi_\alpha$ que preservam a estrutura linear das fibras.
		\item \textbf{Unicidade da Estrutura (Questão 1, Lista 01):} Qualquer atlas suave está contido em um único atlas suave maximal, o que define unicamente a estrutura suave da variedade.
	\end{enumerate}
	
	\textbf{Desenvolvimento Formal}
	
	\textbf{1. Definição da Estrutura de Variedade Suave via Atlas}
	
	Adotamos a estratégia da \textbf{Questão 11 da Lista 01}. Para cada $\alpha \in \Gamma$, definimos o conjunto $W_\alpha = \pi^{-1}(U_\alpha)$. Visto que $\{U_\alpha\}$ cobre $M^n$, a coleção $\{W_\alpha\}$ cobre $E$. Definimos as cartas coordenadas como as trivializações fornecidas:
	\begin{equation*}
		\Psi_\alpha = \Phi_\alpha : W_\alpha \to U_\alpha \times \mathbb{R}^k.
	\end{equation*}
	Como $U_\alpha$ é um aberto de $M^n$, ele é homeomorfo a um aberto de $\mathbb{R}^n$. Logo, o produto $U_\alpha \times \mathbb{R}^k$ é homeomorfo a um aberto de $\mathbb{R}^n \times \mathbb{R}^k \cong \mathbb{R}^{n+k}$. A coleção $\mathcal{A} = \{(W_\alpha, \Psi_\alpha)\}_{\alpha \in \Gamma}$ constitui, portanto, um atlas topológico candidato para $E$.
	
	\textbf{2. Suavidade das Funções de Transição}
	
	Sejam $\alpha, \beta \in \Gamma$ tais que $U_\alpha \cap U_\beta \neq \emptyset$. A função de transição $\Theta_{\alpha\beta} = \Psi_\alpha \circ \Psi_\beta^{-1}$ atua sobre o domínio $(U_\alpha \cap U_\beta) \times \mathbb{R}^k$. Pela hipótese (iv), temos:
	\begin{equation*}
		\Theta_{\alpha\beta}(p, v) = (p, \tau_{\alpha\beta}(p)v).
	\end{equation*}
	Analisemos a suavidade desta aplicação em componentes:
	\begin{itemize}
		\item A primeira componente é a aplicação identidade no aberto $U_\alpha \cap U_\beta \subset M^n$, que é suave.
		\item A segunda componente é a aplicação $(p, v) \mapsto \tau_{\alpha\beta}(p)v$. Como $\tau_{\alpha\beta} : U_\alpha \cap U_\beta \to M(k, \mathbb{R})$ é suave, cada entrada da matriz $\tau_{\alpha\beta}(p)$ é uma função suave de $p$. O resultado da multiplicação $\tau_{\alpha\beta}(p)v$ é uma combinação linear de $v$ com coeficientes suaves em $p$, o que constitui uma aplicação suave em $(p, v)$.
	\end{itemize}
	Sendo as funções de transição suaves, o atlas $\mathcal{A}$ define, pelo Lema de Construção de Variedades, uma estrutura de variedade suave em $E$ de dimensão $n+k$.
	
	\textbf{3. Verificação da Estrutura de Fibrado Vetorial}
	
	\begin{itemize}
		\item \textbf{Projeção Suave:} A aplicação $\pi : E \to M^n$ em coordenadas locais é $\pi \circ \Psi_\alpha^{-1}(p, v) = p$. Esta é a projeção canônica de $U_\alpha \times \mathbb{R}^k$ sobre $U_\alpha$, que é suave.
		\item \textbf{Linearidade das Fibras:} A hipótese (iii) garante que a restrição $\Phi_\alpha|_{E_p} : E_p \to \{p\} \times \mathbb{R}^k$ é um isomorfismo linear. Portanto, cada fibra $E_p$ é um espaço vetorial de dimensão $k$.
		\item \textbf{Trivializações:} As aplicações $\Phi_\alpha$ são, por construção, as trivializações locais suaves.
	\end{itemize}
	
	\textbf{4. Unicidade}
	
	Se existisse outra estrutura suave $\mathcal{A}'$ tal que as $\Phi_\alpha$ fossem cartas suaves, então $\mathcal{A} \subseteq \mathcal{A}'$. Pela \textbf{Questão 1 da Lista 01}, todo atlas suave está contido em um único atlas suave maximal. Logo, a estrutura suave é única.
}
\solucaoalt{
	
	\textbf{Abordagem Alternativa via Grupos de Lie e Ações Suaves}
	
	Nesta solução, evidenciamos a estrutura intrínseca do fibrado utilizando propriedades de Grupos de Lie e aprofundando o papel analítico da matriz de transição.
	
	\textbf{Fundamentação Teórica}
	\begin{enumerate}
		\item \textbf{Produto de Variedades (Questão 5, Lista 01):} O produto cartesiano $U_\alpha \times \mathbb{R}^k$ herda uma estrutura de variedade suave de dimensão $n+k$.
		\item \textbf{Grupos de Lie (Questão 35, Lista 01):} O grupo linear geral $GL(k, \mathbb{R})$ é um grupo de Lie aberto em $M(k, \mathbb{R})$, cuja ação em $\mathbb{R}^k$ é suave.
		\item \textbf{Atlas Maximal (Questão 1, Lista 01):} A compatibilidade suave de um atlas garante a existência de uma única estrutura suave maximal.
	\end{enumerate}
	
	\textbf{Desenvolvimento Formal}
	
	\textbf{1. O Papel do Grupo de Lie $GL(k, \mathbb{R})$}
	
	Pela hipótese (iii), as aplicações $\Phi_\alpha$ e $\Phi_\beta$ são bijetivas. Deste modo, a aplicação de transição $\Theta_{\alpha\beta} = \Phi_\alpha \circ \Phi_\beta^{-1}$ é necessariamente bijetiva no seu domínio $(U_\alpha \cap U_\beta) \times \mathbb{R}^k$. Sabendo que $\Theta_{\alpha\beta}(p, v) = (p, \tau_{\alpha\beta}(p)v)$, para cada ponto $p$ fixado no domínio, a aplicação puramente linear $v \mapsto \tau_{\alpha\beta}(p)v$ tem de ser um isomorfismo em $\mathbb{R}^k$.
	
	Isto obriga que o determinante $\det(\tau_{\alpha\beta}(p))$ seja não nulo. Logo, a imagem da aplicação $\tau_{\alpha\beta}$ não reside apenas em $M(k, \mathbb{R})$, mas está estritamente contida no grupo linear geral $GL(k, \mathbb{R})$. Como provamos na \textbf{Questão 35 da Lista 01}, $GL(k, \mathbb{R})$ é um grupo de Lie. Em um grupo de Lie, a ação sobre o espaço vetorial subjacente (neste caso, a multiplicação matriz-vetor) é infinitamente diferenciável. Assim, a aplicação $(p, v) \mapsto \tau_{\alpha\beta}(p)v$ é suave por tratar-se da avaliação de uma aplicação suave com valores em um grupo de Lie atuando em $\mathbb{R}^k$.
	
	\textbf{2. Construção do Atlas e Topologia}
	
	Adotamos em $E$ a topologia induzida pelas bijeções $\Phi_\alpha^{-1}$. Formalmente, um subconjunto $W \subset E$ é declarado aberto se $\Phi_\alpha(W \cap \pi^{-1}(U_\alpha))$ é um conjunto aberto na topologia de $U_\alpha \times \mathbb{R}^k$ para todo $\alpha \in \Gamma$.
	
	As cartas de $E$ são compostas por $(W_\alpha, \Psi_\alpha)$, onde $W_\alpha = \pi^{-1}(U_\alpha)$ e $\Psi_\alpha = (\phi_i \times \text{Id}_{\mathbb{R}^k}) \circ \Phi_\alpha$, sendo $\phi_i$ cartas locais da variedade base $M^n$. Baseado na \textbf{Questão 5 da Lista 01}, o espaço de chegada $\mathbb{R}^n \times \mathbb{R}^k$ define rigorosamente a dimensão $n+k$ para o espaço total $E$.
	
	\textbf{3. Suavidade das Transições e Diagrama Lógico}
	
	A viabilidade do atlas reside na comprovação de que $\Theta_{\alpha\beta}(p, v) = (p, \tau_{\alpha\beta}(p)v)$ é um difeomorfismo. Separando em componentes:
	\begin{align*}
		(\Theta_{\alpha\beta})_1(p, v) &= p \quad (\text{Trata-se da identidade suave no aberto de } M^n) \\
		(\Theta_{\alpha\beta})_2(p, v) &= \tau_{\alpha\beta}(p) \cdot v \quad (\text{Ação suave do Grupo de Lie } GL(k, \mathbb{R}))
	\end{align*}
	Sendo ambas as componentes compostas por operações de classe $C^\infty$, $\Theta_{\alpha\beta}$ é um difeomorfismo suave. Visualizamos a geometria destas transições com o diagrama:
	
	\begin{figure}[htbp]
		\centering
		\begin{tikzpicture}[node distance=2.5cm, auto, >=Stealth]
			% Nodos
			\node (E) {$ \pi^{-1}(U_\alpha \cap U_\beta) $};
			\node (P1) [below left of=E, xshift=-1cm] {$ (U_\alpha \cap U_\beta) \times \mathbb{R}^k $};
			\node (P2) [below right of=E, xshift=1cm] {$ (U_\alpha \cap U_\beta) \times \mathbb{R}^k $};
			
			% Setas
			\draw[->, thick] (E) -- node[swap, yshift=2pt] {$\Phi_\beta$} (P1);
			\draw[->, thick] (E) -- node[yshift=2pt] {$\Phi_\alpha$} (P2);
			\draw[->, thick] (P1) -- node[below, yshift=-6pt] {$\Theta_{\alpha\beta}(p,v) = (p, \tau_{\alpha\beta}(p)v)$} (P2);
		\end{tikzpicture}
		\caption{Diagrama de colagem local do Fibrado Vetorial via a matriz de transição do Grupo de Lie.}
	\end{figure}
	
	\textbf{4. Estruturação do Fibrado e Unicidade}
	
	Sob esta estrutura suave estabelecida, a aplicação projeção $\pi$ é suave, pois na restrição dos abertos coincide com a projeção cartesiana $U_\alpha \times \mathbb{R}^k \to U_\alpha$. A integridade da estrutura de fibrado vetorial é asseverada pela restrição linear imposta por $\tau_{\alpha\beta}(p)$. Finalmente, a unicidade da referida estrutura suave em $E$ deriva diretamente da \textbf{Questão 1 da Lista 01}: uma vez que as cartas são suavemente compatíveis, elas pertencem e determinam um único atlas suave maximal para $E$.
}

%%%% QUESTÃO 02 %%%%
\questao{Mostre que a seção zero de qualquer fibrado vetorial suave é suave.
}
\solucao{
	
	\textbf{Análise do Enunciado}
	\begin{itemize}
		\item \textbf{Hipóteses:} Seja $\pi: E \to M^n$ um fibrado vetorial suave de posto $k$ sobre uma variedade suave $M^n$. Definimos a \textbf{seção zero} $\sigma: M^n \to E$ como a aplicação global que associa a cada ponto $p \in M^n$ o elemento neutro (vetor nulo) $0_p$ da sua respectiva fibra vetorial $E_p = \pi^{-1}(p)$.
		\item \textbf{Tese:} Demonstrar que a aplicação $\sigma$ é infinitamente diferenciável (suave) em todo o seu domínio $M^n$.
	\end{itemize}
	
	\textbf{Fundamentação Teórica}
	\begin{enumerate}
		\item \textbf{Suavidade Local (Lee, Cap. 2):} Uma aplicação entre variedades suaves é suave se, e somente se, ela é suave na vizinhança de cada ponto do seu domínio. Isso nos permite verificar a suavidade restringindo a aplicação a abertos coordenados ou trivializantes.
		\item \textbf{Trivialização Local (Lee, Cap. 5):} Por definição de fibrado vetorial suave, para todo ponto $p \in M^n$, existe uma vizinhança aberta $U \subset M^n$ contendo $p$ e um difeomorfismo $\Phi: \pi^{-1}(U) \to U \times \mathbb{R}^k$ que preserva as fibras por isomorfismos lineares.
		\item \textbf{Variedades Produto (Questão 5 da Lista 01):} Como provamos exaustivamente na referida questão (que demonstra que o produto de variedades suaves é uma variedade suave e que as projeções canônicas são suaves), o produto de variedades suaves $U \times \mathbb{R}^k$ é uma variedade suave, e uma aplicação que chega nesse produto é suave se, e somente se, suas componentes projetadas o forem.
		\item \textbf{Unicidade e Atlas Maximal (Questão 1 da Lista 01):} A estrutura suave de $E$ (o espaço total) é aquela determinada pelo atlas suave maximal que contém as trivializações locais, garantindo que $\Phi$ seja, de fato, um difeomorfismo suave.
	\end{enumerate}
	
	\textbf{Desenvolvimento Formal}
	
	A suavidade é uma propriedade eminentemente local. Para provar que a aplicação global $\sigma: M^n \to E$ é suave, é suficiente e necessário demonstrar que, para um ponto arbitrário $p \in M^n$, existe um conjunto aberto $U \ni p$ tal que a restrição $\sigma|_U : U \to \pi^{-1}(U) \subset E$ é uma aplicação suave entre essas variedades.
	
	\textbf{1. Representação na Trivialização Local}
	
	Seja $p \in M^n$ um ponto qualquer e escolha um aberto trivializante $U \ni p$ munido do difeomorfismo de trivialização local garantido pela definição de fibrado:
	\begin{equation*}
		\Phi: \pi^{-1}(U) \to U \times \mathbb{R}^k.
	\end{equation*}
	A condição de estrutura vetorial do fibrado exige que, para cada $q \in U$, a restrição da trivialização à fibra $E_q$, denotada por $\Phi|_{E_q} : E_q \to \{q\} \times \mathbb{R}^k$, atue como um isomorfismo linear entre espaços vetoriais reais de dimensão $k$. 
	
	Sabemos da álgebra linear elementar que todo isomorfismo linear deve mapear, obrigatoriamente, o vetor nulo do domínio no vetor nulo do contradomínio. O vetor nulo do espaço vetorial $E_q$ é, por definição, precisamente a imagem da seção zero no ponto $q$, isto é, $\sigma(q) = 0_q$. Por outro lado, o vetor nulo do espaço vetorial produto $\{q\} \times \mathbb{R}^k$ é o par ordenado $(q, 0)$, onde $0 \in \mathbb{R}^k$ representa o vetor cujas $k$ coordenadas são simultaneamente nulas. Consequentemente, obtemos a seguinte identidade rigorosa para todo $q \in U$:
	\begin{equation*}
		\Phi(\sigma(q)) = (q, 0).
	\end{equation*}
	
	\textbf{2. Análise da Composição Local no Produto de Variedades}
	
	Consideremos agora a aplicação composta $\hat{\sigma} = \Phi \circ \sigma|_U : U \to U \times \mathbb{R}^k$. De acordo com a estrutura de variedade produto validada na \textbf{Questão 5 da Lista 01}, a aplicação $\hat{\sigma}$ pode ser decomposta através de projeções cartesianas em suas duas componentes fundamentais:
	\begin{align*}
		\pi_1 \circ \hat{\sigma}(q) &= q \\
		\pi_2 \circ \hat{\sigma}(q) &= 0
	\end{align*}
	Analisemos a suavidade de cada uma independentemente:
	\begin{itemize}
		\item A primeira componente, $q \mapsto q$, é a aplicação identidade do aberto $U$ nele mesmo ($\text{Id}_U$). A aplicação identidade em qualquer variedade suave é, por definição fundamental (Lee, Cap. 2), infinitamente diferenciável.
		\item A segunda componente, $q \mapsto 0$, é uma aplicação constante nula de $U$ no espaço euclidiano $\mathbb{R}^k$. Toda aplicação constante entre variedades é infinitamente diferenciável, visto que suas derivadas em cartas locais são identicamente nulas de todas as ordens.
	\end{itemize}
	Como ambas as funções componentes de $\hat{\sigma}$ são comprovadamente suaves, segue do Lema das Funções em Produtos (novamente ancorado em nossos resultados da \textbf{Questão 5 da Lista 01}) que a aplicação produto $\hat{\sigma}$ é uma aplicação globalmente suave de $U$ no produto $U \times \mathbb{R}^k$.
	
	\begin{figure}[htbp]
		\centering
		\begin{tikzpicture}[scale=0.7][node distance=3cm, auto, >=Stealth]
			% Nodos
			\node (U) {$U \subset M^n$};
			\node (E) [right of=U, xshift=2cm] {$\pi^{-1}(U) \subset E$};
			\node (R) [below right of=U, xshift=1cm, yshift=-1.5cm] {$U \times \mathbb{R}^k$};
			
			% Setas
			\draw[->, thick] (U) -- node[above] {$\sigma|_U$} (E);
			\draw[->, thick] (E) -- node[right, xshift=3pt] {$\Phi$} (R);
			\draw[->, thick] (U) -- node[left, swap, yshift=-4pt] {$\hat{\sigma}(q) = (q, 0)$} (R);
		\end{tikzpicture}
		\caption{Diagrama comutativo da representação local da seção zero $\sigma$ no fibrado trivializado.}
	\end{figure}
	
	\textbf{3. Conclusão da Suavidade Global da Seção Zero}
	
	Temos estabelecida a relação de composição $\Phi \circ \sigma|_U = \hat{\sigma}$. Uma vez que a trivialização local $\Phi$ é um difeomorfismo entre $\pi^{-1}(U)$ e $U \times \mathbb{R}^k$ (conforme a definição de fibrado vetorial no Capítulo 5 do Lee), ela possui uma inversa $\Phi^{-1} : U \times \mathbb{R}^k \to \pi^{-1}(U)$ que é estritamente suave. 
	
	Podemos então isolar a restrição da seção zero compondo ambos os lados à esquerda por esta inversa suave:
	\begin{equation*}
		\sigma|_U = \Phi^{-1} \circ \hat{\sigma}.
	\end{equation*}
	Pelo Teorema da Composição de Funções Suaves (Lee, Cap. 2), a composição de duas aplicações suaves ($\Phi^{-1}$ e $\hat{\sigma}$) resulta necessariamente em uma aplicação suave. Assim, demonstramos conclusivamente que $\sigma|_U$ é suave. 
	
	Visto que o ponto $p$ foi tomado arbitrariamente em $M^n$, a aplicação $\sigma$ é localmente suave na vizinhança de absolutamente todos os pontos de $M^n$. Logo, a seção zero $\sigma: M^n \to E$ é globalmente suave em todo o fibrado.
}

%%%% QUESTÃO 03 %%%%	
\questao{Seja $\pi:E\rightarrow M^{n}$ um fibrado vetorial suave de posto k. Mostre que todo referencial local suave $(\sigma_{1},...,\sigma_{k})$ para E em U é associado com alguma trivialização local $\Phi:\pi^{-1}(U)\rightarrow U\times\mathbb{R}^{k}$.
}
\solucao{
	
	\textbf{Análise do Enunciado}
	\begin{itemize}
		\item \textbf{Hipóteses:} $\pi: E \to M^n$ é um fibrado vetorial suave de posto $k$. Temos um conjunto aberto $U \subset M^n$ e uma coleção de $k$ seções locais suaves $\sigma_1, \dots, \sigma_k : U \to \pi^{-1}(U)$ que formam um referencial local. Isso significa que, para todo ponto $p \in U$, o conjunto de vetores $\{\sigma_1(p), \dots, \sigma_k(p)\}$ constitui uma base para a fibra vetorial $E_p = \pi^{-1}(p)$.
		\item \textbf{Tese:} Demonstrar a existência de uma trivialização local suave $\Phi: \pi^{-1}(U) \to U \times \mathbb{R}^k$ que corresponda a este referencial.
	\end{itemize}
	
	\textbf{Desenvolvimento Formal}
	
	A demonstração consiste em construir uma parametrização local $\Psi$ a partir do referencial dado, provar que ela é bijetora e suave e, em seguida, mostrar que a sua inversa $\Phi = \Psi^{-1}$ é a trivialização suave procurada.
	
	\textbf{1. Construção da Parametrização Trivializante $\Psi$}
	
	Definimos uma aplicação $\Psi: U \times \mathbb{R}^k \to \pi^{-1}(U)$. Para qualquer par $(p, v) \in U \times \mathbb{R}^k$, onde o vetor $v$ é dado por suas coordenadas $v = (v^1, \dots, v^k) \in \mathbb{R}^k$, definimos $\Psi$ utilizando a \textbf{Notação de Einstein}:
	\begin{equation*}
		\Psi(p, v) = v^i \sigma_i(p).
	\end{equation*}
	
	Para verificar a suavidade de $\Psi$, tomamos um ponto arbitrário $p_0 \in U$. Como $\pi: E \to M^n$ é um fibrado vetorial suave, existe por definição um aberto $W \subset U$ contendo $p_0$ e uma trivialização local suave $\Theta: \pi^{-1}(W) \to W \times \mathbb{R}^k$. Na trivialização $\Theta$, cada seção suave $\sigma_i$ restrita a $W$ é expressa como $\Theta(\sigma_i(p)) = (p, s_i(p))$, onde $s_i: W \to \mathbb{R}^k$ são aplicações suaves. 
	
	Avaliando $\Psi$ nesta trivialização local e usando a linearidade nas fibras:
	\begin{equation*}
		\Theta \circ \Psi(p, v) = \Theta( v^i \sigma_i(p) ) = ( p, v^i s_i(p) ).
	\end{equation*}
	Como as componentes $v^i$ são projeções suaves e as aplicações $s_i(p)$ são suaves, a combinação linear acima é suave em $W \times \mathbb{R}^k$. Logo, $\Psi$ é suave.
	
	\textbf{2. Bijectividade e Definição de $\Phi$}
	
	Para cada $p \in U$, a restrição $\Psi_p : \mathbb{R}^k \to E_p$ é dada por $\Psi_p(v) = v^i \sigma_i(p)$. Como $\{\sigma_i(p)\}$ é uma base de $E_p$, $\Psi_p$ é um isomorfismo linear. Assim, $\Psi$ é uma bijeção e definimos a trivialização como:
	\begin{equation*}
		\Phi = \Psi^{-1} : \pi^{-1}(U) \to U \times \mathbb{R}^k.
	\end{equation*}
	
	\textbf{3. Demonstração da Suavidade da Inversa $\Phi$}
	
	Para provar que $\Phi$ é suave, analisamos sua representação local em $W$. Definimos a matriz $A(p)$ cujas colunas são os vetores $s_1(p), \dots, s_k(p)$. A componente $j$ do vetor $v^i s_i(p)$ é escrita em notação de Einstein como $A^j_i(p) v^i$. Assim, a representação local de $\Phi$ é $(p, w) \mapsto (p, [A(p)]^{-1} w)$. Apresentamos agora duas formas de concluir a suavidade desta inversão:
	
	\textit{\textbf{Via I: Abordagem por Grupos de Lie (Elegância Estrutural)}}
	\solucaoalt{
		A suavidade de $\Phi$ é garantida pois a representação local da aplicação $\Psi$ em uma trivialização pré-existente $\Theta$ é dada por $\Theta \circ \Psi(p, v) = (p, A(p)v)$. Conforme demonstrado na \textbf{Questão 35 da Lista 01}, o conjunto $GL(k, \mathbb{R})$ das matrizes inversíveis é um \textbf{Grupo de Lie}. Como tal, a aplicação de inversão $\text{Inv}: GL(k, \mathbb{R}) \to GL(k, \mathbb{R})$, dada por $A \mapsto A^{-1}$, é comprovadamente suave. 
		
		Visto que as seções $\sigma_i$ são suaves, a aplicação $A: W \to GL(k, \mathbb{R})$ é suave. Portanto, a representação de $\Phi$ via $\Theta$, dada por $(p, w) \mapsto (p, A(p)^{-1}w)$, é a composição de aplicações suaves entre variedades. Esta estrutura assegura que a inversão matricial preserva a suavidade, tornando $\Phi$ um difeomorfismo local.
	}
		
	\textit{\textbf{Via II: Abordagem Analítica (Regra de Cramer)}}
	Pela Regra de Cramer, as entradas da matriz inversa são dadas por:
	\begin{equation*}
		(A^{-1})^j_i(p) = \frac{1}{\det(A(p))} \text{Adj}(A(p))^j_i.
	\end{equation*}
	O determinante e a matriz adjunta são polinômios nas entradas $A^a_b(p)$. Como as entradas são suaves e $\det(A(p)) \neq 0$, o quociente é suave. O produto $[A(p)]^{-1} w$ em notação de Einstein é $(A^{-1})^j_i(p) w^i$, que é uma soma de produtos de funções suaves, logo, suave.
	
	\textbf{4. Conclusão}
	
	Em ambas as abordagens, $\Phi: \pi^{-1}(U) \to U \times \mathbb{R}^k$ é um difeomorfismo suave que preserva as fibras e age como isomorfismo linear nelas. Portanto, $\Phi$ é a trivialização local associada ao referencial $(\sigma_i)$.
}

%%%% QUESTÃO 04 %%%%
\questao{Seja $M^n$ uma $n$-variedade suave. Então a estrutura suave já construída no fibrado tangente $TM$ é a única com respeito a qual $\pi : TM \to M^n$ é um fibrado vetorial suave de posto $n$ sobre $M^n$ e cujos campos vetoriais coordenados são seções locais suaves.}

\solucao{
	
	\textbf{Análise do Enunciado:}
	\begin{itemize}
		\item \textbf{Hipóteses:} $M^n$ é uma $n$-variedade suave; $TM$ é o seu fibrado tangente; $\pi: TM \to M^n$ é a projeção natural.
		\item \textbf{Tese:} A estrutura suave de fibrado vetorial em $TM$ é única sob a condição de que os campos coordenados $\tfrac{\partial}{\partial x^i}$ sejam seções locais suaves.
	\end{itemize}
	
	\textbf{Fundamentação Teórica}
	\begin{enumerate}
		\item \textbf{Referenciais e Trivializações (Questão 3, Lista 03):} Estabelecemos que qualquer referencial local suave $(\sigma_1, \dots, \sigma_k)$ em um fibrado vetorial induz univocamente uma trivialização local suave $\Phi: \pi^{-1}(U) \to U \times \mathbb{R}^k$.
		\item \textbf{Unicidade de Atlas (Questão 1, Lista 01):} A estrutura suave de uma variedade é univocamente determinada por seu atlas maximal. Dois atlas suavemente compatíveis definem a mesma estrutura.
	\end{enumerate}
	
	\textbf{Desenvolvimento Formal}
	
	Seja $\mathcal{E}$ qualquer estrutura de variedade suave em $TM$ que satisfaça as condições do enunciado. Consideremos uma carta coordenada $(U, (x^1, \dots, x^n))$ em $M^n$. 
	
	Por definição de espaço tangente, os campos vetoriais coordenados $\left\{ \tfrac{\partial}{\partial x^1}, \dots, \tfrac{\partial}{\partial x^n} \right\}$ formam uma base de $T_p M$ para todo $p \in U$. Pela hipótese do enunciado, esses campos são seções locais suaves sob a estrutura $\mathcal{E}$. Portanto, eles constituem, rigorosamente, um \textbf{referencial local suave} para o fibrado vetorial $(TM, \mathcal{E})$ sobre o aberto $U$.
	
	Aplicando o resultado da \textbf{Questão 3 desta Lista}, a existência deste referencial local suave implica a existência de uma única trivialização local suave $\Phi : \pi^{-1}(U) \to U \times \mathbb{R}^n$ associada a ele, tal que:
	\begin{equation*}
		\Phi\left( \sum_{i=1}^n v^i \tfrac{\partial}{\partial x^i} \right) = (p, (v^1, \dots, v^n)).
	\end{equation*}
	Esta aplicação $\Phi$ é precisamente a trivialização canônica utilizada na construção da estrutura suave padrão de $TM$ (conforme detalhado na Questão 13 da Lista 02 e reafirmado na Questão 1 desta Lista). 
	
	Sendo assim, para qualquer carta de $M^n$, a estrutura $\mathcal{E}$ deve admitir as mesmas trivializações locais que a estrutura canônica. Isso implica que o atlas de trivializações de $\mathcal{E}$ é suavemente compatível com o atlas canônico de $TM$, pois as funções de transição serão idênticas.
	
	Pela \textbf{Questão 1 da Lista 01}, dois atlas suavemente compatíveis pertencem ao mesmo atlas maximal. Portanto, a estrutura suave $\mathcal{E}$ deve coincidir com a estrutura suave canônica de $TM$. Esta prova demonstra a rigidez da estrutura do fibrado tangente: ela é inteiramente determinada pela suavidade dos campos coordenados da base.
}

%%%% QUESTÃO 05 %%%%
\questao{
	Seja $\pi : E \to M$ um fibrado vetorial suave de posto $k$ e suponha que $\sigma_{1}, \dots, \sigma_{m}$ sejam seções locais suaves independentes sobre um conjunto aberto $U \subset M$. Mostre que, para cada $p \in U$, existem seções suaves $\sigma_{m+1}, \dots, \sigma_{k}$ definidas em alguma vizinhança $V$ de $p$ tal que $(\sigma_{1}, \dots, \sigma_{k})$ é um referencial local suave para $E$ em $U \cap V$.
}
\solucao{
	
	\textbf{Análise do Enunciado}
	\begin{itemize}
		\item \textbf{Hipóteses:} $\pi: E \to M$ é um fibrado vetorial suave de posto $k$. As seções locais $\sigma_1, \dots, \sigma_m \in \Gamma(U, E)$ são linearmente independentes em cada fibra $E_q$ para todo $q \in U$.
		\item \textbf{Tese:} Para cada ponto $p \in U$, provar a existência de seções suaves $\sigma_{m+1}, \dots, \sigma_k$ em uma vizinhança $V \ni p$ tais que o conjunto total $\{\sigma_1, \dots, \sigma_k\}$ constitua um referencial local suave (base suave) para $E$ sobre $U \cap V$.
	\end{itemize}
	
	\textbf{Fundamentação Teórica}
	\begin{enumerate}
		\item \textbf{Completamento de Base (Álgebra Linear):} Dado que $E_p$ é um espaço vetorial de dimensão $k$, qualquer conjunto de $m < k$ vetores linearmente independentes pode ser estendido a uma base de $E_p$.
		\item \textbf{Estrutura de Fibrados (Lee, Cap. 5):} A existência de trivializações locais garante que podemos expressar qualquer seção suave como uma combinação linear de seções de um referencial local pré-existente com coeficientes suaves.
		\item \textbf{Estabilidade da Independência Linear:} A condição de que um conjunto de vetores seja linearmente independente é equivalente a a não-anulação do determinante de sua matriz de representação. Como o determinante é uma função contínua, a independência linear é uma propriedade aberta.
	\end{enumerate}
	
	\textbf{Desenvolvimento Formal}
	
	\textbf{1. Extensão Pontual em $p$}
	Seja $p \in U$. Por hipótese, o conjunto $\{\sigma_1(p), \dots, \sigma_m(p)\}$ é linearmente independente no espaço vetorial $E_p$. Pelo Teorema do Completamento de Base, existem vetores $v_{m+1}, \dots, v_k \in E_p$ tais que o conjunto $\{\sigma_1(p), \dots, \sigma_m(p), v_{m+1}, \dots, v_k\}$ seja uma base de $E_p$.
	
	\textbf{2. Construção de Seções Suaves Globais}
	Para propagar esses vetores $v_j$ para a vizinhança de $p$, utilizamos a estrutura de fibrado. Seja $(\tau_1, \dots, \tau_k)$ um referencial local suave definido em um aberto $W \ni p$. Cada vetor $v_j$ pode ser escrito univocamente como:
	\begin{equation*}
		v_j = \sum_{i=1}^k C_j^i \tau_i(p), \quad \text{com } C_j^i \in \mathbb{R}.
	\end{equation*}
	Definimos as seções $\sigma_j : W \to E$ para $j = m+1, \dots, k$ por:
	\begin{equation*}
		\sigma_j(q) = \sum_{i=1}^k C_j^i \tau_i(q), \quad \forall q \in W.
	\end{equation*}
	Visto que cada $\tau_i$ é uma seção suave e os $C_j^i$ são constantes, as seções $\sigma_j$ são, por definição, suaves em $W$. Note que, por construção, $\sigma_j(p) = v_j$.
	
	\textbf{3. Verificação da Independência Linear na Vizinhança}
	Seja $\mathcal{S}(q) = (\sigma_1(q), \dots, \sigma_k(q))$ o conjunto de seções em $U \cap W$. Para analisar a independência linear, expressamos cada $\sigma_\alpha$ em termos do referencial $\tau_\beta$:
	\begin{equation*}
		\sigma_\alpha(q) = A^\beta_\alpha(q) \tau_\beta(q), \quad \alpha, \beta \in \{1, \dots, k\}.
	\end{equation*}
	As funções $A^\beta_\alpha : U \cap W \to \mathbb{R}$ são as componentes coordenadas das seções $\sigma_\alpha$, logo são suaves. A matriz $[A^\beta_\alpha(q)]$ é a matriz de transição entre os referenciais. 
	
	No ponto $p$, o conjunto $\{\sigma_\alpha(p)\}$ é uma base de $E_p$, o que implica que $\det[A^\beta_\alpha(p)] \neq 0$. Sendo a função $\det[A^\beta_\alpha(\cdot)]$ contínua, existe uma vizinhança aberta $V \subset U \cap W$ de $p$ tal que $\det[A^\beta_\alpha(q)] \neq 0$ para todo $q \in V$. 
	
	Nesse aberto $V$, o determinante não nulo garante que as seções $\{\sigma_1(q), \dots, \sigma_k(q)\}$ permanecem linearmente independentes em cada fibra $E_q$. Como são $k$ vetores independentes em um espaço de dimensão $k$, eles formam um referencial local suave para $E$ em $V$.
}
\solucaoalt{
	\textbf{Demonstração via Topologia do Grupo de Lie $GL(k, \mathbb{R})$}
	
	A matriz de transição define uma aplicação suave $\mathcal{A}: U \cap W \to M(k, \mathbb{R})$ dada por $q \mapsto [A^\beta_\alpha(q)]$. O conjunto das matrizes inversíveis $GL(k, \mathbb{R})$ é um Grupo de Lie e, portanto, um subconjunto aberto de $M(k, \mathbb{R})$.
	
	Como $\{\sigma_\alpha(p)\}$ é base, temos $\mathcal{A}(p) \in GL(k, \mathbb{R})$. Pela continuidade de $\mathcal{A}$, a pré-imagem $\mathcal{A}^{-1}(GL(k, \mathbb{R}))$ é um aberto em $U \cap W$. Definindo $V$ como esse aberto, garantimos que $\mathcal{A}(q)$ é invertível para todo $q \in V$, o que assegura que as seções formam um referencial local suave.
}

%%% DEFINIÇÃO %%%
Antes de procedermos à resolução da Questão, convém notar que a análise de certos problemas geométricos pode ser significativamente enriquecida se abordada sob uma perspectiva estrutural. Para viabilizar as \textbf{soluções alternativas} que apresentaremos a seguir, e para evitar a redundância de definições algébricas no corpo das provas, estabelecemos previamente o seguinte resultado fundamental. Este servirá como o suporte teórico necessário para que a transição entre a álgebra e a geometria ocorra de forma natural e rigorosa.

%%%% RESULTADO EXTRA %%%%
\resultadoextra{Álgebra dos Quatérnios ($\mathbb{H}$)}{
	O corpo dos quatérnios $\mathbb{H}$ é uma álgebra de divisão sobre os reais com 4 dimensões. Um quatérnio $q \in \mathbb{H}$ é escrito como:
	\begin{equation*}
		q = a + b\mathbf{i} + c\mathbf{j} + d\mathbf{k}, \quad a, b, c, d \in \mathbb{R}
	\end{equation*}
	
	As unidades imaginárias $\mathbf{i}$, $\mathbf{j}$, $\mathbf{k}$ obedecem às regras fundamentais:
	\begin{equation*}
		\mathbf{i}^2 = \mathbf{j}^2 = \mathbf{k}^2 = \mathbf{i}\mathbf{j}\mathbf{k} = -1
	\end{equation*}
	Desta definição, derivam-se as relações de anticomutatividade: $\mathbf{ij} = \mathbf{k}, \mathbf{jk} = \mathbf{i}, \mathbf{ki} = \mathbf{j}$ e $\mathbf{ji} = -\mathbf{k}, \mathbf{kj} = -\mathbf{i}, \mathbf{ik} = -\mathbf{j}$.
	
	A norma de um quatérnio é definida por $|q|^2 = a^2 + b^2 + c^2 + d^2$. Uma propriedade fundamental é que a norma é multiplicativa: $|q \cdot p| = |q| \cdot |p|$ para quaisquer $q, p \in \mathbb{H}$.
	
	Consequentemente, o conjunto de quatérnios unitários $\mathbb{S}^3 = \{q \in \mathbb{H} : |q| = 1\}$ forma um grupo sob a multiplicação de quatérnios, sendo, portanto, um Grupo de Lie de dimensão 3.
}

%%%% QUESTÃO 06 %%%%
\questao{
	Considere os seguintes campos vetoriais em $\mathbb{R}^{4}$:
	$$X_{1}=-x^{2}t\frac{\partial}{\partial x^{1}}+x^{1}\tfrac{\partial}{\partial x^{2}}+x^{4}\tfrac{\partial}{\partial x^{3}}-x^{3}\tfrac{\partial}{\partial x^{4}}$$
	$$X_{2}=-x^{3}\tfrac{\partial}{\partial x^{1}}-x^{4}\tfrac{\partial}{\partial x^{2}}+x^{1}\tfrac{\partial}{\partial x^{3}}+x^{2}\tfrac{\partial}{\partial x^{4}}$$
	$$X_{3}=-x^{4}\tfrac{\partial}{\partial x^{1}}+x^{3}\tfrac{\partial}{\partial x^{2}}-x^{2}\tfrac{\partial}{\partial x^{3}}+x^{1}\tfrac{\partial}{\partial x^{4}}$$
	Mostre que existem campos vetoriais suaves $V_{1}, V_{2}, V_{3}$ em $\mathbb{S}^{3}$ tais que $V_{j}$ é $\iota$-relacionado a $X_{j}$ para $j=1,2,3,$ onde $\iota:\mathbb{S}^{3}\rightarrow\mathbb{R}^{4}$ é a aplicação inclusão.
	Conclua que $\mathbb{S}^{3}$ é paralelizável.
}
\solucao{
	
	\textbf{Análise do Enunciado}
	\begin{itemize}
		\item \textbf{Hipóteses:} $X_1, X_2, X_3 \in \mathfrak{X}(\mathbb{R}^4)$. $\mathbb{S}^3 \subset \mathbb{R}^4$ é a subvariedade mergulhada definida pela inclusão $\iota$.
		\item \textbf{Tese 1:} Provar que $X_j$ são tangentes a $\mathbb{S}^3$, garantindo a existência de $V_j \in \mathfrak{X}(\mathbb{S}^3)$.
		\item \textbf{Tese 2:} Provar que $\{V_1, V_2, V_3\}$ formam um referencial global suave para $\mathbb{S}^3$.
	\end{itemize}
	
	\textbf{Fundamentação Teórica}
	\begin{enumerate}
		\item \textbf{Tangencialidade:} Um campo $X \in \mathfrak{X}(\mathbb{R}^4)$ é tangente a $\mathbb{S}^3$ se, para todo $p \in \mathbb{S}^3$, $\langle X(p), p \rangle = 0$.
		\item \textbf{Paralelizabilidade:} $\mathbb{S}^3$ é paralelizável se admite 3 campos suaves globais que são linearmente independentes em cada ponto.
	\end{enumerate}
	
	\textbf{Desenvolvimento Formal}
	
	\textbf{1. Verificação da Tangencialidade}
	Seja $p = (x^1, x^2, x^3, x^4) \in \mathbb{S}^3$. O produto escalar entre $p$ e os campos $X_j$ é:
	\begin{align*}
		\langle X_1, p \rangle &= (-x^2)(x^1) + (x^1)(x^2) + (x^4)(x^3) + (-x^3)(x^4) = 0. \\
		\langle X_2, p \rangle &= (-x^3)(x^1) + (-x^4)(x^2) + (x^1)(x^3) + (x^2)(x^4) = 0. \\
		\langle X_3, p \rangle &= (-x^4)(x^1) + (x^3)(x^2) + (-x^2)(x^3) + (x^1)(x^4) = 0.
	\end{align*}
	Como $\langle X_j, p \rangle = 0$ para todo $j \in \{1, 2, 3\}$, os campos $X_j$ são tangentes a $\mathbb{S}^3$. Logo, existem campos $V_j \in \mathfrak{X}(\mathbb{S}^3)$ tais que $\iota_* V_j = X_j$.
	
	\textbf{2. Independência Linear Global}
	Calculamos a ortogonalidade dos campos no ambiente $\mathbb{R}^4$:
	\begin{align*}
		\langle X_1, X_2 \rangle &= (-x^2)(-x^3) + (x^1)(-x^4) + (x^4)(x^1) + (-x^3)(x^2) = 0. \\
		\langle X_1, X_3 \rangle &= (-x^2)(-x^4) + (x^1)(x^3) + (x^4)(-x^2) + (-x^3)(x^1) = 0. \\
		\langle X_2, X_3 \rangle &= (-x^3)(-x^4) + (-x^4)(x^3) + (x^1)(-x^2) + (x^2)(x^1) = 0.
	\end{align*}
	As normas são $\|X_1\|^2 = \|X_2\|^2 = \|X_3\|^2 = \sum (x^i)^2 = 1$ para $p \in \mathbb{S}^3$.
	
	Visto que $\{V_1(p), V_2(p), V_3(p)\}$ formam um conjunto ortonormal em $T_p\mathbb{S}^3$ para todo $p$, eles são linearmente independentes. Como $\dim(\mathbb{S}^3) = 3$, eles constituem um referencial global suave. Portanto, $\mathbb{S}^3$ é paralelizável.
}

\solucaoalt{
	\textbf{Abordagem via Estrutura de Grupo de Lie ($\mathbb{H}$)}
	
	Utilizando a estrutura de Quatérnios definida no \textbf{Resultado Extra anterior}, identificamos $\mathbb{S}^3$ como o Grupo de Lie dos quatérnios unitários.
	
	Os campos $X_j$ são gerados pela translação à direita de $x \in \mathbb{S}^3$ pelas unidades imaginárias $\{\mathbf{i}, \mathbf{j}, \mathbf{k}\}$ da base de $\mathbb{H}$:
	\begin{align*}
		X_1(x) &= x \cdot \mathbf{i} = (-x^2, x^1, x^4, -x^3) \\
		X_2(x) &= x \cdot \mathbf{j} = (-x^3, -x^4, x^1, x^2) \\
		X_3(x) &= x \cdot \mathbf{k} = (-x^4, x^3, -x^2, x^1)
	\end{align*}
	Sendo a multiplicação de quatérnios uma operação suave, os campos $X_j$ são suaves. Como a multiplicação por unidades imaginárias preserva a norma ($|x \cdot \mathbf{i}| = |x| \cdot 1 = 1$), os vetores $X_j(x)$ são tangentes a $\mathbb{S}^3$ (pois a derivada de $\|x\|^2=1$ em qualquer direção de translação de grupo é nula).
	
	Como $\{\mathbf{i}, \mathbf{j}, \mathbf{k}\}$ é uma base para o espaço tangente à identidade $T_1\mathbb{S}^3$, e a translação à direita em um Grupo de Lie é um difeomorfismo que preserva a independência linear, os campos $\{V_1, V_2, V_3\}$ formam obrigatoriamente um referencial global suave para $T\mathbb{S}^3$. Assim, a paralelizabilidade de $\mathbb{S}^3$ decorre da sua estrutura de Grupo de Lie.
}

%%%% QUESTÃO 07 %%%%
\questao{Seja $M^{n}$ uma n-variedade suave e $T^{*}M$ seu fibrado cotangente. Com sua projeção natural e com a estrutura vetorial de cada fibra, mostre que TM é o único fibrado vetorial suave de posto n sobre $M^{n}$ para o qual todos os campos de covetores coordenados são seções locais suaves.}

\solucao{
	
	\textbf{Análise do Enunciado}
	\begin{itemize}
		\item \textbf{Hipóteses:} $M^n$ é uma variedade suave. O conjunto $T^*M = \coprod_{p \in M} T_p^*M$ é o fibrado cotangente munido de sua projeção natural $\pi: T^*M \to M^n$. (Nota analítica: o enunciado original grafa "TM" na segunda proposição, mas o contexto de "campos de covetores coordenados" e a premissa inicial definem inequivocamente que a tese se refere ao espaço total $T^*M$).
		\item \textbf{Tese:} Provar que a estrutura suave padrão de $T^*M$ é a única estrutura de variedade que torna $\pi$ um fibrado vetorial suave sob a condição de que, para qualquer carta, os diferenciais coordenados $(dx^1, \dots, dx^n)$ atuem como seções locais suaves.
	\end{itemize}
	
	\textbf{Fundamentação Teórica}
	\begin{enumerate}
		\item \textbf{Atlas do Fibrado Cotangente:} A estrutura suave padrão de $T^*M$ é fundamentada nas cartas naturais $(T^*U, \tilde{\phi})$. Para uma carta $(U, \phi)$ de $M^n$ com coordenadas $(x^1, \dots, x^n)$, um covetor $\omega_p \in T_p^*M$ escreve-se, em notação de Einstein, como $\omega_p = \omega_i dx^i|_p$. A carta natural é dada por $\tilde{\phi}(\omega_p) = (x^1(p), \dots, x^n(p), \omega_1, \dots, \omega_n)$.
		\item \textbf{Referenciais e Trivializações (Questão 3):} Todo referencial local suave determina univocamente uma trivialização local suave para o fibrado.
		\item \textbf{Atlas Maximal (Questão 1 da Lista 01):} Atlas suavemente compatíveis definem uma única estrutura suave.
	\end{enumerate}
	
	\textbf{Desenvolvimento Formal}
	
	Suponha que exista uma estrutura de variedade suave $\mathcal{E}$ no conjunto $T^*M$ com respeito à qual $\pi: T^*M \to M^n$ seja um fibrado vetorial suave e tal que, para qualquer carta suave $(U, (x^1, \dots, x^n))$ da variedade base $M^n$, os campos de covetores coordenados $(dx^1, \dots, dx^n)$ sejam seções locais estritamente suaves.
	
	Para cada ponto $p \in U$, os covetores $dx^1|_p, \dots, dx^n|_p$ formam a base dual do espaço cotangente $T_p^*M$. Como assumimos por hipótese que estas seções são suaves na estrutura $\mathcal{E}$, a coleção ordenada $(dx^1, \dots, dx^n)$ constitui um referencial local suave para $T^*M$ sobre o aberto $U$.
	
	Pela equivalência demonstrada rigorosamente na \textbf{Questão 3}, este referencial local suave induz uma única trivialização local suave na estrutura $\mathcal{E}$. Denotemos esta trivialização por $\Phi: \pi^{-1}(U) \to U \times \mathbb{R}^n$. Atuando sobre um covetor genérico $\omega_p = \omega_i dx^i|_p$, a trivialização extrai as coordenadas:
	\begin{equation*}
		\Phi(\omega_i dx^i|_p) = (p, (\omega_1, \dots, \omega_n)).
	\end{equation*}
	
	Sendo $\mathcal{E}$ uma estrutura de fibrado vetorial, $\Phi$ é um difeomorfismo local em relação a $\mathcal{E}$.
	
	Analisemos a relação desta trivialização $\Phi$ com a carta padrão $\tilde{\phi}$ do fibrado cotangente. A aplicação de transição é dada por:
	\begin{equation*}
		\tilde{\phi} \circ \Phi^{-1}(p, (\omega_1, \dots, \omega_n)) = \tilde{\phi}(\omega_i dx^i|_p) = (\phi(p), \omega_1, \dots, \omega_n).
	\end{equation*}
	Esta composição é essencialmente o produto direto $\phi \times \text{Id}_{\mathbb{R}^n}$. Sendo $\phi: U \to \mathbb{R}^n$ um difeomorfismo suave da variedade base, a aplicação produto preserva o difeomorfismo.
	
	Logo, as cartas geradas pelas trivializações na estrutura $\mathcal{E}$ são suavemente compatíveis com as cartas naturais que definem a estrutura padrão de $T^*M$. Conclui-se, pelo Lema do Atlas Maximal, que as estruturas são idênticas. A estrutura suave padrão é, portanto, única.
}
\solucaoalt{
	
	\textbf{Demonstração via Transições no Grupo de Lie $GL(n, \mathbb{R})$}
	
	A unicidade da estrutura do fibrado pode ser atestada analisando o comportamento das funções de transição no Grupo de Lie associado ao fibrado.
	
	Sejam $(U, (x^i))$ e $(V, (y^j))$ duas cartas coordenadas suaves superpostas em $M^n$. Na interseção $U \cap V$, a mudança de base no espaço cotangente é regida pela matriz Jacobiana transposta inversa. Em notação de Einstein, a transição entre os referenciais é:
	\begin{equation*}
		dy^j = \tfrac{\partial y^j}{\partial x^i} dx^i.
	\end{equation*}
	A matriz de transição $\mathcal{A}(p)$ cujas entradas são $A^j_i(p) = \tfrac{\partial y^j}{\partial x^i}(p)$ é, em cada ponto, não-singular. Portanto, a aplicação $p \mapsto \mathcal{A}(p)$ mapeia pontos da variedade no Grupo de Lie $GL(n, \mathbb{R})$.
	
	Como as cartas da base são difeomorfismos, as derivadas parciais $\tfrac{\partial y^j}{\partial x^i}$ são funções de classe $C^\infty$. Assim, a aplicação de transição para $GL(n, \mathbb{R})$ é estritamente suave.
	
	Qualquer estrutura de fibrado vetorial $\mathcal{E}$ em $T^*M$ que torne as seções coordenadas $(dx^i)$ suaves será obrigatoriamente \textbf{colada} pelas mesmas funções de transição suaves avaliadas em $GL(n, \mathbb{R})$. Como provamos na \textbf{Questão 1}, a identidade dessas funções de transição dita uma única estrutura de fibrado vetorial suave. Portanto, a estrutura $\mathcal{E}$ coincide topológica e algebricamente com a estrutura padrão.
}

%%%% QUESTÃO 08 %%%%
\questao{Sejam $M^{n}$ uma $n$-variedade suave e $\omega:M^{n}\rightarrow T^{*}M$ uma seção áspera.
	\begin{enumerate}[label=(\alph*)]
		\item Se $\omega=\sum_{i=1}^{n}\omega_{i}dx^{i}$ é a representação coordenada para $\omega$ em uma carta coordenada suave arbitrária $(U,(x^{i}))$, então mostre que $\omega$ é suave em $U$ se, e somente se, suas funções componentes são suaves.
		\item Mostre que $\omega$ é suave se, e somente se, para qualquer campo de vetores suave $X$ sobre um conjunto aberto $U\subset M^{n}$, a função $\omega(X): U \to \mathbb{R}$, definida por $\omega(X)(p)=\omega_{p}(X_{p})$ para todo $p\in U$, é suave.
	\end{enumerate}
}
\solucao{
	
	\textbf{Análise Detalhada do Enunciado}
	\begin{itemize}
		\item \textbf{Configuração Geométrica:} $M^n$ é uma variedade suave e $T^*M$ é o seu fibrado cotangente. Uma seção áspera $\omega$ associa a cada ponto $p \in M$ um covetor $\omega_p \in T_p^*M$, sem assumir a priori a suavidade da aplicação $\omega: M \to T^*M$.
		\item \textbf{Objetivo:} Estabelecer a equivalência entre a suavidade da seção (como mapa entre variedades) e a suavidade de suas componentes locais (a) e a suavidade de sua contração com campos vetoriais (b).
	\end{itemize}
	
	\textbf{Fundamentação Teórica}
	\begin{enumerate}
		\item \textbf{Trivializações do Fibrado Cotangente:} A estrutura suave de $T^*M$ é definida por cartas $\tilde{\phi}$ que mapeia $\omega_p = \omega_i dx^i|_p$ para $(x(p), \omega_1, \dots, \omega_n)$. Uma seção é suave se sua representação local $\hat{\omega} = \tilde{\phi} \circ \omega$ for suave (Lee, Cap. 5).
		\item \textbf{Referenciais Suaves (Questão 3, Lista 03):} a coleção $\{dx^i\}$ constitui um referencial local suave para $T^*M$. Qualquer seção $\omega$ pode ser escrita como combinação linear deste referencial.
		\item \textbf{Dualidade e Contrações:} A função $\omega(X)$ é a avaliação do covetor $\omega_p$ no vetor $X_p$, denotada pelo pareamento dual $\langle \omega_p, X_p \rangle$.
	\end{enumerate}
	
	\textbf{Desenvolvimento Formal - Parte (a)}
	
	Seja $(U, (x^i))$ uma carta coordenada suave. A trivialização canônica $\Phi: T^*U \to \varphi(U) \times \mathbb{R}^n$ associa a cada covetor $\xi_p = \xi_i dx^i|_p$ o par $(\phi(p), (\xi_1, \dots, \xi_n))$.
	
	A seção $\omega$ é suave se, e somente se, a aplicação $\hat{\omega} = \Phi \circ \omega : U \to \varphi(U) \times \mathbb{R}^n$ for suave. A representação de $\hat{\omega}$ é:
	\begin{equation*}
		\hat{\omega}(p) = (\phi(p), \omega_1(p), \dots, \omega_n(p)).
	\end{equation*}
	As primeiras $n$ componentes são as coordenadas da própria variedade $M$, que são suaves por definição. Portanto, a suavidade de $\hat{\omega}$ depende exclusivamente da suavidade das funções componentes $\omega_i : U \to \mathbb{R}$. Assim, $\omega$ é suave em $U$ se, e somente se, cada $\omega_i$ é de classe $C^\infty$.
	
	\textbf{Desenvolvimento Formal - Parte (b)}
	
	\textbf{Condição Necessária ($\implies$):} Assumimos que $\omega$ é uma seção suave. Seja $X$ um campo de vetores suave em $U$, expresso na base coordenada como $X = X^j \tfrac{\partial}{\partial x^j}$. A função $\omega(X): U \to \mathbb{R}$ é definida pela avaliação do covetor $\omega_p$ sobre o vetor $X_p$ para cada $p \in U$:
	
	\begin{equation*}
		\omega(X)(p) = \omega_p(X_p) = \left( \omega_i(p) dx^i|_p \right) \left( X^j(p) \tfrac{\partial}{\partial x^j} \Big|_p \right) = \omega_i(p) X^j(p) \delta^i_j = \omega_i(p) X^i(p),
	\end{equation*}
	
	onde utilizamos a convenção de soma de Einstein e o fato de que $\{dx^i\}$ e $\left\{\tfrac{\partial}{\partial x^j}\right\}$ são bases duais, satisfazendo $dx^i \left( \tfrac{\partial}{\partial x^j} \right) = \delta^i_j$. Como as funções componentes $\omega_i$ e $X^i$ são suaves em $U$, o produto $\omega_i X^i$ também é uma função suave, provando que $\omega(X)$ é suave.
	
	\textbf{Condição Suficiente ($\impliedby$):} Assumimos que, para qualquer campo de vetores suave $X$ em $U$, a função $\omega(X): U \to \mathbb{R}$ é suave. Para extrair a suavidade das funções componentes $\omega_j$, consideramos os campos de vetores coordenados $\tfrac{\partial}{\partial x^j}$ para cada $j \in \{1, \dots, n\}$. Como estes são campos suaves por construção, a hipótese garante que a função $\omega\left(\tfrac{\partial}{\partial x^j}\right)$ é suave em $U$. Calculando a sua representação em cada ponto $p \in U$, temos:
	
	\begin{equation*}
		\omega\left(\tfrac{\partial}{\partial x^j}\right)(p) = \omega_p \left( \left. \tfrac{\partial}{\partial x^j} \right|_p \right) = \left( \omega_i(p) dx^i|_p \right) \left( \left. \tfrac{\partial}{\partial x^j} \right|_p \right) = \omega_i(p) \delta^i_j = \omega_j(p).
	\end{equation*}
	
	Como $\omega\left(\tfrac{\partial}{\partial x^j}\right)$ é uma função suave, conclui-se que cada função componente $\omega_j$ é suave em $U$. Pela fundamentação estabelecida na \textbf{Parte (a)}, a suavidade de todas as funções componentes $\omega_j$ é a condição necessária e suficiente para que a seção $\omega$ seja suave.
}
\solucaoalt{
	\textbf{Justificativa Estrutural via Dualidade de Fibrados}
	
	A suavidade de $\omega$ pode ser interpretada como uma propriedade da estrutura de fibrados. Como $M$ é suave, o fibrado tangente $TM$ possui funções de transição suaves (Jacobianas). O fibrado $T^*M$ é o dual de $TM$, e suas funções de transição são as transpostas das inversas das Jacobianas de $M$.
	
	A equivalência $\omega \text{ suave} \iff \omega(X) \text{ suave}$ reflete o fato de que a dualidade entre $TM$ e $T^*M$ é um isomorfismo suave de feixes. Em termos concretos, a avaliação de uma seção do fibrado dual em seções do fibrado original recupera a topologia da variedade base, permitindo a caracterização da suavidade de $\omega$ através da sua ação em campos vetoriais.
}

%%%% QUESTÃO 09 %%%%
\questao{Se $M^{n}$ é uma $n$-variedade suave, mostre que $T^*M$ é trivial se, e somente se, $TM$ é trivial.}
\solucao{
	
	\textbf{Análise do Enunciado}
	\begin{itemize}
		\item \textbf{Hipóteses:} $M^n$ é uma variedade suave. $TM$ é o fibrado tangente e $T^*M$ é o seu fibrado dual (cotangente).
		\item \textbf{Tese:} Provar a equivalência $\text{Trivialidade}(TM) \iff \text{Trivialidade}(T^*M)$.
	\end{itemize}
	
	\textbf{Fundamentação Teórica}
	\begin{enumerate}
		\item \textbf{Trivialidade de Fibrados (Lee, Cap. 5):} Um fibrado vetorial de posto $n$ é trivial se, e somente se, admite um referencial global suave, isto é, um conjunto de $n$ seções suaves $\{\sigma_1, \dots, \sigma_n\}$ que formam uma base de cada fibra em todo o domínio.
		\item \textbf{Base Dual (Álgebra Linear):} Para qualquer base $\{X_i\}$ de um espaço vetorial $V$, existe uma única base dual $\{\omega^j\}$ de $V^*$ tal que $\omega^j(X_i) = \delta^j_i$.
		\item \textbf{Reflexividade:} Para espaços vetoriais de dimensão finita, existe um isomorfismo canônico $\Psi : V \to V^{**}$.
	\end{enumerate}
	
	\textbf{Desenvolvimento Formal}
	
	\textit{Ida ($\impliedby$):} Assuma que $TM$ é trivial. Então, existe um referencial global suave $\{X_1, \dots, X_n\}$ para $TM$. Para cada $p \in M$, definimos a base dual $\{\omega^1(p), \dots, \omega^n(p)\}$ em $T_p^*M$ tal que $\omega^j(p)(X_i(p)) = \delta^j_i$.
	
	Para provar que cada $\omega^j$ é uma seção suave de $T^*M$, tomemos um campo vetorial suave qualquer $Y \in \mathfrak{X}(M)$. Podemos escrever $Y$ na base $\{X_i\}$ como $Y = \sum_{k=1}^n Y^k X_k$, onde as funções $Y^k: M \to \mathbb{R}$ são suaves, pois $X_i$ formam um referencial suave. A ação de $\omega^j$ sobre $Y$ é:
	\begin{equation*}
		\omega^j(Y) = \omega^j(\sum_{k=1}^n Y^k X_k) = \sum_{k=1}^n Y^k \omega^j(X_k) = \sum_{k=1}^n Y^k \delta^j_k = Y^j.
	\end{equation*}
	Visto que $\omega^j(Y) = Y^j$ e $Y^j$ é uma função suave para qualquer $Y \in \mathfrak{X}(M)$, segue da caracterização de suavidade de seções do fibrado cotangente (\textbf{Questão 8b da Lista 03}) que $\omega^j$ é uma seção suave. Logo, $\{\omega^j\}$ é um referencial global suave para $T^*M$, tornando-o trivial.
	
	\textit{Volta ($\implies$):} Assuma que $T^*M$ é trivial com referencial global suave $\{\theta^1, \dots, \theta^n\}$. Pela reflexividade de espaços de dimensão finita, cada $\theta^i(p) \in T_p^*M$ pode ser vista como um elemento de $(T_p^*M)^* \cong T_p M$.
	
	Definimos os campos vetoriais $V_j$ tais que $\theta^i(V_j) = \delta^i_j$. De forma análoga à prova anterior, para qualquer 1-forma suave $\omega = \sum \omega_k \theta^k$, temos que a contração $\omega(V_j) = \omega_j$. Como as componentes $\omega_j$ são suaves por hipótese de suavidade de $\omega$, a aplicação $p \mapsto V_j(p)$ é uma seção suave de $TM$. Assim, $\{V_j\}$ é um referencial global suave para $TM$, tornando-o trivial.
}
\solucaoalt{
	\textbf{Demonstração via Funções de Transição e Grupos de Lie}
	
	Um fibrado vetorial é trivial se, e somente se, suas funções de transição podem ser reduzidas à identidade $I \in GL(n, \mathbb{R})$ através de uma mudança de referencial.
	
	Sejam $\tau_{\alpha\beta}: U_\alpha \cap U_\beta \to GL(n, \mathbb{R})$ as funções de transição de $TM$. As funções de transição $\tilde{\tau}_{\alpha\beta}$ de $T^*M$ são dadas pela transposta da inversa:
	\begin{equation*}
		\tilde{\tau}_{\alpha\beta}(p) = (\tau_{\alpha\beta}(p)^{-1})^T.
	\end{equation*}
	Se $TM$ é trivial, existe um atlas cujas transições são $\tau_{\alpha\beta} = I$. Substituindo na fórmula acima:
	\begin{equation*}
		\tilde{\tau}_{\alpha\beta} = (I^{-1})^T = I^T = I.
	\end{equation*}
	A identidade das matrizes de transição implica que $T^*M$ também é trivial. Como a operação de inversão e transposição em $GL(n, \mathbb{R})$ é um difeomorfismo (Questão 35, Lista 01), a trivialidade de um fibrado implica a do seu dual.
}

%%%% QUESTÃO 10 %%%%
\questao{
	Considere $f : \mathbb{R}^3 \to \mathbb{R}$, $f(x,y,z) = x^2 + y^2 + z^2$ e a aplicação $F : \mathbb{R}^2 \to \mathbb{R}^3$ (inversa da projeção estereográfica). Calcule $F^* df$ e $d(f \circ F)$ separadamente, e verifique que são iguais.
}
\solucao{
	
	\textbf{Análise do Enunciado}
	\begin{itemize}
		\item \textbf{Hipóteses:} $f$ é uma função suave em $\mathbb{R}^3$ definida por $f(x,y,z) = x^2 + y^2 + z^2$. $F$ é a aplicação inversa da projeção estereográfica, que mapeia o plano $\mathbb{R}^2$ na esfera unitária $\mathbb{S}^2 \subset \mathbb{R}^3$.
		\item \textbf{Tese:} Calcular analiticamente a diferencial da composição $d(f \circ F)$ e o pullback $F^* df$, demonstrando a igualdade entre ambos.
	\end{itemize}
	
	\textbf{Fundamentação Teórica}
	Conforme a teoria de variedades suaves (Lee, Cap. 3 e 8), utilizamos:
	\begin{enumerate}
		\item \textbf{Diferencial de uma função (Lista 02):} Para $f \in C^\infty(M)$, a diferencial $df$ é a 1-forma única que, em cada ponto $p \in M$, atua sobre um vetor tangente $X \in T_p M$ como a derivação de $f$ na direção de $X$, isto é, $df_p(X) = X(f)$.
		\item \textbf{Pullback de diferenciais:} Seja $F : M \to N$ uma aplicação suave. O pullback $F^* df$ é definido tal que, para qualquer vetor $X \in T_p M$, temos $(F^* df)_p(X) = df_{F(p)}(dF_p X)$.
		\item \textbf{Diferencial como Derivação (Lista 02):} A ação da diferencial de $F$ no ponto $p$, $dF_p X$, sobre uma função $f$ é dada por $(dF_p X)(f) = X(f \circ F)$.
		\item \textbf{Inversa da Aplicação $F$ (Questão 3, Lista 01):} A inversa da projeção estereográfica é dada por:
		\begin{equation*}
			F(u, v) = \left( \frac{2u}{u^2+v^2+1}, \frac{2v}{u^2+v^2+1}, \frac{u^2+v^2-1}{u^2+v^2+1} \right).
		\end{equation*}
	\end{enumerate}
	\textbf{Desenvolvimento Formal:}
	
	Primeiro, calculamos a função composta $f \circ F$. Definindo $S = u^2 + v^2 + 1$, temos:
	\begin{align*}
		(f \circ F)(u,v) &=\left( \tfrac{2u}{S} \right)^2 + \left( \tfrac{2v}{S} \right)^2 + \left( \tfrac{u^2+v^2-1 \color{red}{+1-1}}{S} \right)^2 = \left( \tfrac{2u}{S} \right)^2 + \left( \tfrac{2v}{S} \right)^2 + \left( \tfrac{S-2}{S} \right)^2 \\
		& = \tfrac{4(u^2+v^2) + S^2 - 4S + 4}{S^2} = \tfrac{4(S-1) + S^2 - 4S + 4}{S^2} = \tfrac{S^2}{S^2} \\ 
		& = 1.
	\end{align*}
	Como a composta é a função constante $1$, sua diferencial é $d(f \circ F) = 0$.
	
	Para provar que $F^* df = 0$, tomamos um vetor tangente qualquer $X \in T_p \mathbb{R}^2$. Utilizando as definições de pullback e push-forward (estudadas na \textbf{Lista 02}), temos:
	\begin{align*}
		(F^* df)_p(X) = df_{F(p)}(dF_p X) = (dF_p X)(f) = X(f \circ F) = X(1) = 0.
	\end{align*}
	Como a igualdade $(F^* df)_p(X) = 0$ vale para qualquer ponto $p$ e qualquer vetor $X$, concluímos que $F^* df = 0$. Portanto, $F^* df = d(f \circ F)$.
}

%%%% RESULTADO EXTRA %%%%
\resultadoextra{Propriedades da Diferencial e Funções com Diferencial Nula}{
	
	\textbf{Proposição 6.9 (Propriedades da Diferencial).} Sejam $f, g \in C^\infty(M)$. Para provar as identidades abaixo, tomamos um ponto $p \in M$ e um vetor tangente $X \in T_p M$ qualquer. Utilizaremos a definição de diferencial: $df_p(X) = X(f)$.
	
	\begin{enumerate}[label=(\alph*)]
		\item \textbf{Linearidade:} $d(af + bg) = a \, df + b \, dg$.
		\begin{align*}
			(d(af + bg))_p(X) &= X(af + bg) \\
			&= a X(f) + b X(g) \quad \text{(pela linearidade de $X$ como derivação)} \\
			&= a \, df_p(X) + b \, dg_p(X) = (a \, df + b \, dg)_p(X).
		\end{align*}
		
		\item \textbf{Regra do Produto (Leibniz):} $d(fg) = f \, dg + g \, df$.
		\begin{align*}
			(d(fg))_p(X) &= X(fg) \\
			&= f(p) X(g) + g(p) X(f) \quad \text{(pela regra do produto para derivações)} \\
			&= f(p) dg_p(X) + g(p) df_p(X) = (f \, dg + g \, df)_p(X).
		\end{align*}
		
		\item \textbf{Regra do Quociente:} $d(f/g) = (g \, df - f \, dg)/g^2$ onde $g \neq 0$.
		\begin{align*}
			(d(f/g))_p(X) &= X(f/g) \\
			&= \frac{g(p) X(f) - f(p) X(g)}{(g(p))^2} \quad \text{(pela regra do quociente para derivações)} \\
			&= \frac{g(p) df_p(X) - f(p) dg_p(X)}{(g(p))^2} = \left( \frac{g \, df - f \, dg}{g^2} \right)_p(X).
		\end{align*}
		
		\item \textbf{Regra da Cadeia:} $d(h \circ f) = (h' \circ f) \, df$.
		\begin{align*}
			(d(h \circ f))_p(X) &= X(h \circ f) \\
			&= h'(f(p)) X(f) \quad \text{(pela regra da cadeia para derivações)} \\
			&= (h' \circ f)(p) \, df_p(X) = ((h' \circ f) \, df)_p(X).
		\end{align*}
		
		\item \textbf{Função Constante:} Se $f$ é constante, então $df = 0$.
		
		Como $f$ é constante, para qualquer curva $\gamma$ tal que $\gamma(0)=p$ e $\gamma'(0)=X$, a função $(f \circ \gamma)(t)$ é constante. Logo, $X(f) = \tfrac{d}{dt}(f \circ \gamma)(t) \big|_{t=0} = 0$. Assim, $df_p(X) = 0$ para todo $X$, logo $df = 0$.
	\end{enumerate}
	\textbf{Proposição 6.10 (Funções com Diferenciais Nulas).} Se $f$ é uma função suave em $M$, então $df = 0$ se, e somente se, $f$ é constante em cada componente conexa de $M$.
	
	\textbf{Prova:}
	
	$(\impliedby)$ Se $f$ é constante em uma componente conexa, então pela Proposição 6.9(e), $df = 0$ nessa componente. Como a diferencial é definida pontualmente, se $f$ é constante em todas as componentes, $df = 0$ em todo $M$.
	
	$(\implies)$ Suponha que $df = 0$. Seja $C$ uma componente conexa de $M$. Queremos mostrar que $f$ é constante em $C$. 
	Sejam $p, q \in C$. Como as componentes conexas de uma variedade suave são conexas por caminhos (\textbf{Questão 02 da Lista 01}), existe uma curva suave $\gamma : [0,1] \to C$ tal que $\gamma(0)=p$ e $\gamma(1)=q$.
	
	Considere a função real $h(t) = f(\gamma(t))$. Pela regra da cadeia (Proposição 6.9(d)), temos que	$\tfrac{dh}{dt} = df_{\gamma(t)}(\gamma'(t))$.
	
	Como, por hipótese, $df = 0$, temos $\tfrac{dh}{dt} = 0$ para todo $t \in [0,1]$. Do cálculo elementar, sabemos que uma função com derivada nula em um intervalo é constante. Logo, $h(0) = h(1)$, ou seja, $f(p) = f(q)$. Como $p$ e $q$ foram escolhidos arbitrariamente em $C$, $f$ é constante em $C$.\hfill\cqd
	
	\vspace{0.3cm}
	\textbf{\underline{Observação}:} Vale ressaltar que qualquer variedade suave $M$ pode ser decomposta como a união disjunta de suas componentes conexas, $M = \sqcup_{i \in I} C_i$. Como $M$ é localmente conexa, cada componente $C_i$ é, simultaneamente, um subconjunto aberto e fechado de $M$, herdando a estrutura de variedade suave. Portanto, a conclusão de que $f$ é constante em cada componente conexa $C_i$ estende a validade do resultado para a totalidade de $M$.
}


%%%% QUESTÃO 11 %%%%
\questao{Sejam $[a, b] \subset \mathbb{R}$ e $\omega$ um campo de covetores suave em $[a, b]$. Se $t$ denota a coordenada padrão em $[a, b]$, então $\omega = f(t)dt$ para alguma função suave $f : [a, b] \to \mathbb{R}$. Definimos a integral de $\omega$ em $[a, b]$ por $\int_{[a,b]} \omega = \int_a^b f(t)dt$. Mostre que:
	\begin{enumerate}[label=(\alph*)]
		\item Se $\varphi : [c, d] \to [a, b]$ é um difeomorfismo crescente, então $\int_{[c,d]} \varphi^* \omega = \int_{[a,b]} \omega$.
		\item Se $\varphi : [c, d] \to [a, b]$ é um difeomorfismo decrescente, então $\int_{[c,d]} \varphi^* \omega = -\int_{[a,b]} \omega$.
\end{enumerate}}

\solucao{
	
	\textbf{Análise do Enunciado:}
	\begin{itemize}
		\item \textbf{Hipóteses:} $\omega = f(t)dt$ é um campo de covetores suave em $[a,b]$. $\varphi : [c,d] \to [a,b]$ é um difeomorfismo (aplicação suave, bijetiva, com inversa suave). 
		\item \textbf{Tese:} Provar que a integral do pullback $\varphi^* \omega$ sobre $[c,d]$ coincide com a integral de $\omega$ sobre $[a,b]$ se $\varphi$ for crescente, e é a negativa desta se $\varphi$ for decrescente.
	\end{itemize}
	
	\textbf{Fundamentação Teórica:}
	\begin{enumerate}
		\item \textbf{Pullback de campos de covetores (Lee, Cap. 6):} Se $\omega = f dt$ é um campo de covetores em $\mathbb{R}$ e $\varphi : M \to \mathbb{R}$ é uma aplicação suave, o pullback é dado por $\varphi^* \omega = (f \circ \varphi) d\varphi$. Em coordenadas locais $t$ para $[c,d]$, temos $d\varphi = \tfrac{d\varphi}{dt} dt$.
		\item \textbf{Teorema de Mudança de Variável (Cálculo):} Para funções integráveis, $\int_{\varphi(c)}^{\varphi(d)} f(u) du = \int_c^d f(\varphi(t)) \varphi'(t) dt$.
	\end{enumerate}
	
	\textbf{Desenvolvimento Formal:}
	
	Primeiramente, determinamos a expressão do pullback $\varphi^* \omega$. Como $\omega = f(t)dt$, temos:
	\begin{equation*}
		\varphi^* \omega = \varphi^*(f(t)dt) = (f \circ \varphi)(t) d(\varphi(t)) = f(\varphi(t)) \varphi'(t) dt.
	\end{equation*}
	A integral de $\varphi^* \omega$ sobre o intervalo $[c,d]$ é, por definição:
	\begin{equation*}
		\int_{[c,d]} \varphi^* \omega = \int_c^d f(\varphi(t)) \varphi'(t) dt.
	\end{equation*}
	
	\textit{Análise do caso (a): $\varphi$ é um difeomorfismo crescente.}
	Como $\varphi$ é crescente e bijetiva de $[c,d]$ para $[a,b]$, temos obrigatoriamente $\varphi(c) = a$ e $\varphi(d) = b$. Aplicando o Teorema de Mudança de Variável com a substituição $u = \varphi(t)$:
	\begin{equation*}
		\int_c^d f(\varphi(t)) \varphi'(t) dt = \int_{\varphi(c)}^{\varphi(d)} f(u) du = \int_a^b f(u) du = \int_{[a,b]} \omega.
	\end{equation*}
	
	\textit{Análise do caso (b): $\varphi$ é um difeomorfismo decrescente.}
	Como $\varphi$ é decrescente e bijetiva de $[c,d]$ para $[a,b]$, temos obrigatoriamente $\varphi(c) = b$ e $\varphi(d) = a$. Aplicando novamente o Teorema de Mudança de Variável:
	\begin{equation*}
		\int_c^d f(\varphi(t)) \varphi'(t) dt = \int_{\varphi(c)}^{\varphi(d)} f(u) du = \int_b^a f(u) du.
	\end{equation*}
	Utilizando a propriedade de inversão dos limites de integração $\int_b^a = -\int_a^b$, obtemos:
	\begin{equation*}
		\int_b^a f(u) du = - \int_a^b f(u) du = - \int_{[a,b]} \omega.
	\end{equation*}
	Ambos os casos foram demonstrados exaustivamente.
}

%%%% QUESTÃO 12 %%%%
\questao{Seja $M^n$ uma variedade suave. Se $\gamma : [a, b] \to M^n$ é uma curva suave em $M^n$ e $\omega$ é um campo de covetores em $M^n$, então definimos a integral de linha de $\omega$ sobre $\gamma$ por $\int_\gamma \omega = \int_{[a,b]} \gamma^* \omega$. Mostre que $\int_\gamma \omega = \int_a^b \omega_{\gamma(t)}(\gamma'(t)) dt$.}

\solucao{
	
	\textbf{Análise do Enunciado:}
	\begin{itemize}
		\item \textbf{Hipóteses:} $M^n$ é uma $n$-variedade suave; $\gamma : [a, b] \to M^n$ é uma curva suave; $\omega \in \mathfrak{X}^*(M^n)$ é um campo de covetores suave.
		\item \textbf{Tese:} Provar que a definição de integral de linha via pullback, $\int_{[a,b]} \gamma^* \omega$, coincide com a integral de Riemann da função $t \mapsto \omega_{\gamma(t)}(\gamma'(t))$ sobre o intervalo $[a,b]$.
	\end{itemize}
	
	\textbf{Fundamentação Teórica:}
	\begin{enumerate}
		\item \textbf{Definição de Integral em Intervalos (Questão 11, Lista 03):} Estabelecemos que, para qualquer campo de covetores $\eta$ em $[a,b]$, se $\eta = f(t) dt$, então $\int_{[a,b]} \eta = \int_a^b f(t) dt$.
		\item \textbf{Pullback de campos de covetores (Lee, Cap. 6):} O pullback $\gamma^* \omega$ é o campo de covetores em $[a,b]$ definido por $(\gamma^* \omega)_t(X) = \omega_{\gamma(t)}(d\gamma_t X)$ para todo $X \in T_t [a,b]$.
		\item \textbf{Vetor Velocidade (Lee, Cap. 6):} Para uma curva suave $\gamma$, o vetor velocidade $\gamma'(t)$ é definido como a imagem do vetor base $\tfrac{\partial}{\partial t} \in T_t [a,b]$ pela diferencial de $\gamma$, ou seja, $\gamma'(t) = d\gamma_t(\tfrac{\partial}{\partial t})$.
	\end{enumerate}
	
	\textbf{Desenvolvimento Formal:}
	
	Seja $\eta = \gamma^* \omega$. Como $\eta$ é um campo de covetores em $[a,b]$, ele pode ser escrito na forma $\eta = f(t) dt$, onde $f(t)$ é a função obtida ao avaliar $\eta$ no campo de vetores coordenado $\tfrac{\partial}{\partial t}$. Calculamos então essa função:
	\begin{align*}
		f(t) &= \eta_t \left( \left. \tfrac{\partial}{\partial t} \right|_t \right) \\
		&= (\gamma^* \omega)_t \left( \left. \tfrac{\partial}{\partial t} \right|_t \right) \\
		&= \omega_{\gamma(t)} \left( d\gamma_t \left( \left. \tfrac{\partial}{\partial t} \right|_t \right) \right).
	\end{align*}
	
	Pela definição de vetor velocidade $\gamma'(t)$ citada na fundamentação, temos que $d\gamma_t(\tfrac{\partial}{\partial t}) = \gamma'(t)$. Substituindo este resultado na expressão acima, obtemos:
	\begin{equation*}
		f(t) = \omega_{\gamma(t)}(\gamma'(t)).
	\end{equation*}
	
	Portanto, o campo de covetores $\gamma^* \omega$ em $[a,b]$ é dado por:
	\begin{equation*}
		\gamma^* \omega = \omega_{\gamma(t)}(\gamma'(t)) dt.
	\end{equation*}
	
	Agora, aplicando a definição de integral em intervalos estabelecida na \textbf{Questão 11}, temos:
	\begin{align*}
		\int_\gamma \omega &= \int_{[a,b]} \gamma^* \omega \\
		&= \int_{[a,b]} \omega_{\gamma(t)}(\gamma'(t)) dt \\
		&= \int_a^b \omega_{\gamma(t)}(\gamma'(t)) dt.
	\end{align*}
	A igualdade está demonstrada.
}

%%%% QUESTÃO 13 %%%%
\questao{Sejam $M = \mathbb{R}^2 \setminus \{0\}$, $\omega = \frac{xdy - ydx}{x^2 + y^2} \in \mathfrak{X}^*(M)$ e $\gamma(t) = (\cos t, \sin t)$, $t \in [0, 2\pi]$. Calcule $\int_\gamma \omega$.}

\solucao{
	\textbf{Análise do Enunciado:}
	\begin{itemize}
		\item \textbf{Hipóteses:} A variedade é o plano menos a origem $M = \mathbb{R}^2 \setminus \{0\}$; o campo de covetores é $$\omega = \frac{x}{x^2+y^2}dy - \frac{y}{x^2+y^2}dx;$$ a curva é o círculo unitário $\gamma : [0, 2\pi] \to M$ dada por $\gamma(t) = (\cos t, \sin t)$.
		\item \textbf{Tese:} Calcular o valor real da integral de linha $\int_\gamma \omega$.
	\end{itemize}
	
	\textbf{Fundamentação Teórica:}
	\begin{enumerate}
		\item \textbf{Cálculo da Integral de Linha (Questão 12, Lista 03):} Para calcular a integral de linha de um campo de covetores $\omega$ sobre uma curva $\gamma$, utilizamos a fórmula:
		\begin{equation*}
			\int_\gamma \omega = \int_a^b \omega_{\gamma(t)}(\gamma'(t)) dt.
		\end{equation*}
		\item \textbf{Avaliação de Covetores em $\mathbb{R}^n$:} Se $\omega = A(x,y)dx + B(x,y)dy$ e $X = X^1 \frac{\partial}{\partial x} + X^2 \frac{\partial}{\partial y}$, então $$\omega_p(X_p) = A(p)X^1(p) + B(p)X^2(p).$$
	\end{enumerate}
	
	\textbf{Desenvolvimento Formal:}
	
	Primeiramente, identificamos as componentes do campo de covetores $\omega$ e a parametrização da curva $\gamma$. Temos:
	\begin{equation*}
		\omega = \underbrace{\left( \frac{-y}{x^2+y^2} \right)}_{A(x,y)} dx + \underbrace{\left( \frac{x}{x^2+y^2} \right)}_{B(x,y)} dy.
	\end{equation*}
	
	Para a curva $\gamma(t) = (\cos t, \sin t)$, as coordenadas são $x(t) = \cos t$ e $y(t) = \sin t$. O vetor velocidade $\gamma'(t)$ é dado por:
	\begin{equation*}
		\gamma'(t) = \frac{d}{dt}(\cos t) \frac{\partial}{\partial x} + \frac{d}{dt}(\sin t) \frac{\partial}{\partial y} = -\sin t \frac{\partial}{\partial x} + \cos t \frac{\partial}{\partial y}.
	\end{equation*}
	
	Agora, avaliamos o campo de covetores $\omega$ no ponto $\gamma(t)$ e o aplicamos ao vetor $\gamma'(t)$. Note que, sobre a curva $\gamma$, temos $x^2+y^2 = \cos^2 t + \sin^2 t = 1$. Assim:
	\begin{align*}
		\omega_{\gamma(t)}(\gamma'(t)) &= A(x(t), y(t)) \cdot x'(t) + B(x(t), y(t)) \cdot y'(t) \\
		&= \left( \frac{-\sin t}{1} \right)(-\sin t) + \left( \frac{\cos t}{1} \right)(\cos t) \\
		&= \sin^2 t + \cos^2 t \\
		&= 1.
	\end{align*}
	
	Finalmente, substituímos este resultado na fórmula da integral de linha:
	\begin{align*}
		\int_\gamma \omega &= \int_0^{2\pi} \omega_{\gamma(t)}(\gamma'(t)) dt \\
		&= \int_0^{2\pi} 1 \, dt \\
		&= \left[ t \right]_0^{2\pi} \\
		&= 2\pi.
	\end{align*}
	O valor da integral de linha é $2\pi$.
}

%%%% QUESTÃO 14 %%%%
\questao{Sejam $\gamma : [a, b] \to M^n$ e $\tilde{\gamma} : [c, d] \to M^n$ segmentos de curvas suaves em $M^n$. Dizemos que $\tilde{\gamma}$ é uma reparametrização de $\gamma$ se $\tilde{\gamma} = \gamma \circ \varphi$, para algum difeomorfismo $\varphi : [c, d] \to [a, b]$. Se $\varphi$ é crescente (decrescente), dizemos que $\tilde{\gamma}$ é uma reparametrização para frente (para trás). Mostre que, se $\omega \in \mathfrak{X}^*(M)$, então para qualquer reparametrização $\tilde{\gamma}$ de $\gamma$ temos:
	\begin{enumerate}[label=(\alph*)]
		\item $\int_{\tilde{\gamma}} \omega = \int_{\gamma} \omega$, se $\varphi$ é reparametrização para frente.
		\item $\int_{\tilde{\gamma}} \omega = -\int_{\gamma} \omega$, se $\varphi$ é reparametrização para trás.
\end{enumerate}}

\solucao{
	\textbf{Análise do Enunciado:}
	\begin{itemize}
		\item \textbf{Hipóteses:} $\gamma$ e $\tilde{\gamma}$ são curvas suaves; $\omega$ é um campo de covetores suave em $M^n$; $\tilde{\gamma}$ é definida pela composição $\gamma \circ \varphi$, onde $\varphi : [c, d] \to [a, b]$ é um difeomorfismo.
		\item \textbf{Tese:} Provar a invariância da integral de linha sob reparametrizações, distinguindo os casos de preservação (crescente) e inversão (decrescente) de orientação.
	\end{itemize}
	
	\textbf{Fundamentação Teórica:}
	\begin{enumerate}
		\item \textbf{Definição de Integral de Linha (Questão 12, Lista 03):} $\int_{\gamma} \omega = \int_{[a,b]} \gamma^* \omega$.
		\item \textbf{Composição de Pullbacks (Lee, Cap. 6):} Para aplicações suaves $F$ e $G$, temos $(F \circ G)^* \omega = G^* (F^* \omega)$. No nosso caso, $\tilde{\gamma}^* \omega = (\gamma \circ \varphi)^* \omega = \varphi^* (\gamma^* \omega)$.
		\item \textbf{Invariância da Integral em Intervalos (Questão 11, Lista 03):} Se $\eta$ é um campo de covetores em $[a,b]$ e $\varphi : [c,d] \to [a,b]$ é um difeomorfismo, então $\int_{[c,d]} \varphi^* \eta = \int_{[a,b]} \eta$ se $\varphi$ for crescente, e $\int_{[c,d]} \varphi^* \eta = -\int_{[a,b]} \eta$ se $\varphi$ for decrescente.
	\end{enumerate}
	
	\textbf{Desenvolvimento Formal:}
	
	Iniciamos aplicando a definição de integral de linha para a curva $\tilde{\gamma}$:
	\begin{equation*}
		\int_{\tilde{\gamma}} \omega = \int_{[c,d]} \tilde{\gamma}^* \omega.
	\end{equation*}
	Utilizando a hipótese de que $\tilde{\gamma} = \gamma \circ \varphi$ e a propriedade de composição de pullbacks citada na fundamentação, temos:
	\begin{equation*}
		\tilde{\gamma}^* \omega = (\gamma \circ \varphi)^* \omega = \varphi^* (\gamma^* \omega).
	\end{equation*}
	Seja $\eta = \gamma^* \omega$ o campo de covetores resultante do pullback de $\omega$ por $\gamma$. Note que $\eta$ é um campo de covetores suave definido no intervalo $[a,b]$. Substituindo $\eta$ na expressão da integral, obtemos:
	\begin{equation*}
		\int_{\tilde{\gamma}} \omega = \int_{[c,d]} \varphi^* \eta.
	\end{equation*}
	
	Agora, aplicamos o resultado provado na \textbf{Questão 11} para a integral de $\eta$ sob o difeomorfismo $\varphi$:
	
	\textit{Caso (a): $\varphi$ é uma reparametrização para frente (crescente).}
	Pela Questão 11(a), sabemos que $\int_{[c,d]} \varphi^* \eta = \int_{[a,b]} \eta$. Logo:
	\begin{equation*}
		\int_{\tilde{\gamma}} \omega = \int_{[a,b]} \gamma^* \omega = \int_{\gamma} \omega.
	\end{equation*}
	
	\textit{Caso (b): $\varphi$ é uma reparametrização para trás (decrescente).}
	Pela Questão 11(b), sabemos que $\int_{[c,d]} \varphi^* \eta = -\int_{[a,b]} \eta$. Logo:
	\begin{equation*}
		\int_{\tilde{\gamma}} \omega = -\int_{[a,b]} \gamma^* \omega = -\int_{\gamma} \omega.
	\end{equation*}
	A demonstração está completa.
}

%%%% QUESTÃO 15 %%%%
\questao{(Teorema Fundamental das integrais de linha) Sejam $M^n$ uma variedade suave, $f \in C^\infty(M)$ e $\gamma : [a, b] \to M^n$ um segmento de curva suave em $M^n$. Mostre que $\int_\gamma df = f(\gamma(b)) - f(\gamma(a))$.}

\solucao{
	\textbf{Análise do Enunciado:}
	\begin{itemize}
		\item \textbf{Hipóteses:} $M^n$ é uma $n$-variedade suave; $f$ é uma função suave em $M^n$ (logo, $df \in \mathfrak{X}^*(M)$ é um campo de covetores suave); $\gamma : [a, b] \to M^n$ é um segmento de curva suave.
		\item \textbf{Tese:} Provar que a integral de linha do campo de covetores $df$ ao longo de $\gamma$ é igual à diferença dos valores de $f$ nos pontos extremidades da curva.
	\end{itemize}
	
	\textbf{Fundamentação Teórica:}
	\begin{enumerate}
		\item \textbf{Definição de Integral de Linha (Questão 12, Lista 03):} Para um campo de covetores $\omega$, temos $\int_\gamma \omega = \int_a^b \omega_{\gamma(t)}(\gamma'(t)) dt$. Aplicando ao caso $\omega = df$, temos $\int_\gamma df = \int_a^b df_{\gamma(t)}(\gamma'(t)) dt$.
		\item \textbf{A Diferencial como Derivação (Lista 02):} Por definição, a diferencial $df$ atua sobre um vetor tangente $X \in T_p M$ como a derivada de $f$ na direção de $X$. Assim, para o vetor velocidade $\gamma'(t) \in T_{\gamma(t)} M$, temos:
		\begin{equation*}
			df_{\gamma(t)}(\gamma'(t)) = \gamma'(t)(f).
		\end{equation*}
		\item \textbf{Regra da Cadeia para Curvas (Lee, Cap. 3):} A ação do vetor velocidade $\gamma'(t)$ sobre uma função suave $f$ é exatamente a derivada da composição $f \circ \gamma$ em relação ao parâmetro $t$:
		\begin{equation*}
			\gamma'(t)(f) = \frac{d}{dt}(f \circ \gamma)(t).
		\end{equation*}
		\item \textbf{Teorema Fundamental do Cálculo (TFC):} Para qualquer função $h : [a,b] \to \mathbb{R}$ com derivada contínua, $\int_a^b h'(t) dt = h(b) - h(a)$.
	\end{enumerate}
	
	\textbf{Desenvolvimento Formal:}
	
	Partimos da definição de integral de linha para o campo de covetores $df$:
	\begin{equation*}
		\int_\gamma df = \int_a^b df_{\gamma(t)}(\gamma'(t)) dt.
	\end{equation*}
	
	Utilizando a fundamentação teórica (item 2), substituímos a avaliação do covetor $df$ pelo operador de derivação $\gamma'(t)$:
	\begin{equation*}
		\int_\gamma df = \int_a^b \gamma'(t)(f) dt.
	\end{equation*}
	
	Agora, aplicamos a Regra da Cadeia (item 3), reconhecendo que a expressão $\gamma'(t)(f)$ é a derivada da função escalar $h(t) = (f \circ \gamma)(t)$ definida no intervalo $[a,b]$:
	\begin{equation*}
		\int_\gamma df = \int_a^b \frac{d}{dt}(f \circ \gamma)(t) dt.
	\end{equation*}
	
	Finalmente, aplicando o Teorema Fundamental do Cálculo à função $h(t) = f(\gamma(t))$, obtemos:
	\begin{align*}
		\int_\gamma df &= (f \circ \gamma)(b) - (f \circ \gamma)(a) \\
		&= f(\gamma(b)) - f(\gamma(a)).
	\end{align*}
	A tese está provada.
}

%%%% QUESTÃO 16 %%%%
\questao{Sejam $M^n$ e $N^m$ variedades suaves, $F : M^n \to N^m$ uma aplicação suave, $\omega \in \mathfrak{X}^*(N)$ e $\gamma : [a, b] \to M^n$ um segmento de curva suave por partes em $M^n$. Mostre que $\int_{\gamma} F^* \omega = \int_{F \circ \gamma} \omega$.}

\solucao{
	\textbf{Análise do Enunciado:}
	\begin{itemize}
		\item \textbf{Hipóteses:} $M^n$ e $N^m$ são variedades suaves; $F : M^n \to N^m$ é uma aplicação suave; $\omega \in \mathfrak{X}^*(N)$ é um campo de covetores suave em $N$; $\gamma : [a, b] \to M^n$ é um segmento de curva suave por partes.
		\item \textbf{Tese:} Provar que a integral de linha do pullback $F^* \omega$ ao longo de $\gamma$ é igual à integral de linha de $\omega$ ao longo da curva composta $\sigma = F \circ \gamma$ em $N$.
	\end{itemize}
	
	\textbf{Fundamentação Teórica:}
	\begin{enumerate}
		\item \textbf{Definição de Integral de Linha (Questão 12, Lista 03):} Para qualquer campo de covetores $\eta$ e curva $\gamma$, $\int_{\gamma} \eta = \int_a^b \eta_{\gamma(t)}(\gamma'(t)) dt$.
		\item \textbf{Definição de Pullback de Campos de Covetores (Lee, Cap. 6):} O pullback $F^* \omega$ é o campo de covetores em $M$ definido pontualmente por $(F^* \omega)_p(X) = \omega_{F(p)}(dF_p X)$ para todo $X \in T_p M$.
		\item \textbf{Diferencial de Composição (Regra da Cadeia):} Seja $\sigma = F \circ \gamma$ a imagem de $\gamma$ em $N$. O vetor velocidade de $\sigma$ no ponto $t$ é dado por:
		\begin{equation*}
			\sigma'(t) = dF_{\gamma(t)}(\gamma'(t)).
		\end{equation*}
		\item \textbf{Aditividade da Integral:} Como a integral de linha é a soma das integrais sobre cada segmento suave de $\gamma$, basta provar a igualdade para um único segmento suave $\gamma : [a,b] \to M^n$.
	\end{enumerate}
	
	\textbf{Desenvolvimento Formal:}
	
	Seja $\gamma : [a,b] \to M^n$ um segmento suave. Aplicamos a definição de integral de linha ao campo de covetores $F^* \omega$ ao longo de $\gamma$:
	\begin{equation*}
		\int_{\gamma} F^* \omega = \int_a^b (F^* \omega)_{\gamma(t)}(\gamma'(t)) dt.
	\end{equation*}
	
	Agora, utilizamos a definição de pullback citada na fundamentação (item 2), avaliando o covetor $(F^* \omega)$ no ponto $p = \gamma(t)$ e aplicando-o ao vetor $X = \gamma'(t)$:
	\begin{equation*}
		(F^* \omega)_{\gamma(t)}(\gamma'(t)) = \omega_{F(\gamma(t))} \left( dF_{\gamma(t)} (\gamma'(t)) \right).
	\end{equation*}
	
	Definimos a curva $\sigma : [a,b] \to N^m$ como $\sigma = F \circ \gamma$. Pela regra da cadeia (item 3), o vetor velocidade de $\sigma$ é precisamente $\sigma'(t) = dF_{\gamma(t)}(\gamma'(t))$. Substituindo estas expressões na integral, obtemos:
	\begin{equation*}
		\int_{\gamma} F^* \omega = \int_a^b \omega_{\sigma(t)}(\sigma'(t)) dt.
	\end{equation*}
	
	Reconhecemos que a expressão à direita é, por definição (Questão 12), a integral de linha do campo de covetores $\omega$ ao longo da curva $\sigma = F \circ \gamma$:
	\begin{equation*}
		\int_a^b \omega_{\sigma(t)}(\sigma'(t)) dt = \int_{\sigma} \omega = \int_{F \circ \gamma} \omega.
	\end{equation*}
	
	Assim, provamos que $\int_{\gamma} F^* \omega = \int_{F \circ \gamma} \omega$ para um segmento suave. Pela aditividade da integral, o resultado estende-se a qualquer curva suave por partes.
}

%%%% QUESTÃO 17 %%%%
\questao{Mostre que a composição de submersões é uma submersão, a composição de imersões é uma imersão e que a composição de mergulhos suaves é um mergulho suave.}

\solucao{
	\textbf{Análise do Enunciado:}
	\begin{itemize}
		\item \textbf{Hipóteses:} Sejam $M, N, P$ variedades suaves. Consideremos aplicações suaves $F : M \to N$ e $G : N \to P$. 
		\item \textbf{Tese:} 
		\begin{enumerate}[label=(\roman*)]
			\item Se $F$ e $G$ são submersões, então $G \circ F$ é uma submersão.
			\item Se $F$ e $G$ são imersões, então $G \circ F$ é uma imersão.
			\item Se $F$ e $G$ são mergulhos suaves, então $G \circ F$ é um mergulho suave.
		\end{enumerate}
	\end{itemize}
	
	\textbf{Fundamentação Teórica:}
	\begin{enumerate}
		\item \textbf{Definições (Lee, Cap. 7):}
		\begin{itemize}
			\item Uma aplicação $F: M \to N$ é uma \textbf{imersão} se $dF_p : T_p M \to T_{F(p)} N$ é injetiva para todo $p \in M$.
			\item Uma aplicação $F: M \to N$ é uma \textbf{submersão} se $dF_p : T_p M \to T_{F(p)} N$ é sobrejetiva para todo $p \in M$.
			\item Uma aplicação $F: M \to N$ é um \textbf{mergulho suave} se é uma imersão e um homeomorfismo sobre sua imagem $F(M)$ (munida da topologia de subespaço de $N$).
		\end{itemize}
		\item \textbf{Regra da Cadeia (Lee, Cap. 3):} Para a composição $G \circ F$, a diferencial no ponto $p \in M$ é dada pela composição das diferenciais:
		\begin{equation*}
			d(G \circ F)_p = dG_{F(p)} \circ dF_p.
		\end{equation*}
		\item \textbf{Álgebra Linear:} A composição de duas aplicações lineares injetivas é injetiva; a composição de duas aplicações lineares sobrejetivas é sobrejetiva.
	\end{enumerate}
	
	\textbf{Desenvolvimento Formal:}
	
	\textit{Parte (i): Composição de Submersões.}
	Sejam $F: M \to N$ e $G: N \to P$ submersões. Por definição, para qualquer $p \in M$, a aplicação $dF_p : T_p M \to T_{F(p)} N$ é sobrejetiva, e para qualquer $q \in N$ (incluindo $q = F(p)$), a aplicação $dG_q : T_q N \to T_{G(q)} P$ é sobrejetiva. 
	Pela Regra da Cadeia:
	\begin{equation*}
		d(G \circ F)_p = dG_{F(p)} \circ dF_p.
	\end{equation*}
	Como a composição de duas aplicações lineares sobrejetivas é sobrejetiva, $d(G \circ F)_p$ é sobrejetiva para todo $p \in M$. Portanto, $G \circ F$ é uma submersão.
	
	\textit{Parte (ii): Composição de Imersões.}
	Sejam $F: M \to N$ e $G: N \to P$ imersões. Por definição, $dF_p$ é injetiva para todo $p \in M$ e $dG_{F(p)}$ é injetiva para todo $p \in M$. 
	Novamente, pela Regra da Cadeia:
	\begin{equation*}
		d(G \circ F)_p = dG_{F(p)} \circ dF_p.
	\end{equation*}
	Como a composição de duas aplicações lineares injetivas é injetiva, $d(G \circ F)_p$ é injetiva para todo $p \in M$. Portanto, $G \circ F$ é uma imersão.
	
	\textit{Parte (iii): Composição de Mergulhos Suaves.}
	Sejam $F: M \to N$ e $G: N \to P$ mergulhos suaves.
	1. Pela Parte (ii), como $F$ e $G$ são imersões, a composição $G \circ F$ é uma imersão.
	2. Sobre a topologia: $F$ é um homeomorfismo entre $M$ e $F(M) \subseteq N$, e $G$ é um homeomorfismo entre $N$ e $G(N) \subseteq P$. A restrição de $G$ ao subespaço $F(M)$ é, portanto, um homeomorfismo entre $F(M)$ e $G(F(M))$. 
	A composição de dois homeomorfismos é um homeomorfismo. Logo, $G \circ F$ é um homeomorfismo entre $M$ e a sua imagem $(G \circ F)(M) \subseteq P$.
	
	Sendo $G \circ F$ uma imersão e um homeomorfismo sobre sua imagem, conclui-se que $G \circ F$ é um mergulho suave.
}

%%%% QUESTÃO 18 %%%%
\questao{Seja $F : M \to N$ uma imersão injetiva. Se alguma das seguintes condições é satisfeita:
	\begin{enumerate}[label=(\roman*)]
		\item $M$ é compacta;
		\item $F$ é uma aplicação própria.
	\end{enumerate}
	Mostre que $F$ é um mergulho suave com imagem fechada.}

\solucao{
	\textbf{Análise do Enunciado:}
	\begin{itemize}
		\item \textbf{Hipóteses:} $F: M \to N$ é uma imersão injetiva. Temos duas condições alternativas: (i) $M$ é compacta ou (ii) $F$ é uma aplicação própria.
		\item \textbf{Tese:} Provar que, sob qualquer uma dessas condições, $F$ é um mergulho suave (ou seja, $F$ é um homeomorfismo sobre sua imagem $F(M)$) e que $F(M)$ é um subconjunto fechado de $N$.
	\end{itemize}
	
	\textbf{Fundamentação Teórica:}
	\begin{enumerate}
		\item \textbf{Definição de Mergulho Suave (Lee, Cap. 7):} Uma imersão injetiva $F: M \to N$ é um mergulho suave se a aplicação $F : M \to F(M)$ é um homeomorfismo, onde $F(M)$ possui a topologia de subespaço de $N$.
		\item \textbf{Propriedade de Espaços Compactos:} Se $X$ é um espaço compacto e $Y$ é um espaço Hausdorff, qualquer aplicação contínua bijetiva $f: X \to Y$ é, necessariamente, um homeomorfismo.
		\item \textbf{Definição de Aplicação Própria (Lee, Cap. 2):} Uma aplicação $F: M \to N$ é própria se, para todo subconjunto compacto $K \subset N$, a imagem inversa $F^{-1}(K)$ é compacta em $M$.
		\item \textbf{Propriedade de Aplicações Próprias:} Toda aplicação própria entre variedades suaves é uma aplicação fechada (ou seja, leva conjuntos fechados em $M$ a conjuntos fechados em $N$). Além disso, a imagem de uma aplicação própria é sempre um subconjunto fechado do contradomínio.
	\end{enumerate}
	
	\textbf{Desenvolvimento Formal:}
	
	\textit{Análise do caso (i): $M$ é compacta.}
	Como $F$ é uma imersão injetiva, ela é, por definição, uma aplicação contínua e bijetiva entre $M$ e sua imagem $F(M)$. 
	Como $M$ é compacta e $N$ é uma variedade suave (logo, é um espaço Hausdorff), a imagem $F(M)$ também é compacta. Em qualquer espaço Hausdorff, um subconjunto compacto é obrigatoriamente fechado. Portanto, $F(M)$ é fechado em $N$.
	
	Quanto à topologia, aplicamos a propriedade fundamental dos espaços compactos: como $F : M \to F(M)$ é uma bijecão contínua de um espaço compacto para um espaço Hausdorff, ela é automaticamente um homeomorfismo. Sendo $F$ uma imersão e um homeomorfismo sobre sua imagem, conclui-se que $F$ é um mergulho suave com imagem fechada.
	
	\textit{Análise do caso (ii): $F$ é uma aplicação própria.}
	Suponhamos que $F$ seja uma aplicação própria. Como $M$ é uma variedade suave, $M$ é fechado em si mesma. Como $F$ é uma aplicação própria, ela é uma aplicação fechada. Portanto, a imagem $F(M)$ é um subconjunto fechado de $N$.
	
	Para provar que $F$ é um homeomorfismo sobre sua imagem, devemos mostrar que $F$ é uma aplicação aberta de $M$ para $F(M)$, ou equivalentemente, que $F$ leva conjuntos fechados de $M$ a conjuntos fechados em $F(M)$. Como $F$ é uma aplicação própria, ela é uma aplicação fechada em relação a $N$. Se $C \subset M$ é fechado, então $F(C)$ é fechado em $N$. Como $F(C) \subset F(M)$, então $F(C)$ também é fechado na topologia de subespaço de $F(M)$.
	
	Sendo assim, $F : M \to F(M)$ é uma bijecão contínua e fechada, o que a caracteriza como um homeomorfismo. Como $F$ é também uma imersão, conclui-se que $F$ é um mergulho suave com imagem fechada.
}

%%%% QUESTÃO 19 %%%%
\questao{Mostre que a aplicação $f : (1, +\infty) \to \mathbb{R}^2$ dada por $f(t) = \left( \tfrac{1}{t} \cos(2\pi t), \tfrac{1}{t} \sin(2\pi t) \right)$ é um mergulho suave.}

\solucao{
	\textbf{Análise do Enunciado:}
	\begin{itemize}
		\item \textbf{Hipóteses:} O domínio é o intervalo aberto $I = (1, +\infty) \subset \mathbb{R}$; o contradomínio é $\mathbb{R}^2$; a aplicação $f$ é definida por componentes trigonométricas e racionais.
		\item \textbf{Tese:} Provar que $f$ é um mergulho suave, o que exige demonstrar que: (i) $f$ é suave, (ii) $f$ é uma imersão, (iii) $f$ é injetiva e (iv) $f$ é um homeomorfismo sobre sua imagem $f(I) \subset \mathbb{R}^2$.
	\end{itemize}
	
	\textbf{Fundamentação Teórica:}
	\begin{enumerate}
		\item \textbf{Definição de Mergulho Suave (Lee, Cap. 7):} Uma aplicação $f: M \to N$ é um mergulho suave se é uma imersão injetiva e se $f : M \to f(M)$ é um homeomorfismo, onde $f(M)$ possui a topologia de subespaço de $N$.
		\item \textbf{Imersão:} $f$ é uma imersão se a diferencial $df_t : T_t I \to T_{f(t)} \mathbb{R}^2$ é injetiva para todo $t \in I$. Como $\dim(I) = 1$, isso equivale a dizer que o vetor velocidade $f'(t) \neq (0,0)$.
		\item \textbf{Topologia de Subespaço:} Para provar que $f$ é um homeomorfismo sobre sua imagem, basta mostrar que a aplicação inversa $f^{-1} : f(I) \to I$ é contínua.
	\end{enumerate}
	
	\textbf{Desenvolvimento Formal:}
	
	\textit{1. Suavidade:} As componentes de $f$ são compostas por funções suaves em $(1, +\infty)$. Logo, $f$ é suave.
	
	\textit{2. Imersão:} Calculamos a diferencial $df_t$, representada pelo vetor velocidade $f'(t)$:
	\begin{align*}
		f'(t) &= \left( -\tfrac{1}{t^2} \cos(2\pi t) - \tfrac{2\pi}{t} \sin(2\pi t), -\tfrac{1}{t^2} \sin(2\pi t) + \tfrac{2\pi}{t} \cos(2\pi t) \right).
	\end{align*}
	Analisando a norma ao quadrado:
	\begin{equation*}
		\|f'(t)\|^2 = \tfrac{1}{t^4} + \tfrac{4\pi^2}{t^2}.
	\end{equation*}
	Como $t \in (1, +\infty)$, $\|f'(t)\|^2 > 0$, logo $f$ é uma imersão.
	
	\textit{3. Injetividade:} Se $f(t_1) = f(t_2)$, então a norma $\|f(t)\|=1/t$ implica $\tfrac{1}{t_1} = \tfrac{1}{t_2} \implies t_1 = t_2$. Logo, $f$ é injetiva.
	
	\textit{4. Homeomorfismo sobre a imagem:} Para provar que $f : I \to f(I)$ é um homeomorfismo, devemos mostrar que a aplicação inversa $f^{-1} : f(I) \to I$ é contínua. 
	Seja $p = (x,y) \in f(I)$. Pela definição de $f(t)$, temos que a norma do vetor é $\|f(t)\| = \sqrt{\tfrac{1}{t^2}(\cos^2(2\pi t) + \sin^2(2\pi t))} = \tfrac{1}{t}$. 
	Portanto, para qualquer ponto $p \in f(I)$, o valor do parâmetro $t$ é univocamente determinado por:
	\begin{equation*}
		t = \frac{1}{\|p\|} = \frac{1}{\sqrt{x^2+y^2}}.
	\end{equation*}
	Assim, a aplicação inversa $f^{-1} : f(I) \to I$ é dada explicitamente por $f^{-1}(x,y) = (x^2+y^2)^{-1/2}$. 
	Como a função norma $\|p\| = \sqrt{x^2+y^2}$ é contínua em $\mathbb{R}^2 \setminus \{0\}$ e a função inversão $z \mapsto 1/z$ é contínua em $(0, \infty)$, a aplicação $f^{-1}$ é uma composição de funções contínuas, sendo, portanto, contínua em todo o seu domínio $f(I)$.
	Sendo $f$ uma bijecão contínua com inversa contínua, conclui-se que $f$ é um homeomorfismo sobre sua imagem.
	\begin{center}
		\begin{tikzpicture}[
			scale=2.5,
			>=Stealth, % Setas elegantes padrão
			% Seta de direção automática integrada na curva
			direcao/.style={
				postaction={decorate},
				decoration={
					markings,
					mark=at position 0.12 with {\arrow{>}}
				}
			}
			]
			% 1. Eixos finos e pretos (sem o quadriculado)
			\draw[->, thin, black] (-1.2,0) -- (1.2,0) node[right, black] {$x$};
			\draw[->, thin, black] (0,-1.2) -- (0,1.2) node[above, black] {$y$};
			\node[below left, black, font=\footnotesize] at (0,0) {$0$};
			
			% 2. A Espiral em Azul
			\draw[thick, blue!70!black, direcao, samples=250, domain=1:6, variable=\t, smooth] 
			plot ({1/\t * cos(360*\t)}, {1/\t * sin(360*\t)});
			
			% 3. Ponto Inicial Aberto f(1) com contorno em azul
			\filldraw[fill=white, draw=blue!70!black, thick] (1,0) circle (1pt);
			\node[above right, font=\footnotesize, black] at (1,0) {$f(1)$};
			
			% 4. Centro da espiral (indicação de convergência assintótica)
			\fill[blue!20] (0,0) circle (0.02);
			
			% 5. Legenda
				% Legenda
			\node at (0, -1.3) {\small $\gamma(t) = \left( \frac{1}{t}\cos(2\pi t), \frac{1}{t}\sin(2\pi t) \right)$};
			
		\end{tikzpicture}
	\end{center}
	
	Sendo $f$ uma imersão injetiva e um homeomorfismo sobre sua imagem, conclui-se que $f$ é um mergulho suave.
}

%%%% QUESTÃO 20 %%%%
\questao{Mostre que a aplicação $g : \mathbb{R} \to \mathbb{R}^2$ dada por $g(t) = (t^3 - 4t, t^2 - 4)$ é uma imersão, mas não é mergulho suave.}

\solucao{
	\textbf{Análise do Enunciado:}
	\begin{itemize}
		\item \textbf{Hipóteses:} $g$ é uma aplicação de $\mathbb{R}$ em $\mathbb{R}^2$ definida por componentes polinomiais.
		\item \textbf{Tese:} Provar que $g$ satisfaz a condição de imersão (diferencial injetiva), mas não satisfaz as condições de mergulho suave (falha na injetividade e/ou na homeomorfia com a imagem).
	\end{itemize}
	
	\textbf{Fundamentação Teórica:}
	\begin{enumerate}
		\item \textbf{Definição de Imersão (Lee, Cap. 7):} $g$ é uma imersão se $dg_t : T_t \mathbb{R} \to T_{g(t)} \mathbb{R}^2$ é injetiva para todo $t \in \mathbb{R}$. Para curvas, isso equivale a $g'(t) \neq (0,0)$ para todo $t$.
		\item \textbf{Definição de Mergulho Suave (Lee, Cap. 7):} Para ser um mergulho suave, $g$ deve ser, obrigatoriamente, uma imersão injetiva e um homeomorfismo sobre sua imagem $g(\mathbb{R})$.
	\end{enumerate}
	
	\textbf{Desenvolvimento Formal:}
	
	\textit{1. Verificação da Imersão:} Calculamos o vetor velocidade $g'(t)$:
	\begin{equation*}
		g'(t) = \left( \frac{d}{dt}(t^3 - 4t), \frac{d}{dt}(t^2 - 4) \right) = (3t^2 - 4, 2t).
	\end{equation*}
	Para que $g'(t) = (0,0)$, deveríamos ter $2t = 0 \implies t = 0$. No entanto, avaliando a primeira componente em $t=0$, temos $3(0)^2 - 4 = -4 \neq 0$. Portanto, $g'(t)$ nunca é o vetor nulo para qualquer $t \in \mathbb{R}$. Logo, $g$ é uma imersão.
	
	\textit{2. Verificação do Mergulho (Injetividade):} Para que $g$ seja um mergulho, ela deve ser injetiva. Vamos testar se existem $t_1 \neq t_2$ tais que $g(t_1) = g(t_2)$:
	\begin{align*}
		t_1^2 - 4 &= t_2^2 - 4 \implies t_1^2 = t_2^2 \implies t_1 = -t_2 \\
		t_1^3 - 4t_1 &= t_2^3 - 4t_2.
	\end{align*}
	Substituindo $t_1 = -t_2$ na segunda equação:
	\begin{equation*}
		(-t_2)^3 - 4(-t_2) = t_2^3 - 4t_2 \implies -t_2^3 + 4t_2 = t_2^3 - 4t_2 \implies 2t_2^3 - 8t_2 = 0.
	\end{equation*}
	Isso implica $2t_2(t_2^2 - 4) = 0$, resultando em $t_2 = 0, 2, -2$. 
	Se tomarmos $t_1 = -2$ e $t_2 = 2$, observamos que:
	\begin{equation*}
		g(-2) = ((-2)^3 - 4(-2), (-2)^2 - 4) = (0,0) \quad \text{e} \quad g(2) = (2^3 - 4(2), 2^2 - 4) = (0,0).
	\end{equation*}
	Como $g(-2) = g(2)$ mas $-2 \neq 2$, a aplicação $g$ não é injetiva.
	
	\textit{3. Conclusão:} Uma aplicação que não é injetiva não pode ser um homeomorfismo sobre sua imagem. Sendo assim, embora $g$ seja uma imersão, ela não é um mergulho suave.
	\begin{center}
		\begin{tikzpicture}[
			% Escala ideal para caber na página
			scale=0.8, 
			>=Stealth
			]
			% --- 1. Eixos finos e pretos ---
			\draw[->, thin, black] (-4.5,0) -- (4.5,0) node[right, black, font=\footnotesize] {$x$};
			\draw[->, thin, black] (0,-4.5) -- (0,2.5) node[above, black, font=\footnotesize] {$y$};
			\node[below left, black, font=\footnotesize, shift={(0.1,-0.1)}] at (0,0) {$0$};
			
			% --- 2. A Curva g(t) em Azul Estilizado ---
			% Sem a dependência de 'direcao' (decorations), eliminando o erro de tamanho permanentemente
			\draw[thick, blue!70!black, samples=200, domain=-2.3:2.3, variable=\t, smooth] 
			plot ({\t*\t*\t - 4*\t}, {\t*\t - 4});
			
			% --- 3. Seta de Direção Perfeita (Desenhada sobre a própria equação) ---
			% Plota um micro-segmento na subida do lado esquerdo garantindo alinhamento e toque perfeitos
			\draw[->, thick, blue!70!black, domain=-1.25:-1.22, variable=\t] 
			plot ({\t*\t*\t - 4*\t}, {\t*\t - 4});
			
			% --- 4. Pontos Chave para Rigor Matemático e Pedagógico ---
			
			% Ponto de Autointersecção 
			\filldraw[black] (0,0) circle (1.5pt); 
			\node[above, font=\footnotesize, black, shift={(0, 0.3)}] at (0,0) {$g(2)=g(-2)$};
			
			% Ponto mais baixo do loop
			\filldraw[black] (0,-4) circle (1.5pt);
			\node[below right, font=\footnotesize, black] at (0,-4) {$g(0)$};
			
			% --- 5. Legenda ---
			\node at (0, -5.5) {\small $g(t) = (t^3 - 4t, \, t^2 - 4)$};
			
		\end{tikzpicture}
	\end{center}
}

%%%% QUESTÃO 21 %%%%
\questao{Considere a aplicação suave $X : \mathbb{R}^2 \to \mathbb{R}^3$ dada por $X(\varphi, \theta) = ( (2 + \cos \varphi) \cos \theta, (2 + \cos \varphi) \sin \theta, \sin \varphi )$. Mostre que $X$ é uma imersão cuja imagem é a superfície obtida pela rotação do círculo $(y-2)^2 + z^2 = 1$, contido no plano $yz$, ao redor do eixo $z$.}

\solucao{
	\textbf{Análise do Enunciado:}
	\begin{itemize}
		\item \textbf{Hipóteses:} $X$ é uma aplicação de $\mathbb{R}^2$ em $\mathbb{R}^3$; as componentes são funções suaves envolvendo senos e cossenos.
		\item \textbf{Tese:} 
		\begin{enumerate}[label=(\roman*)]
			\item Provar que $X$ é uma imersão (a diferencial $dX_p$ é injetiva para todo $p \in \mathbb{R}^2$).
			\item Provar que a imagem $X(\mathbb{R}^2)$ coincide com a superfície de revolução do círculo $(y-2)^2 + z^2 = 1$ no plano $yz$ em torno do eixo $z$.
		\end{enumerate}
	\end{itemize}
	
	\textbf{Fundamentação Teórica:}
	\begin{enumerate}
		\item \textbf{Definição de Imersão (Lee, Cap. 7):} Uma aplicação $X: M \to N$ é uma imersão se a diferencial $dX_p : T_p M \to T_{X(p)} N$ é injetiva para todo $p \in M$. Para aplicações entre espaços euclidianos, isso equivale a verificar que a matriz Jacobiana possui posto igual à dimensão do domínio ($\text{posto} = 2$).
		\item \textbf{Superfícies de Revolução:} Uma superfície obtida pela rotação de uma curva $\mathcal{C}$ no plano $yz$ ao redor do eixo $z$ é dada por pontos $(x,y,z)$ tais que $\sqrt{x^2+y^2}$ e $z$ satisfaçam a equação da curva $\mathcal{C}$ no plano $yz$ (onde $y$ é substituído por $\sqrt{x^2+y^2}$).
	\end{enumerate}
	
	\textbf{Desenvolvimento Formal:}
	
	\textit{1. Prova da Imersão:}
	Calculamos a matriz Jacobiana de $X(\varphi, \theta)$ em relação às coordenadas $(\varphi, \theta)$:
	\begin{equation*}
		dX = \begin{pmatrix} 
			\tfrac{\partial X_1}{\partial \varphi} & \tfrac{\partial X_1}{\partial \theta} \\ 
			\tfrac{\partial X_2}{\partial \varphi} & \tfrac{\partial X_2}{\partial \theta} \\ 
			\tfrac{\partial X_3}{\partial \varphi} & \tfrac{\partial X_3}{\partial \theta} 
		\end{pmatrix} = 
		\begin{pmatrix} 
			-\sin \varphi \cos \theta & -(2 + \cos \varphi) \sin \theta \\ 
			-\sin \varphi \sin \theta & (2 + \cos \varphi) \cos \theta \\ 
			\cos \varphi & 0 
		\end{pmatrix}.
	\end{equation*}
	Para que $X$ seja uma imersão, as duas colunas da matriz devem ser linearmente independentes para todo $(\varphi, \theta)$. Consideremos o produto vetorial das colunas $\mathbf{v}_1$ e $\mathbf{v}_2$:
	\begin{align*}
		\mathbf{v}_1 \times \mathbf{v}_2 &= 
		\begin{vmatrix} 
			\mathbf{i} & \mathbf{j} & \mathbf{k} \\ 
			-\sin \varphi \cos \theta & -\sin \varphi \sin \theta & \cos \varphi \\ 
			-(2 + \cos \varphi) \sin \theta & (2 + \cos \varphi) \cos \theta & 0 
		\end{vmatrix} \\
		&= \left( -(2 + \cos \varphi)\cos \varphi \cos \theta, -(2 + \cos \varphi)\cos \varphi \sin \theta, -(2 + \cos \varphi)\sin \varphi (\cos^2 \theta + \sin^2 \theta) \right) \\
		&= -(2 + \cos \varphi) \left( \cos \varphi \cos \theta, \cos \varphi \sin \theta, \sin \varphi \right).
	\end{align*}
	A norma desse vetor é $\|\mathbf{v}_1 \times \mathbf{v}_2\| = |2 + \cos \varphi| \sqrt{\cos^2 \varphi (\cos^2 \theta + \sin^2 \theta) + \sin^2 \varphi} = |2 + \cos \varphi|$.
	Como $-1 \leq \cos \varphi \leq 1$, temos que $2 + \cos \varphi \geq 1 > 0$ para todo $\varphi$. Logo, o produto vetorial nunca é nulo, as colunas são linearmente independentes e $X$ é, portanto, uma imersão.
	
	\textit{2. Análise da Imagem:}
	Sejam $(x,y,z)$ as coordenadas de um ponto na imagem de $X$. Temos:
	\begin{equation*}
		x = (2 + \cos \varphi) \cos \theta, \quad y = (2 + \cos \varphi) \sin \theta, \quad z = \sin \varphi.
	\end{equation*}
	Calculamos a distância do ponto ao eixo $z$:
	\begin{equation*}
		\sqrt{x^2 + y^2} = \sqrt{(2 + \cos \varphi)^2 (\cos^2 \theta + \sin^2 \theta)} = |2 + \cos \varphi| = 2 + \cos \varphi.
	\end{equation*}
	Agora, observemos a relação entre $\sqrt{x^2+y^2}$ e $z$:
	\begin{equation*}
		(\sqrt{x^2+y^2} - 2)^2 + z^2 = (\cos \varphi)^2 + (\sin \varphi)^2 = 1.
	\end{equation*}
	Esta equação descreve exatamente a superfície obtida ao rotacionar o círculo $(\rho - 2)^2 + z^2 = 1$ (onde $\rho$ é a distância ao eixo $z$) ao redor do eixo $z$. No plano $yz$ (onde $x=0$), a distância $\rho$ é simplesmente $|y|$. Assim, a curva geradora é $(|y|-2)^2 + z^2 = 1$, que corresponde ao círculo $(y-2)^2 + z^2 = 1$ e sua imagem refletida. Portanto, a imagem de $X$ é o toro de revolução descrito.
}

%%%% QUESTÃO 22 %%%%
\questao{Considere a aplicação suave $F : \mathbb{R}^3 \to \mathbb{R}^3$ dada por $F(\rho, \varphi, \theta) = (\rho \sin \varphi \cos \theta, \rho \sin \varphi \sin \theta, \rho \cos \varphi)$.
	\begin{enumerate}[label=(\alph*)]
		\item Encontre um conjunto aberto $U \subset \mathbb{R}^3$ de tal forma que $F : U \to \mathbb{R}^3$ seja um difeomorfismo local.
		\item Encontre um conjunto aberto $U_0 \subset U$ de tal forma que $F : U_0 \to F(U_0)$ seja um difeomorfismo.
\end{enumerate}}

\solucao{
	\textbf{Análise do Enunciado:}
	\begin{itemize}
		\item \textbf{Hipóteses:} $F$ é a aplicação de coordenadas esféricas de $\mathbb{R}^3$ em $\mathbb{R}^3$. O domínio original é $\mathbb{R}^3$ com coordenadas $(\rho, \varphi, \theta)$.
		\item \textbf{Tese:} 
		\begin{enumerate}[label=(\alph*)]
			\item Determinar um aberto $U$ onde a diferencial $dF_p$ seja inversível para todo $p \in U$.
			\item Determinar um subconjunto $U_0 \subset U$ onde $F$ seja injetiva e a inversa $F^{-1}$ seja suave.
		\end{enumerate}
	\end{itemize}
	
	\textbf{Fundamentação Teórica:}
	\begin{enumerate}
		\item \textbf{Teorema da Função Inversa (Lee, Cap. 7):} Uma aplicação suave $F : M \to N$ entre variedades de mesma dimensão é um difeomorfismo local no ponto $p \in M$ se, e somente se, a diferencial $dF_p : T_p M \to T_{F(p)} N$ é um isomorfismo (ou seja, se a matriz Jacobiana $\det(JF_p) \neq 0$).
		\item \textbf{Difeomorfismo:} Uma aplicação $F : U \to V$ é um difeomorfismo se for suave, bijetiva e possuir inversa suave. Se $F$ é um difeomorfismo local em todo ponto de $U$ e é injetiva em $U$, então $F$ é um difeomorfismo sobre sua imagem.
	\end{enumerate}
	
	\textbf{Desenvolvimento Formal:}
	
	\textit{Parte (a): Condição para Difeomorfismo Local.}
	Para determinar onde $F$ é um difeomorfismo local, calculamos a matriz Jacobiana de $F$ em relação às coordenadas $(\rho, \varphi, \theta)$:
	\begin{equation*}
		JF = \begin{pmatrix} 
			\tfrac{\partial x}{\partial \rho} & \tfrac{\partial x}{\partial \varphi} & \tfrac{\partial x}{\partial \theta} \\ 
			\tfrac{\partial y}{\partial \rho} & \tfrac{\partial y}{\partial \varphi} & \tfrac{\partial y}{\partial \theta} \\ 
			\tfrac{\partial z}{\partial \rho} & \tfrac{\partial z}{\partial \varphi} & \tfrac{\partial z}{\partial \theta} 
		\end{pmatrix} = 
		\begin{pmatrix} 
			\sin \varphi \cos \theta & \rho \cos \varphi \cos \theta & -\rho \sin \varphi \sin \theta \\ 
			\sin \varphi \sin \theta & \rho \cos \varphi \sin \theta & \rho \sin \varphi \cos \theta \\ 
			\cos \varphi & -\rho \sin \varphi & 0 
		\end{pmatrix}.
	\end{equation*}
	Calculamos o determinante $\det(JF)$ expandindo pela terceira linha:
	\begin{align*}
		\det(JF) &= \cos \varphi \left[ (\rho \cos \varphi \cos \theta)(\rho \sin \varphi \cos \theta) - (-\rho \sin \varphi \sin \theta)(\rho \cos \varphi \sin \theta) \right] \\
		&\quad - (-\rho \sin \varphi) \left[ (\sin \varphi \cos \theta)(\rho \sin \varphi \cos \theta) - (-\rho \sin \varphi \sin \theta)(\sin \varphi \sin \theta) \right] \\
		&= \cos \varphi \left[ \rho^2 \sin \varphi \cos \varphi (\cos^2 \theta + \sin^2 \theta) \right] + \rho \sin \varphi \left[ \rho \sin^2 \varphi (\cos^2 \theta + \sin^2 \theta) \right] \\
		&= \rho^2 \sin \varphi \cos^2 \varphi + \rho^2 \sin^3 \varphi \\
		&= \rho^2 \sin \varphi (\cos^2 \varphi + \sin^2 \varphi) = \rho^2 \sin \varphi.
	\end{align*}
	Pelo Teorema da Função Inversa, $F$ é um difeomorfismo local onde $\det(JF) \neq 0$. Isso ocorre quando $\rho \neq 0$ e $\sin \varphi \neq 0$ (ou seja, $\varphi \neq k\pi, k \in \mathbb{Z}$). Definimos, portanto, o conjunto aberto:
	\begin{equation*}
		U = \{ (\rho, \varphi, \theta) \in \mathbb{R}^3 \mid \rho \in (0, \infty), \varphi \in (0, \pi) \cup (\pi, 2\pi), \theta \in \mathbb{R} \}.
	\end{equation*}
	
	\textit{Parte (b): Condição para Difeomorfismo Global.}
	Para que $F$ seja um difeomorfismo sobre sua imagem, ela deve ser, além de um difeomorfismo local, uma aplicação injetiva. Analisamos as redundâncias de $F$:
	
	\begin{enumerate}
		\item \textbf{Periodicidade de $\theta$:} Note que as funções $\cos \theta$ e $\sin \theta$ são $2\pi$-periódicas. Assim, para qualquer $k \in \mathbb{Z}$, temos $F(\rho, \varphi, \theta) = F(\rho, \varphi, \theta + 2\pi k)$. Para garantir a injetividade, devemos restringir $\theta$ a um intervalo de comprimento $2\pi$, como $(0, 2\pi)$.
		
		\item \textbf{Simetria de $\varphi$:} A coordenada $\varphi$ define o ângulo com o eixo $z$. Se tomarmos $\overline{\varphi} = 2\pi - \varphi$, temos $\sin \overline{\varphi} = -\sin \varphi$ e $\cos \overline{\varphi} = \cos \varphi$. Note que:
		\begin{align*}
			F(\rho, 2\pi - \varphi, \theta + \pi) &= (\rho \sin(2\pi - \varphi) \cos(\theta + \pi), \rho \sin(2\pi - \varphi) \sin(\theta + \pi), \rho \cos(2\pi - \varphi)) \\
			&= (\rho (-\sin \varphi) (-\cos \theta), \rho (-\sin \varphi) (-\sin \theta), \rho \cos \varphi) \\
			&= (\rho \sin \varphi \cos \theta, \rho \sin \varphi \sin \theta, \rho \cos \varphi) = F(\rho, \varphi, \theta).
		\end{align*}
		Portanto, valores de $\varphi$ fora do intervalo $(0, \pi)$ podem gerar a mesma imagem que valores dentro desse intervalo, mudando-se $\theta$. Para evitar isso, restringimos $\varphi \in (0, \pi)$.
	\end{enumerate}
	
	Definimos, portanto, o domínio restringido:
	\begin{equation*}
		U_0 = \{ (\rho, \varphi, \theta) \in \mathbb{R}^3 \mid \rho \in (0, \infty), \varphi \in (0, \pi), \theta \in (0, 2\pi) \}.
	\end{equation*}
	
	\textbf{Verificação da Injetividade e Difeomorfismo Global:}
	
	Consideremos a restrição da aplicação ao domínio $U_0$.
	\begin{enumerate}
		\item Como $U_0 \subset U$, e já provamos que $F$ é um difeomorfismo local em todo ponto de $U$, segue que $F$ permanece um difeomorfismo local em cada ponto de $U_0$.
		
		\item Para provar a injetividade, sejam $p = (\rho, \varphi, \theta)$ e $p' = (\rho', \varphi', \theta')$ dois pontos em $U_0$ tais que $F(p) = F(p')$. Denotemos a imagem por $(x,y,z)$. Analisamos as componentes:
		
		Primeiro, calculamos a norma euclidiana da imagem:
		\begin{align*}
			\|F(p)\|^2 &= x^2 + y^2 + z^2 \\
			&= (\rho \sin \varphi \cos \theta)^2 + (\rho \sin \varphi \sin \theta)^2 + (\rho \cos \varphi)^2 \\
			&= \rho^2 \sin^2 \varphi (\cos^2 \theta + \sin^2 \theta) + \rho^2 \cos^2 \varphi \\
			&= \rho^2 (\sin^2 \varphi + \cos^2 \varphi) = \rho^2.
		\end{align*}
		Como $\rho, \rho' > 0$, a igualdade $F(p) = F(p')$ implica $\rho^2 = (\rho')^2$, logo $\rho = \rho'$.
		
		Em seguida, utilizamos a terceira componente $z$:
		\begin{equation*}
			z = \rho \cos \varphi \implies \cos \varphi = \tfrac{z}{\rho}.
		\end{equation*}
		Dado que $\varphi \in (0, \pi)$, a função cosseno é estritamente decrescente e, portanto, injetiva neste intervalo. Assim, $\cos \varphi = \cos \varphi' \implies \varphi = \varphi'$.
		
		Finalmente, analisamos as componentes $x$ e $y$:
		\begin{equation*}
			x = \rho \sin \varphi \cos \theta, \quad y = \rho \sin \varphi \sin \theta.
		\end{equation*}
		Como $\rho > 0$ e $\sin \varphi > 0$ para $\varphi \in (0, \pi)$, a razão entre as componentes é:
		\begin{equation*}
			\frac{y}{x} = \frac{\rho \sin \varphi \sin \theta}{\rho \sin \varphi \cos \theta} = \tan \theta.
		\end{equation*}
		Visto que $\theta \in (0, 2\pi)$ e temos os sinais de $\sin \theta$ e $\cos \theta$ definidos por $y$ e $x$, o valor de $\theta$ é univocamente determinado. Logo, $\theta = \theta'$.
	\end{enumerate}
	
	Sendo $F$ injetiva em $U_0$ e um difeomorfismo local em cada ponto de $U_0$, conclui-se, pelo Teorema da Função Inversa, que $F : U_0 \to F(U_0)$ é um difeomorfismo.
}

%%%% QUESTÃO 23 %%%%
\questao{Seja $\pi : M \to N$ uma submersão.
	\begin{enumerate}[label=(\alph*)]
		\item Mostre que $\pi$ é uma aplicação aberta.
		\item Mostre que qualquer ponto de $M$ está na imagem de uma seção local suave de $\pi$.
\end{enumerate}}

\solucao{
	\textbf{Análise do Enunciado:}
	\begin{itemize}
		\item \textbf{Hipóteses:} $\pi: M \to N$ é uma submersão entre variedades suaves.
		
		Isso implica que a diferencial $d\pi_p : T_p M \to T_{\pi(p)} N$ é sobrejetiva para todo $p \in M$.
		\item \textbf{Tese:} 
		\begin{enumerate}[label=(\alph*)]
			\item Provar que, para qualquer aberto $W \subset M$, a imagem $\pi(W)$ é um aberto em $N$.
			\item Para qualquer $p \in M$, provar a existência de um aberto $U \subset N$ contendo $\pi(p)$ e uma aplicação suave $s : U \to M$ tal que $\pi \circ s = \text{id}_U$ e $s(U)$ contenha $p$.
		\end{enumerate}
	\end{itemize}
	
	\textbf{Fundamentação Teórica:}
	\begin{enumerate}
		\item \textbf{Teorema da Forma Canônica Local para Submersões (Lee, Cap. 7):} Seja $\pi : M \to N$ uma submersão com $\dim M = m$ e $\dim N = n$. Para cada $p \in M$, existem cartas coordenadas $(U, \varphi)$ em torno de $p$ e $(V, \psi)$ em torno de $\pi(p)$ tais que a representação local de $\pi$ é a projeção canônica:
		\begin{equation*}
			\hat{\pi}(x^1, \dots, x^m) = (x^1, \dots, x^n).
		\end{equation*}
		\item \textbf{Propriedades de Projeções:} A projeção canônica $\text{proj} : \mathbb{R}^m \to \mathbb{R}^n$ é uma aplicação aberta.
		\item \textbf{Definição de Seção Local:} Uma aplicação suave $s : U \to M$ é uma seção local de $\pi$ se $\pi \circ s = \text{id}_U$.
	\end{enumerate}
	
	\textbf{Desenvolvimento Formal:}
	
	\textit{Parte (a): $\pi$ é uma aplicação aberta.}
	Seja $W \subset M$ um conjunto aberto. Para mostrar que $\pi(W)$ é aberto em $N$, devemos mostrar que para cada $q \in \pi(W)$, existe uma vizinhança aberta de $q$ contida em $\pi(W)$.
	Seja $p \in W$ tal que $\pi(p) = q$. Pelo Teorema da Forma Canônica Local, existem cartas coordenadas $\varphi : U \to \mathbb{R}^m$ e $\psi : V \to \mathbb{R}^n$ tais que a representação local de $\pi$ é a projeção $\hat{\pi} : \mathbb{R}^m \to \mathbb{R}^n$. 
	
	Podemos escolher $U$ pequeno o suficiente para que $U \subset W$. A imagem de $U$ sob $\pi$ é $\pi(U) = \psi^{-1}(\hat{\pi}(\varphi(U)))$. Como $\varphi(U)$ é um aberto em $\mathbb{R}^m$ e $\hat{\pi}$ é uma projeção canônica (que é uma aplicação aberta), $\hat{\pi}(\varphi(U))$ é um aberto em $\mathbb{R}^n$. Como $\psi^{-1}$ é um difeomorfismo, $\pi(U)$ é, portanto, um aberto em $N$.
	Como $q \in \pi(U) \subset \pi(W)$, provamos que cada ponto de $\pi(W)$ possui uma vizinhança aberta contida em $\pi(W)$. Logo, $\pi(W)$ é aberto.
	
	\textit{Parte (b): Existência de Seções Locais.}
	Seja $p \in M$ e denote $q = \pi(p)$. Utilizando novamente o Teorema da Forma Canônica Local, escolhemos cartas $\varphi : U \to \mathbb{R}^m$ e $\psi : V \to \mathbb{R}^n$ centradas em $p$ e $q$, respectivamente, tais que a representação local de $\pi$ seja:
	\begin{equation*}
		\hat{\pi}(x^1, \dots, x^m) = (x^1, \dots, x^n).
	\end{equation*}
	Definimos a aplicação $\hat{s} : \psi(V) \to \varphi(U)$ no espaço euclidiano da seguinte forma:
	\begin{equation*}
		\hat{s}(y^1, \dots, y^n) = (y^1, \dots, y^n, 0, \dots, 0).
	\end{equation*}
	Esta aplicação é claramente suave e satisfaz $\hat{\pi} \circ \hat{s} = \text{id}_{\mathbb{R}^n}$. Agora, definimos a seção local $s : V \to M$ através da composição:
	\begin{equation*}
		s = \varphi^{-1} \circ \hat{s} \circ \psi.
	\end{equation*}
	Como $s$ é a composição de aplicações suaves, $s$ é suave. Além disso:
	\begin{equation*}
		\pi \circ s = \pi \circ (\varphi^{-1} \circ \hat{s} \circ \psi) = (\pi \circ \varphi^{-1}) \circ \hat{s} \circ \psi = (\psi^{-1} \circ \hat{\pi}) \circ \hat{s} \circ \psi = \psi^{-1} \circ (\hat{\pi} \circ \hat{s}) \circ \psi = \psi^{-1} \circ \text{id} \circ \psi = \text{id}_V.
	\end{equation*}
	Visto que as cartas estão centradas em $p$ e $q$, temos $\psi(q) = 0$ e $\hat{s}(0) = 0$, logo $s(q) = \varphi^{-1}(0) = p$. Assim, $p$ está na imagem da seção local suave $s$.
}

%%%% QUESTÃO 24 %%%%
\questao{Sejam $M, N$ e $P$ variedades suaves, $\pi : M \to N$ uma submersão sobrejetiva e $F : N \to P$ qualquer aplicação. Mostre que $F$ é suave se, e somente se, $F \circ \pi$ é suave.}

\solucao{
	\textbf{Análise do Enunciado:}
	\begin{itemize}
		\item \textbf{Hipóteses:} $M, N, P$ são variedades suaves; $\pi: M \to N$ é uma submersão sobrejetiva; $F: N \to P$ é uma aplicação.
		\item \textbf{Tese:} Provar a equivalência lógica: $F \in C^\infty(N, P) \iff F \circ \pi \in C^\infty(M, P)$.
	\end{itemize}
	
	\textbf{Fundamentação Teórica:}
	\begin{enumerate}
		\item \textbf{Composição de Aplicações Suaves (Lee, Cap. 2):} Se $f: M \to N$ e $g: N \to P$ são aplicações suaves, então a composição $g \circ f: M \to P$ é suave.
		\item \textbf{Teorema da Forma Canônica Local para Submersões (Lee, Cap. 7):} Para qualquer $p \in M$, existem cartas coordenadas $(U, \varphi)$ em torno de $p$ e $(V, \psi)$ em torno de $q = \pi(p)$ tais que a representação local de $\pi$ é a projeção canônica:
		\begin{equation*}
			\hat{\pi}(x^1, \dots, x^m) = (x^1, \dots, x^n).
		\end{equation*}
		\item \textbf{Suavidade Local:} Uma aplicação $F: N \to P$ é suave se, e somente se, para cada $q \in N$, existe uma carta $(V, \psi)$ em torno de $q$ tal que a representação local $\hat{F} = \eta \circ F \circ \psi^{-1}$ seja suave (onde $\eta$ é uma carta em $P$).
	\end{enumerate}
	
	\textbf{Desenvolvimento Formal:}
	
	\textit{Direção ($\implies$):} 
	Suponhamos que $F$ seja suave. Como $\pi$ é uma submersão, ela é, por definição, uma aplicação suave. Pela propriedade da composição de aplicações suaves (item 1 da fundamentação), a composição $F \circ \pi$ é necessariamente suave.
	
	\textit{Direção ($\impliedby$):}
	Suponhamos que $F \circ \pi$ seja suave. Para provar que $F$ é suave, devemos mostrar que ela é suave localmente em cada ponto $q \in N$. 
	Como $\pi$ é sobrejetiva, para qualquer $q \in N$, existe $p \in M$ tal que $\pi(p) = q$. Pelo Teorema da Forma Canônica Local (item 2), existem cartas coordenadas $\varphi: U \to \mathbb{R}^m$ e $\psi: V \to \mathbb{R}^n$ centradas em $p$ e $q$, respectivamente, tais que a representação local de $\pi$ é a projeção $\hat{\pi}(x^1, \dots, x^m) = (x^1, \dots, x^n)$.
	
	Seja $\eta: W \to \mathbb{R}^p$ uma carta coordenada em $P$ centrada em $F(q)$. Consideremos a representação local de $F$ em $V$:
	\begin{equation*}
		\hat{F} = \eta \circ F \circ \psi^{-1} : \psi(V) \to \mathbb{R}^p.
	\end{equation*}
	Agora, analisemos a representação local da composição $F \circ \pi$ em $U$:
	\begin{align*}
		\widehat{F \circ \pi} &= \eta \circ (F \circ \pi) \circ \varphi^{-1} \\
		&= \eta \circ F \circ (\pi \circ \varphi^{-1}) \\
		&= (\eta \circ F \circ \psi^{-1}) \circ (\psi \circ \pi \circ \varphi^{-1}) \\
		&= \hat{F} \circ \hat{\pi}.
	\end{align*}
	Sabemos, por hipótese, que $F \circ \pi$ é suave, portanto $\widehat{F \circ \pi}$ é uma aplicação suave de $\mathbb{R}^m$ em $\mathbb{R}^p$. Explicitamente, para qualquer $(x^1, \dots, x^m) \in \varphi(U)$:
	\begin{equation*}
		\widehat{F \circ \pi}(x^1, \dots, x^m) = \hat{F}(x^1, \dots, x^n).
	\end{equation*}
	Note que o valor da aplicação $\widehat{F \circ \pi}$ depende apenas das primeiras $n$ coordenadas. Como $\widehat{F \circ \pi}$ é suave em $\mathbb{R}^m$, sua restrição ao subespaço $\mathbb{R}^n \times \{0\}^{m-n}$, que coincide exatamente com a aplicação $\hat{F}$, também deve ser suave.
	
	Como $\hat{F}$ é suave para qualquer escolha de $q \in N$, a aplicação $F$ é suave em toda a variedade $N$.
}

%%%% QUESTÃO 25 %%%%
\questao{(Teorema do Posto) Sejam $U \subset \mathbb{R}^m$ e $V \subset \mathbb{R}^n$ subconjuntos abertos e $F : U \to V$ uma aplicação suave com posto constante $k$. Mostre que, para cada $p \in U$, existem cartas coordenadas suaves $(U_0, \varphi)$ para $\mathbb{R}^m$ centrada em $p$ e $(V_0, \psi)$ para $\mathbb{R}^n$, com $U_0 \subset U$ e $F(U_0) \subset V_0 \subset V$, tal que
	\begin{equation*}
		\psi \circ F \circ \varphi^{-1}(x^1, \dots, x^k, x^{k+1}, \dots, x^m) = (x^1, \dots, x^k, 0, \dots, 0).
	\end{equation*}
}

\solucao{
	\textbf{Análise do Enunciado:}
	\begin{itemize}
		\item \textbf{Hipóteses:} $F: U \to V$ é uma aplicação suave entre abertos de espaços euclidianos; o posto da diferencial $dF_q$ é constante e igual a $k$ para todo $q \in U$.
		\item \textbf{Tese:} A existência de coordenadas locais (difeomorfismos $\varphi$ e $\psi$) que transformam $F$ na projeção linear canônica de $\mathbb{R}^m$ em $\mathbb{R}^n$ sobre as primeiras $k$ componentes.
	\end{itemize}
	
	\textbf{Fundamentação Teórica:}
	\begin{enumerate}
		\item \textbf{Teorema da Função Inversa (Lee, Cap. 7):} Se a diferencial de uma aplicação suave em um ponto é um isomorfismo, então a aplicação é um difeomorfismo local em torno desse ponto.
		\item \textbf{Mudança de Base Linear:} Se uma aplicação linear $L: \mathbb{R}^m \to \mathbb{R}^n$ possui posto $k$, existem bases de $\mathbb{R}^m$ e $\mathbb{R}^n$ tais que a representação matricial de $L$ é a matriz bloco $\begin{pmatrix} I_k & 0 \\ 0 & 0 \end{pmatrix}$, onde $I_k$ é a identidade de ordem $k$.
	\end{enumerate}
	
	\textbf{Desenvolvimento Formal:}
	
	Seja $p \in U$ e $q = F(p)$. Sem perda de generalidade, podemos assumir que $p=0$ e $q=0$ através de translações. Pela hipótese de posto constante $k$, a diferencial $dF_0 : \mathbb{R}^m \to \mathbb{R}^n$ possui posto $k$. Pela álgebra linear, podemos escolher bases de $\mathbb{R}^m$ e $\mathbb{R}^n$ tais que a representação matricial de $dF_0$ seja a projeção canônica. Assim, após uma mudança de coordenadas linear, podemos assumir que:
	\begin{equation*}
		dF_0(v^1, \dots, v^m) = (v^1, \dots, v^k, 0, \dots, 0).
	\end{equation*}
	
	Agora, definimos uma nova aplicação $\Phi : U \to \mathbb{R}^m$ para construir a carta $\varphi$. Seja:
	\begin{equation*}
		\Phi(x^1, \dots, x^m) = (F^1(x), \dots, F^k(x), x^{k+1}, \dots, x^m),
	\end{equation*}
	onde $F^i$ são as primeiras $k$ componentes de $F$. Calculamos a matriz Jacobiana de $\Phi$ no ponto $0$:
	\begin{equation*}
		J\Phi_0 = \begin{pmatrix} 
			\tfrac{\partial F^1}{\partial x^1} & \dots & \tfrac{\partial F^1}{\partial x^m} \\ 
			\vdots & \ddots & \vdots \\ 
			\tfrac{\partial F^k}{\partial x^k} & \dots & \tfrac{\partial F^k}{\partial x^m} \\ 
			0 & \dots & 1 & \dots & 0 \\
			\vdots & \ddots & \vdots & \ddots & \vdots \\
			0 & \dots & 0 & \dots & 1
		\end{pmatrix} = \begin{pmatrix} 
			I_k & A \\ 
			0 & I_{m-k} 
		\end{pmatrix}.
	\end{equation*}
	O determinante de $J\Phi_0$ é $1 \neq 0$. Pelo Teorema da Função Inversa, $\Phi$ é um difeomorfismo local em torno de $0$. Definimos, portanto, a carta $\varphi = \Phi$ em um aberto $U_0$ pequeno o suficiente.
	
	Agora, analisamos a composição $F \circ \varphi^{-1}$. Seja $y = (y^1, \dots, y^m)$ a coordenada em $\varphi(U_0)$. Por definição de $\Phi$, temos:
	\begin{equation*}
		F(\varphi^{-1}(y^1, \dots, y^m)) = (y^1, \dots, y^k, F^{k+1}(\varphi^{-1}(y)), \dots, F^n(\varphi^{-1}(y))).
	\end{equation*}
	Seja $G : \varphi(U_0) \to \mathbb{R}^n$ a aplicação $F \circ \varphi^{-1}$. O posto de $G$ é $k$ (pois o posto de $F$ é $k$), mas as primeiras $k$ componentes de $G$ são a identidade. Isso implica que as componentes $G^{k+1}, \dots, G^n$ devem ser constantes. Como $F(0)=0$, essas constantes são todas zero. Assim:
	\begin{equation*}
		G(y^1, \dots, y^m) = (y^1, \dots, y^k, 0, \dots, 0).
	\end{equation*}
	
	Para finalizar, se as componentes de $G$ já estão na forma canônica, a carta $\psi$ em $V$ pode ser a própria identidade $\psi = \text{id}_{\mathbb{R}^n}$ (ou qualquer difeomorfismo que preserve a origem). Logo:
	\begin{equation*}
		\psi \circ F \circ \varphi^{-1}(x^1, \dots, x^m) = (x^1, \dots, x^k, 0, \dots, 0).
	\end{equation*}
	A tese está provada.
}

%%%% QUESTÃO 26 %%%%
\questao{(Teorema do Posto para Variedades) Sejam $M^m$ e $N^n$ variedades suaves e $F : M^m \to N^n$ uma aplicação suave de posto constante $k$. Mostre que, para cada $p \in M$, existem coordenadas suaves $(x^1, \dots, x^m)$ centrada em $p$ e $(y^1, \dots, y^n)$ centrada em $F(p)$ tais que $F$ admite a seguinte representação coordenada:
	\begin{equation*}
		F(x^1, \dots, x^k, x^{k+1}, \dots, x^m) = (x^1, \dots, x^k, 0, \dots, 0).
	\end{equation*}
}

\solucao{
	\textbf{Análise do Enunciado:}
	\begin{itemize}
		\item \textbf{Hipóteses:} $M^m$ e $N^n$ são variedades suaves; $F: M \to N$ é uma aplicação suave; o posto da diferencial $dF_q$ é constante e igual a $k$ para todo $q \in M$.
		\item \textbf{Tese:} Provar a existência de cartas coordenadas locais em $p \in M$ e $F(p) \in N$ tais que a representação local da aplicação seja a projeção canônica sobre as primeiras $k$ coordenadas.
	\end{itemize}
	
	\textbf{Fundamentação Teórica:}
	\begin{enumerate}
		\item \textbf{Definição de Coordenadas Locais (Lee, Cap. 1):} Por definição de variedade suave, para qualquer ponto $p \in M$, existe uma carta coordenada $(U, \varphi)$ tal que $\varphi : U \to \mathbb{R}^m$ é um difeomorfismo sobre um aberto de $\mathbb{R}^m$.
		\item \textbf{Representação Local de uma Aplicação (Lee, Cap. 2):} Dada uma aplicação $F: M \to N$ e cartas coordenadas $\varphi: U \to \mathbb{R}^m$ e $\psi: V \to \mathbb{R}^n$ com $F(U) \subset V$, a representação local de $F$ é a aplicação suave $\hat{F} = \psi \circ F \circ \varphi^{-1} : \varphi(U) \to \psi(V)$.
		\item \textbf{Posto de uma Aplicação (Lee, Cap. 7):} O posto de $F$ no ponto $p$ é definido como o posto da aplicação linear $dF_p$. A representação local $\hat{F}$ preserva o posto, ou seja, $\text{posto}(dF_p) = \text{posto}(d\hat{F}_{\varphi(p)})$.
		\item \textbf{Teorema do Posto para Espaços Euclidianos (Questão 25, Lista 03):} Se $\hat{F} : \mathbb{R}^m \to \mathbb{R}^n$ é uma aplicação suave de posto constante $k$, existem difeomorfismos locais $\alpha: \mathbb{R}^m \to \mathbb{R}^m$ e $\beta: \mathbb{R}^n \to \mathbb{R}^n$ tais que $\beta \circ \hat{F} \circ \alpha^{-1}$ é a projeção canônica.
	\end{enumerate}
	
	\textbf{Desenvolvimento Formal:}
	
	Sejam $p \in M$ e $q = F(p) \in N$. Escolhemos cartas coordenadas arbitrárias $\varphi : U \to \mathbb{R}^m$ centradas em $p$ e $\psi : V \to \mathbb{R}^n$ centradas em $q$, com $F(U) \subset V$. Definimos a representação local de $F$ como:
	\begin{equation*}
		\hat{F} = \psi \circ F \circ \varphi^{-1} : \varphi(U) \to \psi(V).
	\end{equation*}
	Como $F$ possui posto constante $k$ em $M$, a aplicação $\hat{F}$ possui posto constante $k$ em $\varphi(U) \subset \mathbb{R}^m$. Aplicamos agora o resultado da \textbf{Questão 25 desta Lista}. Para o ponto $\varphi(p) \in \mathbb{R}^m$, existem abertos $\tilde{U} \subset \varphi(U)$ e $\tilde{V} \subset \psi(V)$ e difeomorfismos locais $\alpha : \tilde{U} \to \mathbb{R}^m$ e $\beta : \tilde{V} \to \mathbb{R}^n$ tais que:
	\begin{equation*}
		\beta \circ \hat{F} \circ \alpha^{-1}(x^1, \dots, x^m) = (x^1, \dots, x^k, 0, \dots, 0).
	\end{equation*}
	
	Agora, definimos as novas cartas coordenadas para $M$ e $N$. Para $M$, definimos a nova carta $\overline{\varphi} : \overline{U} \to \mathbb{R}^m$ como a composição:
	\begin{equation*}
		\overline{\varphi} = \alpha \circ \varphi, \quad \text{onde } \overline{U} = \varphi^{-1}(\tilde{U}).
	\end{equation*}
	Para $N$, definimos a nova carta $\overline{\psi} : \overline{V} \to \mathbb{R}^n$ como a composição:
	\begin{equation*}
		\overline{\psi} = \beta \circ \psi, \quad \text{onde } \overline{V} = \psi^{-1}(\tilde{V}).
	\end{equation*}
	Ambas são composições de difeomorfismos, logo são cartas coordenadas suaves. Verificamos a representação de $F$ sob estas novas cartas:
	\begin{align*}
		\overline{\psi} \circ F \circ \overline{\varphi}^{-1} &= (\beta \circ \psi) \circ F \circ (\varphi^{-1} \circ \alpha^{-1}) \\
		&= \beta \circ (\psi \circ F \circ \varphi^{-1}) \circ \alpha^{-1} \\
		&= \beta \circ \hat{F} \circ \alpha^{-1}.
	\end{align*}
	
	Pelo resultado da Questão 25, a expressão acima é precisamente a projeção canônica:
	\begin{equation*}
		\overline{\psi} \circ F \circ \overline{\varphi}^{-1}(x^1, \dots, x^m) = (x^1, \dots, x^k, 0, \dots, 0).
	\end{equation*}
	A tese está provada.
}

%%%% QUESTÃO 27 %%%%
\questao{Sejam $\pi : M \to N$ uma submersão e $X \in \mathfrak{X}(N)$. Mostre que existe um campo vetorial suave em $M$ (chamada levantamento de $X$) que é $\pi$-relacionado com $X$. Este campo é único?}

\solucao{
	\textbf{Análise do Enunciado:}
	\begin{itemize}
		\item \textbf{Hipóteses:} $\pi: M \to N$ é uma submersão; $X$ é um campo de vetores suave em $N$.
		\item \textbf{Tese:} 
		\begin{enumerate}[label=(\roman*)]
			\item Provar a existência de um campo $Y \in \mathfrak{X}(M)$ tal que $d\pi_p(Y_p) = X_{\pi(p)}$ para todo $p \in M$ (condição de ser $\pi$-relacionado).
			\item Determinar se tal campo $Y$ é único.
		\end{enumerate}
	\end{itemize}
	
	\textbf{Fundamentação Teórica:}
	\begin{enumerate}
		\item \textbf{Campos $\pi$-relacionados (Lee, Cap. 4):} Um campo $Y \in \mathfrak{X}(M)$ é $\pi$-relacionado a $X \in \mathfrak{X}(N)$ se, para todo $p \in M$, a diferencial $d\pi_p$ mapeia o vetor $Y_p$ no vetor $X_{\pi(p)}$, ou seja, $d\pi_p(Y_p) = X_{\pi(p)}$.
		\item \textbf{Teorema da Forma Canônica Local para Submersões (Lee, Cap. 7):} Para qualquer $p \in M$, existem cartas coordenadas $(U, \varphi)$ e $(V, \psi)$ tais que $\pi$ é representada pela projeção canônica $\hat{\pi}(x^1, \dots, x^m) = (x^1, \dots, x^n)$.
		\item \textbf{Existência de Partição da Unidade (Lista 01):} Conforme provado anteriormente na Lista 01, para qualquer cobertura aberta $\{U_\alpha\}$ de $M$, existe uma família de funções suaves $\rho_\alpha : M \to [0,1]$ tal que $\text{supp}(\rho_\alpha) \subset U_\alpha$ e $\sum \rho_\alpha = 1$.
	\end{enumerate}
	
	\textbf{Desenvolvimento Formal:}
	
	\textit{1. Existência Local:}
	Seja $p \in M$. Pela Forma Canônica Local, existem cartas $\varphi: U \to \mathbb{R}^m$ e $\psi: V \to \mathbb{R}^n$ tais que a representação de $\pi$ é $\hat{\pi}(x^1, \dots, x^m) = (x^1, \dots, x^n)$. No aberto $V$, o campo $X$ pode ser escrito como $X = \sum_{i=1}^n X^i \frac{\partial}{\partial y^i}$. 
	Definimos um campo de vetores local $Y_U$ em $U$ como:
	\begin{equation*}
		Y_U = \sum_{i=1}^n (X^i \circ \pi) \frac{\partial}{\partial x^i}.
	\end{equation*}
	Calculando a ação da diferencial $d\pi$ sobre $Y_U$:
	\begin{equation*}
		d\pi(Y_U) = \sum_{i=1}^n (X^i \circ \pi) d\pi \left( \frac{\partial}{\partial x^i} \right) = \sum_{i=1}^n (X^i \circ \pi) \frac{\partial}{\partial y^i} = X.
	\end{equation*}
	Assim, $Y_U$ é um levantamento local de $X$ em $U$.
	
	\textit{2. Existência Global:}
	Seja $\{U_\alpha\}$ uma cobertura de $M$ por abertos onde existam levantamentos locais $Y_\alpha \in \mathfrak{X}(U_\alpha)$. Seja $\{\rho_\alpha\}$ uma partição da unidade subordinada a $\{U_\alpha\}$. Definimos o campo global $Y$ por:
	\begin{equation*}
		Y = \sum_\alpha \rho_\alpha Y_\alpha.
	\end{equation*}
	Verificamos que $Y$ é $\pi$-relacionado a $X$:
	\begin{align*}
		d\pi_p(Y_p) &= d\pi_p \left( \sum_\alpha \rho_\alpha(p) (Y_\alpha)_p \right) \\
		&= \sum_\alpha \rho_\alpha(p) d\pi_p((Y_\alpha)_p) \\
		&= \sum_\alpha \rho_\alpha(p) X_{\pi(p)} \\
		&= \left( \sum_\alpha \rho_\alpha(p) \right) X_{\pi(p)} = X_{\pi(p)}.
	\end{align*}
	Logo, $Y \in \mathfrak{X}(M)$ é um levantamento suave de $X$.
	
	\textit{3. Análise da Unicidade:}
	O campo $Y$ não é único. Suponha que $Y$ seja um levantamento de $X$. Seja $Z \in \mathfrak{X}(M)$ qualquer campo de vetores tal que $Z_p \in \ker(d\pi_p)$ para todo $p \in M$ (ou seja, $Z$ é um campo vertical). Então, o campo $Y' = Y + Z$ também é um levantamento, pois:
	\begin{equation*}
		d\pi_p(Y'_p) = d\pi_p(Y_p) + d\pi_p(Z_p) = X_{\pi(p)} + 0 = X_{\pi(p)}.
	\end{equation*}
	Como $\pi$ é uma submersão, $\dim(\ker d\pi_p) = m - n$. Se $m > n$, o núcleo é não trivial, e existem infinitos campos vetoriais $Z$, tornando o levantamento não único.
}
%%%%%%%%%%%%%%%%%%%%
%%%% OBSERVAÇÃO %%%%
%%%%%%%%%%%%%%%%%%%%
Diferente do caso de difeomorfismos locais (\textbf{Questão 26, Lista 02}), onde a injetividade da diferencial garante a unicidade do levantamento, aqui a submersão permite a existência de múltiplos levantamentos devido ao núcleo não trivial de $d\pi_p$.

%%%% QUESTÃO 28 %%%%
\questao{Sejam $S_1$ e $S_2$ subvariedades mergulhadas das variedades suaves $M_1$ e $M_2$, respectivamente. Mostre que $S_1 \times S_2$ é uma subvariedade mergulhada da variedade produto $M_1 \times M_2$.}

\solucao{
	\textbf{Análise do Enunciado:}
	\begin{itemize}
		\item \textbf{Hipóteses:} $M_1, M_2$ são variedades suaves; $S_1 \subset M_1$ é uma subvariedade mergulhada de dimensão $k_1$; $S_2 \subset M_2$ é uma subvariedade mergulhada de dimensão $k_2$.
		\item \textbf{Tese:} Provar que o produto $S_1 \times S_2$ é uma subvariedade mergulhada de $M_1 \times M_2$ de dimensão $k_1 + k_2$.
	\end{itemize}
	
	\textbf{Fundamentação Teórica:}
	\begin{enumerate}
		\item \textbf{Variedades Produto (Questão 18, Lista 01):} Se $M_1$ e $M_2$ são variedades suaves de dimensão $m_1$ e $m_2$, então $M_1 \times M_2$ é uma variedade suave de dimensão $m_1 + m_2$. Se $(U, \varphi)$ e $(V, \psi)$ são cartas em $M_1$ e $M_2$, então $(U \times V, \varphi \times \psi)$ é uma carta em $M_1 \times M_2$.
		\item \textbf{Caracterização de Subvariedades Mergulhadas (Lee, Cap. 8):} Um subconjunto $S \subset M$ é uma subvariedade mergulhada de dimensão $k$ se, e somente se, para cada $p \in S$, existe uma carta coordenada $(U, \varphi)$ de $M$ centrada em $p$ tal que:
		\begin{equation*}
			\varphi(S \cap U) = \{ (x^1, \dots, x^m) \in \varphi(U) \mid x^{k+1} = \dots = x^m = 0 \}.
		\end{equation*}
		Tais cartas são chamadas de \textit{cartas adaptadas} ou \textit{fatias} (slices).
	\end{enumerate}
	
	\textbf{Desenvolvimento Formal:}
	
	Seja $(p, q)$ um ponto qualquer de $S_1 \times S_2$, onde $p \in S_1$ e $q \in S_2$. 
	Como $S_1$ é uma subvariedade mergulhada de $M_1$, existe uma carta adaptada (slice) $(U, \varphi)$ para $M_1$ centrada em $p$ tal que $\varphi : U \to \mathbb{R}^{m_1}$ satisfaz:
	\begin{equation*}
		\varphi(S_1 \cap U) = \{ (x^1, \dots, x^{m_1}) \in \varphi(U) \mid x^{k_1+1} = \dots = x^{m_1} = 0 \}.
	\end{equation*}
	Analogamente, como $S_2$ é uma subvariedade mergulhada de $M_2$, existe uma carta adaptada (slice) $(V, \psi)$ para $M_2$ centrada em $q$ tal que $\psi : V \to \mathbb{R}^{m_2}$ satisfaz:
	\begin{equation*}
		\psi(S_2 \cap V) = \{ (y^1, \dots, y^{m_2}) \in \psi(V) \mid y^{k_2+1} = \dots = y^{m_2} = 0 \}.
	\end{equation*}
	
	Consideremos agora a carta produto $(U \times V, \Phi)$ em $M_1 \times M_2$, onde $\Phi = \varphi \times \psi$. O domínio $\Phi(U \times V) = \varphi(U) \times \psi(V)$ é um aberto em $\mathbb{R}^{m_1} \times \mathbb{R}^{m_2} \cong \mathbb{R}^{m_1+m_2}$.
	Analisamos a imagem do conjunto $(S_1 \times S_2) \cap (U \times V)$ sob a aplicação $\Phi$:
	\begin{align*}
		\Phi((S_1 \times S_2) \cap (U \times V)) &= \Phi((S_1 \cap U) \times (S_2 \cap V)) \\
		&= \varphi(S_1 \cap U) \times \psi(S_2 \cap V) \\
		&= \{ (x^1, \dots, x^{m_1}) \mid x^{k_1+1} = \dots = x^{m_1} = 0 \} \times \{ (y^1, \dots, y^{m_2}) \mid y^{k_2+1} = \dots = y^{m_2} = 0 \} \\
		&= \{ (x^1, \dots, x^{m_1}, y^1, \dots, y^{m_2}) \in \Phi(U \times V) \mid x^{k_1+1} = \dots = x^{m_1} = 0, y^{k_2+1} = \dots = y^{m_2} = 0 \}.
	\end{align*}
	
	O conjunto resultante é precisamente a definição de uma \textbf{fatia (slice)} em $\mathbb{R}^{m_1+m_2}$ de dimensão $k_1 + k_2$, onde as coordenadas a partir da posição $k_1+1$ até a posição $m_1$ e a partir de $m_1+k_2+1$ até $m_1+m_2$ são nulas. 
	
	Reordenando as coordenadas (o que constitui um difeomorfismo linear de $\mathbb{R}^{m_1+m_2}$), podemos escrever a imagem como:
	\begin{equation*}
		\{ (z^1, \dots, z^{m_1+m_2}) \in \Phi(U \times V) \mid z^{k_1+k_2+1} = \dots = z^{m_1+m_2} = 0 \}.
	\end{equation*}
	Assim, para qualquer ponto $(p, q) \in S_1 \times S_2$, encontramos uma carta adaptada (slice) em $M_1 \times M_2$. Pela caracterização de subvariedades mergulhadas, concluímos que $S_1 \times S_2$ é uma subvariedade mergulhada de $M_1 \times M_2$.
}

%%%% QUESTÃO 29 %%%%
\questao{Seja $S \subset M^n$ uma $k$-subvariedade mergulhada. Com a topologia induzida, mostre que $S$ é uma variedade topológica de dimensão $k$ e possui uma única estrutura suave tal que a inclusão $\iota : S \hookrightarrow M^n$ seja um mergulho suave.}

\solucao{
	\textbf{Análise do Enunciado:}
	\begin{itemize}
		\item \textbf{Hipóteses:} $M^n$ é uma $n$-variedade suave; $S \subset M^n$ é um subconjunto que satisfaz a condição de subvariedade mergulhada de dimensão $k$. $S$ é munido da topologia de subespaço de $M^n$.
		\item \textbf{Tese:} 
		\begin{enumerate}[label=(\roman*)]
			\item Provar que $S$ é uma variedade topológica de dimensão $k$.
			\item Provar que existe uma única estrutura suave em $S$ tal que a inclusão $\iota : S \to M^n$ seja um mergulho suave.
		\end{enumerate}
	\end{itemize}
	
	\textbf{Fundamentação Teórica:}
	\begin{enumerate}
		\item \textbf{Definição de Variedade Topológica (Lee, Cap. 1):} Um espaço $S$ é uma variedade topológica de dimensão $k$ se for Hausdorff, possuir uma base enumerável e for localmente homeomorfo a $\mathbb{R}^k$.
		\item \textbf{Subvariedade Mergulhada (Lee, Cap. 8):} Um subconjunto $S \subset M^n$ é uma subvariedade mergulhada de dimensão $k$ se, para cada $p \in S$, existe uma carta adaptada (\textit{slice chart}) $(U, \varphi)$ de $M^n$ centrada em $p$ tal que $\varphi(S \cap U) = \{ (x^1, \dots, x^n) \in \varphi(U) \mid x^{k+1} = \dots = x^n = 0 \}$.
		\item \textbf{Mergulho Suave (Lee, Cap. 7):} Uma aplicação $\iota : S \to M$ é um mergulho suave se é uma imersão injetiva e um homeomorfismo sobre sua imagem $\iota(S)$.
	\end{enumerate}
	
	\textbf{Desenvolvimento Formal:}
	
	\textit{1. Prova de que $S$ é uma variedade topológica.}
	Como $S$ é um subespaço de $M^n$ e $M^n$ é Hausdorff e possui base enumerável, $S$ herda automaticamente essas propriedades. Para a homeomorfia local, seja $p \in S$. Pela definição de subvariedade mergulhada, existe uma carta adaptada $(U, \varphi)$ de $M^n$ centrada em $p$. 
	Definimos a aplicação $\psi : S \cap U \to \mathbb{R}^k$ como a composição da carta $\varphi$ com a projeção canônica $\pi : \mathbb{R}^n \to \mathbb{R}^k$:
	\begin{equation*}
		\psi(q) = \pi(\varphi(q)) = (x^1, \dots, x^k).
	\end{equation*}
	Como $\varphi(S \cap U) = \mathbb{R}^k \times \{0\}$, a aplicação $\psi$ é um homeomorfismo entre $S \cap U$ e um aberto de $\mathbb{R}^k$. Assim, $S$ é localmente homeomorfo a $\mathbb{R}^k$, provando que $S$ é uma variedade topológica de dimensão $k$.
	
	\textit{2. Existência de Estrutura Suave.}
	As aplicações $\psi$ definidas acima formam um atlas para $S$. Se $(U_\alpha, \varphi_\alpha)$ e $(U_\beta, \varphi_\beta)$ são duas cartas adaptadas para $M^n$, a função de transição entre as respectivas projeções $\psi_\alpha$ e $\psi_\beta$ em $S$ é a restrição da função de transição de $M^n$ (que é suave) ao subespaço $\mathbb{R}^k \times \{0\}$. Como a restrição de uma função suave a um subespaço linear é suave, as funções de transição de $S$ são suaves. Portanto, $S$ admite uma estrutura suave.
	Sob esta estrutura, a inclusão $\iota : S \to M^n$ é localmente representada por $\iota(x^1, \dots, x^k) = (x^1, \dots, x^k, 0, \dots, 0)$, que é claramente uma imersão injetiva. Como $S$ possui a topologia de subespaço, $\iota$ é por definição um homeomorfismo sobre sua imagem. Logo, $\iota$ é um mergulho suave.
	
	\textit{3. Unicidade da Estrutura Suave.}
	Suponhamos que exista outra estrutura suave $\mathcal{A}'$ em $S$ tal que a inclusão $\iota : (S, \mathcal{A}') \to M^n$ seja um mergulho suave. Seja $\mathcal{A}_{ind}$ a estrutura induzida pelas cartas adaptadas construídas anteriormente.
	A aplicação identidade $\text{id} : (S, \mathcal{A}_{ind}) \to (S, \mathcal{A}')$ pode ser escrita como a composição:
	\begin{equation*}
		\text{id} = \iota^{-1} \circ \iota.
	\end{equation*}
	Sendo $\iota : (S, \mathcal{A}_{ind}) \to M^n$ um mergulho suave e $\iota^{-1} : \iota(S) \to (S, \mathcal{A}')$ a inversa de um mergulho suave, a composição $\text{id}$ é suave. Como a identidade é um difeomorfismo, as estruturas suaves $\mathcal{A}_{ind}$ e $\mathcal{A}'$ coincidem.
}

%%%% QUESTÃO 30 %%%%
\questao{Seja $M(m \times n, \mathbb{R})$ a variedade suave (com sua estrutura suave natural) de todas as matrizes reais de ordem $m \times n$. Considere o subconjunto $M_k(m \times n, \mathbb{R})$ das matrizes de posto igual a $k$. Mostre que $M_k(m \times n, \mathbb{R})$ é uma subvariedade mergulhada de $M(m \times n, \mathbb{R})$ de dimensão $k(m+n-k)$.}

\solucao{
	\textbf{Análise do Enunciado:}
	\begin{itemize}
		\item \textbf{Hipóteses:} $M = M(m \times n, \mathbb{R}) \cong \mathbb{R}^{mn}$ é a variedade de todas as matrizes e $$ = M_k(m \times n, \mathbb{R}) = \{ A \in M \mid \text{posto}(A) = k \}.$$
		\item \textbf{Tese:} Provar que $S$ é uma subvariedade mergulhada de $M$ e que sua dimensão é $k(m+n-k)$ (com codimensão $(m-k)(n-k)$).
	\end{itemize}
	
	\textbf{Fundamentação Teórica:}
	\begin{enumerate}
		\item \textbf{Caracterização de Subvariedades Mergulhadas (Lee, Cap. 8):} Um subconjunto $S \subset M$ é uma subvariedade mergulhada de dimensão $d$ se, para cada $A \in S$, existe uma carta adaptada (\textit{slice chart}) $(U, \varphi)$ tal que $\varphi(S \cap U) = \{ (x^1, \dots, x^{mn}) \in \varphi(U) \mid x^{d+1} = \dots = x^{mn} = 0 \}$.
		\item \textbf{Teorema do Posto (Questão 26, Lista 03):} Para qualquer matriz $A$ de posto $k$, existem mudanças de base (difeomorfismos lineares) nos espaços de domínio e contradomínio que transformam $A$ na forma canônica $\begin{pmatrix} I_k & 0 \\ 0 & 0 \end{pmatrix}$.
		\item \textbf{Complemento de Schur:} Para uma matriz em blocos $\begin{pmatrix} A_{11} & A_{12} \\ A_{21} & A_{22} \end{pmatrix}$, se $A_{11}$ é inversível, o posto da matriz é $k$ se, e somente se, $A_{22} = A_{21} A_{11}^{-1} A_{12}$.
	\end{enumerate}
	
	\textbf{Desenvolvimento Formal:}
	
	Seja $A \in M_k$. Pelo Teorema do Posto, existem matrizes inversíveis $P \in GL(m, \mathbb{R})$ e $Q \in GL(n, \mathbb{R})$ tais que $P A Q = \begin{pmatrix} I_k & 0 \\ 0 & 0 \end{pmatrix}$. A aplicação $B \mapsto PBQ$ é um difeomorfismo de $M(m \times n, \mathbb{R})$ sobre si mesmo. Portanto, basta provar que o conjunto de matrizes de posto $k$ próximas a $\begin{pmatrix} I_k & 0 \\ 0 & 0 \end{pmatrix}$ forma uma subvariedade mergulhada.
	
	Considere a vizinhança $U$ de matrizes $B = \begin{pmatrix} B_{11} & B_{12} \\ B_{21} & B_{22} \end{pmatrix}$ onde $B_{11}$ é uma matriz $k \times k$ inversível. Como a inversão de matrizes é uma operação suave, $U$ é um aberto em $M(m \times n, \mathbb{R})$. 
	Para qualquer $B \in U$, o posto de $B$ é $k$ se, e somente se, o bloco $B_{22}$ (de dimensão $(m-k) \times (n-k)$) satisfaz a condição:
	\begin{equation*}
		B_{22} = B_{21} B_{11}^{-1} B_{12}.
	\end{equation*}
	
	Definimos agora a aplicação $\Phi : U \to \mathbb{R}^{mn}$ dada por:
	\begin{equation*}
		\Phi(B) = \begin{pmatrix} B_{11} & B_{12} \\ B_{21} & B_{22} - B_{21} B_{11}^{-1} B_{12} \end{pmatrix}.
	\end{equation*}
	A aplicação $\Phi$ é suave e possui inversa suave $\Phi^{-1}(C) = \begin{pmatrix} C_{11} & C_{12} \\ C_{21} & C_{22} + C_{21} C_{11}^{-1} C_{12} \end{pmatrix}$, logo $\Phi$ é um difeomorfismo local (uma carta).
	
	Sob a representação de $\Phi$, a condição $\text{posto}(B) = k$ é equivalente a $B_{22} - B_{21} B_{11}^{-1} B_{12} = 0$. Ou seja:
	\begin{equation*}
		\Phi(M_k \cap U) = \{ (B_{11}, B_{12}, B_{21}, 0) \in \mathbb{R}^{mn} \}.
	\end{equation*}
	Esta é precisamente a definição de uma carta adaptada (\textit{slice chart}). A dimensão da subvariedade é a dimensão do espaço dos blocos $(B_{11}, B_{12}, B_{21})$, que é:
	\begin{equation*}
		\dim(M_k) = k^2 + k(n-k) + (m-k)k = k^2 + kn - k^2 + mk - k^2 = mk + nk - k^2 = k(m+n-k).
	\end{equation*}
	A codimensão é $mn - (mk + nk - k^2) = mn - mk - nk + k^2 = (m-k)(n-k)$. 
	Assim, $M_k$ é uma subvariedade mergulhada de dimensão $k(m+n-k)$.
}

%%%% QUESTÃO 31 %%%%
\questao{Considere a aplicação $F : \mathbb{R}^4 \to \mathbb{R}^2$ dada por $F(x,y,s,t) = (x^2 + y, x^2 + y^2 + s^2 + t^2 + y)$. Mostre que $(0, 1)$ é um valor regular de $F$, e que o conjunto de nível $F^{-1}(0, 1)$ é difeomorfo a $\mathbb{S}^2$.}

\solucao{
	\textbf{Análise do Enunciado:}
	\begin{itemize}
		\item \textbf{Hipóteses:} $F: \mathbb{R}^4 \to \mathbb{R}^2$ é uma aplicação suave. O valor alvo é $c = (0,1)$.
		\item \textbf{Tese:} 
		\begin{enumerate}[label=(\roman*)]
			\item Provar que $\text{posto}(dF_p) = 2$ para todo $p \in F^{-1}(0,1)$.
			\item Provar a existência de um difeomorfismo entre o conjunto de nível $S = F^{-1}(0,1)$ e a esfera $\mathbb{S}^2$.
		\end{enumerate}
	\end{itemize}
	
	\textbf{Fundamentação Teórica:}
	\begin{enumerate}
		\item \textbf{Teorema do Valor Regular (Lee, Cap. 7):} Se $c$ é um valor regular, $F^{-1}(c)$ é uma subvariedade mergulhada de codimensão $\dim N$.
		\item \textbf{Posto de Matrizes:} O posto de uma matriz é igual ao tamanho do seu maior menor (determinante de uma submatriz quadrada) não nulo. Para posto 2, basta encontrar um menor $2 \times 2$ diferente de zero.
		\item \textbf{Difeomorfismo:} Uma aplicação é um difeomorfismo se for bijetiva e tanto ela quanto sua inversa forem suaves.
	\end{enumerate}
	
	\textbf{Desenvolvimento Formal:}
	
	\textit{1. Verificação do Valor Regular via Menores:}
	A matriz Jacobiana de $F$ é:
	\begin{equation*}
		JF_p = \begin{pmatrix} 
			2x & 1 & 0 & 0 \\ 
			2x & 2y+1 & 2s & 2t 
		\end{pmatrix}.
	\end{equation*}
	Analisamos os menores de ordem 2 para verificar se o posto é 2 para todo $p \in F^{-1}(0,1)$.
	Seja $M_{ij}$ o determinante da submatriz formada pelas colunas $i$ e $j$:
	\begin{itemize}
		\item $M_{23} = \det \begin{pmatrix} 1 & 0 \\ 2y+1 & 2s \end{pmatrix} = 2s$.
		\item $M_{24} = \det \begin{pmatrix} 1 & 0 \\ 2y+1 & 2t \end{pmatrix} = 2t$.
	\end{itemize}
	Se $s \neq 0$ ou $t \neq 0$, então ao menos um desses menores é não nulo, logo $\text{posto}(JF_p) = 2$.
	Agora, analisamos o caso crítico onde $s=0$ e $t=0$. Se $p \in F^{-1}(0,1)$, as equações do conjunto de nível são:
	\begin{enumerate}
		\item $x^2 + y = 0 \implies y = -x^2$.
		\item $x^2 + y^2 + 0^2 + 0^2 + y = 1$.
	\end{enumerate}
	Substituindo $y = -x^2$ na segunda equação: $x^2 + (-x^2)^2 - x^2 = 1 \implies x^4 = 1 \implies x = \pm 1$. 
	Para $x = \pm 1$, temos $y = -1$. Testamos o menor $M_{12}$ para esses pontos:
	\begin{equation*}
		M_{12} = \det \begin{pmatrix} 2x & 1 \\ 2x & 2y+1 \end{pmatrix} = 2x(2y+1) - 2x = 4xy.
	\end{equation*}
	Para $x=1, y=-1$, temos $M_{12} = 4(1)(-1) = -4 \neq 0$. Para $x=-1, y=-1$, temos $M_{12} = 4(-1)(-1) = 4 \neq 0$. 
	Em todos os casos possíveis, existe um menor não nulo. Portanto, $\text{posto}(JF_p) = 2$ para todo $p \in F^{-1}(0,1)$, e $(0,1)$ é um valor regular.
	
	\textit{2. Difeomorfismo para $\mathbb{S}^2$:}
	Como visto anteriormente, o conjunto de nível é $S = \{ (x, -x^2, s, t) \in \mathbb{R}^4 \mid x^4 + s^2 + t^2 = 1 \}$.
	Definimos a aplicação $\Phi : S \to \mathbb{S}^2$ por:
	\begin{equation*}
		\Phi(x, -x^2, s, t) = \frac{(x, s, t)}{\sqrt{x^2+s^2+t^2}}.
	\end{equation*}
	Esta aplicação é claramente suave, pois o denominador $\sqrt{x^2+s^2+t^2}$ nunca é zero em $S$ (se fosse, teríamos $x=s=t=0$, mas então $x^4+s^2+t^2=0 \neq 1$).
	
	Para provar que $\Phi$ é um difeomorfismo, construímos a inversa $\Psi : \mathbb{S}^2 \to S$. Seja $(u,v,w) \in \mathbb{S}^2$. Buscamos um ponto $(x, -x^2, s, t) \in S$ tal que:
	\begin{equation*}
		(x, s, t) = \lambda(u, v, w) \quad \text{para algum } \lambda > 0.
	\end{equation*}
	Substituindo na equação de $S$ ($x^4 + s^2 + t^2 = 1$):
	\begin{align*}
		(\lambda u)^4 + (\lambda v)^2 + (\lambda w)^2 &= 1 \\
		\lambda^4 u^4 + \lambda^2 (v^2 + w^2) &= 1.
	\end{align*}
	Como $(u,v,w) \in \mathbb{S}^2$, temos $v^2+w^2 = 1-u^2$. Logo:
	\begin{equation*}
		u^4 (\lambda^2)^2 + (1-u^2)\lambda^2 - 1 = 0.
	\end{equation*}
	Esta é uma equação quadrática em $z = \lambda^2$. A solução positiva é:
	\begin{equation*}
		\lambda(u,v,w) = \sqrt{\frac{-(1-u^2) + \sqrt{(1-u^2)^2 + 4u^4}}{2u^4}}.
	\end{equation*}
	Se $u=0$, a equação simplifica para $(1-0)\lambda^2 - 1 = 0 \implies \lambda = 1$. A função $\lambda(u,v,w)$ é suave e positiva em toda a esfera. Assim, definimos $\Psi(u,v,w) = (\lambda u, -(\lambda u)^2, \lambda v, \lambda w)$. 
	Como $\Psi$ é a inversa explícita de $\Phi$ e ambas são suaves, $S$ é difeomorfo a $\mathbb{S}^2$.
}
\solucaoalt{\par
	\textbf{1. Verificação do Valor Regular:}
	Calculamos a matriz diferencial (Jacobiana) de $F$:
	\begin{equation*}
		DF(x,y,s,t) = \begin{pmatrix} 
			2x & 1 & 0 & 0 \\ 
			2x & 2y+1 & 2s & 2t 
		\end{pmatrix}.
	\end{equation*}
	Observamos que esta matriz possui posto máximo (posto 2) sempre que $s \neq 0$ ou $t \neq 0$, pois a submatriz formada pelas colunas 3 ou 4 e a coluna 2 terá determinante não nulo.
	
	Suponhamos agora que $s = t = 0$. Nesse caso, o determinante da submatriz formada pelas duas primeiras colunas é:
	\begin{equation*}
		\det \begin{pmatrix} 2x & 1 \\ 2x & 2y+1 \end{pmatrix} = 2x(2y+1) - 2x = 4xy - 2x + 2x = 4xy.
	\end{equation*}
	Este determinante é zero se, e somente se, $x = 0$ ou $y = 0$. Entretanto, se $p \in F^{-1}(0,1)$, devemos ter $x^2+y=0$ e $x^2+y^2+s^2+t^2+y=1$. 
	Se $x=0$, então $y=0$, e a segunda equação resultaria em $0 = 1$, o que é impossível. Se $y=0$, então $x=0$, resultando no mesmo absurdo.
	Portanto, $DF_p$ é sobrejetiva para todo $p \in F^{-1}(0,1) = S$.
	
	\textbf{2. Construção do Difeomorfismo de $\mathbb{R}^4$:}
	Definimos a aplicação $\varphi : \mathbb{R}^4 \to \mathbb{R}^4$ por:
	\begin{equation*}
		\varphi(x,y,s,t) = (x, x^2+y, s, t).
	\end{equation*}
	Observamos que $\varphi$ é diferenciável e injetiva. Sua inversa é dada por:
	\begin{equation*}
		\varphi^{-1}(u,v,s,t) = (u, v-u^2, s, t),
	\end{equation*}
	a qual também é diferenciável. Logo, $\varphi$ é um difeomorfismo de $\mathbb{R}^4$.
	
	\textbf{3. Análise da Aplicação Composta $\tilde{F}$:}
	Definimos a aplicação composta $\tilde{F} = F \circ \varphi^{-1} : \mathbb{R}^4 \to \mathbb{R}^2$ como:
	\begin{align*}
		\tilde{F}(u,v,s,t) &= F(\varphi^{-1}(u,v,s,t)) \\
		&= F(u, v-u^2, s, t) \\
		&= (u^2 + (v-u^2), u^2 + (v-u^2)^2 + s^2 + t^2 + (v-u^2)) \\
		&= (v, v + (v-u^2)^2 + s^2 + t^2).
	\end{align*}
	
	\textbf{4. Difeomorfismo entre $S$ e $\mathbb{S}^2$:}
	Analisamos a imagem do conjunto de nível $S$ sob a aplicação $\varphi$:
	\begin{equation*}
		\varphi(S) = \varphi(F^{-1}(0,1)) = \{ (x, v, s, t) \in \mathbb{R}^4 \mid v=0 \text{ e } x^4+s^2+t^2=1 \}.
	\end{equation*}
	Este conjunto é uma subvariedade mergulhada de $\mathbb{R}^4$ difeomorfa a $\mathbb{S}^2$ (via projeção radial nos eixos $x,s,t$). Como a restrição $\varphi|_S$ é um difeomorfismo entre $S$ e $\varphi(S)$, concluímos que $S$ é difeomorfo a $\varphi(S)$ e, consequentemente, a $\mathbb{S}^2$.
}

%%%% QUESTÃO 32 %%%%
\questao{Considere a aplicação $F : \mathbb{R}^2 \to \mathbb{R}$ dada por $F(x,y) = x^3 + xy + y^3$. Que conjuntos de nível de $F$ são subvariedades mergulhadas de $\mathbb{R}^2$?}

\solucao{
	\textbf{Análise do Enunciado:}
	\begin{itemize}
		\item \textbf{Hipóteses:} $F: \mathbb{R}^2 \to \mathbb{R}$ é uma aplicação polinomial, portanto suave em todo o seu domínio.
		\item \textbf{Tese:} Determinar para quais valores de $c \in \mathbb{R}$ o conjunto de nível $S_c = F^{-1}(c) = \{ (x,y) \in \mathbb{R}^2 \mid x^3 + xy + y^3 = c \}$ é uma subvariedade mergulhada de $\mathbb{R}^2$.
	\end{itemize}
	
	\textbf{Fundamentação Teórica:}
	\begin{enumerate}
		\item \textbf{Teorema do Valor Regular (Lee, Cap. 7):} Se $c \in \mathbb{R}$ é um valor regular de $F$, então o conjunto de nível $F^{-1}(c)$ é uma subvariedade mergulhada de $\mathbb{R}^2$ de dimensão $\dim(\mathbb{R}^2) - \dim(\mathbb{R}) = 1$.
		\item \textbf{Valor Regular:} Um valor $c$ é regular se, para todo $p \in F^{-1}(c)$, a diferencial $dF_p$ é sobrejetiva. Como o contradomínio é $\mathbb{R}$, isso equivale a dizer que $dF_p \neq 0$.
		\item \textbf{Pontos Críticos:} Um ponto $p$ é crítico se $dF_p = 0$. Os valores de $F$ nos pontos críticos são os valores críticos da aplicação.
	\end{enumerate}
	
	\textbf{Desenvolvimento Formal:}
	
	Primeiramente, calculamos a diferencial de $F$ para encontrar os pontos críticos:
	\begin{equation*}
		dF = \frac{\partial F}{\partial x} dx + \frac{\partial F}{\partial y} dy = (3x^2 + y) dx + (x + 3y^2) dy.
	\end{equation*}
	Os pontos críticos $p = (x,y)$ são aqueles que satisfazem o sistema:
	\begin{align}
		3x^2 + y &= 0 \label{eq:crit1} \\
		x + 3y^2 &= 0 \label{eq:crit2}
	\end{align}
	Da equação \eqref{eq:crit1}, temos $y = -3x^2$. Substituindo na equação \eqref{eq:crit2}:
	\begin{equation*}
		x + 3(-3x^2)^2 = 0 \implies x + 27x^4 = 0 \implies x(1 + 27x^3) = 0.
	\end{equation*}
	Isso nos fornece duas soluções para $x$:
	\begin{enumerate}
		\item Se $x = 0$, então $y = -3(0)^2 = 0$. O ponto crítico é $P_1 = (0,0)$.
		\item Se $1 + 27x^3 = 0$, então $x^3 = -1/27 \implies x = -1/3$. Então $y = -3(-1/3)^2 = -3(1/9) = -1/3$. O ponto crítico é $P_2 = (-1/3, -1/3)$.
	\end{enumerate}
	
	Agora, calculamos os valores críticos de $F$ avaliando a função nestes pontos:
	\begin{itemize}
		\item $F(P_1) = 0^3 + (0)(0) + 0^3 = 0$.
		\item $F(P_2) = (-1/3)^3 + (-1/3)(-1/3) + (-1/3)^3 = -1/27 + 1/9 - 1/27 = 1/27$.
	\end{itemize}
	Os valores críticos de $F$ são $\{0, 1/27\}$. Portanto, para qualquer $c \in \mathbb{R} \setminus \{0, 1/27\}$, $c$ é um valor regular e, pelo Teorema do Valor Regular, $F^{-1}(c)$ é uma subvariedade mergulhada de dimensão 1.
	
	\textit{Análise dos Valores Críticos:}
	
	\begin{itemize}
		\item \textbf{Para $c = 0$:} O conjunto de nível é $F^{-1}(0) = \{ (x,y) \in \mathbb{R}^2 \mid x^3 + xy + y^3 = 0 \}$. 
		Note que no ponto crítico $P_1 = (0,0)$, a diferencial $dF_{(0,0)} = (0,0)$, portanto não podemos aplicar o Teorema do Valor Regular. Para verificar se $F^{-1}(0)$ é uma subvariedade, analisamos o comportamento da equação próximo à origem.
		
		Para um valor fixo de $x$ pequeno, a equação $y^3 + xy + x^3 = 0$ é um polinômio cúbico em $y$. O número de raízes reais de um polinômio cúbico $y^3 + py + q = 0$ é determinado pelo sinal do seu discriminante $\Delta = -4p^3 - 27q^2$. No nosso caso, $p = x$ e $q = x^3$, logo:
		\begin{equation*}
			\Delta = -4x^3 - 27(x^3)^2 = -x^3(4 + 27x^3).
		\end{equation*}
		
		Analisamos o sinal de $\Delta$ para $x$ próximo de zero:
		\begin{enumerate}
			\item Se $x > 0$, então $\Delta < 0$, o que implica que a equação possui apenas \textbf{uma} raiz real para $y$.
			\item Se $x < 0$, então $\Delta > 0$, o que implica que a equação possui \textbf{três} raízes reais distintas para $y$.
		\end{enumerate}
		
		Isso significa que, para $x$ ligeiramente negativo, existem três \textbf{ramos} da curva passando próximos à origem, enquanto para $x$ positivo existe apenas um. Uma variedade de dimensão 1 deve ser localmente homeomorfa a $\mathbb{R}^1$, o que implica que, em qualquer ponto, a curva deve ter exatamente duas direções de saída. A existência de três ramos para $x < 0$ que convergem para $(0,0)$ rompe a homeomorfia local com a reta. Portanto, $F^{-1}(0)$ não é uma subvariedade mergulhada.
		
		\item \textbf{Para $c = 1/27$:} O ponto $P_2 = (-1/3, -1/3)$ é um ponto crítico. Para analisar a natureza deste ponto, calculamos as derivadas de segunda ordem de $F$:
		\begin{align*}
			\tfrac{\partial^2 F}{\partial x^2} &= \tfrac{\partial}{\partial x}(3x^2 + y) = 6x, & \tfrac{\partial^2 F}{\partial x \partial y} &= \tfrac{\partial}{\partial y}(3x^2 + y) = 1, \\
			\tfrac{\partial^2 F}{\partial y \partial x} &= \tfrac{\partial}{\partial x}(x + 3y^2) = 1, & \tfrac{\partial^2 F}{\partial y^2} &= \tfrac{\partial}{\partial y}(x + 3y^2) = 6y.
		\end{align*}
		A matriz Hessiana $H(F)$ avaliada em $P_2 = (-1/3, -1/3)$ é:
		\begin{equation*}
			H(F)_{P_2} = \begin{pmatrix} 6(-1/3) & 1 \\ 1 & 6(-1/3) \end{pmatrix} = \begin{pmatrix} -2 & 1 \\ 1 & -2 \end{pmatrix}.
		\end{equation*}
		Analisamos os autovalores de $H(F)_{P_2}$ ou seu determinante:
		\begin{equation*}
			\det(H(F)_{P_2}) = (-2)(-2) - (1)^2 = 4 - 1 = 3.
		\end{equation*}
		Como $\det(H) > 0$ e a componente $F_{xx} = -2 < 0$, o ponto $P_2$ é um \textbf{máximo local estrito} de $F$. 
		
		Isso implica que existe uma vizinhança $U$ de $P_2$ tal que para todo $q \in U \setminus \{P_2\}$, temos $F(q) < F(P_2) = 1/27$. Portanto, o conjunto de nível $F^{-1}(1/27) \cap U$ consiste unicamente no ponto isolado $\{P_2\}$. Como um ponto isolado é uma subvariedade mergulhada de dimensão 0, $F^{-1}(1/27)$ é uma subvariedade mergulhada.
	\end{itemize}
	
	Concluímos que os conjuntos de nível $F^{-1}(c)$ são subvariedades mergulhadas de $\mathbb{R}^2$ para todo $c \in \mathbb{R} \setminus \{0\}$.
	\begin{center}
		\begin{tikzpicture}[
			scale=2.5, 
			>=Stealth
			]
			% --- 1. Eixos finos e pretos ---
			\draw[->, thin, black] (-1.2,0) -- (1.2,0) node[right, black, font=\footnotesize] {$x$};
			\draw[->, thin, black] (0,-1.2) -- (0,1.2) node[above, black, font=\footnotesize] {$y$};
			
			% --- 2. O Conjunto de Nível c = 0 (Folium de Descartes) ---
			% Parte A: Entrada pelo 4º quadrante
			\draw[thick, blue!70!black, samples=100, domain=-0.78:0, variable=\t, smooth] 
			plot ({- \t / (1 + \t*\t*\t)}, {- \t*\t / (1 + \t*\t*\t)});
			% Parte B: Metade inferior do laço
			\draw[thick, blue!70!black, samples=100, domain=0:1, variable=\t, smooth] 
			plot ({- \t / (1 + \t*\t*\t)}, {- \t*\t / (1 + \t*\t*\t)});
			% Parte C: Metade superior do laço
			\draw[thick, blue!70!black, samples=100, domain=0:1, variable=\u, smooth] 
			plot ({- \u*\u / (1 + \u*\u*\u)}, {- \u / (1 + \u*\u*\u)});
			% Parte D: Saída pelo 2º quadrante
			\draw[thick, blue!70!black, samples=100, domain=-0.78:0, variable=\u, smooth] 
			plot ({- \u*\u / (1 + \u*\u*\u)}, {- \u / (1 + \u*\u*\u)});
			
			% --- 3. Setas de Direção Reais (Caminhada contínua sobre a curva) ---
			% Seta 1: Entrando pelo 4º quadrante em direção à origem
			\draw[->, thick, blue!70!black, domain=-0.6:-0.4, variable=\t]
			plot ({- \t / (1 + \t*\t*\t)}, {- \t*\t / (1 + \t*\t*\t)});
			
			% Seta 2: No laço, movendo-se para a ponta (metade inferior)
			\draw[->, thick, blue!70!black, domain=0.4:0.6, variable=\t]
			plot ({- \t / (1 + \t*\t*\t)}, {- \t*\t / (1 + \t*\t*\t)});
			
			% Seta 3: No laço, voltando para a origem (metade superior)
			\draw[->, thick, blue!70!black, domain=0.6:0.4, variable=\u]
			plot ({- \u*\u / (1 + \u*\u*\u)}, {- \u / (1 + \u*\u*\u)});
			
			% Seta 4: Saindo da origem em direção ao 2º quadrante
			\draw[->, thick, blue!70!black, domain=-0.4:-0.6, variable=\u]
			plot ({- \u*\u / (1 + \u*\u*\u)}, {- \u / (1 + \u*\u*\u)});
			
			% --- 4. Retas Tracejadas e Projeções (Apenas nos semi-eixos negativos) ---
			% Projeção no eixo x negativo
			\draw[dashed, gray!70, thin] (-0.3333, -0.3333) -- (-0.3333, 0);
			\node[above, font=\tiny, fill=white, inner sep=1pt] at (-0.3333, 0) {$-\frac{1}{3}$};
			
			% Projeção no eixo y negative
			\draw[dashed, gray!70, thin] (-0.3333, -0.3333) -- (0, -0.3333);
			\node[right, font=\tiny, fill=white, inner sep=1pt] at (0, -0.3333) {$-\frac{1}{3}$};
			
			% --- 5. Pontos Críticos e Rótulos ---
			% Ponto P1 (0,0) - Autointersecção
			\filldraw[black] (0,0) circle (1pt);
			\node[above right, font=\tiny, black, yshift=2pt] at (0,0) {$P_1(0,0)=F^{-1}(0)$ (Autointersecção)};
			
			% Ponto P2 (-1/3, -1/3) - Máximo local isolado com texto bem posicionado
			\filldraw[red!80!black] (-0.3333, -0.3333) circle (1pt);
			\node[below left, font=\tiny, red!80!black, shift={(0.02,0.02)}] at (-0.3333, -0.2333) {Ponto isolado de $F^{-1}\left(\frac{1}{27}\right)=P_2$};
			
			% Legenda inferior
			\node[anchor=north, font=\small] at (0, -1.3) {$F(x,y) = x^3 + xy + y^3$};
			
		\end{tikzpicture}
	\end{center}
}

%%%% QUESTÃO 33 %%%%
\questao{Considere a aplicação $\Phi : \mathbb{R}^2 \to \mathbb{R}$ dada por $\Phi(x,y) = x^3 - y^2$. Mostre que $\Phi^{-1}(0)$ não é uma subvariedade mergulhada de $\mathbb{R}^2$.}

\solucao{
	\textbf{Análise do Enunciado:}
	\begin{itemize}
		\item \textbf{Hipóteses:} $\Phi: \mathbb{R}^2 \to \mathbb{R}$ é uma aplicação suave; o conjunto de interesse é $$S = \Phi^{-1}(0) = \{ (x,y) \in \mathbb{R}^2 \mid y^2 = x^3 \}.$$
		\item \textbf{Tese:} Provar que $S$ não admite a estrutura de subvariedade mergulhada de $\mathbb{R}^2$.
	\end{itemize}
	
	\textbf{Fundamentação Teórica:}
	\begin{enumerate}
		\item \textbf{Definição de Subvariedade Mergulhada (Lee, Cap. 8):} Se $S$ é uma subvariedade mergulhada de dimensão 1, então para cada $p \in S$, deve existir uma parametrização suave $\gamma : (-\varepsilon, \varepsilon) \to S$ tal que $\gamma(0)=p$ e $\gamma'(0) \neq (0,0)$.
		\item \textbf{Propriedade de Funções Suaves (Cálculo):} Se $x(t)$ é uma função suave tal que $x(0)=0$ e $x(t) \geq 0$ para todo $t$, então $t=0$ é um ponto de mínimo local, o que implica obrigatoriamente que $x'(0) = 0$.
	\end{enumerate}
	
	\textbf{Desenvolvimento Formal:}
	
	A equação $y^2 = x^3$ implica a restrição fundamental $x^3 \geq 0$, o que significa que $x \geq 0$ para qualquer ponto $(x,y) \in S$.
	
	Suponhamos, por absurdo, que $S$ seja uma subvariedade mergulhada. Então, no ponto $p = (0,0)$, deve existir uma curva suave $\gamma(t) = (x(t), y(t))$ tal que $\gamma(0) = (0,0)$ e $\gamma'(0) \neq (0,0)$.
	
	Como a curva $\gamma$ deve estar contida em $S$, a condição $x(t) \geq 0$ deve ser satisfeita para todo $t \in (-\varepsilon, \varepsilon)$. Dado que $x(0)=0$, o ponto $t=0$ é um mínimo local da função $x(t)$. Pela teoria do cálculo, isso implica que:
	\begin{equation*}
		x'(0) = 0.
	\end{equation*}
	
	Agora, analisamos a componente $y(t)$. Como $\gamma(t) \in S$, temos $y(t)^2 = x(t)^3$. Diferenciando ambos os lados em relação a $t$:
	\begin{equation*}
		2y(t)y'(t) = 3x(t)^2 x'(t).
	\end{equation*}
	Avaliando no ponto $t=0$, onde $x(0)=0$ e $y(0)=0$, obtemos $0 = 0$, o que não nos fornece a derivada de $y$. No entanto, podemos analisar a derivada de segunda ordem ou, mais elegantemente, observar que:
	\begin{equation*}
		(y'(0))^2 = \lim_{t \to 0} \frac{(y(t)-y(0))^2}{t^2} = \lim_{t \to 0} \frac{x(t)^3}{t^2}.
	\end{equation*}
	Como $x'(0)=0$, a expansão de $x(t)$ começa em $t^2$ (ou ordem superior), ou seja, $x(t) = O(t^2)$. Assim, $x(t)^3 = O(t^6)$. Portanto:
	\begin{equation*}
		(y'(0))^2 = \lim_{t \to 0} \frac{O(t^6)}{t^2} = 0 \implies y'(0) = 0.
	\end{equation*}
	
	Concluímos que $\gamma'(0) = (x'(0), y'(0)) = (0,0)$. Isso contradiz a hipótese de que $\gamma$ é uma parametrização de subvariedade (que exige vetor velocidade não nulo). Logo, $S$ não pode ser uma subvariedade mergulhada em $(0,0)$.
	\begin{center}
		\begin{tikzpicture}[
			% Escala equilibrada para visualização nítida da cúspide
			scale=1.5, 
			>=Stealth
			]
			% --- 1. Eixos finos e pretos ---
			% Como x >= 0 obrigatoriamente, o eixo x se estende mais para a direita
			\draw[->, thin, black] (-1,0) -- (3,0) node[right, black, font=\footnotesize] {$x$};
			\draw[->, thin, black] (0,-2.5) -- (0,2.5) node[above, black, font=\footnotesize] {$y$};
			
			% --- 2. A Curva S (Parábola Semicúbica) em Azul Estilizado ---
			% Parametrização x = t^2 e y = t^3 garante o comportamento y^2 = x^3 sem falhas de precisão
			\draw[thick, blue!70!black, samples=200, domain=-1.35:1.35, variable=\t, smooth] 
			plot ({\t*\t}, {\t*\t*\t});
			
			% --- 3. Setas de Orientação Perfeitas (Baseadas no crescimento de t) ---
			% Seta 1: Ramo inferior (t se aproximando de zero)
			\draw[->, thick, blue!70!black, domain=-0.95:-0.90, variable=\t] 
			plot ({\t*\t}, {\t*\t*\t});
			
			% Seta 2: Ramo superior (t se afastando de zero)
			\draw[->, thick, blue!70!black, domain=0.90:0.95, variable=\t] 
			plot ({\t*\t}, {\t*\t*\t});
			
			% --- 4. O Ponto Singular (Cúspide) ---
			% Destacando o ponto crítico discutido no seu desenvolvimento formal
			\filldraw[black] (0,0) circle (1.5pt); 
			\node[below left, font=\footnotesize, black, align=right] at (-0.1,-0.1) {Cúspide\\$(0,0)$};
			
			% --- 5. Legenda Identificadora do Conjunto de Nível ---
			\node[anchor=west, font=\small] at (1.5, -1.8) {$S = \Phi^{-1}(0) = \{(x,y) \mid x^3 = y^2\}$};
			
		\end{tikzpicture}
	\end{center}
}

%%%% QUESTÃO 34 %%%%
\questao{Seja $S$ o subconjunto de $\mathbb{R}^3$ definido pelas soluções do sistema 
	\begin{equation*}
		\begin{cases} 
			x^3 + y^3 + z^3 = 1 \\ 
			x + y + z = 0 
		\end{cases}
	\end{equation*}
	$S$ é uma subvariedade mergulhada de $\mathbb{R}^3$?}

\solucao{
	\textbf{Análise do Enunciado:}
	\begin{itemize}
		\item \textbf{Hipóteses:} $S$ é definido como o conjunto de pontos $(x,y,z) \in \mathbb{R}^3$ que satisfazem simultaneamente duas equações polinomiais.
		\item \textbf{Tese:} Verificar se $S$ possui a estrutura de subvariedade mergulhada de $\mathbb{R}^3$.
	\end{itemize}
	
	\textbf{Fundamentação Teórica:}
	\begin{enumerate}
		\item \textbf{Teorema do Valor Regular (Lee, Cap. 7):} Seja $F: M \to N$ uma aplicação suave. Se $c \in N$ é um valor regular de $F$, então o conjunto de nível $F^{-1}(c)$ é uma subvariedade mergulhada de $M$ com dimensão $\dim M - \dim N$.
		\item \textbf{Valor Regular:} Um valor $c$ é regular se, para todo $p \in F^{-1}(c)$, a diferencial $dF_p : T_p M \to T_{F(p)} N$ é sobrejetiva. Para $N = \mathbb{R}^k$, isso equivale a dizer que a matriz Jacobiana de $F$ possui posto $k$ em todos os pontos do conjunto de nível.
	\end{enumerate}
	
	\textbf{Desenvolvimento Formal:}
	
	Definimos a aplicação $F : \mathbb{R}^3 \to \mathbb{R}^2$ tal que o conjunto $S$ seja exatamente o conjunto de nível $F^{-1}(0,1)$. A aplicação é dada por:
	\begin{equation*}
		F(x,y,z) = (x^3 + y^3 + z^3, x + y + z).
	\end{equation*}
	Para verificar se $(0,1)$ é um valor regular, calculamos a matriz Jacobiana $JF$:
	\begin{equation*}
		JF = \begin{pmatrix} 
			\tfrac{\partial F_1}{\partial x} & \tfrac{\partial F_1}{\partial y} & \tfrac{\partial F_1}{\partial z} \\ 
			\tfrac{\partial F_2}{\partial x} & \tfrac{\partial F_2}{\partial y} & \tfrac{\partial F_2}{\partial z} 
		\end{pmatrix} = 
		\begin{pmatrix} 
			3x^2 & 3y^2 & 3z^2 \\ 
			1 & 1 & 1 
		\end{pmatrix}.
	\end{equation*}
	O valor $(0,1)$ é regular se $\text{posto}(JF_p) = 2$ para todo $p \in S$. A matriz $JF_p$ falha em ter posto 2 se, e somente se, suas linhas forem linearmente dependentes. Como a segunda linha $(1,1,1)$ nunca é nula, a dependência linear ocorreria se a primeira linha fosse um múltiplo da segunda:
	\begin{equation*}
		(3x^2, 3y^2, 3z^2) = \lambda(1, 1, 1) \implies 3x^2 = 3y^2 = 3z^2 = \lambda.
	\end{equation*}
	Isso implica que $x^2 = y^2 = z^2$. Portanto, as únicas possibilidades para as coordenadas de $p$ são $x,y,z \in \{ a, -a \}$ para algum $a \in \mathbb{R}$. 
	
	Analisamos agora se algum desses pontos críticos pertence a $S$. Para pertencer a $S$, o ponto deve satisfazer a segunda equação do sistema:
	\begin{equation*}
		x + y + z = 0.
	\end{equation*}
	Se $x^2 = y^2 = z^2 = a^2$, as únicas combinações de sinais que resultariam em soma zero seriam aquelas onde dois valores são iguais e o terceiro é o oposto (ex: $a + a - a = 0$). Contudo, isso implicaria $a = 0$, logo $x=y=z=0$. 
	
	Testamos o ponto $(0,0,0)$ na primeira equação de $S$:
	\begin{equation*}
		0^3 + 0^3 + 0^3 = 0 \neq 1.
	\end{equation*}
	Assim, o ponto $(0,0,0) \notin S$. Concluímos que não existem pontos em $S$ onde a matriz $JF_p$ tenha posto inferior a 2. Portanto, $(0,1)$ é um valor regular de $F$.
	
	Pelo Teorema do Valor Regular, $S = F^{-1}(0,1)$ é uma subvariedade mergulhada de $\mathbb{R}^3$. A dimensão de $S$ é $\dim \mathbb{R}^3 - \dim \mathbb{R}^2 = 3 - 2 = 1$.
}

%%%% QUESTÃO 35 %%%%
\questao{Seja $GL(n, \mathbb{R}) = \{ A \in M(n, \mathbb{R}) \mid \det A \neq 0 \}$ o grupo linear geral das matrizes reais de ordem $n$. O grupo linear especial, denotado por $SL(n, \mathbb{R})$, é formado por todas as matrizes de $GL(n, \mathbb{R})$ cujo determinante é igual a 1. Considere a aplicação $f : GL(n, \mathbb{R}) \to \mathbb{R}$ dada por $f(A) = \det A$. Mostre que 1 é um valor regular de $f$ e conclua que $SL(n, \mathbb{R})$ é uma $(n^2-1)$-subvariedade mergulhada de $GL(n, \mathbb{R})$.}

\solucao{
	\textbf{Análise do Enunciado:}
	\begin{itemize}
		\item \textbf{Hipóteses:} $GL(n, \mathbb{R})$ é um aberto de $M(n, \mathbb{R}) \cong \mathbb{R}^{n^2}$, logo é uma variedade suave de dimensão $n^2$. A aplicação $f(A) = \det A$ é um polinômio nas entradas da matriz, sendo, portanto, suave.
		\item \textbf{Tese:} Provar que $\text{posto}(df_A) = 1$ para todo $A \in SL(n, \mathbb{R})$, e consequentemente que $SL(n, \mathbb{R}) = f^{-1}(1)$ é uma subvariedade mergulhada de dimensão $n^2-1$.
	\end{itemize}
	
	\textbf{Fundamentação Teórica:}
	\begin{enumerate}
		\item \textbf{Teorema do Valor Regular (Lee, Cap. 7):} Se $c$ é um valor regular de $f: M \to N$, então $f^{-1}(c)$ é uma subvariedade mergulhada de $M$ com codimensão $\dim N$.
		\item \textbf{Diferencial do Determinante (Jacobi's Formula):} Para qualquer matriz $A \in GL(n, \mathbb{R})$, a diferencial $df_A : T_A GL(n, \mathbb{R}) \to T_{\det A} \mathbb{R}$ avaliada em um vetor tangente $H \in M(n, \mathbb{R})$ é dada por:
		\begin{equation*}
			df_A(H) = \text{tr}(\text{adj}(A)H),
		\end{equation*}
		onde $\text{adj}(A)$ é a matriz adjunta de $A$.
		\item \textbf{Relação com a Inversa:} Para $A \in GL(n, \mathbb{R})$, sabemos que $\text{adj}(A) = \det(A) A^{-1}$. Assim, a diferencial pode ser escrita como:
		\begin{equation*}
			df_A(H) = \det(A) \text{tr}(A^{-1}H).
		\end{equation*}
	\end{enumerate}
		
	\textbf{Desenvolvimento Formal:}
	
	Seja $A \in SL(n, \mathbb{R})$. Por definição, $\det A = 1$. Substituindo este valor na fórmula da diferencial, temos que para qualquer matriz $H \in M(n, \mathbb{R})$:
	\begin{equation*}
		df_A(H) = (1) \cdot \text{tr}(A^{-1}H) = \text{tr}(A^{-1}H).
	\end{equation*}
	
	Para provar que $1$ é um valor regular, devemos mostrar que a aplicação linear $df_A : M(n, \mathbb{R}) \to \mathbb{R}$ é sobrejetiva. Como o contradomínio é $\mathbb{R}$ (dimensão 1), a sobrejetividade é garantida se $df_A$ não for a aplicação nula. Ou seja, basta encontrar uma única matriz $H \in M(n, \mathbb{R})$ tal que $df_A(H) \neq 0$.
	
	Escolhemos a matriz $H = A$. Como $A$ é inversível, temos:
	\begin{equation*}
		df_A(A) = \text{tr}(A^{-1}A) = \text{tr}(I_n) = n.
	\end{equation*}
	Visto que $n \geq 1$, temos $df_A(A) \neq 0$. Portanto, a diferencial $df_A$ é sobrejetiva para todo $A \in SL(n, \mathbb{R})$, o que prova que $1$ é um valor regular de $f$.
	
	Pelo Teorema do Valor Regular, o conjunto de nível $f^{-1}(1) = SL(n, \mathbb{R})$ é uma subvariedade mergulhada de $GL(n, \mathbb{R})$. A dimensão desta subvariedade é:
	\begin{equation*}
		\dim(SL(n, \mathbb{R})) = \dim(GL(n, \mathbb{R})) - \dim(\mathbb{R}) = n^2 - 1.
	\end{equation*}
}

%%%% QUESTÃO 36 %%%%
\questao{Sejam $S \subset M^n$ uma subvariedade mergulhada e $\Phi : U \subset M^n \to N$ uma aplicação que define localmente $S$. Mostre que $T_p S = \ker(\Phi_* : T_p M \to T_{\Phi(p)} N)$, para todo $p \in S \cap U$.}

\solucao{
	\textbf{Análise do Enunciado:}
	\begin{itemize}
		\item \textbf{Hipóteses:} $S$ é uma subvariedade mergulhada de $M^n$; $\Phi: U \to N$ é uma aplicação suave tal que $S \cap U$ é um conjunto de nível regular de $\Phi$.
		\item \textbf{Tese:} Provar a igualdade de subespaços $T_p S = \ker(\Phi_*)$ para todo $p \in S \cap U$.
	\end{itemize}
	
	\textbf{Fundamentação Teórica:}
	\begin{enumerate}
		\item \textbf{Definição de $T_p S$ (Lee, Cap. 8):} Para uma subvariedade mergulhada $S$, o espaço tangente $T_p S$ é o subespaço de $T_p M$ consistindo nos vetores $\iota_* v$, onde $\iota : S \hookrightarrow M$ é a inclusão e $v \in T_p S$.
		\item \textbf{Aplicação que Define Localmente $S$ (Lee, Cap. 7):} $\Phi$ define localmente $S$ se $S \cap U = \Phi^{-1}(c)$, implicando que $\Phi_*$ é sobrejetiva em $T_p M$.
		\item \textbf{Regra da Cadeia (Lee, Cap. 3):} Se $\gamma$ é uma curva em $S$, a derivada da composição $\Phi \circ \gamma$ é dada pelo push-forward: $(\Phi \circ \gamma)'(0) = \Phi_*(\gamma'(0))$.
	\end{enumerate}
	
	\textbf{Desenvolvimento Formal:}
	
	Para provar a igualdade $T_p S = \ker(\Phi_*)$, provaremos as duas inclusões.
	
	\textit{1. Inclusão $T_p S \subseteq \ker(\Phi_*)$:}
	Seja $v \in T_p S$. Por definição, existe uma curva suave $\gamma : (-\varepsilon, \varepsilon) \to S$ tal que $\gamma(0) = p$ e $\iota_*(\gamma'(0)) = v$. Como a imagem de $\gamma$ está contida em $S = \Phi^{-1}(c)$, temos que $\Phi(\gamma(t)) = c$ para todo $t$. 
	Diferenciando em relação a $t$ no ponto $t=0$:
	\begin{equation*}
		\Phi_*(\iota_* \gamma'(0)) = \Phi_*(v) = 0.
	\end{equation*}
	Portanto, $v \in \ker(\Phi_*)$.
	
	\textit{2. Inclusão $\ker(\Phi_*) \subseteq T_p S$:}
	Utilizamos a comparação de dimensões. Como $\Phi$ define localmente $S$, a dimensão de $S$ é $\dim S = \dim M - \dim N$. Logo, $\dim T_p S = n - \dim N$.
	Pelo Teorema do Núcleo e da Imagem, como $\Phi_*$ é sobrejetiva:
	\begin{equation*}
		\dim(\ker \Phi_*) = \dim T_p M - \text{posto}(\Phi_*) = n - \dim N.
	\end{equation*}
	Como $T_p S \subseteq \ker(\Phi_*)$ e ambos possuem a mesma dimensão finita, eles coincidem.
}
	
\end{document}